




\chapter{Introducción entre-png-jpg - y -  gif} % {\label{Cap1}}\hypertarget{Cap1}{}  %  \hypertarget{Introd}{}\label{Introd}
% https://www.ecured.cu/Python_Game_Library


%  web animada   https://blogthinkbig.com/diferencia-entre-png-jpg-y-gif


%(Para acceder a los eventos de teclado o mouse {\ttfamily pygame} usa la función {\ttfamily pygame.event.get()}, esta función devuelve una lista de eventos que han tenido lugar desde la última vez que se llamó a esta función). 


GUI frente a CLI

Los programas Python que puedes escribir con las funciones integradas de Python solo tratan con texto.
a través de las funciones print() y input(). 

Su programa puede mostrar texto en la pantalla y permitir que el usuario escriba texto desde el teclado.
Este tipo de programa tiene una interfaz de línea de comando, o CLI (que se pronuncia como la primera
parte de “climb” y rima con “sky”). Estos Los programas son algo limitados porque no pueden mostrar gráficos, 
tener colores o usar el ratón. 

Estos programas CLI solo obtienen información del teclado con la función input() y
incluso entonces el usuario debe presionar Enter antes de que el programa pueda responder a la entrada. Esto significa en tiempo real
(es decir, continuar ejecutando código sin esperar al usuario) los juegos de acción son imposibles de ejecutar.
hacer. Pygame proporciona funciones para crear programas con una interfaz gráfica de usuario o GUI
(pronunciado “pegajoso”). En lugar de una CLI basada en texto, los programas con una GUI basada en gráficos pueden
mostrar una ventana con imágenes y colores.


¡Hurra! ¡Acabas de crear el videojuego más aburrido del mundo! Es solo una ventana en blanco con "Hola".
¡Mundo!‖ en la parte superior de la ventana (en lo que se llama la barra de título de la ventana, que contiene el
texto del título). Pero crear una ventana es el primer paso para crear juegos gráficos. Cuando haces clic
en el botón X en la esquina de la ventana, el programa finalizará y la ventana desaparecerá.
Llamar a la función print() para que aparezca texto en la ventana no funcionará porque
print() es una función para programas CLI. Lo mismo ocurre con input() para obtener la entrada del teclado.
del usuario. Pygame utiliza otras funciones para entrada y salida que se explican más adelante en este
capítulo. Por ahora, veamos cada línea de nuestro programa "Hola mundo" con más detalle.




%   https://www.youtube.com/watch?v=oUdd3u7E5ms
%   https://www.youtube.com/watch?v=aBBmR5Tw9yY
%   https://www.youtube.com/watch?v=ADHzoMTTsDg

%%%%%%%%%%%%%%%%%%%%%%%%%%%%%%%%  Inicio   Capitulo  %%%%%%%%%%%%%%%%%%%%%%%%%%



% El codigo ASCII sigla en inglés de American Standard Code for Information Interchange ( Código Estadounidense Estándar para el Intercambio de Información )


%  https://elcodigoascii.com.ar/




%	\addcontentsline{toc}{chapter}{Resumen }	

\newpage

\pagenumbering{arabic}
\setcounter{page}{1}
\pagestyle{plain}


%     https://www.freeconvert.com/es/epub-to-pdf/download


\chapter*{Primera Parte:  FUNDAMENTOS DE PYGAME}%{\label{Cap1}}\hypertarget{Cap1}{}  %  \hypertarget{Introd}{}\label{Introd}

%  Conceptos y módulos básicos 
En proceso....
\\
Comentar el trabajo con IDLE y en Windows...\\





Proyectos  Python, un Juego: Invasión Alienígena
En el primer Capitulo de la Segunda Parte del libro, desarrolláremos  un proyecto de un juego llamado Alienígena Invasión.
Para desarrollar este juego se usará la librería {\ttfamily Pygame}, una colección módulos de Python que facilita la creación juegos,  con sofisticados gráficos, animación e incluso sonido, para que puedas gestión y desarrollar su proyectos.\\
 
%La biblioteca de activos se utiliza junto con Python para llevr a cabo el proceso de programación. Cualquier persona que no tenga conocimientos de los lenguajes de programación generales puede convertir sus ideas de juego en demos jugables utilizando esta aplicación. Además de los juegos, los usuarios también pueden utilizar Pygame para desarrollar aplicaciones multimedia.

Este libro se dirige a todos los makers y entusiastas de la programación que deseen aprender a desarrollar videojuegos con Pygame. También interesará a los principiantes en el lenguaje Python, así como a los informáticos que quieran perfeccionar sus conocimientos sobre esta librería, que se utiliza no solo en el campo de los videojuegos, sino también en el de la simulación.

No es necesario dominar el lenguaje Python para leer este libro, porque el autor comienza presentando los conceptos básicos antes de indicar cómo dar sus primeros pasos con Pygame y detallar la estructura de un juego con Pygame. Posteriormente, a lo largo de los capítulos, se estudian los aspectos principales del desarrollo de videojuegos en dos dimensiones. El lector también va a estudiar cómo gestionar un bucle de juego, cómo dominar los aspectos gráficos con el módulo pygame.draw, cómo añadir sonido con el módulo pygame.mixer, incluso cómo gestionar el tiempo y especialmente las colisiones entre objetos gráficos gracias al concepto de sprite. El autor también ofrece una introducción al 3D y al concepto de motor de juego.

Finalmente, los dos últimos capítulos brindan al lector una documentación precisa de los principales módulos de Pygame que se utilizan en el libro, así como de los que se utilizan con frecuencia en el desarrollo usando Pygame.

Hay elementos adicionales para su descarga en el sitio web www.ediciones-eni.com.


%%%%%%%%%%%%%%    %%%%%%%%%%%%%%%%%   Capitulo 1  %%%%%%%%%%%%%%%%    %%%%%%%%%%%%%%%%%%%%%

%%%%%%%%%%%%%%    %%%%%%%%%%%%%%%%%    %%%%%%%%%%%%%%%%%%    %%%%%%%%%%%%%%%%    %%%%%%%%%%%%%%%%%%%%%

¿Alguna vez has abierto tu computadora y echado un vistazo al interior de la carcasa? No es necesario que lo hagas ahora, pero lo harás.
Descubra que está construido a partir de una serie de piezas necesarias para ofrecer su experiencia informática. la tarjeta de video
genera una imagen y envía una señal a su monitor. La tarjeta de sonido mezcla el sonido y envía
audio a sus parlantes. Luego están los dispositivos de entrada, como el teclado, el mouse y los joysticks, y un
variedad de otros artilugios electrónicos, todos los cuales son esenciales para crear un juego.
En los primeros días de las computadoras domésticas, los programadores con malos cortes de pelo y gafas de montura gruesa tenían
enfrentarse a cada uno de los componentes del ordenador. El programador del juego tuvo que leer la información técnica.
manual de cada dispositivo para escribir el código de computadora para comunicarse con él, todo antes de trabajar
en el juego real. La situación empeoró cuando los fabricantes lanzaron diferentes dispositivos.
y versiones de dispositivos existentes con nuevas capacidades. Los programadores querían soportar tantos dispositivos como fuera posible.
posible, por lo que había un mercado más grande para sus juegos, pero se encontraron empantanados en los detalles
de trabajar con estas nuevas tarjetas gráficas y de sonido. También fue una molestia para el público comprador de juegos, que había
para marcar cuidadosamente la casilla para ver si tenían la combinación correcta de dispositivos para que el juego funcionara.
Las cosas se volvieron un poco más fáciles con la introducción de sistemas operativos gráficos como Microsoft.
Ventanas. Le dieron al programador del juego una forma única de comunicarse con los dispositivos. Significaba que
el programador podría tirar los manuales técnicos porque los fabricantes suministraban drivers, pequeños
Programas que manejan la comunicación entre el sistema operativo y el hardware.
Avancemos rápidamente hasta tiempos más recientes, cuando los programadores todavía tienen malos cortes de pelo pero llantas más delgadas.
sus gafas. La vida de un programador de juegos todavía no es fácil. Aunque existe una forma común de
Al comunicarse con gráficos, audio y entradas, aún puede resultar complicado escribir juegos debido a la variedad
de hardware en el mercado. La PC familiar barata que mamá compró en el supermercado local es muy diferente
de la máquina de alta gama adquirida por un directivo de la empresa. Es esta variedad la que lo hace tan
un esfuerzo para inicializar el hardware y prepararlo para su uso en el juego. Afortunadamente, ahora que Pygame está aquí,
tienen una forma de crear juegos sin tener que preocuparse por estos detalles (y los programadores de juegos tienen tiempo
salir y cortarse el pelo decentemente).
En este capítulo le presentaremos Pygame y le explicaremos cómo usarlo para crear una pantalla gráfica.
y leer el estado de los dispositivos de entrada.

Una palanca de mando o joystick (del inglés joy, alegría, y stick, palo) es un periférico de entrada que consiste en una palanca que gira sobre una base e informa su ángulo o dirección al dispositivo que está controlando.  \\





Blitting

El método de superficie de objetos que probablemente utilizarás con más frecuencia es blit, que es un acrónimo de bit
transferencia en bloque. Blitting simplemente significa copiar datos de imágenes de una superficie a otra. Lo usarás para
¡Dibujar fondos, fuentes, personajes y casi todo lo demás en un juego!

También puede borrar solo una parte de la superficie agregando un objeto de estilo Rect al parámetro que define la
región de origen. Aquí hay dos formas de utilizar el método blit:
\begin{lstlisting}
	pantalla.blit(fondo, (0,0))
\end{lstlisting}
que proyecta una superficie llamada fondo en la esquina superior izquierda de la pantalla. Si el fondo tiene el mismo
dimensiones como pantalla, no necesitaremos llenar la pantalla con un color sólido; l otra manera es
\begin{lstlisting}
	pantalla.blit(ogro, (300, 200), (100*frame_no, 0, 100, 100))
\end{lstlisting}	




\subsection{CREAR IMÁGENES CON UN CANAL ALFA}

Si tomó una fotografía con una cámara digital o la dibujó con algún software de gráficos, probablemente no 
tenga un canal alfa. Agregar un canal alfa a una imagen generalmente implica un poco de trabajo con el software de gráficos.

Para agregar un canal alfa a la imagen de un pez, usé Photoshop, pero también puedes usar otro software, 
como GIMP (www.gimp.org) para hacer esto. 

Para obtener una introducción a GIMP, consulte Beginning GIMP de Akkana Peck:
De principiante a profesional (Apress, 2006). Una alternativa a agregar un canal alfa a una imagen existente es 
crear una imagen con una representación 3D, paquete como 3ds Max de Autodesk o la alternativa gratuita, Blender (www.blender.org). 
Con este tipo de software puede generar directamente una imagen con un fondo invisible (también puede crear varios fotogramas de
animación o vistas desde diferentes ángulos). Esto puede producir los mejores resultados para un juego de apariencia elegante, 
pero puede hacer mucho con la técnica manual del canal alfa. 
Intente tomar una fotografía de su gato, perro o pez dorado y ¡Haciendo un juego con esto!








\subsubsection{Almacenamiento de imágenes}

Hay varias formas de almacenar una imagen en su disco duro. A lo largo de los años, muchos formatos
para archivos de imagen, cada uno con sus pros y sus contras. Afortunadamente un pequeño número
han surgido como los más útiles, dos en particular: JPEG y PNG. Ambos están bien apoyados
en software de edición de imágenes y probablemente no tendrá que utilizar otros formatos para
almacenar imágenes en un juego.\\

• JPEG (Grupo conjunto de expertos en fotografía): los archivos de imágenes JPEG suelen tener la extensión
.jpg o, a veces, .jpeg. Si utiliza una cámara digital, los archivos que produce probablemente
ser archivos JPEG, porque fueron diseñados específicamente para almacenar fotografías. Ellos usan un
proceso conocido como compresión con pérdida, que es muy bueno para reducir el tamaño del archivo. 
La desventaja de la compresión con pérdida es que puede reducir la calidad de la imagen, pero generalmente
es tan sutil que no notarás la diferencia. 
La cantidad de compresión puede también se puede ajustar para lograr un equilibrio entre calidad visual y compresión. Pueden ser
excelentes para fotografías, pero los archivos JPEG son malos para cualquier cosa con bordes duros, como fuentes o diagramas,
porque la compresión con pérdida tiende a distorsionar este tipo de imágenes. Si usted tiene
para cualquiera de estos tipos de imágenes, {\bf PNG probablemente sea una mejor opción}.\\


• PNG (Portable Network Graphics): los archivos PNG son probablemente las imágenes más versátiles, porque  el
formato pueden almacenar una amplia variedad de tipos de imágenes y aún así comprimirse muy bien.
También admiten canales alfa, lo cual es una verdadera ayuda para los desarrolladores de juegos. 
La compresion que utilizan los archivos PNG no tiene pérdidas, lo que significa que las imágenes almacenadas como archivos PNG se
exactamente igual que los originales. La desventaja es que incluso con una buena compresión probablemente será más grande que los archivos JPEG. Además de JPEG y PNG, Pygame su


\chapter{Introducción a la Librería Pygame} % {\label{Cap1}}\hypertarget{Cap1}{}  %  \hypertarget{Introd}{}\label{Introd}
                    % https://www.ecured.cu/Python_Game_Library
\vspace{-1.2cm}                    
%\begin{wrapfigure}{l}{0.5\textwidth} 
\begin{center} 
	\includegraphics[width=0.99\textwidth]{../A Logos de Python/pygame04}
\end{center} %\caption{Python 3} 
%\end{wrapfigure}               
Pygame es una librería o paquete  de Python para crear juegos, se basa en otra librería  llamada Simple DirectMedia Layer (SDL), que  fue escrita por Sam Lantinga mientras trabajaba para Loki Software  o Loki Games (una empresa de juego  desaparecida en 2002) para simplificar la tarea de portar juegos de una plataforma a otra. Proporcionó una forma común de 
crear una visualización en múltiples plataformas, así como trabajar con gráficos y  dispositivos de entradas. 

Debido a su sencillez para trabajar con él, se hizo muy popular entre los desarrolladores de juegos. 
Cuando se lanzó en 1998 y desde entonces se ha utilizado para muchos juegos comerciales y de pasatiempo.
SDL fue escrito en C, un lenguaje comúnmente utilizado para juegos debido a su velocidad y capacidad
para trabajar con el hardware. Pero desarrollar en C, o su sucesor C++, puede ser lento y propenso a errores. 
Entonces los programadores produjeron enlaces a sus lenguajes favoritos, y SDL ahora se puede utilizar desde casi
cualquier lenguaje que exista. Uno de esos enlaces es Pygame, que permite a los programadores de Python
 utilizar la potente biblioteca SDL.

Ambos de código abierto tanto Pygame como SDL, y han estado en desarrollo activo, una gran cantidad de programadores
han trabajado para refinar y mejorar esta excelente herramienta para crear juegos.


\section{Pygame}             % funciones integradas de Python


Pygame es un conjunto de módulos de Python para desarrollar videojuegos. 
Este programa de desarrollo e informática proporciona a los usuarios una biblioteca Simple DirectMedia Layer que contiene 
gráficos de ordenador y clips de sonido.  La principal ventaja de este software es su portabilidad y su soporte
multiplataforma. Con esta aplicación, los usuarios pueden distribuir sus juegos basados en Python y probar los trabajos de otros.\\







En las siguientes secciones se presenta un  breve tutorial o guía sobre  las principales características y funciones de Pygame.
En caso de que no hayas instalado Python ni Pygame en tu computadora, en la próxima sección se dan las instrucciones de instalación, si
 ya ha instalado ambos entonces puedes saltara y continuar en la siguiente.\\

La librería  Pygame al igual que Python viene con varios módulos  que proporcionan
funciones adicionales para sus programas; incluye varios módulos con
funciones para dibujar gráficos, reproducir sonidos, manejar la entrada del mouse y otros.

Este capítulo se describen algunos los módulos y funciones básicos que Pygame proporciona, y se asume que conoce un poco de la programación 
básica de Python. Si tiene problemas con algunos conceptos de programación, puedes leer el libro "Inventa tus propios juegos de computadora con Python", en línea en http://invpy.com/book. Este libro está dirigido a principiantes totales en programación; o revisare le Libro Python 3 de Linis Guerrero, una compilación de varios textos.  

Para obtener información más detallada sobre Pygame, sobre los módulos que proporciona Pygame se recomienda revisar 
la documentación en línea en la web \href{www.pygame.org/docs/}{\color{vino}www.pygame.org/docs/}. 



\subsection{Instalación de la Librería PyGame}

Para descargar e instalar la librería PyGame en su sistema operativo puede visitar la web  https://www.pygame.org/news. En esta web encontrara el método para descargar la versión que desee instalar en su sistema. Además en esta página podrás conocer más sobre pygame sus módulos  y las nuevas noticias acerca de sus versiones.\\


Por supuesto que PyGame requiere de Python (Python 3); Si aún no lo tienes, puedes descargarlo desde https://www.python.org/. Se recomienda ejecutar la última versión de Python, porque suele ser más rápida y tiene mejores funciones que las anteriores.\\

La mejor manera de instalar pygame es con la herramienta pip (que es la que usa Python para instalar paquetes), en la consola o terminal mcd (símbolo del sistema). Tenga en cuenta que esta herramienta viene con Python en versiones recientes. Usamos el indicador --user para indicarle que se instale en el directorio de inicio, en lugar de hacerlo globalmente.

\begin{lstlisting}
python -m pip install -U pygame==2.5.2 --user	
\end{lstlisting}

Una vez que haya instalado el paquete apropiado, puede probarlo abriendo el intérprete de Python
e ingresando las siguiente  línea de código:

\begin{lstlisting}
>>> import pygame
\end{lstlisting}

Si Pygame se instaló correctamente, debería ver la versión mostrada:
{\ttfamily\scriptsize\scriptsize 
	\hspace{-0.65cm} \vspace{-0.3cm}	
	{\color{blue} 	
\begin{verbatim}
pygame 2.5.2 (SDL 2.28.3, Python 3.12.0)
Hello from the pygame community. https://www.pygame.org/contribute.html
\end{verbatim}
}}
 
Ademas puede ver la versión de pygame con {\ttfamily print(pygame.ver)}:

\begin{lstlisting}
>>> print(pygame.ver)
2.5.2	
\end{lstlisting}

Al momento de escribir este artículo, la versión 2.5.2 es la más reciente, pero es posible que encuentre una versión posterior,
que seguro estará disponible cuando estés  leyendo este libro. El código de ejemplo seguirá funcionando, ya que las versiones más nuevas
de Pygame tienden a ser compatibles con versiones anteriores. (Tenga en cuenta que pygame ha dejado de admitir Python 2).


\subsection{Módulos de la Librería PyGame}   %  duende  sprite      & Gestiona imágenes en movimiento


El paquete  PyGame contiene una serie de módulos que se pueden utilizar de forma independiente. Hay
un módulo para cada uno de los dispositivos que puedes usar en un juego, y muchos otros para crear juegos
de manera muy sencilla. 

Siempre tendrás algún tipo de exhibición, por lo tanto, el módulo de visualización es esencial y 
definitivamente necesitará algún tipo de información, ya sea del teclado, joystick o mouse. Otros módulos se utilizan con menos frecuencia, pero en combinación te brindan una de las herramientas de creación de juegos más poderosas que existen.

\noindent En la Tabla \ref{mdulos} se presentan algunos módulos de Pygame:\vspace{0.2cm}


{  }Tabla \ref{mdulos}. Módulos de la librería Pygame:
\vspace{-0.4cm}	
\begin{center}	
	{\footnotesize{{%\scriptsize 
				\hspace{-0.7cm}	
				\begin{tabular}{l l l }	% white
				\multicolumn{1}{|>{\columncolor{vino}} l }{\color{white}\bf Nombre}& 
				\multicolumn{1}{|>{\columncolor{vino}} l }{\color{white}\bf Propósito del módulo}\\ 
				pygame.cdrom     & Accede y controla unidades de CD    \\
				pygame.cursors   & Carga imágenes del cursor       \\
				pygame.display	 &  Accede a la pantalla   \\
				pygame.draw      & Dibuja formas, líneas y puntos  \\
				pygame.event     & Gestiona eventos externos    \\
				pygame.font      &Utiliza fuentes del sistema  \\
				pygame.image     & Carga y guarda una imagen   \\ 
				pygame.joystick  &  Utiliza joysticks y dispositivos similares   \\ 
				pygame.key       & Lee las pulsaciones de teclas del teclado   \\    
				pygame.mixer     & Carga y reproduce sonidos   \\
				pygame.mouse     &  Administra el mouse  \\ 
				pygame.movie     & archivos de película   \\    
				pygame.music     &  Funciona con música y transmisión de audio   \\  
				pygame.overlay    &  Accede a superposiciones de vídeo avanzadas  \\ 
				                   &Contiene funciones de Pygame de alto nivel\\
				pygame.rect       &  Gestiona áreas rectangulares   \\   
				pygame.sndarray    &  manipula datos de sonido   \\ 
				pygame.sprite      & Gestiona imágenes en movimiento   \\    
				pygame.surface     & Gestiona las imágenes y la pantalla   \\   
				pygame.surfarray   &  Manipula los datos de píxeles de la imagen   \\   
				pygame.time        & Gestiona el tiempo y la velocidad de fotogramas  \\ 
				pygame.transform   & Cambia el tamaño y mueve imágenes  \\ 				                  
				\arrayrulecolor{vino}\hline 
				\label{mdulos} 	   
				\end{tabular}		
				%  	\vspace{0.2cm} 
	}}} 
\end{center}


Sin embrago, no hay garantía que todos los módulos de la Tabla \ref{mdulos} estén presentes para todas las plataformas. Es
posible que el hardware en el que se ejecuta el juego no tenga ciertas capacidades, o que no estén instalados los controladores (drivers) necesarios. Si este es el caso, Pygame configurará el módulo como {\ttfamily None}, lo que facilita la prueba. \\

El siguiente fragmento de código detectará si el módulo {\ttfamily pygame.font} está disponible o no:
\begin{lstlisting}
if pygame.font is None:
    print("El módulo font (fuentes) no esta disponible!")
    exit()	
\end{lstlisting}

A través de el espacio de nombres de {\ttfamily pygame} se accede a estos módulos; por ejemplo, {\ttfamily pygame.display} se refiere al módulo de visualización.

\section{Mi primer Juego}

Siguiendo la tradición de usar el popular "¡Hola, mundo!"{} al iniciarse en un nuevo lenguaje. En nuestro primer juego o programa con {\ttfamily pygame} que vamos a crear, usaremos el texto "¡Hola, Mundo Pygame!"{}. En este caso escribimos el texto "¡Hola, Mundo Pygame!"{} en la pantalla como nombre del juego. No se desanime, a  pesar de ser una línea de texto simplemente, y no  muy interesante para empezar a crear juegos, iniciamos con una ventana ventana gráfica en su escritorio y luego seguiremos con  imágenes atractivas. \\

Vamos a crear nuestro primer archivo o script {\ttfamily mi\_primer\_pygame.py},  que abre una ventana gráfica en su escritorio
con el nombre "¡Hola, Mundo Pygame!"{}, al inicio con un fondo de pantalla negro, pero le iremos agregando elementos para hacerlo más interesante.
 
 
\hspace{-2.9cm}{\ttfamily \scriptsize mi\_primer}
 
\vspace{-0.2cm}
\hspace{-2.5cm}{\ttfamily \scriptsize \_pygame.py}\vspace{-0.84cm}
\begin{lstlisting}	
import pygame
from pygame.locals import *
from sys import exit
 	
pygame.init()
screen = pygame.display.set_mode((640, 480), 0, 32)
pygame.display.set_caption("¡Hola, Mundo Pygame!")
 	
while True:   # bucle principal del juego
    for in pygame.evento.get():  # bucle de eventos
        if evento.type == QUIT:
            pygame.quit()
            exit()
    
    pygame.display.update()	
\end{lstlisting}



% screen.blit(background, (0,0))   


Cuando ejecutes este programa, aparecerá una ventana como la ventana que se muestra en la figura \ref{Hola_mundo}, pero con fondo negro.
Hasta ahora, muy bien, ¡Acabas de crear un videojuego muy simple!. Es solo una ventana con fondo negro, con el nombre  ¡Hola, Mundo Pygame! en la parte superior (se le llama la barra de título de la ventana,  por defecto {\ttfamily windows pygame}).\\
 
El primer paso para crear juegos gráficos es crear una ventana; y  finalizar el juego haciendo clic en el botón {\ttfamily X} en la esquina superior derecha de la ventana, el programa finalizará y la ventana desaparecerá.

\vspace{0.2cm}
\begin{minipage}[t]{0.2\textwidth}
	\centering\raisebox{\dimexpr\topskip-\height}{%
		\includegraphics[scale=0.20]{../A Logos de Python/Advertencia1}}  %[width=\textwidth]
	
\end{minipage}%\hfill
\begin{minipage}[t]{0.70\textwidth} 
	{\scriptsize Llamar a la función {\ttfamily print()} para que aparezca texto en la ventana no funcionará porque
		print() es una función para programas CLI (interfaz de línea de comando); lo mismo sucede  con {\ttfamily input()} para obtener la entrada del teclado del usuario. 
		
		Pygame utiliza otras funciones para entrada y salida que se explican más adelante.}
\end{minipage}
\vspace{0.2cm} 

Ahora continuamos describiendo  con mayor detalle cada línea del código fuente de nuestro programa.\\

% Configurar un programa Pygame   \ding{51}

 • La primera línea es una declaración simple que importa la librería {\ttfamily pygame} para que nuestro
programa pueda utilizar los módulos y las funciones que contiene. Tenga en cuenta que cuando importa el módulo pygame, se importa automáticamente todos los módulos están  Pygame, como pygame.images y pygame.mixer.music. Todas las funciones de Pygame relacionadas con gráficos, sonido, y otras características que proporciona Pygame se encuentran en el módulo pygame.
No es necesario importar estos módulos dentro de módulos con declaraciones de importación adicionales.


 • La segunda línea también es una declaración de importación. Sin embargo, en lugar del formato de nombre del módulo de importación,
utiliza el formato  {\ttfamily from modulename import *}:  {\ttfamily from pygame.locals import *}.

Normalmente, si desea llamar a una función que está en un módulo, debe usar el formato {\ttfamily modulename.functionname()}
después de importar el módulo.  Sin embargo, con {\ttfamily from modulename import *}, puede omitir  parte del nombre del módulo y simplemente usar el  {\ttfamily functionname()}, semejante con las funciones integradas de Python.

La razón por la que utilizamos esta forma de declaración de importación para {\ttfamily pygame.locals} es porque
contiene varias variables constantes que son fáciles de identificar como parte del módulo {\ttfamily pygame.locals} sin {\ttfamily pygame.locals} en frente de ellos.

Este módulo contiene varias constantes usadas por {\ttfamily pygame}{\footnote{Para obtener documentación completa sobre los módulos de Pygame, consulte \href{www.pygame.org/docs/}{\color{vino}www.pygame.org/docs/}.}}. Su contenido se coloca automáticamente en el espacio de nombres (namespace) del módulo {\ttfamily pygame}.  Sin embargo, en una aplicación se puede usar {\ttfamily pygame.locals} para  incluir solo las constantes de {\ttfamily pygame} con la sentencia:
\begin{lstlisting}
 from pygame.locals import *	
\end{lstlisting}
y usar constantes como: mejorar y completar referencias
{\ttfamily QUIT} del módulo  {\ttfamily pygame.event}; el módulo {\ttfamily pygame.key}  enumera las constantes y
 modificadores del teclado ( K\_* y MOD\_*) relacionados con los atributos {\ttfamily key} y {\ttfamily mod} de los
  eventos {\ttfamily KEYDOWN} y {\ttfamily KEYUP};  el módulo	 {\ttfamily pygame.display}  contiene indicadores como
   {\ttfamily FULLSCREEN} utilizados por {\ttfamily pygame.display.set\_mode()}; el módulo	{\ttfamily pygame.time}, 
 define {\ttfamily TIMER\_RESOLUTION}.\\





Las  dos primeras líneas de código en el programa  son líneas que comenzarán casi en todo programa que escribes con Pygame.\ding{51}

• La tercera línea se importa  {\ttfamily exit} de {\ttfamily sys} (un módulo en la biblioteca estándar),  para cerrar la ventana pygame, luego de cerrar  pygame, al recibir el objeto evento {\ttfamily QUIT} (al llamarlo hará que la ventana de Pygame desaparezca). Si usa {\ttfamily import sys} en lugar de  {\ttfamily from sys import *} en su programa, tendrías que llamar a {\ttfamily sys.exit()}   en lugar de {\ttfamily exit()}.  Sin embrago,  la mayoría de las veces es mejor usar el nombre completo del módulo y la función para tener claro en qué módulo se usa y evitar confusiones.

• La cuarta línea es el llamado a la función {\ttfamily pygame.init()}, que siempre debe llamarse después de importar
el módulo {\ttfamily pygame} y antes de llamar a cualquier otra función de{\ttfamily Pygame}.

Inicializa el modulo {\ttfamily pygame.init()},  debe llamarse primero para que muchas funciones de Pygame
trabajen.   Si alguna vez ve un mensaje de error como {\ttfamily pygame.error}: fuente no
inicializado, verifique que seguro olvidó llamar a {\ttfamily pygame.init()} al comienzo de su programa.


\vspace{0.2cm}
\begin{minipage}[t]{0.2\textwidth}
	\centering\raisebox{\dimexpr\topskip-\height}{%
		\includegraphics[scale=0.20]{../A Logos de Python/Advertencia}}  %[width=\textwidth]
	
\end{minipage}%\hfill
\begin{minipage}[t]{0.70\textwidth} 
	{\scriptsize El script llamará a salir cuando el usuario haga clic en el botón cerrar; de lo contrario, el
		¡El usuario no tendría forma de cerrar la ventana!
		Si te encuentras en una situación en la que no puedes cerrar la ventana de Pygame, 
		es posible que puedas detener Python en seco presionando {\ttfamily Ctrl+C}. 
		Esta línea bastante simple de código Python en realidad hace un  trabajo muy importante}
\end{minipage}
\vspace{0.2cm}


• La quinta línea es una llamada a la función {\ttfamily pygame.display.set\_mode()}, que devuelve el
objeto {\ttfamily pygame.Surface} para la ventana (los objetos de superficie se describen más adelante en la sección \ref{Surface}). 

Primer parámetro: 

Observe que pasamos un valor de tupla de dos números enteros (640, 480) a la función, que le indica a la 
el ancho y el alto de la ventana en píxeles. En esta caso creará una ventana con 640 píxeles 
de ancho  y 640 píxeles de alto. 

Una llamada a función como {\ttfamily pygame.display.set\_mode(640, 480)} provocará un error que
se ve así: {\ttfamily TypeError}: el argumento debe ser una tupla de 2 elementos, no entero.
El objeto {\ttfamily pygame.Surface} (los llamaremos simplemente objetos Surface para abreviar) devuelto es
almacenado en una variable llamada {\ttfamily DISPLAYSURF}. Lo primero que hacemos con nuestra superficie recién creada es agregarle el nombre con el método {\ttfamily set\_caption} en el módulo de visualización para configurar la barra de título del ventana de Pygame. Configuramos el título como {}"¡Hola, Mundo Pygame!"{}, solo para que sea un guión válido del juego.  \\


Segundo parámetro: 

El parámetro {\ttfamily 32}: Esta es la profundidad de color de la superficie de visualización. Especifica el número de bits utilizados para representar cada píxel en la superficie de visualización, utilizados para almacenar
colores en la pantalla. En este caso, 32 bits se utilizan, lo que significa que cada píxel puede tener 2exp32 (o 4.294.967.296) colores diferentes. 

Un bit, o dígito binario, es la unidad de almacenamiento  fundamental en una computadora. Los bits tienen exactamente dos valores potenciales, 1 o 0, y están organizados en la memoria como grupos de 8. Un grupo de 8 bits se llama byte. 
No te preocupes si esto te suena a tecno-charla.  Python tiende a ocultar este tipo de cosas al programador.
Usaremos el valor 32 para nuestra profundidad de bits porque nos da la mayor cantidad de colores.


En la tabla \ref{tabla_bits} se resume los valores de profundidad de color en bits, puede consulte  para conocer otros posibles valores de profundidad en bits.
\vspace{0.2cm}






Tabla \ref{tabla_bits}. Valores de profundidad en bits:
\vspace{-0.65cm}
\begin{center}	
	{\footnotesize{{%\scriptsize 
				\hspace{-0.7cm}	
				\begin{tabular}{l l l }
					\multicolumn{1}{|>{\columncolor{vino}} l }{\color{white}\bf Profundidad}& 
					\multicolumn{1}{|>{\columncolor{vino}} l }{\color{white}\bf Número de colores}  \\ 
					8 bits  & 256 colores, elegidos entre una paleta de colores más amplia  \\				
					15 bits   &  32,768 colores,  con un   bit libre \\
					16 bits 	&  65.536 colores \\				
					24 bits  	& 132.768 colores, con un  a spare bit\\
					24 bits & 16.7 millones de colores \\
					32 bits  &  16,7 millones de colores, con un  8  bit libre \\										                  
					\arrayrulecolor{vino}\hline 
					\label{tabla_bits} 
				\end{tabular}		
				%  	\vspace{0.2cm} 
	}}} 
\end{center} Es posible tener otras profundidades de bits, pero estas son las más común.


\vspace{0.2cm}
\begin{minipage}[t]{0.2\textwidth}
	\centering\raisebox{\dimexpr\topskip-\height}{%
		\includegraphics[scale=0.20]{../A Logos de Python/Advertencia}}  %[width=\textwidth]
	
\end{minipage}%\hfill

\begin{minipage}[t]{0.70\textwidth} 
	{\scriptsize Algunas  veces, Pygame no puede generar la visualización exacta que configuramos, tal vez se debe a que la tarjeta gráfica
		no admite esa configuración.  Afortunadamente, Pygame elegirá una pantalla que sea compatible 
		con el hardware y emula la pantalla que realmente hemos configurado.}
\end{minipage}
\vspace{0.2cm}


Hay más parámetros indicadores para la función {\ttfamily .set\_mode()}. Los  parámetros indicadores (flags) de  {\ttfamily pygame.display.set\_mode()}, que activar o desactivar una función de la ventana Pygame; y se puede combinar varios parámetros indicadores  con el operador bitwise OR (|) (bit a bit).

Pero por ahora no los trataremos, solo se dará el valor de 0, que el valor por defecto. En la sección  \ref{pantalla} se gestiona los parámetros {\ttfamily  NOFRAME}, {\ttfamily  RESIZABLE}, {\ttfamily  DOUBLEBUF} y  {\ttfamily  HWSURFACE}, para tener una idea puede consultar la tabla \ref{tabla_Flags}, que resume los parámetros indicadores de visualización de la función {\ttfamily pygame.display.set\_mode()} que puede utilizar.  \vspace{0.2cm}



Table \ref{tabla_Flags}. Parámetros indicador de visualización:
\begin{center}\vspace{-0.2cm}	
	{\footnotesize{{%\scriptsize 
				\hspace{-0.7cm}	
				\begin{tabular}{l l l }
					\multicolumn{1}{|>{\columncolor{vino}} l }{\color{white}\bf Parámetro indicador }& 
					\multicolumn{1}{|>{\columncolor{vino}} l }{\color{white}\bf Propósito}  \\ 
					FULLSCREEN  & Crea una visualización que ocupa toda la pantalla  \\				
					DOUBLEBUF   &  Crea una visualización de "{}doble búfer"{}  \\					
					&  Recomendado para HWSURFACE u OPENGL \\
					HWSURFACE 	& Crea una pantalla acelerada por hardware \\
					& Debe combinarse con  FULLSCREEN\\
					OPENGL      &  Crea una visualización renderizable OpenGL \\
					RESIZABLE   & Crea una visualización redimensionable \\
					NOFRAME     & Elimina el borde y la barra de título de la ventana\\							                  
					\arrayrulecolor{vino}\hline 
					\label{tabla_Flags} 
				\end{tabular}		
				%  	\vspace{0.2cm} 
	}}} 
\end{center}


• La línea 6 establece el texto del título que aparecerá en la parte superior de la ventana al llamando de la función {\ttfamily pygame.display.set\_caption()}. El valor de cadena {}"¡Hola, Mundo Pygame!"{}  es pasado en el llamado de función para que el texto aparezca como título de la ventana.
 
•  La línea 7 es un bucle {\ttfamily while} que tiene como condición simplemente el valor {\ttfamily True}, es el bucle principal del juego. El trabajo de este bucle es repetir continuamente el bloque; es decir, nunca sale del bucle, pues su condición para salir se evalúa como False. La única forma de salir del bucle es  ejecutando una instrucción {\ttfamily break} (que mueve la ejecución a la primera línea después de la bucle), o forzarlo a salir de alguna otra manera, en este caso  sale con {\ttfamily  exit()} ( {\ttfamily sys.exit()}) que finaliza el programa al hacer clic en el botón {\ttfamily X} en la esquina superior derecha de la ventana, el programa finalizará y la ventana desaparecerá. 

Todos los juegos tendrán un bucle similar a este, que normalmente se repite una vez por actualización de pantalla.
Dentro del bucle principal del juego tenemos otro bucle: el bucle de eventos, que la mayoría de los juegos ejecutarán. Todos los juegos de este libro tienen estos bucles {\ttfamily while True} junto con un comentario que llama
es el “bucle principal del juego”. 

Un bucle de juego, también llamado bucle principal es un bucle donde el código hace tres cosas:

1. Maneja eventos.

2. Actualiza el estado del juego.

3. Dibuja el estado del juego en la pantalla.\\

El estado del juego es simplemente una forma de referirse a un conjunto de valores para todas las variables de un juego.
En muchos juegos, el estado incluye los valores de las variables que rastrean el la salud y la posición del jugador, 
la salud y la posición de los enemigos, qué marcas se han realizado en un tablero, la puntuación o a quién le toca. 
Cada vez que sucede algo como el jugador recibe un daño (lo que reduce su valor de salud), o un enemigo se mueve a algún lugar, o algo
sucede en el mundo del juego, decimos que el estado del juego ha cambiado.

Si alguna vez has jugado un juego que te permite guardar, el {}"{}estado de guardado"{} es el estado del juego en el punto en el que
lo has guardado. En la mayoría de los juegos, pausar el juego evitará que cambie el estado del juego.
Dado que el estado del juego generalmente se actualiza en respuesta a eventos (como clics del mouse o pulsaciones del teclado)
o el paso del tiempo, el bucle del juego está constantemente comprobando y volviendo a comprobar muchas veces
un segundo para cualquier evento nuevo que haya ocurrido. 

\subsection*{Gestión de eventos externos}

Dentro del bucle principal {\ttfamily while} del juego, hay  otro bucle: el bucle de eventos ({\ttfamily for event in pygame.event.get()}), que analiza si han creado eventos y se hace con  función {\ttfamily pygame.event.get()} (mayoría de los juegos ejecutarán).
\begin{lstlisting}
while True:   # bucle principal del juego
	for evento in pygame.event.get(): # bucle de eventos
		if evento.type == QUIT:
			pygame.quit()
			exit()	
\end{lstlisting}

Un evento es la forma en que {\ttfamily Pygame}  informa que algo sucedió fuera de su código. Los eventos son
creados para muchas cosas, desde presionar teclas hasta recibir información de Internet, y están
en cola para que  Usted  los manejes. La función {\ttfamily get()} en el módulo {\ttfamily pygame.event} devuelve
cualquier evento que haya ocurrido, y que se puede recuperar con un bucle {\ttfamily for}.\\ 

En este {\ttfamily script}, solo estamos interesado en el evento {\ttfamily QUIT}, que es generado por Pygame cuando el usuario hace clic en el botón cerrar en la ventana de Pygame (si no se incluye tendremos un juego inmortal, una ventana que no cierra). Entonces, si el tipo de evento es {\ttfamily QUIT}, se llama a {\ttfamily exit()} para cerrar, y todos los demás eventos son ignorados. 

En un juego, por supuesto, tendríamos que gestionar una mayor cantidad de eventos. Pygame crea otros eventos para informarle sobre cosas como el movimiento del mouse y la pulsación de teclas (en la sección \ref{eventos} se estudian con mas detalles los eventos).



•  La siguiente línea, el bloque del condicional {\ttfamily if evento.type == QUIT}, que evalúa si el evento es la salida,  en caso afirmativo cierran el juego: Es importante llamar a {\ttfamily pygame.quit()} y {\ttfamily exit()} para cerrar completamente {\ttfamily Pygame}. La función {\ttfamily pygame.quit()}  es algo así como lo opuesto a la función {\ttfamily pygame.init()}, al ejecutarla desactiva la librería  {\ttfamily  pygame} (cierra {\ttfamily pygame}), y  {\ttfamily exit()} cierra de la ventana.\\

\subsection*{Bucle principal del juego: {\ttfamily while}}

•  En la última línea del bloque {\ttfamily while}, tenemos el código que evita los destellos de pantalla cuando creas una image: {\ttfamily pygame.display.update()}.  Esto es porque Pygame primero construye una 
imagen en un buffer de respaldo, que es una visualización invisible en la memoria, antes de ser desplegado.

Si no tuviéramos este paso, el usuario vería destellos individuales a medida que ocurren, parpadearían
de lo más desagradable. Para los programadores de juegos, ¡el parpadeo es el enemigo!. Afortunadamente, una llamada a {\ttfamily pygame.display.update()} es todo lo que necesitamos para asegurarnos de que la imagen que tenemos creado en la memoria se muestra al usuario sin parpadeo.

% \section{Mi primer Juego}

\section{Actualizando Mi primer Juego}



Es hora de hacer un juego mas atractivo.  Ahora le agregaremos dos imágenes, una para usar como fondo de pantalla {\ttfamily pygame.jpg},
y la otra {\ttfamily cubo.jpg} que usaremos puntero nuestro el mouse (cursor del ratón).
   
Puede descargar estos archivos para este y otros ejemplos desde código fuente o crearlos para a su necesidad. Si no tienes acceso a Internet en este momento, puedes puede usar archivos de imágenes que tenga en su disco duro o crearlos con cualquier edición gráfica o fotográfica software. Cualquier imagen está bien para el fondo, siempre que tenga un tamaño mínimo de 640 por 480, que es el tamaño de pantalla que hemos creado (si es más grande, se recortará el exceso). Para el cursor del mouse, necesitará una imagen más pequeña, que encaje cómodamente dentro del fondo; un buen tamaño es 80 por 80. 


\hspace{-2.9cm}{\ttfamily \scriptsize mi\_primer}

\vspace{-0.2cm}
\hspace{-2.5cm}{\ttfamily \scriptsize \_pygame.py}\vspace{-0.84cm}
\begin{lstlisting}	
background_image_filename = 'pygame.jpg'
mouse_image_filename = 'cubo.jpg'

import pygame
from pygame.locals import *
from sys import exit

pygame.init()
screen = pygame.display.set_mode((640, 480), 0, 32)
pygame.display.set_caption("¡Hola, Mundo Pygame!")

background = pygame.image.load(background_image_filename).convert()
mouse_cursor = pygame.image.load(mouse_image_filename).convert_alpha()

while True:
	for evento in pygame.event.get():
		if evento.type == QUIT:
		exit()
	
	screen.blit(background,(0,0))
	x, y = pygame.mouse.get_pos()
	x-= mouse_cursor.get_width()/2
	y-= mouse_cursor.get_height()/2
	screen.blit(mouse_cursor, (x, y))
	pygame.display.update()	
\end{lstlisting}




Al ejecutar el archivo, obtendrá una ventana similar a la que se muestra en la Figura \ref{Hola_mundo}.

\begin{figure}[h!]
	\centering
	\includegraphics[scale=0.40]{../A Logos de Python/hola_pygame}
	\caption{Hola, Mundo en Pygame}
	\label{Hola_mundo}
\end{figure} 

\vspace{0.2cm}


• Las dos primeras líneas establecen los nombres de los archivos de imágenes; si está utilizando imágenes diferentes, debe reemplazar los nombres de los archivos con la ubicación de mismo.


• El siguiente par de líneas nuevas en el programa se usa la función {\ttfamily pygame.image.load()} para cargar las dos imágenes,
para el fondo y para el cursor del ratón. Pasamos los nombres de los archivos de las imágenes almacenadas al inicio del {\ttfamily script}:

\begin{lstlisting}	
background = pygame.image.load(background_image_filename).convert()
mouse_cursor = pygame.image.load(mouse_image_filename).convert_alpha()
\end{lstlisting}

La función de carga lee un archivo de su disco duro y devuelve una superficie que contiene datos de la imagen.
Estos son el mismo tipo de objetos que nuestra pantalla, pero representan imágenes almacenadas en la memoria. 

La primera llamada a {\ttfamily pygame.image.load()} lee la imagen de fondo y luego llama inmediatamente a la 
función {\ttfamily convert()}, que es una función miembro para objetos de superficie. Esta función convierte la imagen
al mismo formato que nuestra pantalla, por ser más rápido dibujar imágenes si la pantalla tiene la misma profundidad.

El cursor del mouse se carga de manera similar, pero llamamos la funcion {\ttfamily convert\_alpha()} en lugar de {\ttfamily convert()}.
Esto se debe a que la imagen del cursor del mouse contiene información alfa, lo que significa que partes de la imagen pueden 
ser translúcidas o completamente invisible. Sin información alfa en la imagen de nuestro mouse, estamos limitados a un
Cuadrado o rectángulo antiestético como el cursor de nuestro mouse! (en la sección \ref{formato} se  cubrirá alfa y los 
formatos de imagen con más detalle).\\








Siguiendo con la descripción del código fuente del programa: en el bloque del {\ttfamily while} (bucle principal) y después del bucle de eventos, se agrega una línea de codigo :  {\ttfamily  screen.blit(background, (0,0))}, esta función transfiere la imagen de fondo a la pantalla ("{}blit"{} significa copiar  una imagen sobre otra).  

• Esta línea utiliza la función  {\ttfamily blit()}  miembro el objeto Surface de pantalla, que toma una 
imagen (en este caso, nuestro fondo de 640 × 480) y una tupla que contiene la posición de destino. El fondo nunca se moverá; solo queremos que cubra toda la ventana de Pygame, así que nos dirigimos a la coordenada (0, 0), que está en la parte superior izquierda de la pantalla.



\vspace{0.2cm}
\begin{minipage}[t]{0.2\textwidth}
	\centering\raisebox{\dimexpr\topskip-\height}{%
		\includegraphics[scale=0.20]{../A Logos de Python/Advertencia1}}  %[width=\textwidth]
	
\end{minipage}%\hfill
\begin{minipage}[t]{0.70\textwidth} 
	{\scriptsize Es importante que la imagen abarque cada parte de la pantalla. Si no lo hace, pueden producirse efectos visuales extraños
		cuando presenta animaciones, y el juego puede verse diferente en cada computadora en la que se ejecuta. }
\end{minipage}
\vspace{0.2cm} 



•  En las siguientes líneas código son para dibujar {\ttfamily mouse\_cursor} debajo del puntero habitual del mouse:
\begin{lstlisting}
screen.blit(background, (0,0))
x, y = pygame.mouse.get_pos()
x-= mouse_cursor.get_width()/2
y-= mouse_cursor.get_height()/2
screen.blit(mouse_cursor, (x, y))
pygame.display.update()
\end{lstlisting}

Obtener la posición del mouse es muy sencillo; el módulo {\ttfamily pygame.mouse} contiene todo lo que necesitamos
para trabajar con el mouse. La función {\ttfamily pygame.mouse.get\_pos()}  devuelve una tupla que contiene las coordenadas del mouse.

•  La primera línea descomprime esta tupla en dos valores por conveniencia: x e y. Podríamos usar estos dos valores como
coordenadas cuando desplazamos el cursor del mouse, pero eso colocaría la esquina superior izquierda de la imagen debajo
el mouse, y queremos que el centro de la imagen esté debajo del mouse. Así que hacemos un poco de cálculo (¡no temas!) para
ajuste x e y para que la imagen del mouse se mueva hacia arriba la mitad de su altura y hacia la izquierda la mitad de su ancho. 
Usando estas coordenadas coloca el centro de la imagen justo debajo del puntero del mouse, lo que se ve mejor. Al menos lo hace
para la imagen del cubo de rubik; si desea utilizar una imagen de puntero más típica, ajuste las coordenadas para que la punta
se encuentra debajo de las coordenadas reales del mouse.
 
•  La línea de código {\ttfamily screen.blit(mouse\_cursor,(x, y))}, se utiliza nuevamente la función  {\ttfamily blit()}, como el caso de la imagen de fondo; toma  una imagen y una tupla que contiene la posición de destino.\ding{51}


\subsection{Coordenadas de píxeles}

La ventana que crea el programa {\ttfamily mi\_primer\_pygame} está compuesta simplemente por pequeños puntos cuadrados en tu
pantalla llamada píxeles. Cada píxel comienza en negro, pero se puede configurar en un color diferente. Nuestro objeto de Superficie  tiene 640 píxeles de ancho y  480 píxeles de alto. Para ilustrar suponga que tenemos un objeto de Superficie
con 8 píxeles de ancho por 8 píxeles de alto. Si esa pequeña superficie de 8x8 se amplia para que cada píxel pareciera un
cuadrado en una cuadrícula, y agregamos números para los ejes X e Y, luego una buena representación del mismo
podría verse algo como esto:
\begin{figure}[h!]
	\centering
	\includegraphics[scale=0.20]{../A Logos de Python/coordenadas}
	%\caption{Hola, Mundo en Pygame}
	\label{coordenadas2}
\end{figure} 

Podemos referirnos a un píxel específico utilizando un sistema de coordenadas cartesianas. Cada columna del eje X
y cada fila del eje Y tendrá una "dirección"{} que es un número entero de 0 a 7, de modo que
puede localizar cualquier píxel especificando los números enteros de los ejes X e Y.
Por ejemplo, en la imagen de 8x8 de arriba, podemos ver que los píxeles en las coordenadas XY (4, 0), 
(2,2), (0, 5) y (5, 6) se han pintado de negro, el píxel en (2, 4) se ha pintado de gris, mientras que todos
los otros píxeles están pintados de blanco. Las coordenadas XY también se llaman puntos. Si has tomado una 
clase de matemática y aprendió sobre las coordenadas cartesianas, puede notar que el eje Y comienza en 0  en el origen 
y luego aumenta hacia abajo, en lugar de aumentar a medida que sube. Así   funcionan coordenadas en Pygame (y en casi todos los lenguajes de programación).


El marco de Pygame a generalmente representa las coordenadas cartesianas como una tupla de dos números enteros, como
(4, 0) o (2, 2). El primer número entero es la coordenada X y el segundo es la coordenada Y. (Las coordenadas cartesianose tratan con más detalle en el capítulo 12 de "Inventa tus propios juegos de computadora)


%Por ejemplo, pygame.Rect() y pygame.Surface() son funciones constructoras dentro del módulo pygame que devuelve nuevos objetos Rect y Surface. (Estos objetos son se describe más adelante.)






\subsection{Objetos de superficie y la ventana}

Los objetos de superficie son objetos que representan una imagen 2D rectangular. Los píxeles del objeto Superficie
se puede cambiar llamando a las funciones para dibujo de Pygame (descritas más adelante en este seccion \ref{drawing}) y luego
mostrado en la pantalla. El borde de la ventana, la barra de título y los botones no forman parte de objeto de superficie la pantalla.

En particular, el objeto Surface devuelto por {\ttfamily pygame.display.set\_mode()} se llama
superficie de visualización. Todo lo que se dibuje en el objeto de superficie de visualización se mostrará en la
ventana cuando se llama a la función {\ttfamily pygame.display.update()}. Es mucho más rápido dibujar en un objeto Superficie (que sólo existe en la memoria de la computadora) que dibujar una Superficie de pantalla de la computadora. La memoria de la computadora es mucho más rápida de cambiar que los píxeles de una monitor.

A menudo su programa dibujará varias cosas diferentes en un objeto Surface. Una vez que hayas terminado
dibujando todo en el objeto  superficie de la pantalla para esta iteración del bucle del juego (llamado
marco, tal como se llama una imagen fija en un DVD en pausa) en un objeto de superficie, se puede dibujar en
la pantalla. La computadora puede dibujar cuadros muy rápidamente y nuestros programas a menudo se ejecutarán de un lado a otro
30 fotogramas por segundo (es decir, 30 FPS). Esto se llama "{}velocidad de fotogramas"{} y se explica más adelante en \ref{FPS}.

El dibujo sobre objetos de superficie se tratará en las “Funciones de dibujo primitivas” y “Funciones de dibujo”.
Imágenes” más adelante en este capítulo.


\subsection{Colores}

Hay tres colores primarios de luz: rojo, verde y azul (El rojo, el azul y el amarillo son los
colores primarios para pinturas y pigmentos, pero el monitor de la computadora usa luz, no pintura), 
combinando diferentes cantidades de estos tres colores puedes formar cualquier otro color. En Pygame, 
se representar colores con tuplas de tres números enteros. El primer valor de la tupla es el color rojo. Un valor entero de 0 significa que no hay rojo en este color y un valor de 255 significa que tiene la cantidad máxima de rojo ese el color. El segundo valor es para verde y el tercer valor es para azul.

Estas tuplas de tres números enteros que se utilizan para representar un color suelen denominarse valores RGB, debido a que puedes usar cualquier combinación de 0 a 255 para cada uno de los tres colores primarios, esto significa que Pygame puede dibujar 16.777.216 colores diferentes (es decir, 256 x 256 x 256 colores). Sin embargo, si intenta utilizar un número mayor que 255 o un número negativo, obtendrá un error parecido a "ValueError: argumento de color no válido". Por ejemplo, crearemos la tupla (0, 0, 0) y la almacenaremos en una variable llamada NEGRO. Con sin cantidad de color (luz) rojo , verde o azul, el color resultante es completamente negro. El color negro es el
ausencia de cualquier color. La tupla (255, 255, 255) para una cantidad máxima de rojo, verde y azul para dar como resultado blanco. El color blanco es la combinación completa de rojo, verde y azul.  La tupla (255, 0, 0) representa la cantidad máxima de rojo pero no la cantidad de verde y azul, por lo que el el color resultante es rojo. De manera similar, (0, 255, 0) es verde y (0, 0, 255) es azul.

Puedes mezclar la cantidad de rojo, verde y azul para formar otros colores. Aquí están 


En la Tabla \ref{Tabla03} se presentan algunos valores de RGB para los
colores más comunes:

\begin{center}
	\noindent Table \ref{Tabla03}. Algunos colores más comunes:\\
	{\footnotesize{{%\scriptsize 
				\hspace{-0.7cm}	 {\scriptsize
				\begin{tabular}{l l l }	% white
					\multicolumn{1}{|>{\columncolor{vino}} l }{\color{white}\bf Color}& 
					\multicolumn{1}{|>{\columncolor{vino}} l }{\color{white}\bf  Valores de RGB}\\ 
					Agua &  ( 0, 255, 255)    \\
					Negro   & (0, 0, 0)      \\
					Azul &(0, 0, 255)\\
					Fucsia &(255, 0, 255)\\
					Gris &(128, 128, 128)\\
					Verde &( 0, 128, 0)\\
					Lima &( 0, 255, 0)\\
					Granate& (128, 0, 0)\\
					Azul marino& ( 0, 0, 128)\\
					Oliva &(128, 128, 0)\\
					Púrpura &(128, 0, 128)\\
					Rojo &(255, 0, 0)\\
					Plata& (192, 192, 192)\\
					Verde azulado &( 0, 128, 128)\\
					Blanco &(255, 255, 255)\\
					Amarillo& (255, 255, 0)\\			                  
					\arrayrulecolor{vino}\hline    
				\end{tabular}\label{Tabla03} 		
				%  	\vspace{0.2cm} 
}	}}} 
\end{center}


\subsection{Colores transparentes}

Cuando miras a través de una ventana de vidrio que tiene un tinte rojo intenso, todos los colores detrás de ella tienen un
un tono rojo. Puede imitar este efecto agregando un cuarto valor entero de 0 a 255 a tus valores de color.
Este valor se conoce como valor {\ttfamily alfa}. Es una medida de qué tan opaco es un color (es decir, no transparente).
Normalmente, cuando dibujas un píxel en un objeto superficie, el nuevo color reemplaza por completo a cualquier color que ya estuviera allí. Pero con colores que tienen un valor alfa, puedes en su lugar, simplemente agregue un tinte de color al color que ya está allí.

Por ejemplo, esta tupla  de tres números enteros para el color verde es (0, 255, 0). Pero si añadimos un
cuarto entero para el valor {\ttfamily alfa}, podemos convertirlo en un color verde medio transparente: (0, 255,
0, 128). Un valor {\ttfamily alfa} de 255 es completamente opaco (es decir, no es transparente en absoluto). 
Los colores (0, 255, 0) y (0, 255, 0, 255) se ven exactamente iguales. Un valor {\ttfamily alfa} de 0
significa que el color es completamente transparente. Si dibujas cualquier color que tenga un valor {\ttfamily alfa} de 0 en
un objeto Superficie, no tendrá ningún efecto, porque este color es completamente transparente e invisible.

Para dibujar usando colores transparentes, debes crear un objeto Superficie con el
método {\ttfamily convert\_alpha()}. Por ejemplo, el siguiente código crea un objeto Superficie que
se pueden dibujar colores transparentes en:
\begin{lstlisting}
anotherSurface =  DISPLAYSURF.convert\_alpha()
\end{lstlisting}
Una vez que se han dibujado cosas en el objeto Superficie almacenado en otra Superficie, entonces
{\ttfamily anotherSurface} se puede "{}borrar"{} (es decir, copiar) en {\ttfamily DISPLAYSURF} para que aparezca en la
pantalla (Consulte la sección "{}Dibujar imágenes con pygame.image.load() y blit()"{} más adelante \ref{Drawing}).

Es importante tener en cuenta que no puede utilizar colores transparentes en objetos Surface que no hayan sido devueltos por
una llamada {\ttfamily convert\_alpha()}, incluida la superficie de visualización que se devolvió
{\ttfamily pygame.display.set\_mode()}.

Si tuviéramos que crear una tupla de colores para dibujar el legendario Unicornio Rosa Invisible, usaríamos
(255, 192, 192, 0), que termina pareciendo completamente invisible como cualquier otro color que tiene un 0 como valor {\ttfamily alfa}., después de todo, es invisible.

%legendary Invisible Pink Unicorn,



\subsubsection{Objetos pygame.Color}

Necesitas saber cómo representar un color porque las funciones de dibujo de Pygame necesitan una manera de saber qué color quieres usar. 
Una tupla de tres o cuatro números enteros es unidireccional. Otro forma es como un objeto {\ttfamily pygame.Color}. 
Puede crear objetos de color llamando a la  función {\ttfamily pygame.Color} constructora y pasando tres o cuatro números enteros, puedes almacenar este objeto Color en variables del mismo modo que puedes almacenar tuplas en variables. Intente escribir lo siguiente
en el shell interactivo:
\begin{lstlisting}		
>>> importar pygame
>>> pygame.Color(255, 0, 0)
(255, 0, 0, 255)
>>> miColor = pygame.Color(255, 0, 0, 128)
>>> miColor == (255, 0, 0, 128)
True
\end{lstlisting}


Cualquier función de dibujo en Pygame que tenga un parámetro para
el color, tiene la forma de tupla o la forma de objeto de color para pasar el argumento. A pesar de
que son tipos de datos diferentes, un objeto Color es igual a una tupla de cuatro números enteros si ambos
representan el mismo color (al igual que 42 == 42.0 se evaluará como Verdadero).

Ahora que sabes cómo representar colores (como un objeto {\ttfamily pygame.Color} o una tupla de tres o
cuatro números enteros para rojo, verde, azul y, opcionalmente {\ttfamily alfa}) y coordenadas (como una tupla de dos
 números enteros para X e Y), aprendamos sobre los objetos {\ttfamily pygame.Rect} para que podamos comenzar a usar 
 las funciones para  dibujar de Pygame.

\section{Módulo pygame.Rect y objetos Rect}\label{objeto_Rect}

Pygame tiene dos formas de representar áreas rectangulares (al igual que hay dos formas de representar
colores). La primera es una tupla de cuatro números enteros:

1. La coordenada X de la esquina superior izquierda.

2. La coordenada Y de la esquina superior izquierda.

3. El ancho (en píxeles) del rectángulo.

4. Luego la altura (en píxeles) del rectángulo.\\

La segunda forma es como un objeto {\ttfamily pygame.Rect}, al que llamaremos objetos  {\ttfamily Rect} para abreviar. 

Por ejemplo, el siguiente código crea un objeto  {\ttfamily Rect} con una esquina superior izquierda en (10, 20) que tiene 200 píxeles 
de ancho y 300 píxeles de alto:
\begin{lstlisting}	
>>> import pygame
>>> rectangulo = pygame.Rect(10, 20, 200, 300)
>>> rectangulo == (10, 20, 200, 300)
True
\end{lstlisting}
Lo útil de esto es que el objeto {\ttfamily Rect} calcula automáticamente las coordenadas para
las otras características del rectángulo. Por ejemplo, si necesita saber la coordenada X del 
borde derecha del objeto{\ttfamily pygame.Rect} que almacenamos en la variable {\ttfamily spamRect}, simplemente puede acceder al
atributo derecho del objeto {\ttfamily Rect}:
\begin{lstlisting}
>>> rectangulo.right
210
\end{lstlisting}

El código Pygame para el objeto {\ttfamily Rect} calculó automáticamente que si el borde izquierdo está en la X
coordenada 10 y el rectángulo tiene 200 píxeles de ancho, entonces el borde derecho debe estar en la X
coordenada 210. Si reasigna el atributo {\ttfamily right}, todos los demás atributos se eliminan automáticamente, 
recalculado:
\begin{lstlisting}
>>> rectangulo.right= 350
>>> rectangulo.right
150
\end{lstlisting}

En la Tabla \ref{Tabla04} se muestra una lista de todos los atributos que proporcionan los objetos {\ttfamily pygame.Rect} (en el ejemplo, la variable que almacena el objeto {\ttfamily Rect} es {\ttfamily rectangulo}):

%rect = pygame.draw.rect(DISPLAYSURF,BLACK,(360,360,200,60),1)
%pygame.draw.ellipse(DISPLAYSURF,GREEN,rect,2)

\begin{center}
	\noindent Table \ref{Tabla04}. Lista de los atributos del objeto {\ttfamily Rect}:\\
	{\footnotesize{{%\scriptsize 
				\hspace{-0.7cm}	{\scriptsize
					\begin{tabular}{l l l }	% white
						\multicolumn{1}{|>{\columncolor{vino}} l }{\color{white}\bf Atributo }& 
						\multicolumn{1}{|>{\columncolor{vino}} l }{\color{white}\bf  Descripción}\\ 
						rectangulo.left &El valor int de la coordenada X del lado izquierdo del rectángulo\\
						rectangulo.right &El valor int de la coordenada X del lado derecho del rectángulo\\
						rectangulo.top& El valor int de la coordenada Y del lado superior del rectángulo\\
						rectangulo.bottom &El valor int de la coordenada Y del lado inferior\\
						rectangulo.centerx &El valor int de la coordenada X del centro del rectángulo\\
						rectangulo.centery &El valor int de la coordenada Y del centro del rectángulo\\
						rectangulo.width &El valor int del ancho del rectángulo\\
						rectangulo.height& El valor int de la altura del rectángulo\\
						rectangulo.size& Una tupla de dos enteros: (ancho, alto)\\
						rectangulo.topleft &Una tupla de dos enteros: (izquierda, arriba)\\
						rectangulo.topright& Una tupla de dos enteros: (derecha, arriba)\\
						rectangulo.bottomleft &Una tupla de dos enteros: (izquierda, abajo)\\
						rectangulo.bottomright &Una tupla de dos enteros: (derecha, abajo)\\
						rectangulo.midleft &Una tupla de dos enteros: (izquierda, centro)\\
						rectangulo.midright &Una tupla de dos enteros: (derecha, central)\\
						rectangulo.midtop &Una tupla de dos enteros: (centerx, top)\\
						rectangulo.midbottom &Una tupla de dos enteros: (centrox, abajo)\\			                  
						\arrayrulecolor{vino}\hline    
					\end{tabular}\label{Tabla04} 		
					%  	\vspace{0.2cm} 
	}	}}} 
\end{center}

Desarrollaremos las prácticas de este módulo el la siguiente sección.

\section{Módulo pygame.draw y  funciones primitivas para dibujo}

Pygame proporciona varias funciones para dibujar diferentes formas en un objeto superficie.
Estas formas, son como rectángulos, círculos, elipses, líneas o píxeles individuales, a menudo se denominan
primitivas para dibujo (drawing primitives). 


Abra el editor de archivos de IDLE, escriba el siguiente programa y guárdelo como {\ttfamily dibujos.py}

\hspace{-2.5cm}{\ttfamily \scriptsize dibujos.py}  \vspace{-0.5cm}
\begin{lstlisting}
import pygame, sys
from pygame.locals import *
pygame.init()

# Configuarando  ventana
DISPLAYSURF = pygame.display.set_mode((640, 480), 0, 32)
pygame.display.set_caption('Dibujando Polígonos, Rectas, 
                            Rectángulos, Círculos')

# Configurando colores
BLACK = ( 0, 0, 0)
WHITE = (255, 255, 255)
RED = (255, 0, 0,0)
GREEN = ( 0, 255, 0,0)
BLUE = ( 0, 0, 255)
SUAVE = (153,217,235)
AZUL = (153,217,234)

# dibujando el objeto Superficie (surface object)
DISPLAYSURF.fill(SUAVE)

pygame.draw.polygon(DISPLAYSURF,BLACK,((40, 50),(300,50),
                                       (300, 300),(160, 440),
                                       (40,300)),4)
#  Recta
pygame.draw.line(DISPLAYSURF, BLUE, (80, 80), (250, 80), 4)
pygame.draw.line(DISPLAYSURF, WHITE, (250, 80), (80, 300),2)
pygame.draw.line(DISPLAYSURF, RED, (80, 300), (250, 300), 4)

# Rectas y Polígono
pygame.draw.lines(DISPLAYSURF,GREEN,True,((440, 60),(470, 240),
                                          (340,90)),4)

pygame.draw.circle(DISPLAYSURF,WHITE, (550, 150), 70, 4)
pygame.draw.rect(DISPLAYSURF, BLUE,(320,260,200,60),2)
pygame.draw.ellipse(DISPLAYSURF,RED,(360,360,200,60),2)

# Bucle del juego
while True:
    for evento in pygame.event.get():
        if evento.type == QUIT:
            pygame.quit()
            sys.exit()
    pygame.display.update()	
\end{lstlisting}	


%#rect = pygame.draw.rect(DISPLAYSURF,BLACK,(360,360,200,60),1)
%#pygame.draw.ellipse(DISPLAYSURF,GREEN,rect,2)

Al ejecutar este archivo, tendremos una ventana semejante a la siguiente figura \ref{dibujos}:
\begin{figure}[h!]
	\centering
	\includegraphics[scale=0.32]{../A Logos de Python/dibujos}
	\caption{Ventana dibujos del módulo {\ttfamily pygame.draw}}
	\label{dibujos}
\end{figure}
%\ding{51} 
El juego estará ejecutándose hasta que el usuario cierre la ventana (haciendo clic en {\ttfamily X}).\\




Observe que usamos variables constantes para cada uno de los colores. Esto hacer que nuestro código sea más
legible, porque ver VERDE en el código fuente es mucho más fácil de entender que representar  el color verde como  una tupla (0, 255, 0).

Las funciones de dibujo reciben el nombre de las formas que dibujan. Los parámetros que pasas estas 
funciones les dicen en qué parte del objeto Superficie dibujar la forma, el tamaño, el color y qué tan gruesas deben ser las líneas.


Puedes ver cómo se llaman estas funciones en el archivo de programa {\ttfamily dibujos.py}, pero aquí hay una breve descripción de cada función:\\ 


• {\ttfamily fill(color)} es un método,  no es una función sino un método de {\ttfamily pygame.Surface}.  Rellenará completamente todo el objeto Superficie con cualquier valor de color que le pase como parámetro de color (actúa soble el objeto 
 {\ttfamily DISPLAYSURF.fill(WHITE)}.\\

• {\ttfamily pygame.draw.polygon (superficie, color, lista de puntos, an-\\ cho)}: dibuja un polígono, que  es una figura formada por
sólo segmentos de rectas (lados planos). Los parámetros de superficie y color le indican a la función en qué superficie dibuja el polígono y de qué color hacerlo. El parámetro lista de puntos es una tupla o lista de puntos (es decir, una tupla o lista de dos números enteros de tuplas para coordenadas XY). El polígono se dibuja trazando líneas entre cada punto, y el punto que viene
después en la tupla, luego se traza una línea desde el último punto hasta el primero punto. También puedes pasar una lista de puntos en lugar de una tupla de puntos. El parámetro de ancho es opcional. Si lo omites, el polígono que se dibuje se rellenará adentro, al igual que nuestro polígono verde en la pantalla se llena de color. Si pasas un número entero  para el parámetro ancho, solo se dibujará el contorno del polígono, el numero entero representa cuántos píxeles de ancho tendrá el contorno del polígono.

Pasando 1 para el ancho el parámetro creará un polígono delgado, mientras que pasar 4, 10 o 20 lo hará más grueso.
Si pasa el número entero 0 para el parámetro de ancho, el polígono se completará (como si omitieras el parámetro de ancho por completo).

Todas las funciones de dibujo de {\ttfamily pygame.draw} tienen parámetros de ancho opcionales al final,
y funcionan de la misma manera que el parámetro de ancho de {\ttfamily pygame.draw.polygon()}.
Probablemente un mejor nombre para el parámetro ancho hubiera sido espesor (thickness), ya que ese
El parámetro controla el grosor (thick) de las líneas que dibujas.\\

• {\ttfamily pygame.draw.line(superficie, color, punto\_inicial, punto\_- \\ final, ancho)}: esta función dibuja un
línea entre los parámetros {\ttfamily punto\_ini- \\cial (x,y), punto\_final (x,y)}. \\

• {\ttfamily pygame.draw.lines(superficie, color, cerrado, lista de pun- \\ tos, ancho)}: esta función dibuja una serie
de líneas de un punto al siguiente, muy parecido a  {\ttfamily pygame.draw.polygon()}. La única diferencia es que si pasa {\ttfamily False} para el parámetro cerrado, no dibuja línea del último punto al primer punto en la lista de puntos del parámetro. Si pasa {\ttfamily True}, entonces dibuja una línea desde el último punto hasta el primero.\\


• {\ttfamily pygame.draw.circle(superficie, color, punto\_centro, radio,\\ ancho)}: esta función dibuja un
círculo. El centro del círculo está en el parámetro center\_point. El número entero pasado por
el parámetro de radio establece el tamaño del círculo. El radio de un círculo es la distancia desde el centro hasta el borde. 
(El radio de un círculo es siempre la mitad del diámetro). Al pasar el valor de 70  para el parámetro de radio se dibujará un círculo que
tiene un radio de 70 píxeles.\\


• {\ttfamily pygame.draw.rect(superficie, color, rectángulo\_tuple, an- \\ cho)}: esta función dibuja un
rectángulo. La tupla\_rectángulo es una tupla de cuatro números enteros (para las coordenadas XY
de la esquina superior izquierda, y el ancho y alto) o se puede pasar un objeto pygame.Rect
en cambio. Si la tupla\_rectángulo tiene el mismo tamaño en ancho y alto, se formará un cuadrado.
ser dibujado.\\

• {\ttfamily pygame.draw.ellipse(superficie, color, rectángulo\_limita-\\ dor,ancho)}: esta función dibuja un
elipse (que es como un círculo aplastado o estirado). Esta función tiene todas las características habituales, sin embargo para indicarle   a la función dónde dibujar la elipse y el tamaño, debe especifique el rectángulo delimitador de la elipse. Un rectángulo delimitador es el rectángulo más pequeño que se puede dibujar alrededor de una forma. Aquí hay un ejemplo de una elipse y su rectángulo delimitador:
 

El parámetro {\ttfamily bounding\_rectangle} puede ser un objeto {\ttfamily pygame.Rect} o una tupla de cuatro
números enteros. Tenga en cuenta que no especifica el punto central de la elipse como lo hace para la
función {\ttfamily pygame.draw.circle()}.


%  pygame.draw.arc()   draw an elliptical arc
% arc(surface, color, rect, start_angle, stop_angle) -> Rect
% arc(surface, color, rect, start_angle, stop_angle, width=1) -> Rect
% Draws an elliptical arc on the given surface.

% The two angle arguments are given in radians and indicate the start and stop positions of the arc. The arc is drawn in a
%  counterclockwise  direction from the start_angle to the stop_angle.
%  https://www.pygame.org/docs/ref/draw.html#pygame.draw.arc

\subsection*{Finalizando el juego}

• La línea de llama a la función {\ttfamily pygame.display.update()} actualiza el objeto Surperficie
 en la pantalla del monitor del usuario.
%(recuerde que almacenamos este objeto en la variable  {\ttfamily DISPLAYSURF})
 
Una vez que haya terminado de llamar las funciones de dibujo para que el objeto Superficie de visualización 
tenga el aspecto deseado, debe llamar a {\ttfamily pygame.display.update()}. Dado que el objeto Superficie no ha
cambiado, la imagen es la misma azul claro, se vuelve a dibujar en la pantalla cada vez que se llama a
   {\ttfamily pygame.display.update()}. 
Después de terminar la línea, el bucle {\ttfamily  while} infinito, comienza nuevamente desde el inicio. 
Esto es todo el programa. \ding{51}\\


Lo único que debes recordar es que {\ttfamily pygame.display.update()} solo hará que superficie de visualización aparecen en la pantalla 
(es decir, el objeto superficie que devolvió al llamada a  {\ttfamily pygame.display.set\_mode()}). Si quieres que las imágenes de otros objetos Surface aparezcan en la pantalla, debe usar en método {\ttfamily blit()} para copiarlos en el objeto Superficie de visualización. En la sección \ref{blit} se explica con más detalles este método.


\section{Animación}


Ahora que sabemos cómo hacer que Pygame  dibuje en la pantalla figuras o dibujos, veamos cómo
hacer imágenes animadas en la pantalla o ventana de juego.  

Las imágenes animadas son el resultado de dibujar una imagen en la pantalla, y una fracción de segundo después dibujar una 
segunda imagen ligeramente diferente a la primara en la pantalla, que aparente un movimiento.\\

Considere las  siguientes tres  imágenes en la figura \ref{animacion}, y suponga que se  muestran rápidamente una 
\begin{figure}[h!]
	\centering
	\includegraphics[height=2cm]{../A Logos de Python/animacion}
	\caption{La ardilla mueve hacia el gato}
	\label{animacion}
\end{figure}después de la otra en la pantalla, para el usuario, parecería que el gato se mueve hacia la ardilla.  Es solo  ilusión visual. Para la computadora, son sólo un montón de píxeles. 

\subsection{Nuestra primera animación}

A continuación  prestamos un programa de ejemplo que muestra una animación simple. Se trata en una imagen de un pero 
que aparenta moverse en la pantalla. 

Escriba este código en su editor y guárdelo como {\ttfamily animacion\_perro.py}, también necesita  que los dos archivos de imagen {\ttfamily  perro.png} y {\ttfamily  perro\_regre-\\ sa.png} se encuentre en la misma carpeta que el programa. 

% Puede descargar esta imagen desde http://invpy.com/cat.png.   Este código está disponible en http://invpy.com/catanimation.py.

\hspace{-2.9cm}{\ttfamily \scriptsize animacion}

\vspace{-0.2cm}
\hspace{-2.5cm}{\ttfamily \scriptsize \_perro.py}\vspace{-0.84cm}
\begin{lstlisting}
# -*- coding: utf-8 -*-
"""
Created on Mon Apr  3 17:07:13 2023
@author: Linis
"""
#  El perro regresa

import pygame, sys
from pygame.locals import *
pygame.init()

# Configuración de Ventana
DISPLAYSURF = pygame.display.set_mode((450, 370), 0, 32)
pygame.display.set_caption('Animaciones y algo mas')

FPS = 30 # Configurando fotograma por segundos
fpsReloj = pygame.time.Clock() # velocidad de FPS

# Configuración de colores
WHITE =(255,255,255)
CELESTE = (153,217,235)
RED =(225,0,0)
GREEN = (0,255,0,180)

# Cargando imagenes y posición inicial 
perro = pygame.image.load('perro.png') #  Tamaño 85x76 
perrox = 20
perroy = 20
direction = 'derecho'
perroR = pygame.image.load('perro_regresa.png')
perroRx =365
perroRy =20 

while True: # Bucle principal del juego
    DISPLAYSURF.fill(WHITE)

    if direction == 'derecho':
        perrox  += 5
        if perrox  == 365:
            direction = 'izquierdo' 
            perroRx= 365


    if direction == 'izquierdo':
        DISPLAYSURF.blit(perroR, (perroRx, perroRy))                    
        if direction == 'izquierdo':
            perroRx  -= 5
            if perroRx == 10:
                direction = 'derecho'
                perrox = perroRx
    else:
        DISPLAYSURF.blit(perro, (perrox-5, perroy-5)) 
    
    for evento in pygame.event.get():
        if evento.type == QUIT:
            pygame.quit()
            sys.exit()

    pygame.draw.line(DISPLAYSURF, GREEN, (10, 90), (440, 90), 4)
    pygame.draw.circle(DISPLAYSURF,CELESTE,(300, 50),40, 4)
    pygame.draw.ellipse(DISPLAYSURF,RED,(20,15,200,73),2)
    
    pygame.display.update()
    fpsReloj.tick(FPS)
\end{lstlisting}
%  imagen  = {"perro":pygame.image.load('perro2.png'),"perroR":pygame.image.load('perro_regresa.png')}      
% DISPLAYSURF.blit(imagen["perro"], (perrox-5, perroy-5))
% DISPLAYSURF.blit(imagen["perroR"], (perroRx, perroRy))

Aunque no es el mejor ejemplo de animación, se observa que el perro se mueve en la pantalla, simulando movimiento. \\
 

Al ejecutar este archivo, tendremos una ventana semejante a la siguiente en la figura \ref{movimiento_perrro}:
\begin{figure}[h!]
	\centering
	\includegraphics[scale=0.48]{../A Logos de Python/perrotex}
	\caption{Simulando movimiento del perro}
	\label{movimiento_perrro}
\end{figure}
%\ding{51}
A continuación describimos en programa en detalle:

• En la línea 7 se define una constante igual a 30 FPS, es la configuración de fotogramas por segundo (FPS), en la próxima 
sección se da una referencia teórica básica al respecto.
  
\subsection{Velocidad de Fotogramas}   %  Cuadros Frames

 %y objetos {\ttfamily pygame.time.Clock()} 
 
La velocidad de fotogramas o frecuencia de actualización es la cantidad de imágenes que el programa dibuja por segundo, y
se mide en fotogramas por segundo (FPS). En los monitores de computadora el nombre común para FPS es hertz, 
muchos monitores tienen una velocidad  60 hertz o 60 fotogramas por segundo. Una velocidad de fotogramas baja en
videojuegos pueden hacer que el juego parezca entrecortado o nervioso. 

Si el programa tiene demasiado código para ejecutar dibujos en  pantalla con suficiente frecuencia,
entonces el FPS baja. Pero los juegos de este libro son lo suficientemente simple como para que esto no 
sea un problema ni siquiera en computadoras antiguas.\\


• En la siguiente línea  tenemos un objeto del módulo {\ttfamily pygame.time}, {\ttfamily fpsReloj = pygame.time.Clock()}; un  objeto con la finalidad de hacer que nuestro programa se ejecute en un máximo de  FPS.  Este objeto {\ttfamily fpsReloj} asegura que nuestros programas de juego no se ejecuten demasiado rápido al instalar pequeñas pausas en cada iteración del bucle del juego. 

Si no tuviéramos estas pausas, nuestro juego (programa) se ejecutaría tan rápido como la computadora pudiera ejecutarlo. Esto suele ser demasiado rápido para el jugador y a medida que las computadoras se vuelvan más rápidas, también ejecutarán el juego más rápido. 

Una llamada al método {\ttfamily tick(FPS)} del objeto {\ttfamily fpsReloj} en el bucle del juego  garantiza que el juego se ejecute a la misma velocidad sin importar qué tan rápido sea la computadora en la que se ejecuta.

El método  {\ttfamily fpsReloj.tick(FPS)}   debe llamarse al final del ciclo del juego, después de
llamar a {\ttfamily pygame.display.update()}. La duración de la pausa se calcula en función de cuánto tiempo ha pasado desde la llamada anterior a {\ttfamily tick()}, que habría tenido lugar al final de la iteración anterior del bucle del juego. La primera vez que se llama al método  {\ttfamily tick()}, no hay pausar en absoluto. 


Intente modificar la variable constante FPS para ejecutar el mismo programa a diferentes velocidades de fotogramas.
Configurarlo  con un valor más bajo haría que el programa se ejecutara más lento, y a un valor más alto haría
el programa se ejecuta más rápido.


% {fpsClock = pygame.time.Clock(), velocidad de FPS}  
• Luego tenemos las líneas de código: {\ttfamily  pygame.image.load("{}Imagen.png"{})}, esta función  carga los archivos {\ttfamily perro.png} y {\ttfamily perro\_regreresa.png} (formato de imagen),  y devuelve un objeto Surface al llamarlo.


\subsection{Dibujar y cargar imágenes con pygame.image.load() y blit()}\label{blit}
%  https://www.fsansberro.com/post/4-manejo-de-im%C3%A1genes

 
Para cargar  las imágenes se le pasa la  cadena  {\ttfamily "{}perro.png"{}} a la función {\ttfamily pygame.image.load()}, 
luego al llamado de la función {\ttfamily pygame.image.load()} devolverá un objeto Surface con la imagen dibujada. 

Este objeto Surface será un objeto  separado del objeto Surface de visualización, por lo que  el objeto Surface de la imagen
debe ser blit (es decir, copiar) sobre el objeto objeto Surface de visualización (display Surface object).

• Ahora en el cuerpo del {\ttfamily while}, tenemos la línea de código {\ttfamily DISPLAYSURF.-\\ blit(perro, (perrox, perroy))} un método del objeto {\ttfamily  DISPLAYSURF}. El método {\ttfamily blit()} se usa para copiar el objecto Surface de la imagen en el objeto Surface de visualización,  {\ttfamily perro} en   {\ttfamily DISPLAYSURF}.  Blitting consiste en dibujar el contenido de una superficie sobre otra (copiar una imagen en una nueva ubicación). 

El método {\ttfamily DISPLAYSURF.blit(perro, (perrox, perroy))}, tiene dos parámetros, el primero es el objeto Superficie de la imagen,
que se debe copiar en el objeto Surface DISPLAYSURF. El segundo parámetro es una tupla de dos números enteros para los valores X e Y de la esquina superior izquierda donde se debe copiar la imagen.

Si perrox  y  perroy estuvieran configurados en 100 y 200, y el ancho de la imagen  fuera 125 y la altura
era 79, esta llamada  {\ttfamily blit()} copiaría esta imagen en {\ttfamily DISPLAYSURF} de tal forma que la esquina superior izquierda
del objeto queda con coordenada XY (100, 200) y la esquina inferior derecha con  coordenada XY  (225, 279).

% Tenga en cuenta que no puede hacer blit a una superficie que esté actualmente "bloqueada" (como cuando un objeto PixelArray se ha elaborado a partir de él y aún no se ha eliminado).

           

El resto del ciclo del juego consiste simplemente en cambiar las variables perrox, perroy, la dirección para que
el perro se mueve en linea horizontal  en la ventana a la altura de la coordenada Y 20, y retorne el perro con el objeto {\ttfamily perroR} de la segunda imagen. 

Se agregaron  tres dibujos, una línea recta color verde, un círculo y una elipse, opcional. Y el código finaliza con las  dos líneas de código:
\begin{lstlisting}
pygame.display.update()
fpsReloj.tick(FPS)	
\end{lstlisting} Esto es todo el programa,  \ding{51} \\


Puede descargar el archivo movimientos de formas, para ver movimiento de dibujos círculos y elipses...


%%%%%%%%%%%%%%%%%%%%%%
%  font = pygame.font.Font('freesansbold.ttf', 32)

%   Displaying text: Since we made an image move around the edge
%%%%%%%%%%%%%%%%%%%%%%




\subsection{Crear texto con el módulo pygame.font}
 
Pygame tiene un modulo para gestionar fuentes {\ttfamily pygame.font}, este  módulo  proporciona otras funciones además
de objetos de fuente, que ocasionalmente puede necesitar. Es posible que lo uses para mostrar instrucciones del juego,
opciones de menú, etc. 

El módulo {\ttfamily pygame.font},  usa el tipo de fuente TrueType fonts (TTFs), que es usado en  mayoría de los sistemas 
para representar  texto fluido y de alta calidad. Hay muchas fuentes de este tipo instalado en su computadora que puede 
ser utilizado por el módulo de {\ttfamily font}.

%render = representar    % pygame.font.Font()

En {\ttfamily pygame} se necesitan seis pasos para hacer que aparezca un texto en la pantalla con  el módulo {\ttfamily pygame.font},  a continuación los describimos e ilustramos con el texto {\ttfamily texto = "{}¡Cuidado!"{ } perro bravo}:\\

\noindent 1. Crear un objeto fuente en {\ttfamily pygame}  con {\ttfamily pygame.font}. Ejemplo:  {\ttfamily  fuente = pygame.font.Font('freesansbold.ttf',40)}  (o  {\ttfamily fuente = pygame.font.-\\SysFont("{}arial"{}, 50)} una fuente común que es fácil de leer, Pygame buscará una fuente con el nombre "{}arial"{} en las fuentes instaladas; si  no la encuentra, se devolverá una fuente predeterminada). 

Los parámetros de la función constructora {\ttfamily pygame.font.Font()} son dos, el nombre de la fuente que desea crear  (tipo font),  y el tamaño de la fuente (medido en píxeles, como los procesadores lo miden tamaño de la fuente).

{\ttfamily  freesansbold.ttf} es una fuente que viene con Pygame (archivo de fuente TrueType extensión .ttf).  Y  para {\ttfamily pygame.font.SysFont} puede utilizar cualquiera de las fuentes que tienes instaladas en tu computadora. Puede obtener una lista de las fuentes instaladas en su sistema llamando a {\ttfamily pygame.font.get\_fonts()}.   Las fuentes también se pueden crear directamente a partir de archivos {\ttfamily  .ttf } mediante llamando a {\ttfamily  pygame.font.Font()}, que toma un nombre de archivo. 

Una vez que haya creado un objeto fuente, puede usarlo para construir un objeto Suface del texto.\\

\noindent 2. Crear un objeto Surface del texto con  el método {\ttfamily render()}. Ejemplo: 
{\ttfamily textoSurface = fuente.render("{}¡Cuidado!{} perro bravo", True,\\ YELLOW)} (el objeto superficie que contiene el texto). 

Los parámetros del método {\ttfamily  render()} son la cadena  texto a renderizar (representar), un valor booleano
para especificar si queremos suavizado\footnote{El suavizado es una técnica gráfica para hacer que el texto y las formas parezcan menos bloques, se logra haciendo un poco borroso sus bordes. Se necesita un poco más de tiempo de cálculo para dibujar con suavizado,
aunque los gráficos se vean mejor, su programa puede funcionar más lento (pero sólo un poco).
Para hacer que el texto de Pygame use anti-aliasing, simplemente se pase True al segundo parámetro del
método render().}, el color del texto y el
color del fondo. El fondo es opcional, si desea un fondo transparente, simplemente omita el parámetro de color de fondo en la llamada al método.\\


\noindent 3. Crear un objeto  {\ttfamily  Rect} a partir del objeto Surface del texto, mediante el método {\ttfamily get\_rect()}. 
Ejemplo: {\ttfamily textoRect = textoSurface.get\_rect()}. 

Este objeto  {\ttfamily Rect} tendrá el ancho y alto configurado 
correctamente para el texto que se representó, pero los atributos para las coordenadas del rectángulo de la esquina superior 
izquierdo serán (0,0).\\

\noindent 4. Establecer la posición del objeto {\ttfamily  Rect} cambiando uno de sus atributos. Por ejemplo,  configuramos el centro del objeto {\ttfamily   textoRect.center = (225, 150)}.\\


\noindent 5. Ahora solo hace falta copiar (Blit) el objeto Surface del texto en el objeto Superficie de visualización devuelto por
{\ttfamily  pygame.display.set\_mode()}.  Ejemplo: {\ttfamily DISPLAYSURF.blit(textoSurface, textRect)}

\noindent 6. Finalmente se llamar  a {\ttfamily pygame.display.update()} para actualizar la superficie de visualización aparecerá la pantalla.\\

%Consulte http://invpy.com/usingotherfonts para obtener más información sobre el uso de otras fuentes.

Todos estos pasos y códigos los agregamos al archivo {\ttfamily animacion\_perro.py}:

\hspace{-2.9cm}{\ttfamily \scriptsize animacion}

\vspace{-0.2cm}
\hspace{-2.5cm}{\ttfamily \scriptsize \_perro.py}\vspace{-0.84cm}
\begin{lstlisting}		
.
recorte
.
texto = '¡Cuidado! Perro bravo'
fuente = pygame.font.SysFont("arial",50)
textoSurface = fuente.render(texto, True,YELLOW) # omitir el color de fondo
textoRect = textoSurface.get_rect()
textoRect.center = (225, 150)
# salvar archivo de imagen de textoSurface en formato .pnf
pygame.image.save(textoSurface, "textoSurface.png")

while True: # Bucle principal del juego
    DISPLAYSURF.fill(WHITE)
    DISPLAYSURF.blit(textoSurface, textoRect)
    .
    recorte
    .
    pygame.display.update()
    fpsReloj.tick(FPS)
\end{lstlisting}

Al ejecutar el archivo se obtendrá una ventana de juego semejante a la siguiente figura \ref{perroLadra}:
\begin{figure}[h!]
	\centering
	\includegraphics[scale=0.45]{../A Logos de Python/texto}
    \caption{Ventana agregando texto con  módulo {\ttfamily font.SysFont()}}
   \label{perroLadra} 
   \end{figure}  
   
   
En este archivo se agregó una línea de código: 
\begin{lstlisting}
pygame.image.save(textoSurface,"textoSurface.png")	
\end{lstlisting}antes del bucle del {\ttfamily while}, este código permite guardar el archivo en formato PNF en el directorio en que se encuentra el programa, y además  se puede ver con cualquier  visualizador  en su computador, verifique.  Descarga de archivo en ..... \\

Para obtener detalles completos sobre el módulo de fuente, consulte la documentación en \href{www.pygame.org/docs/ref/font.html}{\color{vino} www.pygame.org/docs/ref/font.html}.
 .
 
	
\vspace{0.2cm}
\begin{minipage}[t]{0.2\textwidth}
	\centering\raisebox{\dimexpr\topskip-\height}{%
		\includegraphics[scale=0.20]{../A Logos de Python/Advertencia}}  %[width=\textwidth]
	
\end{minipage}%\hfill
\begin{minipage}[t]{0.70\textwidth} 
	{\scriptsize Las fuentes instaladas varían de una computadora a otra y no siempre es posible confiar en una fuente en particular.
	Si Pygame no encuentra la fuente que ha configurado, usará una fuente predeterminada, y tal vez  puede verse
	un poco diferente. La solución es distribuir el juego con archivos extensión .ttf, pero asegúrate de tener permiso o que sean libres.}
\end{minipage}
\vspace{0.2cm}



\subsection{Reproducir sonidos}

Ya hemos recorrido una buena parte de los conceptos básicos de Pygame, ahora estudiaremos como reproducir sonidos
a partir de archivos de sonido almacenados; pero no se preocupe, este trabajo es muy fácil incluso es más sencillo 
que mostrar imágenes a partir de archivos de imágenes.\\

• Primero, debes crear un objeto de sonido con la función constructora  {\ttfamily pygame.mixer.Sound()}.
Pygame puede cargar archivos WAV, MP3 u OGG.  Ejemplo: {\ttfamily sonidoObj = pygame.mixer.Sound("perro\_ladra.wav")} 
%La diferencia entre estos formatos de archivos de audio se explica en http://invpy.com/formats.   Página en  construcción 
el  parámetro de esta función es una cadena,  el nombre del archivo de sonido. \\

• Segundo, reproducir el sonido del objeto de sonido, esto se logra  llamado al método {\ttfamily  play()}. Ejemplo: {\ttfamily sonidoObj.play()}.
La ejecución del programa continúa inmediatamente después de llamar a play(). El sonido se reproduce inmediatamente,  no espera el
pasar a la siguiente línea de código.\\

• Tercero, para detener el objeto de sonido que se reproduce se llama al método  {\ttfamily stop()}. Ejemplo: {\ttfamily sonidoObj.stop()}
Este método  no tiene argumentos.

Aquí hay un código de muestra:

\hspace{-2.2cm}{\ttfamily \scriptsize sonido.py} \vspace{-0.7cm}
\begin{lstlisting}		
import pygame
pygame.init()
sonidoObj = pygame.mixer.Sound("perro_ladra.wav")  #  
sonidoObj.play()
import time
time.sleep(3) # Espera durante 3 segundo y deja que el sonido se reproduzca 
sonidoObj.stop()
\end{lstlisting}
El método {\ttfamily time.sleep()} espera y permite la ejecución del programa durante un número específico de segundos dados.

\subsubsection{Sonido o música de fondo}

Los objetos de sonido son buenos para que se reproduzcan efectos de sonido cuando el jugador recibe premio, corta un
espada, o compra o colecciona una moneda. Pero los juegos podrían ser mejores si tuvieran música de fondo, 
independientemente de lo que estuviera pasando en el juego.

Pygame solo puede cargar un archivo de música para reproducir en segundo plano a la vez.\\ 

La función {\ttfamily pygame.mixer.music.load()} se usa para cargar un archivo de música de fondo, 
y le pasa una  cadena  como  argumento, el nombre del archivo de sonido a carga. Igualmente este 
archivo puede tener formato WAV, MP3 o MIDI.

• Primero, para comenzar a reproducir el archivo de sonido cargado como música de fondo, se llame a la función 
{\ttfamily pygame.mixer.music.play(-1, 0.0)}. El argumento -1 hace que la música de fondo se repita
siempre cuando llega al final del archivo de sonido. Si lo configuras como un número entero
0 o mayor, entonces la música solo se repetirá esa cantidad de veces en lugar de repetirse para siempre. 

• El segundo argumento  0.0 significa tiempo para inicia a reproducirse el sonido. Si pasa un número entero más grande o flotante,
la música comenzará a reproducirse después de ese número de segundos. Por ejemplo, si pasas 13.5 para el segundo parámetro, 
el archivo de sonido comenzará a reproducirse en el punto 13.5 segundos en desde el principio.

• Para detener la reproducción de la música de fondo inmediatamente, llame a pygame.mixer.music.stop()
función. Esta función no tiene argumentos.\\

Aquí hay un código de ejemplo de los métodos y funciones de sonido:

\hspace{-2.9cm}{\ttfamily \scriptsize musica}

\vspace{-0.2cm}
\hspace{-2.5cm}{\ttfamily \scriptsize \_fondo.py}\vspace{-0.84cm}
\begin{lstlisting}		
import pygame, sys
pygame.init()

# Cargando y reproduciendo un efecto de sonido:
sonidoObj = pygame.mixer.Sound("perro_ladra.wav")    
sonidoObj.play()
import time
time.sleep(2) # espera y deja que el sonido se reproduzca durante 2 segundo
sonidoObj.stop()
\end{lstlisting}

Pero para continuar con la ilustración didáctica, agregamos este código al archivo  {\ttfamily animacion\_perro.py}:

\hspace{-2.9cm}{\ttfamily \scriptsize animacion}

\vspace{-0.2cm}
\hspace{-2.5cm}{\ttfamily \scriptsize \_perro.py}\vspace{-0.84cm}
\begin{lstlisting}		
	.
	recorte
	.
textoRect = textoSurface.get_rect()
textoRect.center = (225, 150)

# Configuración  de Sonido

pygame.mixer.music.load("perro_ladra.wav")
pygame.mixer.music.play(-1, 0.0)

while True: # Bucle principal del juego
    DISPLAYSURF.fill(WHITE)
    .
    recorte
    .
    pygame.display.update()
    fpsReloj.tick(FPS)
\end{lstlisting}
	
Al ejecutar obtendrás un juego igual al que se muestra en la página  \pageref{perroLadra}, pero ahora escuchará ladrar el perro. \ding{51} \\
	
Descargar archivo .... en proyecto...	


\section{Módulo pygame.event y Cola de eventos}\label{eventos} 

En Pygame el método {\ttfamily pygame.event} nos facilita  todos los eventos que han ocurrido en una cola de eventos. 
Esta cola se puede recorrer con un {\ttfamily for in}, para recuperar los eventos de {\ttfamily pygame.event.get()}, luego  procesar o ignorar estos eventos.  Una vez que un evento se procesa o se ignora, se elimina de la cola. 

%  Un evento es la forma en que Pygame le informa que algo sucedió fuera de su código. Los eventos son creados para muchas cosas, desde presionar teclas hasta recibir información de Internet, y son haciendo cola para ti hasta que los manejes. La función get en el módulo pygame.event devuelve cualquier evento esperándonos, que luego recorremos en un bucle for. Para este script, solo estamos interesado en el evento QUIT, que es generado por Pygame cuando el usuario hace clic en el botón cerrar en la ventana de Pygame. Entonces, si el tipo de evento es SALIR, llamamos a salir para cerrar, y todos los demás eventos son ignorados.

En los juegos (programas) que hemos tratado hasta ahora solo hemos  manejamos el evento {\ttfamily "{}QUIT"{}} (SALIR), que es 
esencial a menos que quieras tener ¡ventanas Pygame que no cierran, ventanas inmortales!  

Pygame crea otros eventos para informarle de cosas que suceden, como el movimiento del ratón y pulsaciones de teclas.  Eventos que se pueden generar en cualquier momento, sin importar el momento actual de ejecución del programa.

Por ejemplo, en un juego de combate, el código podría dibujar un tanque en la pantalla, al mismo tiempo que el usuario presiona un botón ejecutar cierta acción, digamos un disparo. Debido a que no puedes reaccionar a los eventos en el instante en que suceden, Pygame almacena
en una cola hasta que esté listo para manejarlos (normalmente al comienzo de la página principal bucle de juego). 

Puedes idearlo  se puede imaginar una fila de personas esperando para entrar a un edificio, cada uno de ellos con una información específica sobre un evento. Cuando el jugador presiona el botón de disparo, llega el evento del joystick, que lleva información sobre qué tecla se presionó. De manera similar, cuando el jugador suelta el botón de disparo, llega un clon del mismo evento del joystick con información
sobre el botón que se soltó. Podrían ir seguidos de un evento de mouse y un evento del teclado. En la próxima  \ref{recuparar} se trata  la recuperación de  estos eventos con detalles.



\subsection{Recuperar eventos}\label{recuparar}

El módulo {\ttfamily pygame.event} almacena los evento que suceden fuera de su código; recordemos que los eventos son generados por
muchas situaciones externas al código, estas van desde presionar teclas, mover el mouse,  hasta recibir información de internet; estos  se van organizando en cola hasta que el usuario los maneje. 
 
La función {\ttfamily pygame.event.get()} devuelve cualquier evento ocurrido o esperado. Para recuparar todos o cualquier eventos se usa un bucle {\ttfamily for}.  En los programas anteriores sólo se esperaba el evento {\ttfamily QUIT}, que es generado por Pygame cuando el usuario hace clic en el botón cerrar en la ventana de Pygame. Entonces, si el tipo de evento es {\ttfamily QUIT}, llamamos a
{\ttfamily pygame.quit()} para cerrar  Pygame, y todos los demás eventos son ignorados, inmediatamente se cierra ventana con {\ttfamily sys.exit()}. 

Además, se llamó a {\ttfamily pygame.event.get()} para recuperar todos los eventos y sacarlos de la cola, lo cual es
como abrir la puerta y dejar entrar a todos.  Esto probablemente sea la mejor manera de lidiar con los eventos, ya que garantiza que hayamos manejado  todo antes de pasar a dibujar algo en la pantalla, pero hay otras formas de trabajar con la cola de eventos. \\

Pygame tiene una función que hace esperar la ocurrencia del próximo evento antes de regresar, se trata de 
{\ttfamily pygame.event.wait()}. Para ilustrar su funcionamiento, puede imaginarse esperado en la puerta hasta 
que llegue alguien. Generalmente esta función no se utiliza casi en juegos, bebido a que si usted llama a 
{\ttfamily pygame.event.wait()}, Pygame pausa el {\ttfamily script} hasta que suceda algo, pero puede ser útil para
 aplicaciones Pygame donde coopere con otros programas en su sistema, como reproductores multimedia.


Una alternativa es la función  {\ttfamily pygame.event.poll()}, que devuelve un  único evento si hay uno en espera, 
o un evento ficticio (dummy)  {\ttfamily NOEVENT} si no hay eventos en la cola. Cualquiera que sea el método que utilice,
es importante no permitir que se acumulen, porque la cola de eventos tiene un tamaño limitado y los eventos se perderán
si la cola se desborda (overflows).

Es necesario llamar al menos a una de las funciones de manejo de eventos a intervalos regulares para que Pygame puede procesar eventos internamente. 

Si no utiliza ninguna de las funciones de manejo de eventos, puedes llamar a la función  {\ttfamily pygame.event.pump()} 
en lugar de un bucle de eventos. 

Estos objetos de evento contienen pocas variables miembro que describen el evento que ocurrió. La información que contienen
varía dependiendo del evento. Lo único común en todo  objetos evento es el tipo, que es un valor que indica el tipo de evento; y es el primero valor que se consulta para que puedas decidir qué hacer con ella. 

La Tabla \ref{tipo_evento} enumera  algunos eventos estándar que puede recibir.\\

\noindent Table \ref{tipo_evento}. Eventos estándar:
\begin{center}\vspace{-0.35cm}	
	{\footnotesize{{%\scriptsize 
				\hspace{-0.7cm}	
				\begin{tabular}{l l l }
					\multicolumn{1}{|>{\columncolor{vino}} l }{\color{white}\bf Evento}& 
					\multicolumn{1}{|>{\columncolor{vino}} l }{\color{white}\bf Propósito} &
					\multicolumn{1}{|>{\columncolor{vino}} l }{\color{white}\bf Parámetros} \\ 
					QUIT  (SALIR)    & El usuario ha hecho  clic   & no tiene\\
					&  en el botón cerrar         &   \\ 
					ACTIVEEVENT     & Pygame ha sido a & gain, state  \\
					& activado u oculto &  (ganancia, estado)  \\
					KEYDOWN 	    & Se ha presionado la tecla   &  unicode, clave, mod  \\
					KEYUP           & La tecla KEYUP  & key, mod\\
					&   ha sido liberada &  (clave, modo) \\
					MOUSEMOTION     & Se ha movido el ratón        &  pos, rel = distancia , botones \\
					MOUSEBUTTONDOWN & Se presionó el botón  & pos, buttons  \\				
					& del mouse                & (posición, botón) \\
					MOUSEBUTTONUP   & Se libero el botón      & pos, button\\ 
					& del mouse                &(posición, botón)  \\
					JOYAXISMOTION   & Se movió el joystick o pad   &  joy, axis =  eje, valor\\ 
					JOYBALLMOTION   & La bola de Joy se  movió  & joy,  ball = esférica, rel  \\    
					JOYHATMOTION    & Se movió el  hat       & joy, hat , valor  \\
					&  del joystick          &   \\   
					JOYBUTTONDOWN   &  Se presionó el botón  & joy, button  \\ 
					& del joystick o pad     &  (botón)  \\
					JOYBUTTONUP     & Se libero el botón     &  joy, button\\ 
					&    del joystick o pad   &  (botón) \\                               
					VIDEORESIZE     & Se cambió el tamaño     & tamaño, ancho, alto \\
					&  de la ventana de Pygame &   \\
					VIDEOEXPOSE     &  Parte o toda la ventana  & no tiene\\ 
					& de Pygame quedó expuesta &   \\
					USEREVENT       & Se ha producido  &   código \\ 
					& un evento de usuario& \\			                  
					\arrayrulecolor{vino}\hline 
					\label{tipo_evento}   
				\end{tabular}		
				%  	\vspace{0.2cm} 
	}}} 
\end{center}


Para ilustrar y ver todos los eventos que se generan, considere el siguiente {\ttfamily script  pygame\_eventos.py} de Pygame
usa {\ttfamily pygame.event.get()} para recuperar los evento. 

% Consejo Si cambia el color de relleno en el Listado 3-2 a (0,0,0) y el color de la fuente a (0, 255, 0), se verá como un poco como el código estilo Matrix. ¡Quizás tengas que usar un poco tu imaginación!

\hspace{-2.4cm}{\ttfamily \scriptsize pygame\_}

\vspace{-0.2cm}
\hspace{-2.4cm}{\ttfamily \scriptsize eventos.py}\vspace{-0.84cm}
\begin{lstlisting}%[frame=single]
import pygame
from pygame.locals import *
from sys import exit

pygame.init()
tamañoVentana = (450, 370)
#CELESTE = (153,217,235)
CELESTE = (170,238,187)
screen = pygame.display.set_mode(tamañoVentana, 0, 32)

while True:
    screen.fill(CELESTE)
    for evento in pygame.event.get():
       print(event)
       if evento.type == QUIT:
            pygame.quit()
            exit()
    pygame.display.update()
\end{lstlisting}

Al  ejecutar archivo {\ttfamily pygame\_eventos.py}, verá una ventana color azul celeste simple. Al mover el ratón sobre la ventana pygame,  {\ttfamily MOUSEMOTION} comienza a transmitir el movimiento del mouse, que se crean cada vez que el mouse cambia de posición
(ver Figura \ref{ver_eventos} ). 

\begin{figure}[h!]
	\centering
	\includegraphics[scale=0.41]{../A Logos de Python/ver_eventos2}
	\caption{Eventos del mouse}
	\label{ver_eventos}
\end{figure}  

Estos eventos especifican la posición actual del mouse, qué tan lejos se ha desplazado el mouse desde el último evento de movimiento y qué botones están presionados actualmente.   

Puedes obtener la  posición actual del mouse con el módulo {\ttfamily pygame.mouse}, pero corres el riesgo de perder información sobre lo que ha estado haciendo el jugador.  Este es un problema particular en computadoras de escritorio que realizan mucho trabajo en segundo plano y ocasionalmente pueden pausar el juego por un breve período de tiempo. 

Para un cursor de mouse, sólo necesita saber dónde está el mouse al comienzo de cada cuadro, por lo que es razonable usar {\ttfamily pygame.mouse.get\_pos()}. Si estuvieras usando el movimiento del mouse para conducir un tanque y los botones para disparar, sería mejor
trabajar con eventos para que el juego pueda monitorear más de cerca lo que el jugador ha estado haciendo.

% Tan pronto como consigue uno, lo convierte en un {\ttfamily string}  con  {\ttfamily str()} y lo agrega a una lista. El resto del código muestra el nuevo evento junto con los eventos anteriores. Utiliza el módulo de fuente para mostrar texto. 

\subsection{Eventos de movimiento del mouse}

Como ha visto, los eventos {\ttfamily MOUSEMOTION} se emiten cada vez que mueve el mouse sobre la ventana de Pygame. Contienen estos tres valores:\\

• pos: una tupla que contiene la posición del mouse cuando se generó el evento.

• rel: una tupla que contiene la distancia que se ha movido el mouse desde el último evento de movimiento del mouse  (a veces es llamado mouse mickies).

• buttons (botones): una tupla de tres números que corresponden a los botones del mouse.
Entonces {\ttfamily buttons[0]} es el botón izquierdo del mouse, {\ttfamily buttons[1]} es el botón central y {\ttfamily buttons[2]}
es el botón derecho. Si se presiona el botón, su valor se establece en 1; si no se presiona, el valor será 0. Se pueden presionar varios botones a la vez.


\subsection{Tipos de eventos del mouse}

Además de los eventos {\ttfamily MOUSEMOTION}, el mouse genera dos tipos de eventos más {\ttfamily MOUSEBUTTONDOWN} y 
{\ttfamily MOUSEBUTTONUP}.  Si hace clic en el mouse, primero verá en el script de la cola de mensajes, el evento inactivo (down event), seguido de un evento suelta (up event) cuando quita el dedo del botón. Entonces, ¿por qué tienen los dos eventos?; si usted usa el botón del mouse como disparador para disparar un cohete, solo necesitarías uno de los eventos, pero es posible que tengas un tipo diferente de arma, como una pistola  automática  que dispara continuamente mientras el se mantiene pulsado el botón. En este caso, comenzarías acelerar la pistola  automática con evento  presionar y disparar hasta que obtengas el evento liberar.

Ambos tipos de eventos del botón del mouse  tienen los  siguientes  dos valores:\\

• buttons: el número del botón que se presionó. Un valor de 1 indica si presionó el botón de la izquierda  del mouse, un valor 2 indica que se presionó el botón central y 3 indica que presionó el botón derecho.

• pos: una tupla que contiene la posición del mouse cuando se generó el evento.

\subsection{Manejo de eventos de teclado}

El teclado y el joystick tienen eventos similares a los de {\ttfamily KEYDOWN} y  {\ttfamily KEYUP}, que se transmiten 
cuando se pulsa una tecla  (down) y  cuando se suelta (up) la tecla. 

En el siguiente archivo se ilustra el proceso cómo podría responder a los eventos {\ttfamily KEYDOWN} y  {\ttfamily KEYUP}  para mover algo en la pantalla con las teclas del cursor. 

\subsubsection{Usar eventos de teclado para mover un fondo}

Al ejecutar este archivo, verá una ventana que contiene una imagen de fondo simple. Presiona arriba, abajo,
izquierda o derecha y el fondo se deslizará en esa dirección una unidad. 
% Quita el dedo de la tecla del cursor y el fondo dejará de moverse.

\hspace{-2.9cm}{\ttfamily \scriptsize movimientos}

\vspace{-0.2cm}
\hspace{-2.5cm}{\ttfamily \scriptsize \_fondo.py}\vspace{-0.84cm}
\begin{lstlisting}		
archivo_imagen_fondo = 'pygame.jpg'

import pygame
from pygame.locals import *
from sys import exit
pygame.init()

#CELESTE = (153,217,235)   # Color azul Celeste 
CELESTE = (170,238,187)
pantalla = pygame.display.set_mode((640, 480), 0, 32)
fondo = pygame.image.load(archivo_imagen_fondo).convert()
x, y = 0, 0
mover_x, mover_y = 0, 0

while True:
    for evento in pygame.event.get():
        if evento.type == QUIT:
            pygame.quit()
            exit()
        if evento.type == KEYDOWN:
            if evento.key == K_LEFT:
                mover_x = -1
            elif evento.key == K_RIGHT:
                mover_x = +1
            elif evento.key == K_UP:
                mover_y = -1
            elif evento.key == K_DOWN:
                mover_y = +1

        elif evento.type == KEYUP:
            if evento.key == K_LEFT:
                mover_x = 0
            elif evento.key == K_RIGHT:
                mover_x = 0
            elif evento.key == K_UP:
                mover_y = 0
            elif evento.key == K_DOWN:
                mover_y = 0
              
        x+= mover_x
        y+= mover_y
        pantalla.fill(CELESTE)
        pantalla.blit(fondo, (x, y))
   
        pygame.display.update()
\end{lstlisting}

El programa empieza igual que los anteriores juegos, importa, inicializa Pygame, y carga un
imagen de fondo. El bucle de eventos en este  {\ttfamily script} es diferente porque maneja los eventos  {\ttfamily KEYDOWN} y
 {\ttfamily KEYUP}. 
 
Estos eventos de teclas contienen los mismos tres valores:\\
 
• key (Tecla): este es un número que representa la tecla que se presionó o liberó. Cada tecla (física) del teclado tiene
una constante que comienza con {\ttfamily  K\_}. Las teclas del alfabeto son {\ttfamily  K\_a} hasta la 
{\ttfamily  K\_z}, pero también hay constantes para todas las demás  teclas (key), como {\ttfamily  K\_SPACE} y
{\ttfamily K\_RETURN}.  Para obtener una lista completa de las constantes del teclado que puede consultar en
 \href{https://www.pygame.org/docs/ref/key.html}{\color{vino}https://www.pygame.org/docs/ref/key.html}\\

• mod: este valor representa las teclas que se utilizan en combinación con otras teclas, como
 {\ttfamily SHIFT} (Mayúscula),  {\ttfamily Alt} y  {\ttfamily Ctrl}. Cada una de estas teclas modificadoras está representada por una constante que comienza con {\ttfamily  KMOD\_}, como {\ttfamily  KMOD\_SHIFT}, {\ttfamily  KMOD\_ALT} y {\ttfamily  KMOD\_CTRL}. 
La información del modificador se puede decodificar usando el operador bitwise {\ttfamily  AND} (bit a bit). Por ejemplo,  {\ttfamily mod} \&  {\ttfamily KMOD\_CTRL} se evaluará como  {\ttfamily True} si la tecla  {\ttfamily Ctrl} se presiona.  Puede encontrar una lista completa de las teclas modificadoras en  \href{www.pygame.org/docs/ref/key.html}{\color{vino}www.pygame.org/docs/ref/key.html}.\\

La información del modificador se puede decodificar usando un operador  {\ttfamily  AND} (excepto KMOD\_NONE, que debe compararse usando iguales ==). Por ejemplo:


• Unicode: este es el valor Unicode de la tecla que se presionó. Se produce combinando
la tecla presionada con cualquiera de las teclas modificadoras que se presionaron. Hay un valor  Unicode
para cada símbolo en el alfabeto inglés y otros idiomas. 

Generalmente no usarás estos valores en un juego, porque las teclas tienden a usarse más como interruptores que para ingresar
texto. Una excepción sería ingresar a una tabla de puntuación, donde se quiere que el jugador  escriba
letras que no estén en inglés, así como mezclar mayúsculas y minúsculas.\\
 
 
Dentro del controlador de {\ttfamily KEYDOWN} verificamos las cuatro  constantes  de teclas que corresponden al
cursor. Si se presiona {\ttfamily K\_LEFT}, entonces el valor de {\ttfamily mover\_x} se establece en –1; 
si se presiona {\ttfamily K\_RIGHT}, está establecido en +1. Este valor se agrega luego a la coordenada x del fondo 
para moverlo hacia la izquierda o hacia la derecha. También hay un valor {\ttfamily mover\_y}, que se establece si se presiona {\ttfamily K\_UP} o {\ttfamily K\_DOWN}, que se moverá el fondo verticalmente.

 
También es necesario manejar el evento {\ttfamily KEYUP}, porque queremos que el fondo deje de moverse cuando el usuario 
suelta la tecla del cursor. El código dentro del controlador para eventos {\ttfamily  KEYUP} es similar al {\ttfamily down}, 
pero establece  {\ttfamily mover\_x} o {\ttfamily mover\_y} nuevamente a cero para evitar que el fondo se mueva.\\


Después del ciclo de eventos, todo lo que tenemos que hacer es agregar los valores {\ttfamily mover\_x} y 
{\ttfamily mover\_y} a x e y, luego dibuja el fondo en (x, y). Y se usa  {\ttfamily fondo.fill(CELESTE)},
 para borrar la pantalla y dejarla en azul celeste. Esta linea es necesaria porque si movemos la imagen de fondo ya no cubre toda la pantalla, lo cual supongo que técnicamente significaría que ya no es un fondo.

%   <Event(32867-UserEvent {'mensaje': '¡acción  de tecleado!'})>
%  <Event(769-KeyUp {'unicode': 'a', 'key': 97, 'mod': 4096, 'scancode': 4, 'window': None})>

\subsubsection{Bloquear y permitir eventos}  %  https://www.code4tomorrow.org/courses/python/games/ch.3-events/3.3-blocking-allowing-events

En muchas ocasiones no  es necesario gestionar todos los eventos en el juegos, y es sano suspender los eventos que no se necesitan.
Existen formas alternativas  gestionar eventos. Por ejemplo, si estás usando {\ttfamily pygame.mouse.get\_-\\pos()}, no necesitará  al evento {\ttfamily MOUSEMOTION}, con la función {\ttfamily pygame.event.-\\set\_blocked(pygame.MOUSEMOTION)} o 
 {\ttfamily pygame.event.set\_blocked([MOU-\\SEMOTION])} se bloquea. \\

• {\ttfamily set\_blocked()}, si pasas un evento de pygame o una lista de eventos, todos esos eventos se bloquearán.  Si le pasa {\ttfamily None} como entrada, todos los eventos de la cola se bloquearán.

Por ejemplo, la siguiente línea deshabilitará todas las entradas del teclado, bloqueando los eventos 
{\ttfamily KEYDOWN} y  {\ttfamily KEYUP}:
\begin{lstlisting}		
pygame.event.set_blocked([KEYDOWN,KEYUP])
\end{lstlisting}

•  {\ttfamily set\_allowed()}, hace lo contrario que el método anterior,  desbloquear los eventos.  Requiere un único tipo de evento o una lista de tipos de eventos pygame en la cola que se le permitirá.  Además  si le pasa {\ttfamily None}, permitirán todos los eventos. 

\noindent Ejemplo:  {\ttfamily pygame.event.set\_allowed(pygame.QUIT)}.\\
 
•  {\ttfamily get\_blocked()}, con esta función puedes preguntarle a Pygame si un evento o una lista de eventos están bloqueados. 
Devolverá un valor booleano (True o False) si un evento está bloqueado. 

\noindent Ejemplo: {\ttfamily print(pygame.event.get\_blocked([KEYUP,QUIT])) } o  \\ {\ttfamily print(pygame.event.get\_blocked(pygame.QUIT))}.
% print(pygame.QUIT == QUIT) es  True

\subsection{Eventos personalizados con pygame.event.Event}

Generalmente es Pygame el que crea todos los eventos, pero  puedes crear sus propios eventos.
Podría usar esta capacidad para reproducir demostraciones (replicando la entrada del reproductor) o simular los
efectos de un gato caminando sobre el teclado. 

Para enviar un evento a la cola, primero construye un objeto de evento con {\ttfamily pygame.event.Event} y luego se enlista
con {\ttfamily pygame.event.post}. El evento se colocará al final de la cola listo para ser recuperado
en el bucle de eventos. 

A continuación se ilustra cómo simular al jugador presionando la barra espaciadora:

\begin{lstlisting}		
mi_evento = pygame.event.Event(KEYDOWN, tecla = K_SPACE, mod=0, unicode=u'')
pygame.event.post(mi_evento)	

\end{lstlisting}
El constructor de eventos toma el tipo de evento, como uno de los eventos en la Tabla \ref{tipo_evento},
seguido de los valores que debe contener el evento. Como estamos simulando el evento  {\ttfamily KEYDOWN},
debe proporcionar todos los valores (parámetros) que el controlador de eventos esperaría que estuvieran allí.
Si lo prefieres, puede proporcionar estos valores como un diccionario.

Esta línea creará el mismo objeto de evento:

\begin{lstlisting}
mi_evento = pygame.event.Event(KEYDOWN,{"key":K_SPACE,
                               "mod":0,"unicode":u' '})
\end{lstlisting}

Además de simular eventos generados por Pygame, puedes crear eventos completamente nuevos. Todo lo que tienes que hacer
es usar un valor por encima de  {\ttfamily USEREVENT} para el evento.   {\ttfamily USEREVENT} es el valor  máximo  
que Pygame utilizará para sus propios eventos IDs. 

%   USEREVENT   User Event
%   type(pygame.event.post(mi_evento))  <class 'bool'>

A veces esto puede resultar útil si quiere hacer un bucle de eventos antes de pasar a dibujar 
en la pantalla. Aquí hay un ejemplo de un evento personalizado para responder a una acción de tecleado: 
%\vspace{-0.5cm}
\begin{lstlisting}      %  Accion_Teclado2   =	Accion_Teclado + 1, para un segundo evento
Acción_Teclado = USEREVENT+1  
mi_evento = pygame.event.Event(Accion_Teclado,
                               mensaje= "¡Acción de tecleado!")
pygame.event.post(mi_evento)
\end{lstlisting}

%  {\ttfamily User Event} es una biblioteca complementaria para Testing Library que simula las interacciones del usuario al emitir los eventos que ocurrirían si la interacción se llevara a cabo en un navegador.

El manejo de eventos de usuario se realiza de la misma manera que los eventos habituales que genera Pygame,
simplemente verifique el tipo de evento para ver si coincide con su evento personalizado. Así es como podrías manejar
un evento {\ttfamily Acción\_Teclado}:

\begin{lstlisting}		
for event in pygame.event.get():
    if event.type == Acción_Teclado:
       print(event.mensaje)	
\end{lstlisting}


\section{Opciones para pantalla} 

La resolución de pantalla es una decisión que también depende de cuánta acción tendrás en el juego, 
cuántos objetos estarán moviéndose en pantalla al mismo tiempo, cuanto mas objetos más lento será la ejecutará el juego. Es posible que tengas que compensar seleccionando una resolución más baja que permita acelerará las cosas para lograr que funcionen como se quiere.

Generalmente es más fácil utilizar una resolución de pantalla fija (tamaño de pantalla) porque puede simplificar el código. 
Sin embargo la mejor solución podría  ser dejar que el jugador decida en qué resolución quiere ejecutar, de modo que
pueda ajustar la pantalla hasta que tengan una calidad aceptable entre lo visual y la fluidez con la que se ejecuta el juego. 
Si sigues esta ruta, tendrás que asegurarte de que tu juego se vea bien en todas las resoluciones potenciales.  Pero este trabajo 
no se inicia hasta que llegue el momento de escribir tu juego. Por ahora simplemente seleccione una resolución que funcione mientras experimentas.

\subsection{Pantalla completa}

El primer parámetro de la ventana a crear es el tamaño.  Un tamaño de 640x480 píxeles crea una pequeña ventana 
que cabe cómodamente en la mayoría de los computadoras de escritorios, y para ejecutar y  depurar el juego en esa ventana 
es excelente, sin embargo la mayoría de los juegos ocupan toda la pantalla, y no tienen los bordes ni la barra de título habituales. 
\vspace{0.2cm}
\begin{minipage}[t]{0.2\textwidth}
	\centering\raisebox{\dimexpr\topskip-\height}{%
		\includegraphics[scale=0.20]{../A Logos de Python/Advertencia1}}  %[width=\textwidth]
	
\end{minipage}%\hfill
\begin{minipage}[t]{0.70\textwidth} 
	{\scriptsize Si algo sale mal con el {\ttfamily script} con modo de pantalla completa, a veces puede resultar difícil
		solucionarlo, vuelve a su escritorio. Es mejor probarlo primero en modo ventana. También debe proporcionar una
		forma alternativa de salir del script porque el botón de cerrar no está visible en el modo de pantalla completa.
		Puede agregar un evento auxiliar para salir, por ejemplo al pulsar la letra {\ttfamily q}
\begin{lstlisting}	
if evento.type ==KEYDOWN:
    if evento.key==K_q:
        pygame.quit()
        exit()		
\end{lstlisting} }

\end{minipage}
\vspace{0.2cm}


El modo de pantalla completa suele ser más rápido porque su {\ttfamily cript} no tiene que cooperar con otras ventanas 
en su escritorio.  Para configurar el modo de pantalla completa, puede usar el parámetro indicador {\ttfamily FULLSCREEN}  
el segundo parámetro de {\ttfamily set\_mode()}.

\begin{lstlisting}		
pantalla =pygame.display.set_mode((640, 480), FULLSCREEN, 32)
\end{lstlisting}

Cuando trabaje con pantalla completa, su tarjeta de vídeo probablemente cambiará a un modo de vídeo diferente,
que cambiará el ancho y el alto de la pantalla, y potencialmente cuántos colores puede mostrar al mismo tiempo.
Las tarjetas de vídeo solo admiten algunas combinaciones de tamaño y cantidad de colores, pero Pygame te ayudará
si intentas seleccionar un modo de vídeo que la tarjeta no admite directamente. 

Si el tamaño de pantalla que configurado no es compatible, Pygame seleccionará el siguiente tamaño y
copia la pantalla en el centro, lo que puede generar bordes negros en la parte superior e inferior de la misma. 
Para evitar estos límites, se selecciona una de las resoluciones estándar que prácticamente todas las tarjetas de vídeo soportan:
(640, 480), (800, 600) o (1024, 768). 

Para ver exactamente qué resoluciones  de visualización  soporta, puede utilizar {\ttfamily pygame.display.list\_modes()},
que devuelve una lista de tuplas que contienen resoluciones soportadas. 

Probemos esto desde el intérprete interactivo:
\begin{lstlisting}	
>>> importar pygame
>>> pygame.init()
>>> pygame.display.list_modes()
[(1366, 768), (1360, 768), (1280, 768), (1280, 720), 
(1280, 600), (1024, 768), (800, 600), (640, 480)]

\end{lstlisting}

Si la tarjeta de video no puede brindarle la cantidad de colores que configuro, Pygame ajustara
los colores en la superficie de la pantalla automáticamente, lo que puede resultar en una ligera caída en la calidad de la imagen.

En el siguiente {\ttfamily script} se muestra  cómo pasar del modo de ventana al modo de pantalla completa. 

Si presiona la tecla F, la pantalla ocupará toda la pantalla (puede haber un retraso de algunos
segundos mientras esto sucede). Presione F por segunda vez y la pantalla volverá a una ventana.

\hspace{-2.65cm}{\ttfamily \scriptsize pantalla}

\vspace{-0.2cm}
\hspace{-2.7cm}{\ttfamily \scriptsize \_completa.py}\vspace{-0.84cm}
\begin{lstlisting}
archivo_imagen_fondo = 'pygame.jpg'

import pygame
from pygame.locals import *
from sys import exit
pygame.init()
#CELESTE = (153,217,235)   # Color azul Celeste 
CELESTE = (170,238,187)
#pantalla = pygame.display.set_mode((640, 480), FULLSCREEN, 32)

pantalla = pygame.display.set_mode((640, 480), 0, 32)

fondo = pygame.image.load(archivo_imagen_fondo).convert()

Fulls_pantalla = False

while True:
    for evento in pygame.event.get():
        if evento.type ==KEYDOWN:
            if evento.key==K_q:
                pygame.quit()
                exit()	
        if evento.type == QUIT:
             pygame.quit() 
             exit()
        if evento.type == KEYDOWN:
            if evento.key == K_f:
                Fulls_pantalla = not Fulls_pantalla        
                if Fulls_pantalla:
                    pantalla = pygame.display.set_mode((640, 480), FULLSCREEN, 32)
                else :
                    screen = pygame.display.set_mode((640, 480), 0, 32)      


    pantalla.fill(CELESTE)
    pantalla.blit(fondo, (0,0))

    pygame.display.update()

\end{lstlisting}


\subsection{Redimensionar ventanas Pygame}

Ocasionalmente, es posible que el usuario quiera cambiar el tamaño de una ventana de Pygame, lo que normalmente
logra haciendo clic en la esquina de la ventana y arrastrando con el mouse. Es esta tarea bastante fácil,
usando el parámetro indicador {\ttfamily RESIZABLE} cuando llamas a {\ttfamily set\_mode()}. Pygame informa a su 
código que el usuario ha cambiado el tamaño de la ventana enviando un evento {\ttfamily VIDEORESIZE} que 
contiene el nuevo ancho y alto de la ventana. Cuando obtenga uno de estos eventos, debe llamar
a la función {\ttfamily pygame.display.set\_mode()} nuevamente para configurar las nuevas dimensiones del tamaño de ventana.

En el siguiente {\ttfamily script} se muestra cómo tratar los eventos {\ttfamily VIDEORESIZE} para redimensionar una ventana:


\hspace{-2.65cm}{\ttfamily \scriptsize redimens}

\vspace{-0.2cm}
\hspace{-2.7cm}{\ttfamily \scriptsize \_ventana.py}\vspace{-0.84cm}
\begin{lstlisting}		
archivo_imagen_fondo = 'pygame.jpg'

import pygame
from pygame.locals import *
from sys import exit
pygame.init()
tamana = (640, 480)
CELESTE = (170,238,187)   # Color azul Celeste 

# Inicializar la pantalla
pantalla = pygame.display.set_mode(tamana, RESIZABLE, 32)
fondo = pygame.image.load(archivo_imagen_fondo).convert()
pygame.display.set_caption("Redimensión de Ventana")
# Inicializar la variable para manejar el tamaño de la pantalla
pantalla_ancho, pantalla_alto = tamana

# Bucle principal
while True:
    # Manejar eventos
    for evento in pygame.event.get():
        if evento.type == QUIT:
            pygame.quit()
            exit()    
    if evento.type == KEYDOWN: 
        if evento.key==K_q:
            pygame.quit()
            exit()
            
    # redimensionar ventana e imagen de fondo
    if evento.type == VIDEORESIZE:
        tamana = evento.size
        pantalla = pygame.display.set_mode(tamana, RESIZABLE, 32)
        pantalla_ancho, pantalla_alto = tamana
        (wp, hp) = tamana
        imagen_redimensionada = pygame.transform.scale(fondo, (wp, hp))
        fondo = imagen_redimensionada
        pygame.display.set_caption("Redimensión de Ventana" +str(evento.size))

    # Dibujar el fondo en la pantalla
    pantalla.fill(CELESTE)    
    for y in range(0, pantalla_alto, fondo.get_height()):       
        for x in range(0, pantalla_ancho, fondo.get_width()):            
            pantalla.blit(fondo,(x, y)) 

    pygame.display.update()
\end{lstlisting}	


Al ejecutar este {\ttfamily script}, mostrará una ventana  Pygame con  
imagen de fondo {\ttfamily pygame.jpg}. Si hace clic en la esquina o el borde de la ventana y arrastra con el mouse, 
el {\ttfamily script} genera un evento {\ttfamily VIDEORESIZE}.  
En el controlador de ese evento se llamada la función {\ttfamily set\_mode()}, y se crea
una nueva superficie de pantalla que coincide con las nuevas dimensiones. 

El evento {\ttfamily VIDEORESIZE}, contiene los siguientes valores:\\

• {\ttfamily size} (tamaño): es una tupla que contiene las nuevas dimensiones de la ventana; {\ttfamily size[0]} es el ancho
y el {\ttfamily size[1]} es la altura.\\

• {\ttfamily w}: este valor contiene el nuevo ancho de la ventana. Es el mismo valor que el {\ttfamily size[0]}, pero
puede ser más conveniente.\\

• {\ttfamily h}: este valor contiene la nueva altura de la ventana. Es el mismo valor que {\ttfamily  size[1]}, pero
puede ser más conveniente.\\

Además, se ha configurado el objeto superficie de imagen de fondo, con la función  {\ttfamily pygame.transform.scale(fondo, (w, h))}, 
a las nuevas dimensiones.  Es importante tener presente que la imagen se puede ver un poco distorsionada y tal vez no se verá con una buena calidad, así es recomendable modificar el tamaño de imagen con un editor de imágenes en lugar de hacerlo con PyGame, pero al final es decisión tuya. También se pueden manejas como opción repetir el fondo en cuadriculas, múltiples copias de fondo para completar la nueva ventana, el siguiente código genera las cuadrículas  para la imagen de fondo: 
\begin{lstlisting}		
for y in range(0, pantalla\_alto, fondo.get\_height()):
    for x in range(0, pantalla\_ancho, fondo.get\_width()):      
        pantalla.blit(fondo,(x, y))
\end{lstlisting}         

Las dos llamadas al rango producen las coordenadas necesarias para colocar la imagen de fondo en cuadrículas (múltiples).


Generalmente los juegos se ejecutan en pantalla completa, por lo que las pantallas redimensionables quizás no
sean una característica que utilices. ¡Pero está ahí en tu caja de herramientas si lo necesitas! \ding{51} 


\subsection{Ventanas sin bordes}\label{sin_bordes}

Generalmente, cuando creas una ventana de Pygame estándar con barras de título y frontera. 
Sin embargo, es posible crear una ventana que no tenga estas características para que el usuario no pueda mover 
ni cambiar el tamaño de la ventana, ni cerrarla mediante el botón de cerrar. Un ejemplo de tal uso es la ventana
 utilizada para las pantallas de presentación. 

Algunos juegos pueden tardar un poco en cargar porque contienen muchos archivos de imagen y sonido. 
Si no hay nada visible en la pantalla mientras esto sucede, el jugador puede sentir que el juego no funciona e intentar reiniciarlo.
 Para configurar una visualización sin bordes, se usa el indicador {\ttfamily NOFRAME} a llamar  la función 
 {\ttfamily  pygame.display.set\_mode()}. 
 
 Por ejemplo, el siguiente {\ttfamily scritp} muestra una ventana sin títulos y sin bordes: 
 
\hspace{-2.3cm}{\ttfamily \scriptsize sin}

\vspace{-0.2cm}
\hspace{-2.5cm}{\ttfamily \scriptsize \_bordes.py}\vspace{-0.84cm}
\begin{lstlisting}
from sys import exit
pygame.init()

# Inicializar  el tamaño de la pantalla
tamana = (640, 480)
CELESTE = (170,238,187)   # Color azul Celeste 

# Inicializar la pantalla sin marcos 
pantalla = pygame.display.set_mode(tamana, NOFRAME, 32)

fondo = pygame.image.load(archivo_imagen_fondo).convert()

# Bucle principal
    while True:
        # Manejar eventos
        for evento in pygame.event.get():       
            if evento.type == KEYDOWN: 
                if evento.key==K_q:
                    pygame.quit()
                    exit()

    # Dibujar el fondo en la pantalla
    pantalla.fill(CELESTE)     
    pantalla.blit(fondo,(0,0))                

    pygame.display.update()
\end{lstlisting}


Al  ejecutar archivo {\ttfamily sin\_bordes.py}, verá una ventana con el fondo de imagen {\ttfamily pygame.jpg}, sin bordes, y sin títulos, como una presentación, semejante a la siguiente figura \ref{sin_marco}:

\begin{figure}[h!]
	\centering
	\includegraphics[scale=0.41]{../A Logos de Python/sin_marco_ventana}
	\caption{Ventana sin borde y sin títulos}
	\label{sin_marco}
\end{figure}  
Se manejan los objetos eventos ({\ttfamily pygame.event.get()}) para cerrar el juego, creando una salida al presiona la tecla {\ttfamily q} ({\ttfamily K\_q}). \\




%Podemos averiguar qué eventos han ocurrido llamando a pygame.event.get() función, que devuelve una lista de objetos pygame.event.Event (que simplemente llamaremos Event objetos para abreviar).


\subsection{Parámetros adicionales para visualización} \label{pantalla}  %  Display 

% Tienes razón, "bandera" o "flags" no es la mejor traducción al español. En este contexto, los términos más apropiados serían:
% Parámetros indicadores  Indicadores de modo de visualización  Opciones de modo de visualización
% Por ejemplo, se podría decir: "Estas son las opciones de modo de visualización que se combinan utilizando el operador bitwise OR (|):"


La función  {\ttfamily pygame.display.set\_mode( )} tiene otras opciones importante que  puedes usar,  pero deben ser usadas 
con más cuidado, debido a que pueden causar problemas de compatibilidad en algunas plataformas y afectar el rendimiento si se usan incorrectamente. Por ejemplo, hay un parámetro indicador de modo de visualización, el {\ttfamily hardware} es que no es  compatible con todas las plataformas. Por lo general es mejor usar el valor 0 para ventanas en  pantallas y  {\ttfamily FULLSCREEN} para pantalla completa, debido a que garantizan el buen funcionamiento en todas las plataformas.

Sin embargo, si sabes lo que está haciendo, puede configurar algunas opciones avanzadas para obtener un rendimiento adicional.\\

Por ejemplo, si configura el parámetro indicador {\ttfamily HWSURFACE}, creará lo que se llama una superficie de {\ttfamily hardware}\footnote{En Pygame, una superficie "{}hardware"{} se refiere a una superficie de visualización que se beneficia de la aceleración por hardware. Esto significa que la superficie de visualización se gestiona y procesa directamente por la tarjeta gráfica de tu computadora, lo que puede resultar en un rendimiento más rápido y eficiente en comparación con una superficie de visualización que se maneja exclusivamente por software.


Al utilizar una superficie "{}hardware"{} en Pygame, puedes aprovechar la capacidad de la tarjeta gráfica para realizar operaciones de dibujo y renderizado de manera más eficiente, lo que puede ser especialmente útil para aplicaciones gráficamente intensivas como juegos o simulaciones.}. Este es un tipo especial  de superficie de visualización que se almacena en la memoria de su tarjeta gráfica. Sólo se puede utilizar en combinación con el parámetro indicador {\ttfamily FULLSCREEN}, de la siguiente forma:
\begin{lstlisting}		
pantalla = pygame.display.set_mode(tamaño, HWSURFACE | FULLSCREEN, 32)
\end{lstlisting}
Esta línea crea una superficie de {\ttfamily hardware} con doble {\ttfamily  búfer}\footnote{En informática, un búfer (del inglés, buffer) es un espacio de memoria, en el que se almacenan datos de manera temporal, normalmente para un único uso (generalmente ocupan un sistema de cola FIFO); su principal función es evitar que el programa o recurso que los requiere, ya sea hardware o software, se quede sin datos durante una transferencia (entrada/salida) de datos irregular o por la velocidad del proceso.}. \\   
% https://es.wikipedia.org/wiki/B%C3%BAfer_de_datos


Las superficies de {\ttfamily hardware} también benefician al parámetro indicador {\ttfamily DOUBLEBUF}. Esto crea efectivamente dos superficies de {\ttfamily hardware}, pero solo una es visible en cualquier momento. 
Las superficies de  {\ttfamily hardware} pueden ser más rápidas que las superficies creadas en la memoria (normal) del sistema,
porque pueden aprovechar más funciones de su tarjeta gráfica para acelerar el blitting.
La desventaja de las superficies de  {\ttfamily hardware} es que no son  compatibles con todas las plataformas.
Suelen funcionar en plataformas Windows pero no tan bien en otras. \\

La siguiente línea crea un doble buffer superficie del hardware:
\begin{lstlisting}		
pantalla = pygame.display.set_mode(SCREEN_SIZE, 
                                   DOUBLEBUF | HWSURFACE | FULLSCREEN, 32)
\end{lstlisting}

A continuación  se describe una  explicación detallada de cada parámetro indicador de visualización:

{\ttfamily pygame.display.set\_mode()}: Esta función inicializa una ventana o pantalla para la visualización. 
Recibe tres parámetros: el tamaño de la ventana, los parámetros indicadores de visualización y la profundidad de color.
{\ttfamily DOUBLEBUF | HWSURFACE | FULLSCREEN}:  son los parámetros indicadores de visualización, que se combinan utilizando el operador {\ttfamily bitwise OR (|)}. \\

%  El operador bitwise OR (|) en Python se utiliza para realizar una operación a nivel de bits donde el resultado es 1 si al menos uno de los bits en la misma posición es 1. Esto significa que si cualquiera de los bits es 1, el bit resultante será 1; de lo contrario, será 0. Es una operación fundamental en programación para manipular datos a nivel de bits de una manera eficiente y precisa.

• {\ttfamily DOUBLEBUF}: Este parámetro indicador configura un doble buffer de visualización. Esto significa que la superficie de visualización se divide en dos partes: un buffer frontal donde se dibuja el juego y un buffer posterior donde se dibuja el juego antes de ser intercambiado con el buffer frontal. Esto ayuda a prevenir el {\ttfamily flicker}\footnote{El "{}flicker"{} { } se refiere a un efecto visual donde los objetos en la pantalla parecen parpadear o titilar. Esto puede ocurrir cuando la pantalla no se actualiza de manera suave y consistente, lo que hace que los objetos se vean intermitentes o inestables.}  y el  {\ttfamily tearing}{\footnote{El "{}tearing"{} { } se refiere a un efecto visual donde la imagen en la pantalla se ve dividida o partida en dos o más partes. Esto ocurre cuando la pantalla se actualiza mientras el juego está dibujando un nuevo cuadro, lo que resulta en que se muestren partes de dos cuadros diferentes al mismo tiempo.} de la visualización. En el contexto de Pygame, "{}flicker"{} y "{}tearing"{} se refieren a problemas de visualización que pueden ocurrir cuando se mueven o actualizan objetos en la pantalla.\\

• {\ttfamily HWSURFACE}: Este parámetro indicador configura una superficie de visualización acelerada por hardware. Esto significa que la superficie de visualización se maneja por el hardware, lo que puede mejorar el rendimiento.\\

• {\ttfamily FULLSCREEN}: Este parámetro indicador configura una visualización en pantalla completa. Esto significa que la ventana cubrirá toda la pantalla y no tendrá una borde ni título.\\



% Algunas posibles causas del flicker en Pygame incluyen:

% Problemas de sincronización entre el juego y la tasa de refresco de la pantalla.
% Actualización incompleta o ineficiente de la pantalla.
% Problemas con la aceleración de hardware.

% Tearing: El "tearing" se refiere a un efecto visual donde la imagen en la pantalla se ve dividida o partida en dos o más partes. Esto ocurre cuando la pantalla se actualiza mientras el juego está dibujando un nuevo cuadro, lo que resulta en que se muestren partes de dos cuadros diferentes al mismo tiempo.

% Algunas posibles causas del tearing en Pygame incluyen:

% Falta de sincronización entre el juego y la tasa de refresco de la pantalla. Problemas con la aceleración de hardware.
% Uso de una superficie de visualización de software en lugar de una superficie acelerada por hardware.

% Tanto el flicker como el tearing pueden ser problemas comunes en juegos y aplicaciones gráficas, y pueden requerir técnicas como el doble búfer, la sincronización vertical (vsync) o el uso de aceleración de hardware para resolverlos.

Tanto el flicker como el tearing pueden ser problemas comunes en juegos y aplicaciones gráficas, y pueden requerir técnicas como el doble búfer, la sincronización vertical  o el uso de aceleración de hardware para resolverlos.\\

%(vsync)
Normalmente, cuando llamas a {\ttfamily pygame.display.update()}, se copia una pantalla completa en la memoria, lo que lleva un poco de tiempo. Las superficies con doble buffered  le permiten cambiar a la nueva pantalla al instante y así hace que su juego se ejecute un poco más rápido.\\

Por último el parámetro indicador de visualización que puede utilizar es {\ttfamily  OPENGL} (\href{http://www.opengl.org/}{\color{vino}http://www.opengl.org/ }) y que lo desarrollaremos con más detalles en la sección \ref{OPENGL}. Se trata de una biblioteca de gráficos
que utiliza el acelerador de gráficos 3D,  y que se encuentra en casi todas las tarjetas gráficas.  La desventaja de usar el  parámetro indicador  {\ttfamily  OPENGL}  es que ya no podrás usar las funciones de gráficos 2D de Pygame. 



Además,  si utiliza una pantalla con doble {\ttfamily buffer}, debe llamar a
\begin{lstlisting}		
pygame.display.flip()
\end{lstlisting}
en lugar de {\ttfamily pygame.display.update()}. Esto hace que la visualización instantánea cambie en lugar de copiar los datos de la pantalla. \\

En la tabla \ref{tabla_Flags} se presenta un resumen de los  parámetros indicadores utilizadas en la creación de la ventana de la función {\ttfamily pygame.display.set\_mode()}. 




%%%%%%%%%%%%%%%%%%%%%%%%%%%%%%%%%%%

% Cualquier imagen está bien para el fondo, siempre que tenga un tamaño mínimo de 640 por 480 (Si es más grande, se recortará el exceso). Para el cursor del mouse, necesitará una imagen más pequeña que encaje cómodamente dentro del fondo; un buen tamaño es 80 por 80. A seguir con el  tema del primer capítulo, el fondo oficial será una imagen de tazones y palillos, y una imagen de un  muy crudo para el cursor del mouse.  

% En la parte superior del script importamos el módulos externos, clases, funciones, etc. que necesitaremos al ejecutar el ejemplo: importar pygame desde pygame.locals importar * la primera línea importa el paquete pygame, que nos da acceso a todos sus submódulos, como como pygame.image y pygame.sound. 

% La segunda línea importa una serie de funciones y constantes (valores que no cambian) en el espacio de nombres de nivel superior. No es imprescindible hacer esto en para usar Pygame, pero es conveniente porque no tenemos que preceder los utilizados con frecuencia valores con el espacio de nombres pygame.

% La última declaración de importación importa una sola función desde sys (un módulo en la biblioteca estándar). Como habrás adivinado, el propósito de la salida es Termine inmediatamente con el guión. Llamarlo hará que la ventana de Pygame desaparezca y Python para cerrar. El script llamará a salir cuando el usuario haga clic en el botón cerrar; de lo contrario, el ¡El usuario no tendría forma de cerrar la ventana!



% pygame.init() Inicializa cada uno de los submódulos en el paquete pygame, que puede cargar controladores y consultar hardware para que Pygame esté listo para usar todos los dispositivos en su computadora. Puedes inicializar sólo los módulos que desea utilizar llamando a la función init en cada submódulo individualmente; por ejemplo, pygame.sound.init() inicializará el módulo de sonido. Esto puede hacer que el


%  La pantalla podría ser una ventana en su escritorio o podría ser la pantalla completa, pero siempre acceda a él a través de un objeto Pygame Surface. La llamada a pygame.display.set\_mode en nuestro script regresa el objeto Superficie que representa la ventana en su escritorio. Se necesitan tres parámetros; sólo se requiere el primero, que debe ser una tupla que contenga el ancho y el alto de la pantalla queremos crear. Nuestra ventana será de 640 × 480 píxeles, que es lo suficientemente grande como para que podamos ver lo que está sucediendo, pero no tan grande como para ocultar gran parte del escritorio. 

%%%%%%%%%%%%%%%%%%%%%%%%%%%%%%%%%%%


% El siguiente parámetro Le damos a set\_mode un valor que contiene las etiquetas parámetros indicadores utilizadas en la creación de la pantalla. 


% parámetros indicadores utilizadas en la creación de la pantalla.
















%%%%%%%%%%%%%%%%%%%%%%%%%%%%%%%%%%%%%%%%%%%
% http://inventwithpython.com/pygame

% al@inventwithpython.com

\section{Resumen}

Estos temas  cubren los conceptos básicos para creación de juegos gráficos con la librería  Pygame. 
Por supuesto,  conocer y manipular estas funciones probablemente no sea suficiente para ayudarte a crear juegos. En los siguientes
capítulos del libro se centran en analizar código fuente de juegos completo. Esto le dará más  ideas formación  de cómo "{}lucen"{} los 
programas de juego completos, para que pueda iniciar  sus propios programas de juegos.

A asumimos que conoce los conceptos básicos de la programación en Python. Si tiene problemas para recordar esos conceptos puede revisar el libro un curso intensivo de Linis Guerrero, descarga gratis....\\

% Pero si todavía dudas, puedes leer el libro "Inventar con Python" (Invent with Python book)  es dirigidos a  personas que son completamente nuevas en la programación, para gratis en línea en http://inventwithpython.com.








%%%%%%%%%%%%%%%%%%%%%%%%%%%%%%%%%%%%%%%%%%%%%%%%%%%%%%%%%%%%%%%%%%%%%%%%%%%%%%%%%%%%%

\chapter{Movimientos, vectores y colisiones}  Introducción a Python y PyGame
% https://docs.w3cub.com/pygame/ref/math#pygame.math.Vector2
módulo pygame para clases vectoriales

En el mundo real, los objetos se mueven de diferentes maneras y en cualquier dirección, dependiendo de lo que estén modelando,
y en nuestro caso que son los juegos debe aproximar esos movimientos para crear una representación virtual convincente y agradable.

Algunos juegos pueden escapar con movimiento poco realista:  por ejemplo  Pac-Man (figura \ref{Pac_Man}), se mueve en línea recta con una velocidad  constante y puede cambiar de dirección en un instante, pero si aplicaras ese tipo de movimiento a un automóvil en un juego 
de conducción, sería destruir la ilusión.
\begin{figure}[h!]
	\centering
	\includegraphics[scale=0.55]{../A Logos de Python/Pac_Man}
	\caption[Juego Pac-Man]{Juego Pac-Man}
	\label{Pac_Man}
\end{figure}
 Después de todo, en un juego de conducción es de esperar que el coche tarde algún
tiempo en alcanzar su máxima capacidad velocidad y definitivamente no debería poder girar 180 grados en un instante!

Para juegos con un toque de realismo, el programador del juego debe tener en cuenta como se  mueven los personajes(vehículos, animales, maquinas, etcétera). 

Veamos un juego de  típico conducción. Independientemente de si el vehículo es una bicicleta, un coche de rally o
un semirremolque, siempre hay una fuerza del motor que lo impulsa hacia adelante. También hay otras fuerzas que actúan 
sobre él. La resistencia de las ruedas que varía según la superficie sobre la que conduzca, por lo que la resistencia del vehículo 
será  diferente para  el barro y para el  asfalto.  Y por supuesto también está la gravedad, que tira constantemente del coche 
hacia la tierra, generando fuerza de rose. Hay muchos factores que que generan cientos de otras fuerzas que se combinan para crear el movimiento de un vehículo, afortunadamente para nosotros, los programadores de juegos, sólo tenemos que simular algunas de estas fuerzas para crear una  ilusión  convincente del movimiento. 


Una vez escrito nuestro código de simulación, podemos aplicarlo a muchos objetos en el juego. 
La gravedad, por ejemplo, afectará a todo (a menos que el juego esté ambientado en el espacio), por lo que podemos aplicar el
código relacionado con la gravedad a cualquier objeto, ya sea una granada de mano arrojada, un tanque que cae de un acantilado o un hacha
volando por el aire.


En este capítulo se describe cómo mover objetos por la pantalla de forma predecible y cómo hacer
ese movimiento consistente en las computadoras de otras personas.

En la próxima sección \ref{fotograma} trataremos la  velocidad de fotogramas, muy importante para crear la ilusión convincente de la animación.


\section{Velocidad de fotogramas}\label{fotograma}

Lo primero que debemos saber sobre el movimiento en un juego de computadora es que nada se mueve realmente, al menos
no en ningún sentido físico. Una pantalla de computadora o un televisor nos presenta una secuencia de imágenes, y
cuando el tiempo entre las imágenes es lo suficientemente corto, nuestro cerebro combina las imágenes para crear la ilusión de fluido movimiento, como un libro animado. 

El número de imágenes, o fotogramas, necesarios para producir un movimiento suave puede variar
de persona a persona. Las películas usan 24 fotogramas (imágenes) por segundo, pero los juegos de computadora tienden a
requerir un cuadro más, una tasa más rápida. 

Treinta fotogramas por segundo es una buena velocidad, pero en términos generales cuanto mayor sea la velocidad de fotogramas, más suave se verá el movimiento. Para un evento deportivo, la escena de un crimen o si hay objetos que se mueven realmente rápido, entonces necesita una mayor velocidad de fotogramas, que puede llegar a 60 FPS.

La velocidad de fotogramas de un juego también está limitada por el número de veces por segundo que el dispositivo de visualización
(como su monitor) se puede actualizar. Por ejemplo, mi monitor LCD tiene una frecuencia de actualización de 60 hercios, lo que
significa que actualiza la pantalla 60 veces por segundo. 

Generar fotogramas más rápido de lo que puede generar la frecuencia de actualización es lo que se conoce como “tearing”, 
donde parte del cuadro siguiente se combina con un cuadro anterior, obtener una buena velocidad de fotogramas generalmente
significa comprometer los efectos visuales, porque cuanto más trabajo  haga su computadora, más lenta será la velocidad de fotogramas.
La buena noticia es que la computadora de escritorio probablemente sea lo suficientemente rápido para generar las imágenes que necesita.



Todos los videos que vemos no son más que fotografías que se muestran a una determinada velocidad para crear el efecto de movimiento. Esas fotografías se llaman fotogramas y el número de fotogramas por segundo de un video es lo que llamamos velocidad de fotogramas.

En general, la velocidad de fotogramas estándar es de 24 a 30 fotogramas por segundo (FPS) para videos como películas y programas de televisión,  y de 60 fotogramas por segundo para grabaciones de eventos deportivos, juegos y otras escenas que implican objetos en rápido movimiento.

Si se dice que un video es de 30 FPS, esto significa que cada segundo  pasan treinta fotogramas consecutivos ante nuestros ojos.


\section{Movimiento en línea recta, vertical y horizontal}


Comencemos examinando el movimiento rectilíneo simple. 
Si movemos una imagen una cantidad fija en la ventana, entonces la imagen aparenta  moverse. Para moverlo horizontalmente, agregaríamos a la coordenada $x$ una velocidad o incremento, de igual forma para moverlo verticalmente el sumaríamos a la coordenada $y$ un incremento. 
Las imágenes en movimiento en 2D a menudo se denominan {\ttfamily sprites}.



El archivo {\ttfamily movimiento\_horizontal.py} presenta un algoritmo para mover una imagen horizontalmente; dibuja una imagen en una la coordenada {\ttfamily x} específica  y luego agrega el valor de 10.0 a cada fotograma, de modo que en el siguiente
marco se habrá desplazado un poco hacia la derecha. Cuando la coordenada {\ttfamily x} llega al borde derecho de la ventana,
se devuelve a 0 para que no desaparezca por completo. 

Algoritmo para el movimiento simple en línea recta (horizontal):

\hspace{-2.85cm}{\ttfamily \scriptsize movimiento}

\vspace{-0.25cm}
\hspace{-2.95cm}{\ttfamily \scriptsize \_horizontal.py}\vspace{-0.87cm}
\begin{lstlisting}		
import pygame
from pygame.locals import *
from sys import exit
archivo_image_fondo = 'pygame.jpg'

perro = 'perro.png'
pygame.init()

pantalla = pygame.display.set_mode((640, 480), 0, 32)
fondo = pygame.image.load(archivo_image_fondo).convert()
perro_sprite = pygame.image.load(perro)

# La coordenada x de nuestro sprite
x = 0.
while True:
    for event in pygame.event.get():
        if event.type == QUIT:
            pygame.quit()
            exit()
    pantalla.blit(fondo, (0,0))
    pantalla.blit(perro_sprite, (x, 100))
    x += 1   # incremento en x, velocidad constante
    # Si la imagen llega al final de la pantalla, se devuelve 
    # inicio de la ventana, en cooredenada x = 0
    if x > 640:
    x -= 640
    pygame.display.update()	
\end{lstlisting}


Si ejecuta, verá la imagen del perro deslizándose de izquierda a derecha. Esto es exactamente el efecto que estábamos 
buscando, pero hay una falla en el diseño. El problema es que no podemos saber exactamente cuánto tiempo llevará dibujar la imagen en la pantalla, y  problema es que el {\ttfamily perro\_sprite} se moverá muy rápido en computadoras  potentes y más lento en máquinas  computadoras menos potentes.

Parece razonablemente fluido porque estamos creando un marco extremadamente simple, pero en un juego el tiempo para 
dibujar un marco variará dependiendo de cuánto actividad que hay en pantalla. Y si queremos un juego que no ralentice (lento un proceso) 
justo cuando se vuelve interesante.  El truco para resolver este problema es hacer que el movimiento se base en el tiempo, empleando el modulo de pygame para el tiempo  {\ttfamily pygame.time.Clock()}.

Para el caso de manipular en forma manual el tiempo, se necesita saber cuanto tiempo tiene ha pasado desde el fotograma anterior
para que podamos colocar todo en la pantalla siguiente. 

\subsection{El tiempo de pygame}

El módulo  {\ttfamily pygame.time.Clock} genera un objeto del tiempo o Reloj, y lo podemos usar para realizar seguimiento  del tiempo.
Para crear un objeto de reloj se llama a su función constructora  {\ttfamily tiempo =  pygame.time.Clock()} 
Una vez creado el objeto reloj, debe llamar al método {\ttfamily tiempo.tick()}, que que devuelve
el tiempo transcurrido en milisegundos (hay 1.000 milisegundos en un segundo)  desde la llamada anterior; ejemplo 
{\ttfamily tiempo\_pasado = tiempo.tick()}, devuelve el tiempo transcurrido desde el ultimo llamado.

La función  {\ttfamily tick()}  también toma un parámetro opcional para la indicar la velocidad máxima de fotogramas, y es
 muy posible que necesite  configurar este parámetro si el juego se ejecuta en el escritorio para que no utilice toda la
  potencia de procesamiento de la computadora:
\begin{lstlisting}
# El juego se ejecutará a un máximo de 30 fotogramas por segundo.
tiempo_FPS= tiempo.tick(30)
\end{lstlisting}

Los milisegundos se utilizan a menudo para cronometrar eventos en los juegos porque puede ser más fácil manejar
valores enteros en lugar de tiempos fraccionarios, y 1000 ticks de reloj por segundo generalmente son lo 
suficientemente precisos para la mayoría de las tareas del juego. 


En lo personal  prefiero trabajar en segundos cuando trato conceptos como la velocidad, porque 250 píxeles por
segundo tiene más sentido para mí que 0,25 píxeles por milisegundo. La conversión de milisegundos a segundos es
simple, dado en milisegundos se divide por 1000, y resulta el valor es segundos:
 {\ttfamily tiempo\_pasado\_segundos = tiempo\_pasado/1000.0}






A continuación tenemos que calcular hasta qué punto el objeto se ha movido en el corto período de tiempo desde el último cuadro, 
y agrega ese valor a la coordenada x. La matemática para esto es sencilla, simplemente multiplica la velocidad del sprite por
{\ttfamily tiempo\_pasado\_segundos}.  Configurado el  movimiento basado en el tiempo y el objeto se moverá  a la misma velocidad independientemente de la velocidad de la computadora en la que lo ejecutas.


El siguiente código muestra el movimiento basado en el tiempo:


\hspace{-2.5cm}{\ttfamily \scriptsize movimiento}

\vspace{-0.24cm}
\hspace{-2.35cm}{\ttfamily \scriptsize \_tiempo.py}
\vspace{-0.8cm}
\begin{lstlisting}
import pygame
from pygame.locals import *
from sys import exit

archivo_image_fondo = 'pygame.jpg'

pygame.init()
pantalla = pygame.display.set_mode((640, 480), 0, 32)
fondo = pygame.image.load(archivo_image_fondo).convert()
sprite = pygame.image.load('perro.png')

# Configurando Reloj  objecto Reloj
tiempo = pygame.time.Clock()

# Velocidad en píxeles por segundo
velocidad = 300

# La coordenada x de nuestro sprite
x = 0.
while True:
    for event in pygame.event.get():
         if event.type == QUIT:
            pygame.quit()
            exit()
    pantalla.blit(fondo, (0,0))
    pantalla.blit(sprite, (x, 100))
    tiempo_pasado = tiempo.tick()

    tiempo_pasado_seconds = tiempo_pasado/1000.0
    distance_moved= tiempo_pasado_seconds * velocidad
    x += distance_moved   

    # Si la imagen llega al final de la pantalla, muévala hacia atrás
    if x > 640:
        x -= 640
    pygame.display.update()
\end{lstlisting}

Con este cálculo la imagen se mueve a una velocidad de 300 fotograma por segundo. El proceso se hizo manual, es la velocidad de un sprite en el juego.

Es importante entender la diferencia entre la velocidad de fotogramas y la velocidad de un sprite en el juego.


\vspace{0.2cm}
\begin{minipage}[t]{0.2\textwidth}
	\centering\raisebox{\dimexpr\topskip-\height}{%
		\includegraphics[scale=0.20]{../A Logos de Python/Advertencia}}  %[width=\textwidth]
	
\end{minipage}%\hfill
\begin{minipage}[t]{0.70\textwidth} 
	{\scriptsize Si no está utilizando Python 3+, asegúrese de dividir por un valor de punto flotante de 1000,0; si 
		no lo hace incluye el punto flotante, el resultado se redondeará hacia abajo al número entero más cercano.
		Entonces, ¿cómo usamos time\_passed\_segundos para mover un objeto? Lo primero que debemos hacer es elegir un
		velocidad para el sprite. Digamos que nuestro sprite se mueve a 300 píxeles por segundo. A esa velocidad, el sprite
		cubrir el ancho de una pantalla de 640 píxeles en 2,13 segundos (640 dividido por 300). }
\end{minipage}
%\vspace{0.2cm}



\subsubsection*{Velocidad de fotogramas (FPS)}
La velocidad de fotogramas se refiere a cuántos fotogramas o frames se muestran por segundo en la pantalla. Esto determina la fluidez de la animación y el movimiento en general del juego.
Por ejemplo, si un juego corre a 30 FPS, significa que se muestran 30 fotogramas cada segundo. A mayor FPS, más fluida se verá la animación, pero también requiere más potencia de procesamiento.
La velocidad de fotogramas es una propiedad global del juego, no de un sprite en particular. Se establece una vez para todo el juego


\subsubsection*{Velocidad de un sprite}
La velocidad de un sprite se refiere a cuán rápido se mueve ese sprite en particular dentro del juego. Esto es independiente de la velocidad de fotogramas.  Por ejemplo, puedes tener un sprite que se mueve a 100 píxeles por segundo, independientemente de si el juego corre a 30 FPS o 60 FPS.  La velocidad de un sprite se puede definir de dos formas en Game Maker Studio:

• Fotogramas por segundo: el sprite avanza 1 fotograma cada cierto tiempo, independiente de los FPS del juego.

• Fotogramas por frame de juego: el sprite avanza 1 fotograma cada frame del juego, por lo que a mayor FPS, más rápido se verá la animación.\\


En resumen, la velocidad de fotogramas (FPS) determina la fluidez general del juego, mientras que la velocidad de un sprite
determina que tan rápido se mueve y anima ese sprite en particular dentro del juego. Son conceptos relacionados pero independientes.\\


En el siguiente código se presenta dos imágenes dos sprites moviéndose en la pantalla, uno con velocidad por píxeles
y el otro por FPS. El de arriba se mueve a 30 fotogramas, por segundo o tan suavemente como lo permita su computadora; 
el otro simula una computadora lenta actualizando sólo cada cinco fotogramas.

Deberías ver que aunque el movimiento es muy entrecortado para el segundo perro\_sprite, en realidad  no  se mueven a la misma velocidad promedio. Entonces, para los juegos que usan movimiento basado en el tiempo, una velocidad de fotogramas lenta será
resultará en una experiencia visual menos placentera, pero en realidad no ralentizará la acción.

El siguiente código presenta dos imágenes en la ventana, uno con velocidad de 30 FPS y el otro con velocidad  
píxeles  por segundos,  basado en el tiempo:

\hspace{-2.5cm}{\ttfamily \scriptsize compara}

\vspace{-0.2cm}
\hspace{-2.7cm}{\ttfamily \scriptsize \_velocidad.py}\vspace{-0.84cm}
\begin{lstlisting}
archivo_image_fondo = 'pygame.jpg'
import pygame
from pygame.locals import *
from sys import exit
pygame.init()
pantalla = pygame.display.set_mode((640, 480), 0, 32)
fondo = pygame.image.load(archivo_image_fondo).convert()
sprite = pygame.image.load('perro.png')
perro = pygame.image.load('perro.png')
perrox = 0
perroy = 200
direction = 'derecho'

# Configurando Reloj  objecto Reloj
tiempo = pygame.time.Clock()
velocidad = 30  # FPS
# Velocidad en pixeles por secondo
velocidad = 300
# La coordenada x de nuestro sprite
y = 300
x = 0.
while True:
    for event in pygame.event.get():
        if event.type == QUIT:
            pygame.quit()
            exit()
    pantalla.blit(fondo, (0,0))
    pantalla.blit(perro, (perrox, 200))         
    pantalla.blit(sprite, (x, 300))

    if direction == 'derecho':
        perrox  += 1
        if perrox  == 640:
            perrox = perrox -640

    tiempo_pasado = tiempo.tick()
    tiempo_pasado_seconds = tiempo_pasado/1000.0
    distance_moved= tiempo_pasado_seconds * velocidad
    x += distance_moved  # distancia en tiempo del ultimo 

    # Si la imagen llega al final de la pantalla, muévala hacia atrás
    if x > 640:
        x -= 640

    pygame.display.update()
    tiempo.tick(velocidad) # todo el juego
\end{lstlisting}

Al ejecutar puede notar la diferencia de movimientos. 

\section{Movimiento en cualquier dirección, y colisión sencilla}

El movimiento en línea recta es útil, pero un juego probablemente se sería bastante aburrido si todo se moviera en forma 
horizontal y  vertical. Una simulación de movimientos de un objeto necesita moverse en cualquier dirección, para lograr este objetivo 
se hacen algunos ajustes en las coordenadas {\ttfamily x} y {\ttfamily y}.\\

El siguiente script simula un objeto que se mueve en dirección diagonal, agregando el movimiento basado en el tiempo en ambas coordenadas. 
Además, se detecta la colisión del sprite con los bordes de la ventana y se hace rebotar en sentido contrario, esta es una  “detección de colisiones”, aunque muy trivial; de esta forma  el sprite lugar de llevarlo de regreso a una posición inicial, cuando llega al borde de la ventana, el sprite rebota en el sentido contrario.\\

Para lograr este rebote, primero tenemos que detectar que el sprite ha tocado un borde, esto se hace con algunos cálculos 
matemáticos simples en las coordenadas. Si la coordenada x es menor o igual que 0, sabemos que se encuentra el lado   
izquierdo de la pantalla porque la coordenada del borde izquierdo es 0.  Si x es mayor que la suma de la posición actual en {\ttfamily x} del objeto más el ancho del sprite, sabemos que el sprite ha tocado el lado derecho de la pantalla. El caso la coordenada {\ttfamily y}  es similar, pero usamos la altura del sprite.

El siguiente archivo {\ttfamily py}, el sprite es un circulo de radio igual a 40 rebotando en la ventana pygame:


\hspace{-2.2cm}{\ttfamily \scriptsize rebote}

\vspace{-0.2cm}
\hspace{-2.4cm}{\ttfamily \scriptsize \_objeto.py}\vspace{-0.84cm}
\begin{lstlisting}
import pygame
import sys
import time
# Inicializar Pygame
pygame.init()

# Definir colores
BLACK = (0, 0, 0)
WHITE = (255, 255, 255)
CELESTE = (170,238,187)

tiempo= pygame.time.Clock()
velocidad = 200

# Tamaño de la pantalla
ancho, alto = 640, 480
SCREEN = pygame.display.set_mode((ancho, alto))
pygame.display.set_caption("Rebote de objeto")

# Tamaño del objeto
Obj_ancho, Obj_alto = 40, 40

# Posición inicial del objeto
#obj_x = ancho//2 - Obj_ancho//2
#obj_y = alto//2 - Obj_alto//2
obj_x = 50
obj_y = 100
# Velocidad del objeto
vel_x = 5
vel_y = 5

# Bucle principal
while True:
    for event in pygame.event.get():
        if event.type == pygame.QUIT:
            pygame.quit()
            sys.exit()

    # Actualizar posición del objeto
    obj_x += vel_x
    obj_y += vel_y

    # Rebotar en los bordes
    if obj_x <= 40 or obj_x + Obj_ancho >= ancho:
       vel_x = -vel_x
    if obj_y <= 40 or obj_y + Obj_alto >= alto:
       vel_y = -vel_y

    # Limpiar pantalla
    SCREEN.fill(CELESTE)

    # Dibujar objeto
    #pygame.draw.rect(SCREEN, WHITE, (obj_x, obj_y, Obj_ancho, Obj_alto))
    pygame.draw.circle(SCREEN,WHITE,(obj_x, obj_y),40)

    # Actualizar pantalla
    pygame.display.update()
    #pygame.display.flip()
    tiempo.tick(velocidad)
\end{lstlisting}

El sprite es un circulo de radio 40, rebotando en la ventana, una  “detección de colisiones” simple.\\

En el siguiente código se agregamos un obstáculo en el centro de la ventana para generar colisión del objeto sprite con el obstáculo, se detecta la colisión y generar el rebote el sprite: 

\hspace{-2.5cm}{\ttfamily \scriptsize obstáculo}

\vspace{-0.2cm}
\hspace{-2.4cm}{\ttfamily \scriptsize \_rebote.py}\vspace{-0.84cm}
\begin{lstlisting}	
# -*- coding: utf-8 -*-
"""
Created on Mon Apr  3 17:07:13 2023

@author: Linis
"""
import pygame
import sys
import time
# Inicializar Pygame
pygame.init()

# Definir colores
BLACK = (0, 0, 0)
WHITE = (255, 255, 255)
CELESTE = (170,238,187)

# Tamaño de la pantalla
WIDTH, HEIGHT = 640, 480
SCREEN = pygame.display.set_mode((WIDTH, HEIGHT))
#pygame.display.set_caption("Rebote de objeto")

# Tamaño del objeto
OBJ_WIDTH, OBJ_HEIGHT = 20, 20

# Posición inicial del objeto
obj_x = WIDTH // 2 - OBJ_WIDTH // 2
obj_y = HEIGHT // 2 - OBJ_HEIGHT // 2

# Velocidad del objeto y timpo
tiempo= pygame.time.Clock()
velocidad = 200
vel_x = 5
vel_y = 5

pygame.display.set_caption("Colisión circular")

# Tamaño y posición del objeto circular
circle_radius = 40
circle_x = WIDTH//2
circle_y = HEIGHT//2

# Tamaño y posición del obstáculo rectangular en el centro
obstacle_width = 100
obstacle_height = 100
obstacle_x = WIDTH//2 - obstacle_width//2
obstacle_y = HEIGHT//2 - obstacle_height//2

# Bucle principal
while True:
    for event in pygame.event.get():
        if event.type == pygame.QUIT:
            pygame.quit()
            sys.exit()

    # Actualizar posición del objeto
    obj_x += vel_x
    obj_y += vel_y

    # Rebotar en los bordes
    if obj_x <= 20 or obj_x + OBJ_WIDTH >= WIDTH:
        vel_x = -vel_x
    if obj_y <= 20 or obj_y + OBJ_HEIGHT >= HEIGHT:
        vel_y = -vel_y

    # Limpiar pantalla
    SCREEN.fill(CELESTE)

    # Dibujar objeto circular
    pygame.draw.circle(SCREEN,WHITE,(obj_x, obj_y),20)
    obstacle_x = obj_x
    obstacle_y = obj_y
    obstacle_width = 2.5
    obstacle_height = 2.5
    pygame.draw.circle(SCREEN, BLACK, (circle_x, circle_y), 5)
   
    # Detectar colisión entre el obstaculo y el pelota
    distance_x = abs(circle_x - obstacle_x - obstacle_width/2)
    distance_y = abs(circle_y - obstacle_y - obstacle_height/2)

    if (distance_x <= (obstacle_width/2 + circle_radius) and
        distance_y <= (obstacle_height/2 + circle_radius)):
        # Colisión detectada
        print("Colisión detectada")
        # Rebotar en los bordes
        if obj_x ==  WIDTH // 2 or obj_x + OBJ_WIDTH >= WIDTH:
            vel_x = -vel_x
        if obj_y == HEIGHT // 2 or obj_y + OBJ_HEIGHT >= HEIGHT:
            vel_y = -vel_y

    # Actualizar pantalla
    pygame.display.flip()
    tiempo.tick(velocidad)
\end{lstlisting}	


Hemos visto que agregar un valor basado en el tiempo a las coordenadas x e y de un objeto crea un
movimiento en cualquier dirección. Sin embargo la mejor manera de trabajar el movimiento es con vectores.
En PyGame los vectores se tratan con clases, una clase {\ttfamily Vector2 } para dos dimensiones y una clase {\ttfamily Vector3 } para tres dimensiones, en la siguiente sección \ref{vectores} se estudia  el tema de clases de vectores pygame y sus métodos.



%%%%%%%%%%%%%%%%%%%%%%%%%%%%%%%%%%%%%%%%%%%%%%%%%%%%%%%%%%%%%%%%%%%%%%%%%%%%%%%%%%%%%




\section{Explorando vectores}   \label{vectores}

Para generar el movimiento diagonal usamos dos valores de velocidad, uno para la coordenada {\ttfamily x} y otra para  {\ttfamily y}. 
Estos dos valores combinados forman un punto, y si tiene dirección es lo que se conoce como vector. Un vector es algo que los desarrolladores de juegos tomaron prestado de las matemáticas y se utilizan con frecuencia en Juegos 2D y 3D.

Los vectores son similares a los puntos en que ambos tienen un valor para {\ttfamily x} y otra para  {\ttfamily y} (en 2D), 
pero se usan para diferentes propósitos. Un punto en la coordenada (10, 20) siempre será el mismo punto en la pantalla, pero un vector de
(10, 20) significa sumar 10 a la coordenada x y 20 a la coordenada y desde la posición inicial. Entonces podrías pensar en un punto como si fuera un vector desde el origen (0, 0).


\subsection*{Creando vectores}

Puedes crear un vector a partir de dos puntos cualesquiera digamos $ A(x_{0},y_{0}) = (x_{0},y_{0})$ y  $B(x_{1},y_{1}) =(x_{1},y_{1}) $, restando a los valores del primer punto $A$ los del segundo punto $B$:   
$ vector = \overrightarrow{AB} = \overrightarrow{AB}(x_{1} - x_{0},y_{1} - y_{0})$.


% https://www.latex.um.es/curso_2000/sesiones/sesion5a.pdf

Ilustremos con un ejemplo de un juego ideal. 
Supongamos un personaje de un juego de batallas, un soldado del futuro llamado Alpha que debe  destruir un centinela droide clase Beta con un rifle de francotirador. Alfa está escondiéndose detrás de un arbusto en la coordenada $A(10, 20)$ y apuntando a Beta en la coordenada $B(30, 35)$.  Para calcular un vector $\overrightarrow{AB}$ dirigido al objetivo, Alfa tiene que restar los componentes de $B$ a las componentes de $A$.

\begin{figure}[h!]
	\centering
	\includegraphics[scale=0.35]{../A Logos de Python/vectorSuma}
	\caption{Vector $\overrightarrow{AB} $}
	\label{vectorS}
\end{figure}



Es decir, el vector  $\overrightarrow{AB}$  es igual a  $(30 - 10, 35 - 20) = (20, 15)$, resultando
 $\overrightarrow{AB}(20, 15)$. Esto nos dice que para llegar de $A$ a $B$ tendríamos que recorrer 20 unidades en la dirección $x$ y
15 unidades en la dirección $y$, como se muestra en la figura \ref{vectorS}. El juego necesitaría esta información para dirigir 
un  proyectil, un arma, o  un rayo láser entre los dos puntos.

\subsection*{Tratamiento de vectores en PyGame}  %  Vectores en Python


En Python, los vectores se pueden representar utilizando listas o matrices. La biblioteca {\ttfamily NumPy} es una forma
común de representar vectores y ofrece operaciones de álgebra lineal. Puedes crear vectores utilizando el método
{\ttfamily numpy.array()} y puedes realizar operaciones básicas con vectores en Python, como sumas y restas.


Sin embargo, en Pygame los vectores se utilizan para representar posiciones y direcciones en el espacio, y la librería  {\ttfamily pygame.math} proporciona dos clases  {\ttfamily Vector2} y {\ttfamily Vector3} para vectores bidimensionales y tridimensionales respectivamente. Estas clases ofrecen operaciones numéricas elementales, como sumas, restas, y operaciones más complejas como productos escalares, productos vectoriales y ángulo entre vectores. 

Para definir un vector, ahora podemos usar objetos generados por la clase {\ttfamily Vector2}. Por ejemplo, una llamada a 
{\ttfamily v1 = pygame.Vector2(10, 20)} produce un objeto Vector2 llamado {\ttfamily v1}.  Veamos algunos ejemplos del uso de la clase Vectorial:
\begin{lstlisting}
import pygame
v1 = pygame.Vector2(10, 20)
v2 = pygame.Vector2(3, 4)
V_sum = v1 + v2
V_diff = v1 - v2
V_produc =  v1*v2  # producto escalar
V_produc_cruz = v1.cross(v2)  # producto cruz
\end{lstlisting}
Además, podemos referirnos a los componentes del vector individualmente como {\ttfamily v1.x} y {\ttfamily v1.y}. \\

Esta es una forma clara y didáctica para tratar vectores en PyGame.  Claro que podemos utilizar {\ttfamily NumPy}, pero lo más didáctico por ahora es usar las clases vectoriales del modulo integrado de pygame  {\ttfamily pygame.math}, {\ttfamily Vector2} y {\ttfamily Vector3}.  
A medida que avancemos en los ejemplos, cubriremos  vectores y la matrices.\\


Los métodos de la clase vectorial de Pygame, como Vector2 y Vector3, proporcionan funcionalidades esenciales 
para operar con vectores en dos y tres dimensiones respectivamente, permitiendo realizar cálculos y manipulaciones 
matemáticas sobre ellos de manera eficiente. Para mas información visitar la página oficial  {\scriptsize
\href{https://www.pygame.org/docs/ref/math.html#pygame.math.Vector3.as_polar}{\color{vino}{https://www.pygame.org/docs/ref/math.html{\ttfamily\#}pygame.math.Vector2.as\_polar}}}.\\
 
En la siguiente tabla \ref{Vector2} se presenta la clase {\ttfamily Vector2}  con los métodos  más comunes y su propósito:\\ 


%\vspace{-0.4cm}
\begin{center}\noindent Table \ref{Vector2}. Método de la clase vectorial Vector2:\vspace{0.1cm}	
	{\footnotesize{{\scriptsize 
				\hspace{-0.7cm}	
				\begin{tabular}{l l l }
					\multicolumn{1}{|>{\columncolor{vino}} l }{\color{white}\bf Método }& 
					\multicolumn{1}{|>{\columncolor{vino}} l }{\color{white}\bf Propósito} \\ 
			     	v1.dot(v2)              & calcula el producto escalar \\
					v1.cross(v2)            & calcula el producto cruzado o vectorial \\ 
					v1.magnitude()          &  devuelve la magnitud euclidiana del vector \\
					v1.magnitude\_squared() &   devuelve la magnitud al cuadrado del vector \\
				    v1.length()             &   devuelve la longitud euclidiana del vector \\
					v.1normalize()            &  devuelve un vector con la \\
					                          &  misma dirección pero de longitud 1 \\
				    v1.normalize\_ip()        &  normaliza el vector en su \\
				                              &  lugar para que su longitud sea 1  \\
					v1.is\_normalized()       &  prueba si el vector está normalizado\\	
					                          &  es decir, tiene una longitud == 1 \\			
					v1.scale\_to\_length(float)  &  escala el vector a una longitud determinada \\
			    	v1.reflect(v2)          &  devuelve un vector reflejado de una normal dada  \\ 
					v1.reflect\_ip(v2)      &  reflejar el vector de una normal dada en su lugar   \\
				    v1.distance\_to(v2)       &  calcula la distancia euclidiana a un vector dado \\ 
					v1.lerp(v2, float)        &  devuelve una interpolación lineal al vector dado \\    
					v1.rotate(angle)          &  gira un vector en un ángulo dado en grados \\
				    v1.rotate\_rad(angle)     &  gira un vector en un ángulo dado en radianes  \\   
				    v1.rotate\_ip(angle)      &  gira el vector en un ángulo \\ 
				                              &  determinado en grados en su lugar  \\
				    v1.angle\_to(v2)          &   Calcula el ángulo entre dos vectores dado en grados  \\				
					\arrayrulecolor{vino}\hline 
					\label{Vector2}  
				\end{tabular}		
				%  	\vspace{0.2cm} 
	}}} 
\end{center}

El siguiente código crear un vector entre dos puntos {\ttfamily A} y {\ttfamily B} dados,  mediante las clases vectoriales de pygame:
\begin{lstlisting}
A = (10.0, 20.0,0)      # Puntos A
B = (30.0, 35.0,0)      # Punto B
A = pygame.Vector3(A)   # Vector PyGame A
B = pygame.Vector3(B)   # Vector PyGame B

AB = B-A                #  Crea un vector de A a B
A.distance_to(B)        # Distancia de A a B

print("Vector AB =", B-A)
print("Distancia de A a B =", A.distance_to(B)," unidades")
print("Coordenadas \n x, y = ", AB.x,",",AB.y)
\end{lstlisting}

La ejecución de este código  produce el siguiente resultado:
{\ttfamily\scriptsize\scriptsize 
	\hspace{-0.65cm} \vspace{-0.3cm}	
	{\color{blue}
\begin{verbatim}	
Vector AB = [20, 15, 0]
Distancia de A a B = 25.0 unidades
Coordenadas 
x, y =  -1.0 , 1.0	
\end{verbatim}
}}


\subsection{Vectores y el movimiento}

Ahora que tenemos el conocimiento básico sobre vectores pygame, podemos usarlos para mover objetos o personajes del juego 
de diversas maneras, e implementar magnitudes físicas simples basadas en la fuerza que simulen más realidad el juego.

\subsubsection{Movimiento diagonal}  %  crear un movimiento diagonal más preciso  .normalize()
	
Para mover un sprite de una posición a otra  en la pantalla, a una velocidad constante.   
El primer paso es crear un vector a partir de la posición actual al destino ({\ttfamily AB = B-A}), necesitamos 
información de la dirección de este vector pero no la magnitud, por lo que lo normalizamos el vector ({\ttfamily AB.normalize()}), para
obtener el rumbo (sentido) del sprite.  
\begin{lstlisting}	
import pygame	
pygame.init()

tiempo = pygame.time.Clock()   # Objeto tiempo, velocidad de FPS	
velocidad = 250	#píxeles
A = (60, 50,0)   # Punto A
B = (100, 100,0)  # Punto B
A = pygame.Vector3(A)  # vector A
B = pygame.Vector3(B)  Vector B
AB = B-A                      # Vector de A a B. Posición
rumbo = AB.normalize()        #  dirección 
\end{lstlisting}

Calculamos también {\ttfamily velocidad*tiempo\_pasado\_segundos} qué indica cuanto se ha movido el spritedel anterior cuadro, y
luego lo  multiplicamos por el vector de rumbo. El vector resultante nos da el cambio desde la posición  anterior en {\ttfamily x} e {\ttfamily y}, luego lo agregamos a la posición actual (recorrido) del sprite.



\begin{lstlisting}
tiempo_pasado = tiempo.tick()
tiempo_pasado_segundos = tiempo_pasado/1000.0
velocidad  = velocidad_I*tiempo_pasado_segundos 
cambio = rumbo*velocidad      #  cambio en x e y
\end{lstlisting}

Ahora dentro del bucle principal del juego calculamos qué tanto se ha movido el sprite con velocidad * tiempo\_pasado\_segundos, 
luego multiplicamos por el vector de rumbo. El vector resultante nos da el cambio en x e y desde el cuadro anterior. Este cambio lo agregamos a la posición (recorrido) del sprite, pero en el bucle principal del juego.



\begin{lstlisting}
# Rebotar en los bordes
if recorrido.x <= 20 or recorrido.x + OBJ_WIDTH >= WIDTH:
    cambio.x = -cambio.x       
if recorrido.y <= 20 or recorrido.y + OBJ_HEIGHT >= HEIGHT:
    cambio.y = -cambio.y
cambio = pygame.Vector3(cambio.x,cambio.y,0)
\end{lstlisting}
En el bucle principal agregamos el cambio a la posición del sprite, y para segir la idea anterior colocamos una “detección de colisión”, para que el sprite rebote en la ventana PyGame.

Ahora presentamos todo el trabajo en el siguiente script {moviniento\_vectorial.py}

\hspace{-2.8cm}{\ttfamily \scriptsize movimiento}

\vspace{-0.2cm}
\hspace{-2.4cm}{\ttfamily \scriptsize \_vectorial.py}\vspace{-0.84cm}
\begin{lstlisting} 
archivo_imagen_fondo = 'pygame.jpg'
import pygame, sys
from pygame.locals import *

tamaño = (640, 480)
CELESTE = (170,238,187)      # Color azul Celeste 
WHITE =(255,255,225)
velocidad_I = 250  #  Configurando fotograma por segundos

pygame.init()
tiempo = pygame.time.Clock() # Objeto tiempo, velocidad de FPS

# Inicializar la pantalla
Obj_ancho, Obj_alto = 20, 20
ancho, alto = 640, 480
pantalla = pygame.display.set_mode((ancho, alto), 0, 32)
pygame.display.set_caption("Movimiento vectorial")
fondo = pygame.image.load(archivo_imagen_fondo).convert()

velocidad = 250  	   #    píxeles  
A = (60, 50,0)         #  Punto A
B = (100, 100,0)       #  Punto B
A = pygame.Vector3(A)  #  vector A
B = pygame.Vector3(B)  #  Vector B
recorrido = B-A        #  AB = B-A,  Vector de A a B.  Posición
rumbo = recorrido.normalize()        #  dirección o rumbo

tiempo_pasado = tiempo.tick()
tiempo_pasado_seconds = tiempo_pasado/1000.0
velocidad  = tiempo_pasado_seconds *velocidad_I
cambio = rumbo*velocidad      

while True:
    for event in pygame.event.get():
        if event.type == QUIT:
            pygame.quit()
            sys.exit()
    pantalla.fill(CELESTE )
    pantalla.blit(fondo, (0,0))
    pygame.draw.circle(pantalla,WHITE,(recorrido.x,recorrido.y),20)

    # Rebotar en los bordes
    if recorrido.x <= 20 or recorrido.x + Obj_ancho >= ancho:
        cambio.x = -cambio.x       
    if recorrido.y <= 20 or recorrido.y + Obj_alto >= alto:
        cambio.y = -cambio.y
    cambio = pygame.Vector3(cambio.x,cambio.y,0)
    recorrido += cambio   
    pygame.display.update()
    tiempo.tick(velocidad_I)
\end{lstlisting}

El script implementa un movimiento basado en el tiempo utilizando vectores. 
Cuando lo ejecutes, verás un sprite moviéndose  en la ventana, pero una vez que haces clic en la pantalla, el sprite inicia el movimiento en la posición indicada por el mouse, con la dirección anterior,  el código calculará un vector para el nueva posición y establece el objeto en movimiento en 250 píxeles por segundo.

%  print(sprite.get_size() )   

\hspace{-2.8cm}{\ttfamily \scriptsize posic\_mous}

\vspace{-0.2cm}
\hspace{-2.8cm}{\ttfamily \scriptsize \_vectorial.py}\vspace{-0.84cm}
\begin{lstlisting}
archivo_imagen_fondo = 'pygame.jpg'
import pygame, sys
from pygame.locals import *
CELESTE = (170,238,187)
tiempo = pygame.time.Clock() # Objeto tiempo, velocidad de FPS
ancho = 640
alto = 480
velocidad_I = 250
pygame.init()
pantalla = pygame.display.set_mode((ancho,alto), 0, 32)
fondo= pygame.image.load(archivo_imagen_fondo).convert()
pygame.display.set_caption("Posición vectorial del mouse")
#nave alinen2 
sprite_image = "cubo03.jpg"
Obj_ancho = 42
Obj_alto = 40.5
sprite = pygame.image.load(sprite_image).convert_alpha()
clock = pygame.time.Clock()
posicion = pygame.Vector2(100.0, 100.0)
recorrido = posicion
rumbo = posicion.normalize()
cambio = rumbo
print(posicion)
speed = 250
while True:
    pantalla.fill(CELESTE)
    pantalla.blit(sprite, (posicion.x, posicion.y))
    for event in pygame.event.get():
        if event.type == QUIT:
            pygame.quit()
            sys.exit()
        if event.type == MOUSEBUTTONDOWN:             
            destination = pygame.Vector2(event.pos)           
            if destination.x == 0 or destination.y ==0:
                destination = (pygame.Vector2(event.pos) + 
                               (pygame.Vector2(sprite.get_size())))
            if (destination.x >= ancho - Obj_ancho  or 
                destination.y >= alto -Obj_alto): 
                destination = (pygame.Vector2(event.pos) - 
                              (pygame.Vector2(sprite.get_size())))
            posicion = destination
            cambio = destination.normalize()
    # Rebotar en los bordes
    if posicion.x <= 0 or posicion.x + 2*Obj_ancho >= ancho:
        cambio.x = -cambio.x       
    if posicion.y <= 0 or posicion.y + 2*Obj_alto >= alto:
       cambio.y = -cambio.y
    cambio = pygame.Vector2(cambio.x,cambio.y)   

    pantalla.blit(fondo, (0,0)) 
    pantalla.blit(sprite, (posicion.x, posicion.y))  

    posicion += cambio
    pygame.display.update()
    tiempo.tick(velocidad_I)    
\end{lstlisting}	


El siguente script {\ttfamily hacia\_el\_mouse.py} implementa un movimiento basado en el tiempo utilizando vectores. Cuando lo ejecutes, verás un sprite moviéndose  en la ventana, pero una vez que haces clic en la pantalla, el código calculará un vector para el nueva posición. 

\hspace{-2.4cm}{\ttfamily \scriptsize hacia\_el}

\vspace{-0.2cm}
\hspace{-2.4cm}{\ttfamily \scriptsize mouse.py}\vspace{-0.84cm}
\begin{lstlisting}		
archivo_imagen_fondo = 'pygame.jpg'
sprite_image = "alien03.jpg"
Obj_ancho = 24
Obj_alto = 20
ancho = 640
alto = 480
CELESTE = (170,238,187)

import pygame, sys
from pygame.locals import *
pygame.init()
tiempo = pygame.time.Clock() # Objeto tiempo, velocidad de FPS

pantalla = pygame.display.set_mode((ancho,alto), 0, 32)
fondo= pygame.image.load(archivo_imagen_fondo).convert()
pygame.display.set_caption("Hacia la posición del mouse")
sprite = pygame.image.load(sprite_image).convert_alpha()

posicion = pygame.Vector2(100.0, 100.0)
recorrido = posicion
rumbo = posicion.normalize()
cambio = rumbo
destination = posicion
while True:
    pantalla.fill(CELESTE)
    pantalla.blit(sprite, (recorrido.x, recorrido.y)) 
    for event in pygame.event.get():
        if event.type == QUIT:
            pygame.quit()
            sys.exit()

        if event.type == MOUSEBUTTONDOWN:             
            destination = pygame.Vector2(event.pos)           
            posicion = destination - recorrido 
            cambio = posicion.normalize()
    recorrido += cambio
    # Rebotar en los bordes
    if recorrido.x <= 0 or recorrido.x + 2*Obj_ancho >= ancho:
        cambio.x = -cambio.x       
    if recorrido.y <= 0 or recorrido.y + 2*Obj_alto >= alto:
        cambio.y = -cambio.y
    pantalla.blit(fondo, (0,0)) 
    pantalla.blit(sprite, (recorrido.x, recorrido.y))  

    pygame.display.update()
    tiempo.tick(200)
\end{lstlisting}
Si vuelves a hacer clic, se calculará un nuevo vector, y el objeto cambiará su dirección hacia la posición del mouse. 


\section{Resumen}

Mover un objeto, o cualquier otra cosa en la pantalla, requiere que agregues pequeños cambios o valores a las coordenadas en cada una, pero si desea que el movimiento sea suave y consistente, debe basarse en el actual tiempo, o más específicamente, el tiempo transcurrido desde el último fotograma. Conocer  movimiento basado en el tiempo también es importante para ejecutar el juego en una gama tan amplia de computadoras como sea posible; las computadoras pueden variar mucho en el número de fotogramas por segundo que pueden generar.

Se ha tratado vectores con la clases {\ttfamily Vector2} y  {\ttfamily Vector3}, que son una parte esencial de la caja de herramientas de cualquier desarrollador de juegos. Los vectores simplifican gran parte de los cálculos los tendrás que hacer al escribir tu juego y los encontrarás notablemente versátiles. 


Las técnicas para moverse en dos dimensiones se extienden fácilmente a las tres dimensiones. Encontrarás que la
clase  {\ttfamily Vector3} contiene muchos de los métodos utilizados en la clase {\ttfamily Vector2} pero con la componente adicional en {\ttfamily z}. Ahora sería un buen momento para empezar a experimentar moviendo elementos en pantalla y mezclando varios
tipos de movimiento. Puedes divertirte mucho creando gráficos de tus amigos y familiares en pantalla y
¡hacer que se deslicen y reboten! 



Pendiente  el teclado y el joystick, a sprites, para que el jugador pueda interactuar con el mundo del juego.\\


En el siguiente capítulo, se estudiara el tema relacionado con gráficos en tres dimensiones, ahora incluimos una conocida como  {\bf{ OpenGL}} (Open Graphics Library). Es la  API de gráficos más adoptada, una librería  multiplataforma, compatible con múltiples sistemas operativos, incluyendo Windows, macOS y Linux; los resultados visuales son consistentes en cualquier sistema operativo compatible. Además, OpenGL se puede utilizar junto con una variedad de lenguajes de alto nivel mediante enlaces: bibliotecas de software que unen dos lenguajes de programación, permitiendo que funciones de un lenguaje se utilicen en otro. 




\chapter{Interacción del Jugador con el juego}   % Capítulo 6  Will  y Harrison

Hay una variedad de formas en que el jugador puede interactuar con un juego, por ejemplo, el teclado, mouse, joysticks y joypads, etc;   el este capítulo se cubren algunos dispositivos de entrada en detalle.


Además de recuperar información de los dispositivos, también exploraremos cómo traducir la actividad del jugador con
eventos significativos en el juego. Esto es extremadamente importante para cualquier juego: independientemente de lo 
bien que se vea y suene un juego, también debe ser fácil interactuar con él.



\section*{Controlando el juego}

En una computadora doméstica típica, podemos confiar en que habrá un teclado y un mouse. Esta es la configuración preferida
por los fanáticos de los juegos de disparos en primera persona a quienes les gusta controlar el movimiento 
de la cabeza (es decir, mirar a su alrededor) con el mouse y mantener una mano en el teclado para controles 
direccionales y disparos. Los teclados se pueden utilizar para el movimiento asignando una tecla para las 
cuatro direcciones básicas: arriba, abajo, izquierda y derecha. Se pueden indicar otras cuatro direcciones
presionando estas teclas en combinación (arriba + derecha, abajo + derecha, abajo + izquierda, arriba + izquierda).
Pero teniendo ocho direcciones sigue siendo muy limitante para la mayoría de los juegos, por lo que no son ideales 
para juegos que necesitan un poco de delicadeza para jugar.

La mayoría de los jugadores prefieren un dispositivo analógico como el mouse. El  estándar es ideal para controlar la 
direccional porque puede detectar con precisión cualquier cosa, desde pequeños ajustes hasta barridos rápidos en 
cualquier dirección.

Si alguna vez has jugado a un juego de disparos en primera persona con solo el control del teclado, apreciarás la diferencia.
Al presionar la tecla "rotar a la izquierda", el jugador gira como un robot con una velocidad constante, lo cual no es
casi tan útil como poder darse la vuelta rápidamente para disparar a un monstruo que se acerca por un lado.

El teclado y el mouse funcionan bien para juegos, lo cual quizás sea un poco irónico porque ninguno de los dos dispositivos
fue diseñado pensando en los juegos. Sin embargo, los joysticks y joypads, fueron diseñados exclusivamente para juegos y han
evolucionado junto con los juegos. Los primeros joysticks se inspiraron en los controles utilizados en
avión y tenía palancas direccionales simples con un solo botón. Eran populares entre las consolas de juegos, pero a los jugadores les resultaba incómodo tener que agarrar la base del joystick con una mano mientras movía el otro. 

Esto llevó al desarrollo de joypads, que podían sostenerse con ambas manos y así daban al jugador un fácil acceso 
a los controles con los dedos y los pulgares. Los primeros joypads tenían un control direccional encendido
por un lado y botones de disparo por el otro. Hoy en día los joypads cuentan con numerosos botones ubicados por todas partes
hay un dedo libre para presionarlos, junto con varias palancas. Los pads direccionales clásicos todavía se encuentran en
joypads, pero la mayoría también tiene palancas analógicas que pueden detectar ajustes finos. 

Muchos también tienen retroalimentación, características que pueden agregar una dimensión adicional a 
los juegos al hacer que el joypad tiemble o retumbe en respuesta a eventos en la pantalla. Sin duda, en el futuro se 
agregarán otras funciones a los joypads que mejorarán aún más mejorar la experiencia de juego.


Hay otros dispositivos que se pueden utilizar para juegos, pero la mayoría imitan los dispositivos de entrada estándar.
Entonces, si su juego se puede jugar con el mouse, también se podrá jugar con estos dispositivos similares a un mouse. 

% Los dispositivos de entrada se mejoran constantemente; considerar el estado actual de la realidad virtual, como por ejemplo con el Oculus Rift, que simula algunos de los que esperaría cuando opera el mouse en “vista libre” en juegos. 

Por ahora discutiremos el manejo de la entrada con teclados, PyGame tiene el módulo {\ttfamily pygame.key} para detecta los eventos de pulsaciones  en el teclado, y lo describimos en la siguiente sección.



\section{Comprender el teclado para controlar el juego}

La mayoría de los teclados que se utilizan hoy en día son teclados QWERTY, llamados así porque las seis letras iniciales 
de la primera fila que deletrea QWERTY. Existen variaciones entre marcas; las teclas pueden tener formas y tamaños ligeramente
diferentes, pero tienden a estar aproximadamente en la misma posición en el teclado. Esto es bueno, ya que en la computadora
los usuarios no tienen que volver a aprender a escribir cada vez que compran un teclado nuevo,  los teclados que he
utilizado han tenido cinco filas para teclas de escritura estándar: una fila para teclas de función F1–F12 y cuatro teclas 
para moviendo un cursor por la pantalla. También tienen un "teclado numérico", que es un bloque de teclas para ingresar números
y hacer algunas operaciones.  El teclado numérico suele omitirse en los teclados de portátiles para ahorrar espacio, porque las teclas numéricas están  duplicadas en la parte principal del teclado. 

Aunque es el más común, QWERT Y no es la única distribución de teclado; hay otros teclados
como AZERT Y y Dvorak; que se pueden utilizar las mismas constantes del teclado, pero las teclas pueden estar en diferentes
ubicaciones en el teclado.\\

Para detectar o identificar la teclas se usa el módulo . En la próxima sección se estudia las pulsaciones de teclas 
como  evento PyGame.
% Además, le puedes dar al jugador la opción de seleccionar sus propias teclas  para usarlas en el juego, puede seleccionar los controles que sean más apropiados para su teclado.


\subsection{Detección de pulsaciones de teclas}

Hay dos formas para detectar la pulsación de una tecla en PyGame. Una forma es manejar eventos {\ttfamily KEYDOWN}, que se emiten
cuando se presiona una tecla, y la otra son eventos  {\ttfamily KEYUP}, que se emiten cuando se suelta la tecla. 

Esto es genial para escribir texto porque siempre obtendremos eventos de teclado incluso si se ha presionado y 
soltado una tecla en el tiempo transcurrido desde el último fotograma. 

Los eventos también capturarán toques muy rápidos de la tecla para los botones de disparo. Pero cuando usamos el teclado
como entrada para movimiento, simplemente necesitamos saber si la tecla está presionada o no antes de dibujar el siguiente cuadro. 
En este caso, podemos utilizar el módulo {\ttfamily pygame.key} de forma más directa. Cada tecla tiene una constante asociada, que es un valor que podemos usar para identificar la tecla en el código. Cada constante comienza con {\ttfamily K\_}. Hay letras ( {\ttfamily K\_a} a   {\ttfamily K\_z}), números ({\ttfamily K\_0}  a {\ttfamily K\_9}) y muchos otras constantes como {\ttfamily K\_f1}, {\ttfamily K\_LEFT} y {\ttfamily K\_RETURN}. Puede consultar en la  documentación de Pygame para obtener una lista completa \href{https://www.pygame.org/docs/ref/key.html}{\color{vino} https://www.pygame.org/docs/ref/key.html}.

Debido a que existen constantes para cada tecla, desde {\ttfamily K\_a} a {\ttfamily K\_z}, tal vez Usted puede suponer que también 
hay versiones equivalentes para el caso de  mayúsculas, pero no es así. La razón de esto es que las letras mayúsculas son
el resultado de combinar teclas (Shift + tecla). Si necesita detectar letras mayúsculas u otras teclas desplazadas, debe  utilizar el
parámetro Unicode en eventos  que contienen el resultado de dichas combinaciones de teclas.\\

La función {\ttfamily pygame.key.get\_pressed} devuelve una lista
de booleanos (valores Verdaderos o Falso) si se presiona una tecla.  Para buscar una clave en particular, debe usar su constante como índice en la lista presionada. Por ejemplo, si usáramos la barra espaciadora como botón de disparo,
podría escribir el código de activación de esta manera:
\begin{lstlisting}
pressed_keys = pygame.key.get_pressed()
if pressed_keys[K_SPACE]:
    # Tecla espaciadora ha sido presionada
    fire()
\end{lstlisting}
Es importante saber que algunos teclados tienen varias reglas con respecto a cuántos pulsaciones de teclas simultáneos
admiten. Si observa los teclados para juegos, encontrará que casi siempre enumeran, como punto de venta, cuántas 
pulsaciones simultáneas de teclas admiten, muchos teclados baratos admiten tan solo
de 3 a 5 pulsaciones de teclas simultáneas, mientras que los teclados para juegos pueden soportar hasta más de veinte.
En muchos juegos, es posible que tengas una clave para agacharte, avanzar, ametrallar un poco y moverte lentamente. Tal vez
mientras haces este ataque lento y agachado en diagonal, debes hacer otra acción. 


En el siguiente script se utiliza el función {\ttfamily get\_pressed}  para detectar cualquier tecla presionada y
muestra una lista de ellas en la pantalla. La función genera una lista de las teclas presionadas como con valores booleanos,
mediante un bucle {\ttfamily for} que itera a través de cada una de ellas valor.  

\hspace{-2.4cm}{\ttfamily \scriptsize demo}

\vspace{-0.2cm}
\hspace{-2.3cm}{\ttfamily \scriptsize \_tecla.py}\vspace{-0.84cm}
\begin{lstlisting}
import pygame
from pygame.locals import *
from sys import exit
pygame.init()
screen = pygame.display.set_mode((640, 480), 0, 32)
font = pygame.font.SysFont("arial", 32);
font_height = font.get_linesize()
while True:
    for event in pygame.event.get():
        if event.type == QUIT:
           pygame.quit()
           exit()

    screen.fill((255, 255, 255))
    pressed_key_text = []
    pressed_keys = pygame.key.get_pressed()
    y = font_height
    for key_constant, pressed in enumerate(pressed_keys):
        if pressed:
            key_name = pygame.key.name(key_constant)
            print( key_name)
            text_surface = font.render("pressed" + key_name, True, (0,0,0))
            screen.blit(text_surface, (8, y))
            y+= font_height
        pygame.display.update()
\end{lstlisting}

El bucle itera sobre las teclas presionadas  sobre {\ttfamily enumerate}, que es una función incorporada que devuelve una tupla del índice (el primer valor es 0, el segundo es 1, etc.) y el valor de la lista iterada. Si el valor iterado (presionado) es Verdadero, significa que esa tecla ha sido presionada, luego con un bloque de código se muestra en la pantalla. Este bloque de código
utiliza otra función en el módulo de teclado,  {\ttfamily pygame.key.name}, que toma una clave constante y devuelve una cadena con una descripción de la clave (es decir, se convierte {\ttfamily K\_SPACE} en "{}espacio"{}).\\

Aquí puedes verificar cuántas teclas admite su teclado. Los  escenarios pueden variar desde admitir hasta 8 o más, pero algunas veces
tan solo 4, o algo así.  Para más información  a cerca del funcionamiento de los teclados en  \href{http://www.sjbaker.org/wiki/index.php?title=Keyboards\_Are\_Evil}{\color{vino} http://www.sjbaker.org/wiki/index.php?title=Keyboards\_Are\_Evil}.\\



Ahora describimos el módulo {\ttfamily pygame.key} con más detalle:\\


•  {\ttfamily key.get\_focused}: una ventana de Pygame solo recibe eventos clave cuando la ventana
está enfocado, generalmente haciendo clic en la barra de título de la ventana. Esto es cierto para todos las ventanas de alto nivel.
La función  {\ttfamily get\_focused} devuelve True si la ventana tiene foco y
puede recibir eventos clave; de lo contrario, devuelve False. Cuando Pygame se está ejecutando en
modo de pantalla completa, siempre tendrá foco porque no tiene forma para  compartir la pantalla con otras aplicaciones.

•  {\ttfamily key.get\_pressed}: devuelve una lista de valores booleanos para cada clave. Si alguno de los valores
están configurados en True, se presiona la tecla para ese índice.

•  {\ttfamily key.get\_mods}: devuelve un valor único que indica cuál de las teclas modificadoras
están presionados. Las teclas modificadoras son teclas como Shift, Alt y Ctrl que se utilizan en
combinación con otras teclas. Para verificar si se presiona una tecla modificadora, debe usar el
Operador  {\ttfamily AND bitwise} (\&) con {\ttfamily KMOD\_constante}. 
Por ejemplo, para comprobar si se presiona la tecla Shift izquierda, usaría  {\ttfamily pygame.key.get\_mods()} \&
 {\ttfamily KMOD\_LSHIFT}.


•  {\ttfamily pygame.key.set\_mods}: también puedes configurar una de las teclas modificadoras para imitar el
efecto de presionar una tecla. Para configurar una o más teclas modificadoras, se combina  {\ttfamily KMOD\_constante}
con el operador {\ttfamily bitwise OR (|)}. Por ejemplo, para configurar las teclas Shift y Alt podrías usar 
{\ttfamily pygame.key.set\_mods(KMOD\_SHIFT | KMOD\_ALT)}.


•  {\ttfamily pygame.key.set\_repeat}: si abres tu editor de texto favorito y mantienes presionada una letra, 
verá que después de un breve retraso la tecla comienza a repetirse, lo cual es útil si desea ingresar un carácter 
varias veces sin presionar y soltar. Puedes pedirle a Pygame que te envíe eventos  {\ttfamily KEY\_DOWN} repetidos con 
la función {\ttfamily set\_repeat}, que toma un valor para el retraso inicial antes de que se repita una tecla y un valor para
el retraso entre teclas repetidas. Ambos valores están en milisegundos (1000 milisegundos
en un segundo). Puede desactivar la repetición de teclas llamando a  {\ttfamily set\_repeat} sin parámetros.


•  {\ttfamily pygame.key.name}: esta función toma una  {\ttfamily KEY\_constante} y devuelve una cadena  descriptiva
para ese valor. Es útil para depurar cuando el código se está ejecutando y sólo podemos ver el valor de una  {\ttfamily KEY\_constante}  y no su nombre. Por ejemplo, si obtengo un evento  {\ttfamily KEY\_DOWN} para una tecla con un valor de 103, puedo usar  {\ttfamily key.name} para imprimir el nombre de la clave (que en este caso es “g”).

\section{Movimiento direccional con teclas}

Se puedes mover un objeto en la pantalla con el teclado asignando una tecla para arriba, abajo, izquierda y derecha. 
Cualquier tecla pueden ser usada  para movimiento direccional, pero las teclas más obvias para usar son las teclas de cursor, porque
están diseñadas para movimiento direccional y están colocados en el lugar correcto para operar con una mano. 
Los aficionados a los shooters en primera persona también están acostumbrados a utilizar las teclas W, A, S y D para desplazarse.
Para convertimos las pulsaciones de teclas en movimientos direccionales, se necesita crear un vector de rumbo que apunte en la
dirección que queremos ir. Si sólo una de las cuatro teclas de dirección está presionado, el vector de rumbo es bastante simple. 
En la siguiente tabla \ref{direcction1}se muestran los cuatro vectores de dirección básicos.


\begin{center}{\noindent  Table \ref{direcction1}. Cuatro vectores de dirección básicos: \\%\vspace{0.1cm}\\
\vspace{0.1cm}
{\footnotesize{{%\scriptsize  %\hspace{-0.7cm}	
\begin{tabular}{l l l }
\multicolumn{1}{|>{\columncolor{vino}} l }{\color{white}\bf Dirección }& 
\multicolumn{1}{|>{\columncolor{vino}} c }{\color{white}\bf Vector} \\ 
					Left     & 	-1, 0    \\
					Right    &	+1, 0    \\ 				
					Up       &  0, -1     \\ 
					Down     &  0, 1      \\				
					\arrayrulecolor{vino}\hline 
					\label{direcction1}  
\end{tabular}\vspace{0.2cm} 	}}} 
}\end{center}



Además del movimiento horizontal y vertical, se necesita que el usuario pueda moverse en dirección  diagonal,
presionando dos teclas al mismo tiempo. Por ejemplo, si se presionan las teclas Arriba y Derecha, el sprite debería
moversb en diagonal hacia la parte superior derecha de la pantalla. Podemos crear este vector diagonal sumando dos de los
vectores simples. Si sumamos Arriba (0.0, –1.0) y Derecha (1.0, 0,0), obtenemos (1.0, –1.0), que apunta hacia arriba, pero no 
podemos usar esto como vector de rumbo porque ya no es un vector unitario (longitud 1). Si tuviéramos que
usarlo como vector de rumbo, encontraríamos que nuestro sprite se mueve más rápido en diagonal que en vertical u horizontalmente, 
lo cual no es tan útil.

Antes de usar nuestro rumbo calculado, debemos convertirlo nuevamente en un vector unitario normalizándolo, lo cual
nos da un rumbo de aproximadamente (0,707, –0,707) (Ver Figura \ref{diagonal1} para obtener una representación visual de la diagonal).



\begin{figure}[h!]
	\centering
	\includegraphics[scale=0.35]{../A Logos de Python/diagonal1}
	\caption{Vectores diagonales combinando vectores simples.}
	\label{diagonal1}
\end{figure}

El siguiente código implementa este  movimiento direccional. Al ejecutarlo  ves un objeto que se  mueve horizontal o verticalmente presionando cualquiera de las teclas del cursor, o  en diagonal presionando dos teclas del cursor en combinación.
Si ejecuta el código  sin modificar, encontrará un error generado por intentar dividir algo por cero. Puedes manejarlo dentro 
con un método de normalización:




\hspace{-2.5cm}{\ttfamily \scriptsize movimiento\_}

\vspace{-0.2cm}
\hspace{-2.3cm}{\ttfamily \scriptsize teclas.py}\vspace{-0.84cm}
\begin{lstlisting}	
import pygame
from pygame.locals import *
from sys import exit


fondo_image_filename = 'pygame.jpg'
sprite_image_filename = 'nave.png'
pygame.init()
screen = pygame.display.set_mode((640, 480), 0, 32)
fondo = pygame.image.load(fondo_image_filename).convert()
sprite = pygame.image.load(sprite_image_filename).convert_alpha()
clock = pygame.time.Clock()
sprite_pos = Vector2(200, 150)
sprite_speed = 300
pygame.display.set_caption('Movimiento con teclas')

while True:
    for event in pygame.event.get():
        if event.type == QUIT:
            pygame.quit()
            exit()
            
    pressed_keys = pygame.key.get_pressed()
    key_direction = pygame.Vector2(1, 1)
    if pressed_keys[K_LEFT]:
        key_direction.x = -1
    elif pressed_keys[K_RIGHT]:
        key_direction.x = +1
    if pressed_keys[K_UP]:
       key_direction.y = -1
    elif pressed_keys[K_DOWN]:
       key_direction.y = +1
    
    # Evitra error al dividir por cero 
    if key_direction ==(0,0):
        key_direction 
    elif key_direction != (0,0):
        key_direction = key_direction.normalize()
    
    screen.blit(fondo, (0,0))
    screen.blit(sprite, (sprite_pos.x,sprite_pos.y))
    time_passed = clock.tick(30)
    time_passed_seconds = time_passed / 1000.0
    sprite_pos += key_direction * sprite_speed * time_passed_seconds
    pygame.display.update()
\end{lstlisting}





En el script se evita dividir por cero al normalizar, lo tratamos con el condicional  {\ttfamily if ... ilif}, poco de habilidad al calcular el vector de dirección. Establece el componente x en –1 o +1 si se presiona  {\ttfamily K\_LEFT} o  {\ttfamily K\_RIGHT} y la componente y a –1 o +1 si se presiona {\ttfamily K\_UP} o  {\ttfamily K\_DOWN}. Esto nos da el mismo resultado que sumar dos de los vectores con dirección  horizontal y vertical simples. Si alguna vez observa  un atajo matemático que te permite hacer algo con menos código, no dudes en probarlo. 

Es posible que hayas notado que solo se utilizan ocho vectores para este movimiento vectorial. Si tuviéramos que
precalcular estos vectores e insertarlos directamente en el código, podríamos reducir la cantidad de trabajo
hecho al ejecutar el script.

Vale la pena hacerlo como ejercicio si te gustan los desafíos, pero dado que calcular
el vector de dirección solo se realiza una vez por cuadro, acelerarlo no hará ninguna diferencia notable en la velocidad 
de fotogramas. Esté atento a situaciones como esta; reduciendo la cantidad de trabajo que el juego debe hacer para
crear un marco que se llama optimización y se vuelve más importante cuando hay mucha acción en marcha.

\subsection{Movimiento rotacional con teclas}

Hasta ahora tenemos ocho direcciones para  el movimiento, sin embargo es un poco artificial porque no se ven muchas cosas moviéndose
así en la vida real. La mayoría de los objetos móviles pueden girar libremente, pero moverse en la dirección a la que apuntan,
o quizás hacia atrás, pero definitivamente en más de ocho direcciones de la brújula. Todavía podemos usar las mismas teclas Arriba, Abajo, Izquierda y Derecha para simule esto, pero tenemos que cambiar lo que controlan las teclas. 
Lo que queremos es tener teclas para izquierda y derecha que controlan la rotación y teclas para el movimiento hacia adelante
y hacia atrás, lo que le daría a nuestro sprite la capacidad de moverse en cualquier dirección. 

En el siguiente script se usa exactamente el mismo conjunto de teclas para el movimiento de un sprite alrededor de la pantalla, pero con movimiento de rotación libre. 


 
\hspace{-2.5cm}{\ttfamily \scriptsize movimiento}

\vspace{-0.2cm}
\hspace{-2.5cm}{\ttfamily \scriptsize \_rotacion.py}\vspace{-0.84cm}
\begin{lstlisting}
import pygame
from pygame.locals import *
from sys import exit
from math import *

fondo_image_filename = 'pygame.jpg'
sprite_image_filename = 'nave.png'
pygame.init()
screen = pygame.display.set_mode((640, 480), 0, 32)
fondo = pygame.image.load(fondo_image_filename).convert()
sprite = pygame.image.load(sprite_image_filename).convert_alpha()
clock = pygame.time.Clock()
sprite_pos = pygame.Vector2(200, 150)
sprite_speed =  300
sprite_rotation = 0
sprite_rotation_speed = 360 # Degrees per second
pygame.display.set_caption('Rotación con teclas')
while True:
    for event in pygame.event.get():
        if event.type == QUIT:
            pygame.quit()
            exit()
    pressed_keys = pygame.key.get_pressed()
    rotation_direction = 0.
    movement_direction = 0.
    if pressed_keys[K_LEFT]:
        rotation_direction = +1.0
    if pressed_keys[K_RIGHT]:
        rotation_direction = -1.0
    if pressed_keys[K_UP]:
        movement_direction = -1.0
    if pressed_keys[K_DOWN]:
        movement_direction = +1.0
   
    screen.blit(fondo, (0,0))
    rotated_sprite = pygame.transform.rotate(sprite, sprite_rotation)
    w, h = rotated_sprite.get_size()
    sprite_draw_pos = pygame.Vector2(sprite_pos.x-w/2, sprite_pos.y-h/2)
    screen.blit(rotated_sprite, (sprite_draw_pos.x,sprite_draw_pos.y))
    time_passed = clock.tick()
    time_passed_seconds = time_passed / 1000.0
    sprite_rotation += rotation_direction * sprite_rotation_speed * time_passed_seconds
    # Calculos
    heading_x = sin(sprite_rotation*pi/180.0)
    heading_y = cos(sprite_rotation*pi/180.0)
    heading = pygame.Vector2(heading_x, heading_y)
    heading *= movement_direction
    sprite_pos+= heading * sprite_speed * time_passed_seconds
    pygame.display.update()
\end{lstlisting}	

El código funciona de manera similar al ejemplo anterior, pero calcula el encabezado de manera diferente.
En lugar de crear un vector de dirección  a partir de la tecla que se presione, creamos la dirección  partir de la 
rotación del sprite (almacenado en la variable  {\ttfamily sprite\_rotation}. Es este valor el que modificamos al presionar la tecla;
cuando se presiona la tecla Derecha, agregamos rotación al sprite, girándolo en esa dirección. Y cuando  se pulsa la tecla izquierda, restamos  rotación al sprite para girarla en sentido contrario.

Para calcular el vector de dirección a partir de la rotación, debemos calcular el seno y el coseno del ángulo,
lo cual podemos hacer con las funciones {\ttfamily sin()}  y {\ttfamily cos()} que se encuentran en el módulo de matemáticas. 
La función {\ttfamily sin()} calcula la componente x, y la función {\ttfamily cos()} calcula el componente y. 
Ambas funciones toman un ángulo en radianes (hay $2 * \pi$  radianes en un círculo, 360 grados), debido a que 
en el script  {\ttfamily movimiento\_rotacion.py} se  usan grados,  deben convertirse la rotación a radianes 
multiplicando por $\pi$ y dividiendo por 180. Después de ensamblar el vector de dirección de los dos componentes, para
usarse como antes para dar el movimiento basado en el tiempo.



Antes de que el código muestre el sprite se gira visualmente para que podamos ver en qué dirección está orientado.
Esto se hace con la función del módulo {\ttfamily pygame.transform} llamada {\ttfamily rotate} ({\ttfamily pygame.transform.rotate(sprite, sprite\_rotation)}), que toma una superficie más un ángulo y devuelve la nueva superficie que contiene el objeto girado.

Un problema con esta función es que el sprite devuelto puede no tener las mismas dimensiones que el original (consulte la Figura \ref{rotando2}, por lo que no podemos convertirlo en el original misma coordenada que el objeto sin rotar o se dibujará en la posición incorrecta en la pantalla; una manera de trabajar  esto es dibujar el sprite en una posición que coloque el centro de la imagen del sprite debajo del Posición del sprite en la pantalla. De esa manera, el sprite estará en la misma posición en la pantalla independientemente del tamaño de la superficie rotada.

\begin{figure}[h!]
	\centering
	\includegraphics[scale=0.35]{../A Logos de Python/rotando03}
	\caption{Al girar una superficie cambia su tamaño.}
	\label{rotando2}
\end{figure}


\section{Implementación del control con el  mouse}

El mouse existe desde hace casi tanto tiempo como el teclado. El diseño del mouse no ha cambiado mucho por días, los dispositivos se volvieron un poco más ergonómicos (con forma de mano), pero el diseño siguió siendo el mismo. Los ratones clásicos tienen una bola recubierta de goma debajo que rueda sobre el escritorio o la alfombrilla del ratón. El movimiento de la pelota es captado por dos rodillos dentro del mouse que están en contacto con la pelota.

Hoy en día, casi todos los mouse son láser, lo que ofrece mayor precisión. Atrás quedaron los días en los
amigos podían quitarnos la bola del mouse como una broma, pero ahora podemos pegar papel sobre el orificio del láser para
hacer la broma divertida.

A medida que pasa el tiempo, cada vez más mouse tienen también varios botones. También hay muchos mouse con muchos botones diferentes.
Los juegos suelen utilizar estas innovaciones de mouses. Los botones adicionales siempre son útiles para un acceso rápido
a los controles del juego, y la rueda del mouse podría haberse diseñado para cambiar de arma en primera persona
¡tirador! Y, por supuesto, la precisión adicional de los mouses láser siempre resulta útil a la hora de eliminar enemigos con un rifle de francotirador.



\subsection{Movimiento rotacional con el mouse}

Ya has visto que dibujar el cursor del ratón en la pantalla es bastante sencillo, sólo necesitas
obtener las coordenadas del mouse de un evento {\ttfamily MOUSEMOTION} o directamente con la función  {\ttfamily pygame.mouse.get\_pos}. Cualquiera de los métodos está bien si sólo desea mostrar el cursor del mouse, pero el movimiento del mouse también puede
usarse para controlar algo  que no sea una posición absoluta, como girar o mirar hacia arriba y hacia abajo
en un juego 3D. En este caso, no podemos usar la posición del mouse directamente porque las coordenadas serían
restringido a los bordes de la pantalla, y no queremos que el jugador esté restringido en cuántas veces
gira a la izquierda o a la derecha. En estas situaciones, necesitamos  obtener el movimiento relativo del mouse, a menudo llamado
{\ttfamily mouse mickeys}, que simplemente significa qué tan lejos se ha movido el mouse desde el cuadro anterior.


En el script {\ttfamily rotacion\_mouse.py} se  agrega movimiento de rotación a mouse al movimiento del sprite.
Además de las teclas del cursor, el objeto rotará cuando el ratón se mueve hacia la izquierda o hacia la derecha.




\hspace{-2.9cm}{\ttfamily \scriptsize rotacion}

\vspace{-0.2cm}
\hspace{-2.5cm}{\ttfamily \scriptsize \_mouse.py}\vspace{-0.84cm}
\begin{lstlisting}		
import pygame
from pygame.locals import *
from sys import exit
from math import *

fondo_image_filename = 'pygame.jpg'
sprite_image_filename = 'nave.png'
pygame.init()
screen = pygame.display.set_mode((640, 480), 0, 32)
fondo = pygame.image.load(fondo_image_filename).convert()
pygame.display.set_caption('Movimiento y rotación con mouse')
sprite = pygame.image.load(sprite_image_filename).convert_alpha()
clock = pygame.time.Clock()

# Desactivar el cursor del mouse
pygame.mouse.set_visible(False)
# PyGame tiene el control total el mouse
pygame.event.set_grab(True)
sprite_pos = pygame.Vector2(200, 150)

sprite_speed = 300.
sprite_rotation = 0.
sprite_rotation_speed = 360. # Degrees per second

while True:
    for evento in pygame.event.get():       
        if evento.type == KEYDOWN: 
            if evento.key==K_q:
                pygame.quit()
                exit()  
            
    pressed_keys = pygame.key.get_pressed()
    pressed_mouse = pygame.mouse.get_pressed()
    rotation_direction = 0.
    movement_direction = 0.
    rotation_direction = pygame.mouse.get_rel()[0] / 3.
    if pressed_keys[K_LEFT]:
        rotation_direction = +1.
    if pressed_keys[K_RIGHT]:
        rotation_direction = -1.
    if pressed_keys[K_UP] or pressed_mouse[0]:
        movement_direction = +1.
    if pressed_keys[K_DOWN] or pressed_mouse[2]:
        movement_direction = -1.
    screen.blit(fondo, (0,0))
    rotated_sprite = pygame.transform.rotate(sprite, sprite_rotation)
    w, h = rotated_sprite.get_size()
    sprite_draw_pos = pygame.Vector2(sprite_pos.x-w/2, sprite_pos.y-h/2)
    screen.blit(rotated_sprite, (sprite_draw_pos.x, sprite_draw_pos.y))
    time_passed = clock.tick()
    time_passed_seconds = time_passed / 1000.0
    sprite_rotation += rotation_direction * sprite_rotation_speed * time_passed_seconds
    heading_x = sin(sprite_rotation*pi/180.)
    heading_y = cos(sprite_rotation*pi/180.)
    heading = pygame.Vector2(heading_x, heading_y)
    heading *= movement_direction
    sprite_pos+= heading * sprite_speed * time_passed_seconds
    pygame.display.update()
\end{lstlisting}

Observe que para salir debe usar la tecla {\ttfamily q}, debido a que el mouse se está usando para el movimiento del sprite; y el área de acción del  mouse se restringe al área física de la ventana (pantalla). Esto se hace con las siguientes dos líneas:
\begin{lstlisting}		
pygame.mouse.set_visible(False)
pygame.event.set_grab(True)
\end{lstlisting}

La llamada a {\ttfamily set\_visible(False)} desactiva el cursor del mouse y la llamada a {\ttfamily set\_grab(True)} le da todo el
control sobre mouse a PyGame. Un efecto secundario de esto es que no puedes usar el mouse para
cerrar la ventana, por lo que debe proporcionar una forma alternativa de cerrar el script.

Usar el mouse de esta manera también dificulta el uso de otras aplicaciones, por lo que es normalmente se utiliza mejor en modo de pantalla completa. 

La siguiente línea obtiene los {\ttfamily mickeys} del mouse para el eje x (en el índice [0]). 
Los divide por 5, porque descubrí que usar el valor directamente rotaba el sprite demasiado rápido. Puede ajustar este valor a
cambie la sensibilidad del mouse (cuanto mayor sea el número, menos sensible será el control del mouse). 
Juegos a menudo permiten que el jugador ajuste la sensibilidad del mouse en un menú de preferencias.

\begin{lstlisting}		
rotation_direction = pygame.mouse.get_rel()[0] / 5.	
\end{lstlisting}	

Además de rotar con el mouse, los botones se usan para mover el sprite hacia adelante (izquierdo del mouse).
botón) y hacia atrás (botón derecho del mouse). La detección de botones presionados del mouse es similar a las teclas. 
La función {\ttfamily pygame.mouse.get\_pressed()} devuelve una tupla de tres valores booleanos para el botón izquierdo del ratón, el
botón central del ratón y el botón derecho del ratón. Si alguno es Verdadero, entonces el botón correspondiente del mouse
se mantuvo presionado en el momento de la llamada.

Una forma alternativa de utilizar esta función es la siguiente, que descomprime la tupla en tres valores:
\begin{lstlisting}		
lmb, mmb, rmb = pygame_mouse.get_pressed()	
\end{lstlisting}


Veamos {\ttfamily pygame.mouse} con más detalle. Es otro módulo sencillo, y tiene sólo ocho funciones:\\

• {\ttfamily pygame.mouse.get\_pressed}: devuelve los botones presionados del mouse como una tupla de tres.
booleanos, uno para los botones izquierdo, central y derecho del mouse.

• {\ttfamily pygame.mouse.get\_rel}: devuelve el movimiento relativo del mouse (o mickeys) como un
tupla con el movimiento relativo x e y.

• {\ttfamily pygame.mouse.get\_pos}: devuelve las coordenadas del mouse  como una tupla de x e y
valores.

• {\ttfamily pygame.mouse.set\_visible}: cambia la visibilidad del cursor del mouse estándar. Si es falso, el cursor será invisible.


• {\ttfamily pygame.mouse.get\_focused}: devuelve True si la ventana de Pygame está recibiendo entrada del mouse. 
Cuando Pygame se ejecuta en una ventana, solo recibirá la entrada del mouse cuando la ventana está seleccionada y 
en el frente de la pantalla.


• {\ttfamily pygame.mouse.set\_cursor}: configura la imagen estándar del cursor. Esto rara vez es necesario
porque se pueden lograr mejores resultados modificando una imagen a la coordenada del mouse.

• {\ttfamily pyGame.mouse.get\_cursor}: obtiene la imagen  estándar del cursor. Ver el punto anterior.



\subsubsection{Jugabilidad del ratón}

El humilde mouse puede permitirte ser bastante creativo en un juego. Incluso un simple cursor puede ser una gran
herramienta de juego para juegos de rompecabezas y estrategia, donde el mouse se usa a menudo para seleccionar y
soltar elementos del juego en el campo de juego. Estos elementos podrían ser cualquier cosa, 
desde espejos que reflejan láser hasta ogros sedientos de sangre, dependiendo del juego. No es necesario que el puntero 
del mouse sea una flecha. Para un juego de DIOS, sería  probablemente  la mano omnipotente de una deidad o quizás 
una misteriosa sonda alienígena. El clásico point 'n' click, los juegos de aventuras también hicieron un excelente
uso del mouse. En estos juegos, el personaje del juego caminaba a lo que sea en lo que el jugador hizo clic y échale un vistazo. 
En estos juegos, el personaje del juego caminaba a lo que el jugador hice clic y écha un vistazo. Si fuera un objeto útil, 
podría almacenarse para su uso más tarde  en el juego.


Para juegos con un control más directo del personaje del jugador, el mouse se usa para rotar o moverlo. 
Un simulador de vuelo puede  usar los ejes x e y del mouse para ajustar el cabeceo y la guiñada del avión de modo que el jugador
puede hacer ajustes finos a la trayectoria. 


Una forma sencilla de hacer esto en Pygame sería agregar al mouse el ángulo del avión. También me gusta usar el mouse para
tiradores (shooters) en primera persona y usar el eje x para girar hacia la izquierda y hacia la derecha, y el eje y para mirar 
hacia arriba y hacia abajo. Esto parece bastante natural porque las cabezas tienden moverse de esa manera; puedo girar la 
cabeza hacia cualquier lado y mirar hacia arriba y hacia abajo. También puedo inclinar la cabeza para de cualquier lado, 
pero no lo hago muy a menudo. 

También puede combinar controles de mouse y teclado. Me gustan especialmente los juegos que usan el teclado para
movimiento y el ratón para apuntar. Un tanque, por ejemplo, podría moverse con las teclas del cursor, pero usa el mouse
para girar la torreta\footnote{Una torreta es usualmente una plataforma rotativa que puede ser montada en un edificio fortificado o estructura tales como baterías de costa, sobre un tanque, buque de guerra, o avión.} y poder dispararle a un enemigo sin mirar en la misma dirección. Por supuesto, no estás limitado a utilizar el mouse del mismo modo que otros juegos. Sea tan creativo como usted
deseo, ¡solo asegúrate de probarlo primero con algunos de tus amigos!





\section{Implementación del control por palanca de mando}
Pendiente...\\

Los joysticks no se han visto limitados por la necesidad de utilizarlos en otros lugares además de los juegos, y han sido completamente
Libertad para innovar, por lo que los joysticks modernos tienen un diseño moldeado suave y aprovechan todos los dígitos sobrantes.
tienen a su disposición. Aunque la palabra joystick se utiliza a menudo para referirse a un controlador de juego de cualquier tipo,
describe más técnicamente una palanca de vuelo utilizada en simulaciones de vuelo.
Hoy en día, el joypad es el más popular entre los jugadores y, a menudo, viene con las consolas. Encajan
cómodamente en dos manos y tienen muchos botones, además de un pad direccional y dos operados con el pulgar
palos analógicos. Piensa en tus controladores de Xbox o PlayStation. Mucha gente conecta su controlador Xbox para
su computadora.

El módulo de joystick de Pygame admite todo tipo de controladores con la misma interfaz y no
distinguir entre los distintos dispositivos disponibles. Esto es bueno porque no queremos tener que hacerlo.
¡Escribe código para cada controlador de juego que existe!


Conceptos básicos del mando


Comencemos mirando el módulo pygame.joystick.




\section{Resumen}

Hay tres módulos de controles en PyGame:  {\ttfamily pygame.keyboard},  {\ttfamily pygame.mouse} y  {\ttfamily pygame.joystick}, 
entre ellos, puedes admitir cualquier dispositivo que el jugador quiera usar en un juego.  Puede  darle al jugador la opción para seleccionar el método de control.

A veces, los controles influyen instantáneamente en el personaje del jugador, de modo que obedecerá exactamente lo que el
el jugador le dice que haga. Los juegos de disparos clásicos eran así; presionaste hacia la izquierda, ella nave instantáneamente 
giró hacia la izquierda y viceversa. Sin embargo, para un poco más de realismo, los controles deberían afectar al personaje 
del jugador indirectamente al aplicando fuerzas como empuje y rotura. 

Ahora que sabes cómo leer información de los controladores de juegos, los siguientes capítulos le enseñarán cómo utilizar 
esta información para manipular objetos de forma más realista los mundos de juego.

Los controles de un juego son la forma en que el jugador interactúa con el mundo del juego, porque no podemos
ingresar al juego, estilo Matrix, los controles deben sentirse lo más naturales posible. 
Al considerar el control como métodos para un juego, es una buena idea observar cómo se juegan juegos similares.
Los controles no varían mucho para juegos de un género similar. Esto no se debe a que a los diseñadores de juegos 
les falte imaginación; es solo que los jugadores esperan que los juegos funcionen de manera similar.
Si utilizas un control más creativo métodos, ¡prepárate para justificárselo a los jugadores u ofrecerles una alternativa más estándar!

En el próximo capítulo exploraremos el tema de la inteligencia artificial y aprenderá cómo agregar
personajes no jugadores a un juego.




%%%%%%%%%%%%%%%%%%%%%%%%%%%%%%%%%%%%%%%%%%%%%%%%%%%%%%%%%%%%%%%%%%%%%%%%%%%%%%%%%%%%%


%%%%%%%%%%%%%%%%%%%%%%%%%%%%%%%%%%%%%%%%%%%%%%%%%%%%%%%%%%%%%%%%%%%%%%%%%%%%%%%%%%%%%












%%%%%%%%%%%%%%%%%%%%%%%%%%%%%%%%%%%%%%%%%%%%%%%%%%%%%%%%%%%%%%%%%%%%%%%%%%%%%%%%%%%%%

%%%%%%%%%%%%%%%%%%%%%%%%%%%%%%%%%%%%%%%%%%%%%%%%%%%%%%%%%%%%%%%%%%%%%%%%%%%%%%%%%%%%%


%%%%%%%%%%%%%%%%%%%%%%%%%%%%%%%%%%%%%%%%%%%%%%%%%%%%%%%%%%%%%%%%%%%%%%%%%%%%%%%%%%%%%


%%%%%%%%%%%%%%%%%%%%%%%%%%%%%%%%%%%%%%%%%%%%%%%%%%%%%%%%%%%%%%%%%%%%%%%%%%%%%%%%%%%%%
% https://www.youtube.com/watch?v=9AlK77POFaU   futbol jugadores
% https://www.youtube.com/watch?v=xEDNEQBIj-w

% https://www.google.com/search?q=dibujo+de+un+jugador+de+futbol&rlz=1C1VDKB_esVE1058VE1058&oq=dibujo+de+un+jugador+de+futbol+&gs_lcrp=EgZjaHJvbWUyBggAEEUYOdIBCjE1OTk4ajBqMTWoAgiwAgE&sourceid=chrome&ie=UTF-8#vhid=CcO7J6Z-p1eyzM&vssid=l&ip=1


%http://pygame.org/docs.
% brief tutorial on all of Pygame’s     
% Pygame provides functions for creating programs with a graphical user interface, graphical user interface, or GUI 
% IDLE stands for Interactive DeveLopment Environment.

%   al@inventwithpython.com

Dos proyectos de juegos los alienigenas y fútbol... 


\chapter*{SEGUNDA PARTE: Mis Proyectos}%{\label{Cap1}}\hypertarget{Cap1}{}  %  \hypertarget{Introd}{}\label{Introd}
\addcontentsline{toc}{section}{SEGUNDA PARTE: Mis Proyectos}

Ahora conocemos  lo suficiente sobre Python y PyGame para comenzar a crear proyectos interactivos y significativos. 
Al crear sus propios proyectos le enseñará nuevas habilidades y solidificará su comprensión de los conceptos 
introducidos en la primera parte.

En la segunda parte presentamos tres  proyectos, que usted puede elegir cualquiera o todos estos proyectos en el orden que desee. 

A continuación se incluye una breve descripción de cada proyecto para ayudarle a decidir cuál iniciara primero.

\section*{Proyecto 1: Un juego con Python y PyGame}
% https://depositphotos.com/es/vector/boy-cartoon-playing-football-151737316.html  futbol 
\noindent{\bf Invasión alienígena: }\\
En el proyecto Invasión alienígena (capítulos 12, 13 y 14), utilizarás el paquete PyGame para desarrollar un juego en 2D. 
El objetivo del juego es derribar una flota de alienígenas a medida que descienden por la pantalla, en niveles que aumentan 
en velocidad y dificultad. Al final del proyecto, habrás aprendido habilidades que te permitirán desarrollar tus propios 
juegos 2D en Pygame.\\

Usaremos Pygame, una colección de divertidos y poderosos Módulos de Python que gestionan gráficos,
animación e incluso sonido, lo que facilita para que puedas crear juegos sofisticados. 
Con Pygame al manejar tareas como dibujar imágenes en la pantalla, usted Puede centrarse en la lógica de nivel
superior de la dinámica del juego.

Mientras creas este juego, también aprenderás a gestionar grandes proyectos que abarcan varios archivos.
Refactorizaremos una gran cantidad de código y administraremos archivos contenidos para organizar el proyecto y hacer 
que el código sea eficiente.


\subsection*{Proyecto 2: Visualización de datos}

Los proyectos de visualización de datos comienzan en el Capítulo \ref{15}, donde aprenderá a generar datos y a crear una serie 
de visualizaciones  y funcionales de esos datos usando la librerías {\ttfamily Matplotlib} y  {\ttfamily Plotly}. 

El Capítulo \ref{16} le enseña a acceder a datos de fuentes en línea e introducirlos en un paquete de visualización para 
crear gráficos de datos meteorológicos y un mapa de la actividad sísmica global. 

Finalmente, el Capítulo 17 le muestra cómo escribir un programa para descargar y visualizar datos automáticamente. 
Aprender a hacer visualizaciones le permite explorar el campo de la ciencia de datos, que es una de las áreas de programación de mayor demanda en la actualidad.

\section*{Proyecto 3: Aplicaciones web}


En el proyecto de aplicación web (capítulos 18, 19 y 20), utilizará el paquete Django para crear una aplicación web sencilla que permita a los usuarios llevar un diario sobre diferentes temas que han estado aprendiendo. Los usuarios crearán una cuenta con un nombre de usuario y contraseña, ingresarán un tema y luego harán entradas sobre lo que están aprendiendo. 
También implementará su aplicación en un servidor remoto para que cualquier persona en el mundo pueda acceder a ella.
Después de completar este proyecto, podrá comenzar a crear sus propias aplicaciones web simples y estará listo para profundizar en recursos más completos sobre la creación de aplicaciones con Django.






%\section*{Planificando tu proyecto}

\vspace{0.2cm}
\begin{minipage}[t]{0.2\textwidth}
	\centering\raisebox{\dimexpr\topskip-\height}{%
		\includegraphics[scale=0.20]{../A Logos de Python/Sugerencia}}  %[width=\textwidth]
	
\end{minipage}%\hfill
\begin{minipage}[t]{0.70\textwidth} 
	{\scriptsize Cuando se va a construir un proyecto grande o pequeño, es recomendable e importante hacer un plan antes de
		comenzar a escribir el código, en el plan debe describir la idea en forma clara y las etapas posibles a seguir. 
		Un plan te mantendrá enfocado en el objetivo planteado y le permitirá ir cumpliendo metas para completar el proyecto. 
		Es importante hacer una descripción general del juego; aunque al principio tal vez la descripción no cubra todos los
		detalles que el juego requiere, pero proporciona una idea clara para empezar a construir el juego.}
\end{minipage}
\vspace{0.2cm}


En el próximo capítulo \ref{JuegoAlien} se  desarrolla el juego Invasión Alienígena, un juego donde el jugador controla una
nave que inicialmente aparece en la parte inferior central de la pantalla. Para la primera fase de desarrollo, construimos una 
nave que pueda moverse correctamente hacia la izquierda y derecha cuando el jugador presiona las teclas de flecha, y que dispare 
balas cuando el jugador presiona la barra espaciadora.  Después de configurar este comportamiento, podemos crear los alienígenas o extraterrestres y perfeccionar el juego.



%%%%%%%%%%%%%%%%%%%%%%%%%%%%%%%%%%%%%%%%%%%%%%%%%%%%%%%%%%%%%%%%%%%%%%%%%%%%%%%%%%%%%%%%%%%%%%%%%%%%%


%%%%%%%%%%%%%%%%%%%%%%%%%%%%%%%%%%   Invasión Alienígena   %%%%%%%%%%%%%%%%%%%%%%%%%%%%%%%%%%%%%%%%%%%

%%%%%%%%%%%%%%%%%%%%%%%%%%%%%%%%%%%%                       %%%%%%%%%%%%%%%%%%%%%%%%%%%%%%%%%%%%%%%%%%%


%%%%%%%%%%%%%%%%%%%%%%%%%%%%%%%%%%%%%%%%%%%%%%%%%%%%%%%%%%%%%%%%%%%%%%%%%%%%%%%%%%%%%%%%%%%%%%%%%%%%%%%




\chapter{Invasión Alienígena}{\label{JuegoAlien}}\hypertarget{Cap1}{}  %  \hypertarget{Introd}{}\label{Introd}

%  https://www.pngwing.com/es/free-png-bemav

\begin{wrapfigure}{l}{0.4\textwidth} 
	\vspace{-25pt}\begin{center} \includegraphics[width=0.41\textwidth]{../A Logos de Python/Invasion03}
	\end{center}\vspace{-30.5pt} %\caption{Python 3} 
\end{wrapfigure}

En este capítulo, se configura su ambiente de trabajo, y se inicia escribiendo el código para crear una nave espacial
que se mueve hacia la derecha y hacia la izquierda,  y que dispara balas en respuesta a la entrada del jugador. 
Luego, en los próximo dos capítulos, se creará una flota de alienígenas que debe destruir el jugador;  luego se continua 
refinando el juego, estableciendo límites en la cantidad de naves que puedes usar y se  añade un marcador.

\vspace{0.2cm}
\begin{minipage}[t]{0.2\textwidth}
	\centering\raisebox{\dimexpr\topskip-\height}{%
		\includegraphics[scale=0.20]{../A Logos de Python/Sugerencia}}  %[width=\textwidth]
	
\end{minipage}%\hfill
\begin{minipage}[t]{0.70\textwidth} 
	{\scriptsize Alienígena Invasión abarca varios archivos diferentes, así que es recomendable crear una carpeta digamos {\ttfamily invasion\_alien} en su sistema. Asegúrese de guardar todos los archivos del proyecto en esta carpeta para que las declaraciones de importación funcionen  correctamente. Además, si se siente cómodo usando el control de versiones, es posible que desee usarlo para este proyecto. Si no ha utilizado el control de versiones antes, consulte el Apéndice D para obtener más información.}
\end{minipage}
\vspace{0.2cm}

InvasionAlien



\section{Una Clase base para el juego}    %  (800,600)
%\addcontentsline{toc}{section}{Crear una ventana }

Iniciamos con la creación de una clase vacía para representar el juego  {\ttfamily class InvasionAlien}. Para esto debe crear en  su editor de texto un script y guárdelo en la carpeta principal para el juego, digamos {\ttfamily invasión\_alien.py}. En este archivo escriba el siguiente código:

%Invasión Alienígena
\hspace{-2.9cm}{\ttfamily \scriptsize invasión}

\vspace{-0.2cm}
\hspace{-2.5cm}{\ttfamily \scriptsize \_alien.py}\vspace{-0.84cm}
\begin{lstlisting}
import sys
import pygame

class InvasionAlien:
    """Clase general para gestionar comportamiento del juego."""
    def __init__(self):
        """Inicializa el juego y crea recursos para el juego."""
        pygame.init()
	
        self.pantalla = pygame.display.set_mode((640, 480))
        pygame.display.set_caption("Invasión Alienígena")

    def run_jugar(self):
        """Inicia el bucle principal del juego."""
        while True:
            # Recuperar los Eventos del teclado y el mouse.
            for event in pygame.event.get():
                if event.type == pygame.QUIT:
                    pygame.quit()                    
                    sys.exit()
            # Hacer visible la pantalla dibujada más recientemente.
            pygame.display.flip()
            
if __name__ == '__main__':
    #  Crea una instancia de juego y ejecuta el juego.
    ia = InvasionAlien()
    ia.run_jugar()
\end{lstlisting}

Observe, primero importamos los módulos {\ttfamily sys} y {\ttfamily pygame}. El módulo  {\ttfamily pygame} que contiene la funcionalidad que necesitamos para hacer un juego.  Y Usamos la herramienta del módulo sistema para salir del juego cuando el jugador lo desee, haciendo clic en el botón de cerrar la ventana del juego. PyGame cierra y desaparece la ventana. \\

Invasión Alienígena comienza como una clase llamada {\ttfamily InvasionAlien}. En el método {\ttfamily\_\_init\_\_()}, se incluye 
la función {\ttfamily  pygame.init()}, que inicializa la configuración de fondo que PyGame necesita funcionar correctamente.
Luego llamamos a  {\ttfamily pygame.display.set\_mode()} para crear una ventana de visualización,    
en la que dibujaremos todos los  elementos gráficos del juego. El argumento (640, 480) es una tupla que define las 
dimensiones  de la ventana del juego, en este caso el tamaño es (640, 480), que tendrá 640 píxeles de ancho por 480 píxeles de alto. Y le damos nombre a la ventana PyGame con el método  {\ttfamily pygame.display.set\_caption("{}Invasión Alienígena"{})} \\
 

Asignamos esta ventana de visualización al atributo {\ttfamily self.pantalla}, por lo que estará disponible en todos los métodos
de la clase. El objeto que asignamos a {\ttfamily self.pantalla} se llama superficie. 

La superficie devuelta por {\ttfamily display.set\_mode()} representa toda la ventana del juego. Cuando activamos el bucle principal ({\ttfamily while}) de animación del juego, esta superficie se volverá a dibujar en cada paso por el bucle, por lo que se puede actualizar con cualquier cambio provocado por la entrada del usuario.

El juego está controlado por el método {\ttfamily run\_jugar()}. Este método contiene un bucle {\ttfamily while}  que se ejecuta continuamente. El bucle {\ttfamily while}   contiene un bucle {\ttfamily for }  (anidado) para detectar y responder a eventos específicos, que por ahora solo interesa el evento {\ttfamily pygame.QUIT}. 
Luego el código {\ttfamily pygame.display.flip()} gestiona las actualizaciones de la pantalla.

 
La llamada a {\ttfamily pygame.display.flip()} le dice a {\ttfamily pygame} que haga visible la pantalla dibujada más recientemente.
En este caso,  simplemente dibuja una pantalla vacía en cada paso por el bucle {\ttfamily while}, borrando la pantalla anterior 
para que solo sea visible la nueva. Cuando movemos los elementos del juego, {\ttfamily pygame.display.flip()} actualiza 
continuamente la pantalla para mostrar las nuevas posiciones de los elementos del juego y ocultar las antiguas, creando la
ilusión de un movimiento suave.

Al final del archivo, creamos una instancia del juego y luego llamamos a {\ttfamily run\_jugar()}. Colocamos {\ttfamily run\_jugar()} en un bloque  {\ttfamily if} que solo se ejecuta si el archivo se llama directamente. 
Al ejecutar el archivo {\ttfamily invasion\_alien.py}, debería ver una ventana de {\ttfamily pygame} vacía.\\
  

Es importante tener en cuenta que si utiliza una pantalla con doble {\ttfamily buffer}, debe llamar a
\begin{lstlisting}		
pygame.display.flip()
\end{lstlisting}
en lugar de {\ttfamily pygame.display.update()}. Esto hace que la visualización instantánea cambie en lugar de copiar los datos de la pantalla. \\


En la tabla \ref{tabla_Flags} se presenta un resumen de la información con las diferencias entre {\ttfamily pygame.display.update()} y 	{\ttfamily pygame.display.flip()}:\\

\noindent Table \ref{tabla_Flags}. Diferencias entre los métodos {\ttfamily .update()} y {}\ttfamily .flip()}\vspace{-0.25cm}
\begin{center}
	{\footnotesize{{%\scriptsize 
				\hspace{-0.7cm}	
				\begin{tabular}{l l l l l}	% white
					\multicolumn{1}{|>{\columncolor{azul2}} l }{\color{black}\bf Método}& 
					\multicolumn{1}{|>{\columncolor{azul2}} l }{\color{black}\bf  {\ttfamily pygame.display.update()}  } &	
				     \multicolumn{1}{|>{\columncolor{azul2}} l }{\color{black}\bf {\ttfamily pygame.display.flip()}  } \\  

Funcionalidad &Actualiza solo las partes de la& Actualiza toda la pantalla, \\
              &pantalla que han cambiado      &  incluyendo áreas que no  \\
              &                               &  han cambiado \\

Rendimiento   &Puede ser más eficiente para   & Puede ser más eficiente para  \\
              &aplicaciones que requieren     & aplicaciones que requieren una \\
              &una actualización rápida y     &actualización completa y  \\
              & precisa de solo ciertas áreas  &sincronizada de toda  \\                               
              & de la pantalla                 & la pantalla       \\  

Uso           &	Se utiliza cuando se necesita   & 	Se utiliza cuando se necesita   \\
              & actualizar solo ciertas partes &  actualizar toda la pantalla,   \\
              &de la pantalla, como en juegos  &  como en aplicaciones que \\   
              &que requieren una actualización  &requieren una actualización \\
              &rápida de solo ciertas áreas	     &  completa y sincronizada de  \\            
               &                                 &toda la pantalla   \\
Efectos visuales & 	Puede causar artefactos     & 	No causa artefactos visuales,  \\
                 & visuales  si no se utiliza  &   ya que actualiza toda la pantalla    \\
                 &correctamente,como por ejemplo,&de manera sincronizada\\
                 & si se actualizan áreas de la&                \\
                 &pantalla que no han cambiado &                 \\
\arrayrulecolor{vino}\hline    
\end{tabular}\label{tabla_Flags} 		
%  	\vspace{0.2cm} 
}}} \\
\end{center}

\subsection{Agregamos la velocidad de fotogramas}

Idealmente, los juegos deberían ejecutarse a la misma velocidad o velocidad de fotogramas en todos los sistemas.
Para controlar la velocidad de fotogramas de un juego que puede ejecutarse en múltiples sistemas es una cuestión compleja, 
pero Pygame ofrece una forma relativamente sencilla de lograr este objetivo, el método  {\ttfamily  pygame.time.Clock()}. 
Lo agregamos en  el método{\ttfamily \_\_init\_\_()} después de {\ttfamily pygame.init()},  y aseguraremos de que el reloj
marque una vez en cada paso por el bucle principal. Cada vez que el bucle se procese más rápido que la  velocidad que definimos, Pygame calculará la cantidad de tiempo correcta para pausar  el juego y se ejecute a  una velocidad constante:

\begin{lstlisting}
class InvasionAlien:
    """Clase general para gestionar comportamiento del juego."""
    def __init__(self):
        """Inicializa el juego y crea recursos para el juego."""
        pygame.init()
        self.clock = pygame.time.Clock()
        --recorte--
\end{lstlisting}

Después de inicializar pygame, creamos un objeto  {\ttfamily self.clock }, con el módulo {\ttfamily pygame.time.Clock()}. Luego dejamos que el reloj avance al final del bucle while en {\ttfamily run\_jugar()}:
\begin{lstlisting}
    def run_jugar(self):
        """Inicia el bucle principal del juego."""
        while True:
           # Recuperar los Eventos del teclado y el mouse.
           --recorte--
           pygame.display.flip()
           self.clock.tick(60)
\end{lstlisting}
El método {\ttfamily tick()} toma  la velocidad de fotogramas del juego como el argumento; en este caso 
usamos un valor de 60, por lo que Pygame  ejecutara el bucle exactamente 60 veces por segundo.

\vspace{0.2cm}
\begin{minipage}[t]{0.2\textwidth}
	\centering\raisebox{\dimexpr\topskip-\height}{%
		\includegraphics[scale=0.20]{../A Logos de Python/Advertencia1}}  %[width=\textwidth]
	
\end{minipage}%\hfill
\begin{minipage}[t]{0.70\textwidth} 
	{\scriptsize  El reloj de Pygame debería ayudar a que el juego se ejecute de manera consistente en la mayoría de los sistemas.
		Si hace el juego se ejecuta de manera menos consistente en tu sistema, puedes probar diferentes valores de velocidad. 
		Si no puede encontrar una buena velocidad de fotogramas en su sistema, puede probar con velocidad de fotograma reloj redactar 
		por completo y ajusta la configuración del juego para que funcione bien en tu sistema.}	
\end{minipage}
\vspace{0.2cm}




\subsection{Configurar el color de fondo}

Pygame crea una pantalla negra de forma predeterminada, pero eso es muy alegre. Podemos establecer un color de fondo
agradable. El valor del color (230, 230, 230) mezcla cantidades iguales de rojo, azul y verde, lo que produce un color de fondo gris claro.  Este procedimiento lo haremos  al final del método {\ttfamily \_\_init\_\_()}, donde asignamos este color a la variable  {\ttfamily self.color\_fondo}:
 
\begin{lstlisting}
class InvasionAlien:
    """Clase general para gestionar comportamiento del juego."""
    def __init__(self):
        """Inicializa el juego y crea recursos para el juego."""
        pygame.init()
        --recorte--
        pygame.display.set_caption("Invasión Alienígena")
        # Configurando color de fondo
        self.color_fondo = (170,238,187)
    def run_game(self):
        --recorte--
    
        # Colorea la pantalla en cada paso del bucle
        self.pantalla.fill(self.color_fondo)
        # Actualiza. Hacer visible la pantalla dibujada más recientemente.
        pygame.display.flip()
        self.clock.tick(60)
\end{lstlisting}
 Y rellenamos la pantalla con el color de fondo usando el método {\ttfamily  fill()} de la forma siguiente {\ttfamily self.pantalla.fill(self.color\_fondo)}, que actúa sobre una superficie y toma un solo argumento el color.

\section{Crear una clase de configuración}

Cada vez que introducimos nuevas funciones en el juego, normalmente también creamos algunas configuraciones. Podemos escribir un módulo llamado configuración que contenga una clase llamada Configuración para almacenar todos esos valores en un solo lugar, en lugar
de agregar configuraciones a lo largo del código. Este enfoque nos permite trabajar con el objeto configuración cada vez que necesitemos acceder a una configuración individual; además  facilita la modificación, apariencia y el comportamiento del juego a medida que nuestro proyecto crece. Ahora para modificar el juego cambiaremos los valores relevantes en {\ttfamily configuración.py}, en lugar de buscar diferentes configuraciones a lo largo del proyecto.

Debemos crear el nuevo archivo llamado {\ttfamily configuración.py} y guardarlo en la carpeta principal {\ttfamily invasión\_alienigena}, con el siguiente código: 

\hspace{-2.75cm}{\ttfamily \scriptsize configurar.py}  \vspace{-0.5cm}
\begin{lstlisting}		
class Configuracion:
    """Almacena todas las confuguraciones para Invasion Alienigena."""
    def __init__(self):
        """Inicializa las configuraciones del juego"""
        # Configuración de pantalla
        self.pantalla_width = 800
        self.pantalla_height = 600
        self.Color_fondo = (170,238,187)
\end{lstlisting}

Para crear una instancia de configuración en el proyecto y usarla para acceder a nuestra configuración, necesitamos modificar el módulo {\ttfamily invasión\_alien.py} de la siguiente manera:



\hspace{-2.4cm}{\ttfamily \scriptsize invasión}

\vspace{-0.2cm}
\hspace{-2.1cm}{\ttfamily \scriptsize \_alien.py}\vspace{-0.84cm}
\begin{lstlisting}
import sys
import pygame
from configurar import Configuracion
class InvasionAlien:
    """Clase general para gestionar comportamiento del juego."""
    def __init__(self):
        """Inicializa el juego y crea recursos para el juego."""
        pygame.init()
        self.clock = pygame.time.Clock()
        self.configuracion = Configuarcion()
        self.pantalla= pygame.display.set_mode(
                         (self.configurar.pantalla_width,
                          self.configurar.pantalla_height))
        pygame.display.set_caption("Invasión Alienígena")

    def run_jugar(self):
        """Inicia el bucle principal del juego."""
        while True:
            --recorte--
            # Dibuja la pantalla durante cada paso del bucle
            self.pantalla.fill(self.color_fondo)
            #self.pantalla.blit(self.fondo, (0,0))

            self.pantalla.fill(self.settings.Color_fondo)
            # Hacer visible la pantalla dibujada más recientemente.
            pygame.display.flip()
           --recorte--
\end{lstlisting}

Del  módulo {\ttfamily configurar.py} se importa la clase {\ttfamily Configuracion()}; y después de la llamada de {\ttfamily self.clock = pygame.time.Clock() }  en método {\ttfamily \_\_init\_\_()}, creamos una instancia y se la asignamos a la variable {\ttfamily  self.configuracion = Configuarcion()}. 


Al ejecutar el script  {\ttfamily invasion\_alien.py}  no verás ningún cambio, 
porque lo único que hemos hecho es mover las configuraciones que ya estábamos usando  al archivo {\ttfamily configuar.py}.\\

Ahora estamos listos para comenzar a agregar nuevos elementos a la pantalla.


\section{Agregar la imagen de la nave}

Para dibujar la nave  del jugador en la pantalla, cargaremos una imagen y luego usaremos el método  {\ttfamily blit()} de Pygame para dibujar la imagen.\\

Puedes usar casi cualquier tipo de archivo de imagen en tu juego, pero es el más fácil
cuando usas un archivo de mapa de bits (.bmp) porque Pygame carga mapas de bits de forma predeterminada.
Sin embargo, puedes configurar PyGame para usar otros tipos de archivos, algunos tipos de archivos
dependen de ciertas bibliotecas de imágenes que deben estar instaladas en su computadora.
La mayoría de las imágenes que encontrarás están en formatos .jpg o .png, pero puedes convertirlas
a mapas de bits usando herramientas como Photoshop, GIMP o Paint.

\vspace{0.2cm}
\begin{minipage}[t]{0.2\textwidth}
	\centering\raisebox{\dimexpr\topskip-\height}{%
		\includegraphics[scale=0.20]{../A Logos de Python/Advertencia1}}  %[width=\textwidth]
	
\end{minipage}%\hfill
\begin{minipage}[t]{0.70\textwidth} 
	{\scriptsize Cuando elijas ilustraciones para tus juegos, asegúrate de prestar atención a la concesión de licencias. 
		La forma más segura y económica de empezar es utilizar  gráficos con licencia libre que puede usar y modificar desde un sitio web como \href{https://opengameart.org}{\color{vino} https://opengameart.org}.}
\end{minipage}
\vspace{0.2cm}
Se recomienda un archivo con un fondo transparente o sólido que se puede reemplazar con cualquier color de fondo, usando un editor de imágenes. Los juegos se verá mejor si el color de fondo de la imagen coincide con el color de fondo de tu juego. Alternativamente, puedes hacer coincidir color de fondo de tu juego con el fondo de la imagen.

\begin{figure}[h!]
	\centering
	\includegraphics[scale=0.70]{../A Logos de Python/nave05.jpg}
	\caption{La nave para Invasión Alienígena.}
	\label{naveReducida}
\end{figure}


En este caso, el color de fondo del archivo coincide con la configuración que estamos usando en este proyecto; y usaremos el  archivo {\ttfamily nave.bmp} (Figura \ref{naveReducida}), que es un archivo manipulado para ilustrar el juego.\\ 
% disponible en los recursos de este libro en https://ehmatthes.github.io/pcc\_3e. 

Ahora debemos crea una carpeta llamada imágenes dentro de la carpeta principal, y guardar el archivo {\ttfamily nave.bmp} en esa carpeta.





\subsection{Una clase Nave()}

Después de seleccionar la imagen para la nava, necesitamos mostrarla a en pantalla. Para ello creamos un módulo llamado {\ttfamily nave.py} que contiene una clase que llamaremos {\ttfamily Nave()}; en esta clase se maneja el comportamiento de la nave en el juego.

Pygame es eficiente porque te permite tratar todos los elementos del juego como rectángulos ({\ttfamily rect}), incluso si no tienen
exactamente la forma de rectángulos. Tratar un elemento como un rectángulo, es eficaz porque los rectángulos son formas 
geométricas simples. Cuando PyGame necesita determinar si dos elementos del juego han chocado, por ejemplo, puede hacerlo más
rápidamente si trata cada objeto como un rectángulo. Este enfoque suele funcionar lo suficientemente bien. 
En este juego notará que no estamos trabajando con la forma exacta de cada elemento del juego.
En esta clase trataremos la nave y la pantalla como rectángulos.


\hspace{-2.0cm}{\ttfamily \scriptsize nave.py}  \vspace{-0.5cm}
\begin{lstlisting}
import pygame
class Nave:
    """Una clase que maneja el comportamiento de la nave"""
    def __init__(self, ia_jugar):
        """Inicia la nave y configura la posicion inicial."""
        self.pantalla = ia_jugar.pantalla
        self.pantalla_rectang = ia_jugar.pantalla.get_rect()
        # Carga la imagen y obtiene su rectangulo
        self.imagen = pygame.imagen.load('imagenes/nave.bmp')
        self.rectang = self.image.get_rect()
               
        # comienza cada nave nueva en  centro de la pantalla
        self.rectang.mi_botom = self.pantalla_rectang.midbottom
       
    def blitme(self):
        """Dibuja la nave en la posicion actual"""
        self.pantalla.blit(self.image, self.rect)
\end{lstlisting}	
En el script importamos el módulo pygame antes de definir la clase. Luego se define la clase {\ttfamily Nave()}, y en el método  {\ttfamily  \_\_init\_\_()} se  definen dos parámetros: {\ttfamily self} y  una referencia a la instancia  {\ttfamily ia\_jugar}  actual de la clase {\ttfamily InvasionAlien}. Esto le dará  la clase {\ttfamily Nave()}  acceso a todos los recursos del juego definidos en InvasionAlien. Cuando asignamos la pantalla a un atributo de {\ttfamily Nave()}(  {\ttfamily self.pantalla = ia\_jugar.pantalla}),  podemos acceder  fácilmente a todos 
los métodos de esta clase. Accedemos al atributo {\ttfamily rect} de la pantalla usando el método {\ttfamily get\_rect()} y se asigna a {\ttfamily self.pantalla\_rectang}. Con esto nos permite colocar la nave en la ubicación correcta en la pantalla.



Para cargar la imagen, llamamos a {\ttfamily  pygame.image.load()}  y le damos la ubicación de la imagen de nuestra nave. 
Esta función devuelve una superficie que representa la nave, que asignamos a la propia imagen. Cuando se carga la imagen, llamamos a   {\ttfamily  get\_rect()} para acceder al atributo  {\ttfamily rect} de la superficie de la nave, para que luego podamos usarlo para colocar la nave.


Cuando se trabaja con objetos {\ttfamily rect}, puedes usar coordenadas $x$ y
$y$ de los bordes superior, inferior, izquierdo y derecho del rectángulo, como así como el centro, para colocar el objeto  (ver en sección \ref{objeto_Rect}). Puede establecer cualquiera de estos valores para ubicar la posición actual del objeto {\ttfamily rect}. 

Cuando estás centrando un elemento del juego, es recomendable trabajar con los atributos {\ttfamily center},
{\ttfamily centerx} o {\ttfamily centery} de un objeto {\ttfamily rect}. Cuando está trabajando en un borde de la pantalla,
puede trabajar  con la parte superior, inferior, izquierda o atributos {\ttfamily rect}. 
También hay atributos que combinan estas propiedades, como medio inferior, medio superior, medio izquierdo y medio derecho. 
Cuando estás ajustando para colocar horizontal o vertical objeto {\ttfamily rect}, puedes usar los  atributos  $x$ e $y$, que son las coordenadas  $x$ e $y$ de su esquina superior izquierda.

Estos atributos te evitan tener que hacer cálculos que los desarrolladores de juegos antes tenías que hacerlo manualmente
y los usarás con frecuencia.

\vspace{0.2cm}
\noindent \begin{minipage}[t]{0.2\textwidth}
	\centering\raisebox{\dimexpr\topskip-\height}{% \includegraphics[scale=0.35]{../Imagen_Machine/Pajaro}}  %[width=\textwidth]
	
\includegraphics[scale=0.20]{../A Logos de Python/Busquedas_R}}
	\end{minipage}%\hfill
	\begin{minipage}[t]{0.75\textwidth}
			{\scriptsize En Pygame, el origen (0, 0) está en la esquina superior izquierda de la pantalla y las coordenadas
		aumenta a medida que bajas y va hacia la derecha. En una pantalla de 1200×800, el origen está en el
		esquina superior izquierda y la esquina inferior derecha tiene las coordenadas (1200, 800). Estos
		coordenadas se refieren a la ventana del juego, no a la pantalla física.}
\end{minipage}
\vspace{0.2cm}
	
La nave la ubicaremos en la parte inferior central de la pantalla. Para ello hacemos que el valor de {\ttfamily self.rectang.midbottom} coincida con el atributo  {\ttfamily  midbottom}  del objeto {\ttfamily rect} en la pantalla (ver en sección \ref{objeto_Rect}): 
\begin{lstlisting}		
self.rectang.midbottom = self.pantalla_rectang.midbottom 	
\end{lstlisting}

Pygame usa los atributos del objeto {\ttfamily rect} para posicionar la imagen de la nave para que esté centrado horizontalmente y alineado con la parte inferior de la pantalla. La otra forma seria hacer los cálculos de forma manual. 

Finalmente, definimos el método {\ttfamily  blitme()}, que dibuja la imagen en el pantalla en la posición especificada por {\ttfamily self.rectang}.


\subsubsection{Dibujando la nave en la pantalla}

Ahora actualicemos en módulo {\ttfamily invasión\_alien.py} para que cree la nave y la dibuje en la pantalla.


\hspace{-2.4cm}{\ttfamily \scriptsize invasión}

\vspace{-0.2cm}
\hspace{-2.1cm}{\ttfamily \scriptsize \_alien.py}\vspace{-0.84cm}
\begin{lstlisting}
import sys
import pygame
fondo_image_filename = 'pygame.jpg'

from configurar import Configuracion
from nave import Nave

class InvasionAlien:
    """Clase general para gestionar comportamiento del juego."""
    def __init__(self):class InvasionAlien:
        --recorte--
       pygame.display.set_caption("Invasión Alienígena")
       self.fondo = pygame.image.load(
                      self.configuracion.fondo_image_filename).convert()
       self.nave = Nave(self)
    
    def run_game(self):
        --recorte--
        self.pantalla.blit(self.fondo,(
             (self.configuracion.pantalla_width-640)/2,
              (self.configuracion.pantalla_height-480)/2))                      
        self.nave.blitme()
        # Actualiza. Hace visible la pantalla dibujada más recientemente
        pygame.display.flip()
        --recorte--	
\end{lstlisting}

Puede ver que importamos la clase {\ttfamily Nave()} de {\ttfamily nave.py}, y luego creamos una instancia de  {\ttfamily Nave()}
después de que se haya creado la pantalla PyGame. La llamada a {\ttfamily Nave()}  requiere un argumento: una instancia de la clase {\ttfamily InvasionAlien}. En este caso el argumento {\ttfamily self}, se refiere a la instancia actual de Invasión alienígena. Este  parámetro  le  da ({\ttfamily Nave(self)}) acceso a los recursos del juego, como el objeto de pantalla.

Asignamos esta instancia de  {\ttfamily Nave()}  a la variable {\ttfamily self.nave = Nave(self)}. 
Después de llenar el fondo, dibujamos la nave en la pantalla llamando a  {\ttfamily self.nave.blitme()}, para que la nave
aparezca encima del fondo. 

Al ejecutar  {\ttfamily invasión\_alien.py}, deberías ver una pantalla de juego con la nave  en la parte inferior central, como se muestra en la Figura \ref{juego01}.



\begin{figure}[h!]
	\centering
	\includegraphics[scale=0.30]{../A Logos de Python/naveReducida05.jpg}
	\caption{La nave en la parte inferior central de la pantalla.}
	\label{juego01}
\end{figure}


\section{Refactorización}
Métodos auxiliares {\ttfamily\_chequear\_eventos()} y {\ttfamily \_actualizar\_pantalla()}\\


En proyectos grandes, a menudo se refactorizara el código que hayas escrito antes de agregar más código.
La refactorización simplifica la estructura del código que ya tienes escrito, lo que facilita la construcción. 

En esta sección, dividimos el método {\ttfamily  run\_jugar()}  que se está haciendo largo, en dos métodos auxiliares.
Un método  ayuda o auxiliar  funciona dentro de una clase pero no debe ser utilizado por código fuera de la clase.

En Python, un único guión bajo al inicio del nombre de un método  indica que es un método auxiliar (helper). 

\subsection{El método \_chequear\_eventos()}

Moveremos el código que administra los eventos a un método auxiliar separado llamado  {\ttfamily\_chequear\_eventos()}. 
Esto simplificará {\ttfamily  run\_jugar()} y aislará el bucle que  gestiona de los  eventos. Aislar el bucle de eventos le permite administrar eventos por separado de otros aspectos del juego, como la actualización de la pantalla. 

En la clase {\ttfamily InvasionAlien} se crea el nuevo  {\ttfamily\_chequear\_eventos()},
que sólo afecta al código del método {\ttfamily run\_jugar()}:

\hspace{-2.4cm}{\ttfamily \scriptsize invasión}

\vspace{-0.2cm}
\hspace{-2.1cm}{\ttfamily \scriptsize \_alien.py}\vspace{-0.84cm}
\begin{lstlisting}
--recorte--	
    def run_jugar(self):
        """Inicia el bucle principal del juego."""
        while True:
            self._check_events()
            --recorte--
        
    def _chequear_eventos(self):
          """Responde a las teclas presionada y a eventos del mouse"""
           for event in pygame.event.get():
               if event.type == pygame.QUIT:
                   pygame.quit()
                   sys.exit()	
\end{lstlisting}

Creamos un nuevo método {\ttfamily \_chequear\_eventos()} y movemos las líneas de código  que verifican
si el jugador ha hecho clic para cerrar la ventana a este nuevo método.

Para llamar a un método dentro de una clase, se usa la notación de puntos con la variable
{\ttfamily self} y el nombre del método ({\ttfamily self.\_chequear\_eventos()}). Llamamos al método dentro del bucle while en el método {\ttfamily run\_jugar()}.

\subsection{El método \_actualizar\_pantalla()}  


Para simplificar aún más {\ttfamily run\_jugar()}, moveremos el código para actualizar la pantalla
a un método separado llamado  {\ttfamily\_actualizar\_pantalla()}:


\begin{lstlisting}
--recorte--		
    def run_game(self):
        """Inicia el bucle principal del juego."""
        while True:
            self._chequear_eventos()
            self._actualizar_pantalla()
            self.clock.tick(60)
    def _chequear_eventos(self):
        --recorte--
    def _actualizar_pantalla(self):
        """Actualiza images sobre la pantalla, y flip actualiza la nuave pantalla."""
        self.pantalla.fill(self.configuracion.color_fondo)
        self.nave.blitme()
        pygame.display.flip()	

\end{lstlisting}	

Hemos movimos el código que dibuja el fondo de pantalla y la nave al método {\ttfamily\_actualizar\_pantalla()}. 

Ahora el cuerpo del bucle principal en {\ttfamily run\_jugar()} es mucho más simple. 
Es fácil ver que estamos buscando nuevos eventos, actualizando la pantalla y  marcando el reloj en cada paso por el bucle.

Si ya has creado varios juegos, probablemente comenzarás dividiendo tu código en métodos como estos. Pero si nunca ha abordado
un proyecto como este, probablemente al principio no sabrá exactamente cómo estructurar su código. Este enfoque le da una idea
de un proceso de desarrollo realista: comienza escribiendo su código de la manera más simple posible y luego lo refactoriza a
medida que su proyecto se vuelve más complejo.

Ahora que hemos reestructurado el código para que sea más fácil agregarlo, ¡podemos trabajar en los aspectos dinámicos del juego!

\section{Ejercicios}  % INTÉNTALO TÚ MISMO

Cielo azul: crea una ventana de Pygame con un fondo azul.\\


Personaje del juego: busque una imagen de mapa de bits de un personaje del juego que le guste o convierta una imagen
a un mapa de bits. \\

Cree una clase que dibuje el personaje en el centro de la pantalla, luego haga coincidir el color de fondo
de la imagen con el color de fondo de la pantalla o viceversa.

\section{Pilotando la nave}

A continuación, le daremos al jugador la posibilidad de mover la nave hacia la derecha y hacia la izquierda. 
Escribiremos un código que responda cuando el jugador presione la tecla de flecha derecha o izquierda. 
Nos centraremos primero en el movimiento  hacia la derecha y luego aplicaremos los mismos principios para controlar el 
movimiento hacia la izquierda. 
A medida que agreguemos este código, aprenderá cómo controlar el movimiento de las imágenes en la pantalla y responder a las entradas del usuario.

\subsubsection{Responder a una pulsación de tecla}

Cada vez que el jugador presiona una tecla, esa pulsación se registra en Pygame como un evento. 
Cada evento es recogido por el método {\ttfamily event.get()}. Necesitamos especificar en nuestro método {\ttfamily \_chequear\_eventos()} qué tipo de eventos queremos que el juego verifique.

Cada pulsación de tecla se registra como un evento {\ttfamily KEYDOWN}. Cuando Pygame detecta un evento {\ttfamily KEYDOWN}, debemos verificar si la tecla que se presionó es la que desencadena determinada acción. Por ejemplo, si el jugador presiona la tecla de flecha derecha, se debe  aumentar el valor {\ttfamily rect.x} de la nave para molerla hacia la derecha:


\hspace{-2.5cm}{\ttfamily \scriptsize curso.py}  \vspace{-0.5cm}
\begin{lstlisting}		
    def _chequear_eventos(self):
        """Responde a las teclas presionada y a eventos del mouse"""
        for event in pygame.event.get():
            if event.type == pygame.QUIT:
            pygame.quit()
            sys.exit()	
       elif event.type == pygame.KEYDOWN:
            if event.key == pygame.K_RIGHT:
                # Mueve la nave a la derecha
                self.nave.rectang.x += 1
\end{lstlisting}	



Dentro de {\ttfamily \_chequear\_eventos()} agregamos un bloque {\ttfamily elif} al bucle de eventos, para responder
cuando PyGame detecta un evento {\ttfamily KEYDOWN}. Comprobamos si la tecla presionada es la tecla de flecha derecha. La tecla de flecha derecha está representada por {\ttfamily pygame.K\_RIGHT}. Si se presionó la tecla de flecha derecha, movemos la nave a la
derecha  aumentando el valor de {\ttfamily self.nave.rect.x} en una unidad.

Al ejecutas {\ttfamily invasion\_alien.py}, la nave debe moverse hacia la derecha una unidad de píxel cada vez que presiona la tecla de flecha derecha. Eso es un comienzo, pero no lo es una manera eficiente de controlar la nave. Mejoremos este control permitiendo
movimiento continuo.


\subsubsection{Permitir el movimiento continuo}

Si el jugador mantiene presionada la tecla de flecha derecha la nave debe moverse continuamente hacia la derecha, hasta que el jugador suelte la tecla. Para esto necesitamos que el juego detecte los eventos {\ttfamily pygame.KEYUP} hasta que se suelta la tecla; usando los eventos {\ttfamily KEYDOWN} y {\ttfamily KEYUP} con un identificador llamado {\ttfamily  movimiento\_derecho}podemos el movimiento continuo.
Cuando el identificador {\ttfamily  movimiento\_derecho} sea {\ttfamily False}, la nave estará inmóvil, y cuando el jugador presiona la tecla de flecha derecha, estableceremos el identificador como {\ttfamily True}, cuando el jugador suelta la tecla, volveremos a establecer la identificador como {\ttfamily False}.


La clase {\ttfamily Nave()} controla todos los atributos la nave, así que le daremos un atributo
llamado  {\ttfamily  self.movimiento\_derecho} y un método  {\ttfamily actualizar()} para verificar el estado del identificador 
{\ttfamily  self.movimiento\_derechoo}. El método  {\ttfamily actualizar()} cambiará la posición la nave si el 
 identificador es  {\ttfamily True}. 
 
Llamaremos a este método una vez en cada paso por el bucle while para actualizar la posición de la nave.
Veamos los cambios en  {\ttfamily nave.py}:

\hspace{-2.0cm}{\ttfamily \scriptsize nave.py}  \vspace{-0.5cm}
\begin{lstlisting}
import pygame
class Nave:
    """Una clase que maneja el comportamiento de la nave"""
    def __init__(self, ia_jugar):
        --recorte--
        # Cada nave nueva empieza al centro inferior de la pantalla
        self.rectang.midbottom = self.pantalla_rectang.midbottom    
        # Señal de identificador; comienza la nave sin movimiento
        self.movimiento_derecho = False

    def actualizar(self):
        """Actualiza posición de la nave basado en la señal del indicador."""
        if self.movimiento_derecho:
            self.rectang.x += 1

    def blitme(self):
         --recorte--	
\end{lstlisting}



Se agrego el atributo {\ttfamily self.movimiento\_derecho} en el método {\ttfamily \_\_init\_\_()}, y lo configuramos inicialmente como {\ttfamily False}. Luego agregamos en el método {\ttfamily actualizar()}   el código:
\begin{lstlisting}		
if self.movimiento_derecho:
    self.rectang.x += 1	
\end{lstlisting}
que permite mover la nave hacia la derecha si el indicador es {\ttfamily True}. El método  {\ttfamily actualizar()} será  llamado desde afuera de la clase, por lo que no se considera un método auxiliar. 

Ahora necesitamos modificar el método {\ttfamily \_chequear\_eventos(self)} para establecer  el indicador  {\ttfamily self.movimiento\_derecho}  en {\ttfamily True} cuando se presiona  la tecla de flecha derecha y en {\ttfamily False} cuando se suelta:

\begin{lstlisting}
    def  _chequear_eventos(self):
        """Responde a las teclas presionada y a eventos del mouse"""
        --recorte--
        elif event.type == pygame.KEYDOWN:
             if event.key == pygame.K_RIGHT:
                 self.nave.movimiento_derecho = True
        elif event.type == pygame.KEYUP:
             if event.key == pygame.K_RIGHT:
                 self.nave.movimiento_derecho = False
\end{lstlisting}
Y eliminamos el código del módulo  {\ttfamily \_chequear\_eventos(self)}:
\begin{lstlisting}
elif event.type == pygame.KEYDOWN:
    if event.key == pygame.K_RIGHT:
        # Mueve la nave a la derecha
        self.nave07.rectang.x += 1
\end{lstlisting}

Aquí, se modifico la forma cómo responde el juego cuando el jugador presiona la tecla de flecha derecha: en lugar de cambiar
la posición de la nave directamente, simplemente configuramos {\ttfamily self.nave.movimiento\_derecho = True}.  Luego agregamos un nuevo bloque  {\ttfamily elif}, que responde a los eventos {\ttfamily KEYUP} cuando el jugador suelta la tecla de flecha derecha ({\ttfamily K\_RIGHT}), y se configura inicialmente en False ({\ttfamily self.nave.movimiento\_derecho = False}.\\


A continuación, modificamos el bucle while en {\ttfamily run\_jugar()} para que llame al método  {\ttfamily actualizar()} de la nave en cada paso por el bucle:

\hspace{-2.4cm}{\ttfamily \scriptsize invasión}

\vspace{-0.2cm}
\hspace{-2.1cm}{\ttfamily \scriptsize \_alien.py}\vspace{-0.84cm}
\begin{lstlisting}
    def run_jugar(self):
        """Inicia el bucle principal del juego."""
        while True:
            self._chequear_eventos()
            self.nave.actualiza()
            self._actualiza_pantalla()
            self.clock.tick(60)
\end{lstlisting}

La posición de la nave se actualizará después de que hayamos verificado los eventos del teclado y antes de actualizar la pantalla.
Esto permite que la posición de la nave se actualice en respuesta a la entrada del jugador y garantiza que la posición actualizada se utilizará al dibujar la nave en la pantalla.


Al  ejecutar {\ttfamily invasion\_alien.py} y mantiene presionada la tecla de flecha derecha, la nave debe moverse continuamente hacia la
derecha hasta que suelte la tecla.


\subsubsection{Moviéndose tanto hacia la izquierda como hacia la derecha}

Ahora que la nave puede moverse continuamente hacia la derecha, agregando el movimiento
hacia la izquierda. Nuevamente, modificaremos la clase {\ttfamily Nave()} y el método  {\ttfamily \_chequear\_eventos()}:. 

Los cambios en {\ttfamily nave.py},  en los métodos  {\ttfamily \_\_init\_\_()}  y {\ttfamily actualizar()} se muestran a continuación:

\hspace{-2.0cm}{\ttfamily \scriptsize nave.py}  \vspace{-0.5cm}
\begin{lstlisting}	
 --recorte--
    def __init__(self, ai_game):
        --recorte--
        # Señal de identificador; comienza la nave sin movimiento
        self.movimiento_derecho = False
        self.movimiento_izquierdo = False

    def actualizar(self):
        """Actualiza posición de la nave basado en la señal del indicador."""
        if self.movimiento_derecho:
            self.rectang.x += 1
        if self.movimiento_izquierdo:
            self.rect.x -= 1	
\end{lstlisting}

En el método {\ttfamily \_\_init\_\_()}, se agregó  el identificador {\ttfamily self.movimiento\_\\izquierdo}. 
Y en el método {\ttfamily actualizar()}, usamos dos bloques {\ttfamily if} separados, en lugar de un {\ttfamily elif}, para permitir que
el valor {\ttfamily rect.x} de la nave aumente y  disminuya cuando se mantienen presionadas ambas teclas de flecha, provoca que la nave se detenga. 

Si usamos {\ttfamily elif} para el movimiento hacia la izquierda, la tecla de flecha derecha siempre tendría prioridad. 
Usar dos bloques {\ttfamily if} hacen los movimientos precisos cuando el jugador  presione momentáneamente
ambas teclas al cambiar de dirección.


Finalizamos estos ajustes haciendo dos cambios adiciones en la clase {\ttfamily class InvasionAlien()}, específicamente   en el método  {\ttfamily  \_chequear\_eventos()}: 



\hspace{-2.4cm}{\ttfamily \scriptsize invasión}

\vspace{-0.2cm}
\hspace{-2.1cm}{\ttfamily \scriptsize \_alien.py}\vspace{-0.84cm}
\begin{lstlisting}
--recorte--		
    def _chequear_eventos(self):
        """Responde a las teclas presionada y a eventos del mouse"""
        --recorte-- 
        
        elif event.type == pygame.KEYDOWN:
            if event.key == pygame.K_RIGHT:
                self.nave07.movimiento_derecho = True
            elif event.key == pygame.K_LEFT:
                self.nave07.movimiento_izquierdo = True

        elif event.type == pygame.KEYUP:
            if event.key == pygame.K_RIGHT:
                self.nave07.movimiento_derecho = False
            elif event.key == pygame.K_LEFT:               
                self.nave07.movimiento_izquierdo = False
 --recorte--                
\end{lstlisting}

Si ocurre un evento {\ttfamily KEYDOWN} generado por la tecla {\ttfamily K\_LEFT}, el indicador  {\ttfamily movimiento\_izquierdo}{\ttfamily True}; si  ocurre  un evento {\ttfamily KEYUP}  generado por la tecla {\ttfamily K\_LEFT}, el indicador  {\ttfamily movimiento\_izquierdo}  es configurado como {\ttfamily False}.  No podemos usar un bloquea  {\ttfamily elif},  porque cada evento está conectado a una sola tecla, y al jugador presionas ambas teclas a la vez, se detectarán dos eventos separados.


Al  ejecutar {\ttfamily invasion\_alien.py}, debe poder mover la nave continuamente a la derecha y a la izquierda. 
Además, si mantienes presionadas ambas teclas, la nave debería dejar de moverse.\\

A continuación, perfeccionaremos aún más el movimiento la nave. Ajustemos la velocidad de la nave y limitemos qué tan
lejos puede moverse para que no desaparezca de los lados de la pantalla.


\subsubsection{Ajustar la velocidad de la nave}

Actualmente, la nave se mueve un píxel por ciclo a través del bucle while, pero podemos tomar un control más preciso de
la velocidad de la nave agregando un atributo  la velocidad de la nave {\ttfamily velocidad\_nave}, en la  clase {\ttfamily Configuracion()}. 

Usaremos este atributo para determinar  qué tan lejos mover ella nave en cada paso por el circuito.

Agregamos el nuevo atributo en  {\ttfamily configurar.py}:

\hspace{-2.75cm}{\ttfamily \scriptsize configurar.py}  \vspace{-0.5cm}
\begin{lstlisting}
class Configuracion:
    """Una clase que almacena todas las confuguarciones 
    para Invasion Alienigena."""
    def __init__(self):
        """Inicializa las configuraciones del juego"""
        # Configuración de pantalla
        self.fondo_image_filename = 'pygame.jpg' 
        self.pantalla_width = 800
        self.pantalla_height = 600
        self.color_fondo = (170,238,187)#(230, 230, 230)
        self.velocidad_nave = 1.5
\end{lstlisting}
Establecemos el valor inicial de {\ttfamily velocidad\_nave} en 1,5. Cuando ella nave se mueva ahora, su posición se ajusta en 1,5 píxeles 
(en lugar de 1 píxel) en cada paso por el bucle. 
Y usamos un flotador para configurar la velocidad para tener un control más preciso de la velocidad de la nave cuando aumentamos
el ritmo del juego más adelante. Sin embargo, los atributos {\ttfamily rect} como $x$ se almacenan solo valores enteros,
por lo que debemos   realizar algunas modificaciones en la clase {\ttfamily Nave()}:

\hspace{-2.0cm}{\ttfamily \scriptsize nave.py}  \vspace{-0.5cm}
\begin{lstlisting}
import pygame
class Nave:
    """Una clase que maneja el comportamiento de la nave"""
    def __init__(self, ia_jugar):
        """Inicia la nave y configura la posicion inicial."""
        self.pantalla = ia_jugar.pantalla
        self.configuracion = ia_jugar.configuracion
        
        self.pantalla_rectang = ia_jugar.pantalla.get_rect()

        # Carga la imagen y obtiene su rectangulo
        #self.imagen = pygame.image.load('../A Logos de Python/nave.bmp')   
        self.imagen = pygame.image.load('imagenes/nave05.bmp') #("nave.bmp")
        self.rectang = self.imagen.get_rect()        
            

        # Cada nave nueva empieza al centro inferior de la pantalla
        self.rectang.midbottom = self.pantalla_rectang.midbottom 
        # Almacena un float para la posicion exacta horizontal de la nave.
        self.x = float(self.rectang.x)           
        # Señal de identificador; comienza la nave sin movimiento
        self.movimiento_derecho = False
        self.movimiento_izquierdo = False

    def actualizar(self):
        """Actualiza posición de la nave basado en la señal del indicador."""
        # Actualiza el valor x de la nave, no en el objeto rect.
        if self.movimiento_derecho:
            self.x += self.configuracion.velocidad_nave
            #self.rectang.x += 1
        if self.movimiento_izquierdo:
             self.x -= self.configuracion.velocidad_nave
             #self.rectang.x -= 1	
        # Actualizael objeto rectang con self.x.
        self.rectang.x = self.x
    def blitme(self):
        """Dibuja la nave en la posicion actual"""
        self.pantalla.blit(self.imagen, self.rectang)
\end{lstlisting}

Creamos un atributo {\ttfamily self.configuracion = ia\_jugar.configuracion} de configuración en clase {\ttfamily Nave()}, 
para usarlo en {\ttfamily actualizar()}. Debido a que estamos ajustando la posición de la nave en fracciones de píxel,
es necesario asignar la posición a una variable a la que se le pueda asignar un flotante. Se puede usar un flotante para establecer 
un atributo del objeto  {\ttfamily rectang}, pero el {\ttfamily rect} solo mantendrá la porción entera de ese valor. 
Para realizar un seguimiento preciso de la posición de la nave,  definimos un nuevo atributo {\ttfamily self.x}; y usamos la función {\ttfamily float()} para convertir el valor de {\ttfamily self.rect.x} en un flotante y luego asignar este valor a  {\ttfamily self.x}. 

Ahora,  cambiamos la posición de la nave en {\ttfamily actualizar()}, el valor de  {\ttfamily self.x} se ajusta según la cantidad 
almacenada en {\ttfamily self.configuracion.velocidad\_nave}. 

Después que  {\ttfamily self.x} se ha actualizado, utilizamos el nuevo valor para actualizar {\ttfamily self.rect.x}, que controla
la posición de la nave. 

Sólo la porción entera de  {\ttfamily self.x} será asignado a {\ttfamily self.rect.x}, pero eso está bien.


Ahora podemos cambiar el valor de {\ttfamily velocidad\_nave} a cualquier valor mayor que 1, y  hará que la nave se mueva más rápido. 
Esto ayudará a que la nave responda lo suficientemente rápido como para derribar alienígenas, y nos permitirá cambiar el ritmo
del juego a medida que el jugador avanza en el juego.

\subsubsection{Limitar el alcance de la nave}

La nave desaparecerá en cualquier borde de la pantalla si mantiene presionada una tecla de flecha el tiempo suficiente. 
Para evitar esta situación haremos que nave se detenga al llega al borde de la pantalla. Esto lo hacemos modificando el
método {\ttfamily actualizar()} en {\ttfamily Nave()}:



\hspace{-2.5cm}{\ttfamily \scriptsize curso.py}  \vspace{-0.5cm}
\begin{lstlisting}		
    --recorte--		
    def _chequear_eventos(self):
    """Responde a las teclas presionada y a eventos del mouse"""
    --recorte-- 

    def actualizar(self):
        """Actualiza posición de la nave basado en la señal del indicador."""
        # Actualiza el valor x de la nave, no en el objeto rect       
        if self.movimiento_derecho and self.rectang.right < self.pantalla_rectang.right:        
            self.x += self.configuracion.velocidad_nave        
        if self.movimiento_izquierdo and self.rectang.left > 0:       
            self.x -= self.configuracion.velocidad_nave      	
        # Actualiza el objeto rectang con self.x.
        self.rectang.x = self.x
       
        def blitme(self):
            --recorte-- 
\end{lstlisting}
% Este código comprueba la posición de la nave antes de cambiar el valor de self.x. 

El código {\ttfamily self.rectang.right} devuelve la coordenada $x$ del borde derecho del rectángulo de la nave. 
Si este valor es menor que el valor devuelto por {\ttfamily self.pantalla\_rect.right}, ella nave no ha llegado al borde derecho de la pantalla.  De igual forma con el borde izquierdo: si el valor del lado izquierdo del rectángulo es mayor que 0,
la  nave no ha llegado al borde izquierdo de la pantalla. Esto asegura que la nave está dentro de estos límites antes de 
ajustar el valor de {\ttfamily self.x}.\\


Al ejecutar  {\ttfamily invasion\_alien.py}, la nave debería dejar de moverse en cualquier borde de la pantalla. 
Esto está muy bien; todo lo que hemos hecho es agregar una prueba condicional en una declaración {\ttfamily if}, 
pero se siente como si la nave golpeara una pared o un campo de fuerza en los borde de la pantalla.


\section{Refactorización \_chequear\_eventos()}


El método {\ttfamily \_chequear\_eventos()}  en  {\ttfamily invasion\_alien.py} aumentará en longitud a medida que avancemos en el desarrollo del juego,  así que dividamos  en dos métodos separados: uno que maneja eventos {\ttfamily KEYDOWN} y otro que maneja eventos{\ttfamily KEYUP}:


\hspace{-2.4cm}{\ttfamily \scriptsize invasión}

\vspace{-0.2cm}
\hspace{-2.1cm}{\ttfamily \scriptsize \_alien.py}\vspace{-0.84cm}
\begin{lstlisting}
--recorte--    
    def  _chequear_eventos(self):
        """Respuesta a teclas presionads y eventos del muuse"""
        for event in pygame.event.get():
            if event.type == pygame.QUIT:
                pygame.quit()
                sys.exit()

            elif event.type == pygame.KEYDOWN:
                self._chequear_keydown_eventos(event)
            elif event.type == pygame.KEYUP:
                self._chequear_keyup_eventos(event)

    def _chequear_keydown_eventos(self, event):
        """Respuesta a teclas presionadas"""
        if event.key == pygame.K_RIGHT:
            self.nave.movimiento_derecho = True
        elif event.key == pygame.K_LEFT:
            self.nave.movimiento_izquierda = True
        if event.key == pygame.K_q:
            pygame.quit()
            sys.exit()            

    def _chequear_keyup_eventos(self, event):
        """respuesta a tecla liberada"""
        if event.key == pygame.K_RIGHT:
            self.nave.movimiento_izquierda = False
        elif event.key == pygame.K_LEFT:
            self.nave.movimiento_izquierda  = False
   --recorte-- 
\end{lstlisting}
Creamos dos métodos  auxiliares para reemplazamos el código antiguo del  método {\ttfamily \_chequear\_eventos()}, el método {\ttfamily \_chequear\_keydown\_eventos()} y el método {\ttfamily \_chequear\_keyup\_eventos}.  Cada uno con un parámetro {\ttfamily self} y un parámetro  para los eventos. Esta estructura de código es más limpia y facilitará el desarrollo de respuestas adicionales a las entradas del jugador.

 
\subsubsection{Presionando Q para salir}

Ahora respondemos a las pulsaciones de teclas de manera eficiente, y podemos agregar otra forma adicional para salir del juego. Muy útil cuando se trabaja con ventanas sin borde (ver sección \ref{sin_bordes}).
Hace clic  en el botón {\ttfamily X} en la esquina superior derecha de la ventana  del juego para finalizar el juego cada vez que pruebas una nueva función algunas veces es tedioso, por lo que agregaremos un atajo de teclado para finalizar el juego cuando el jugador presiona $q$:


En  el  archivo {\ttfamily invasion\_alien.py} agregamos la segunda opcion para salir del juego, específicamente al final del
método auxiliar {\ttfamily \_chequear\_keydown\_eventos()}, agregamos un nuevo bloque{\ttfamily elif} que finaliza el juego 
cuando el jugador presiona {\ttfamily q} o  {\ttfamily Q}. 

\hspace{-2.4cm}{\ttfamily \scriptsize invasión}

\vspace{-0.2cm}
\hspace{-2.1cm}{\ttfamily \scriptsize \_alien.py}\vspace{-0.84cm}
\begin{lstlisting}
--recorte--    
    def  _chequear_eventos(self):
        """Respuesta a teclas presionads y eventos del muuse"""
        for event in pygame.event.get():
            if event.type == pygame.QUIT:
                pygame.quit()
                sys.exit()

    def _chequear_keydown_eventos(self, event):
        """Respuesta a teclas presionadas"""
        if event.key == pygame.K_RIGHT:
            self.nave07.movimiento_derecho = True              
        elif event.key == pygame.K_LEFT:
           self.nave07.movimiento_izquierdo = True 
        elif event.key==pygame.K_q:
            pygame.quit()
            sys.exit()
        --recorte-- 
\end{lstlisting}
	
Ahora, al ejecutar {\ttfamily invasion\_alien.py},  puedes presionar Q para cerrar el juego, como una opción más.

\subsection{Ejecutar el juego en modo de pantalla completa}

Para ejecutar el juego en modo de pantalla debemos hacer unos cambios sencillos en configuración de pantalla en el archivo {\ttfamily invasion\_alien.py}, dentro del método {\ttfamily init}:


\hspace{-2.4cm}{\ttfamily \scriptsize invasión}

\vspace{-0.2cm}
\hspace{-2.1cm}{\ttfamily \scriptsize \_alien.py}\vspace{-0.84cm} 
\begin{lstlisting}
        --recorte-- 
class AlienInvasion:
    """Clase general para gestionar comportamiento del juego."""
    def __init__(self):
        """Inicializa el juego y crea recursos para el juego."""
        pygame.init()
        self.clock = pygame.time.Clock()
        self.configuracion = Configuracion()
        self.pantalla = pygame.display.set_mode((0, 0), pygame.FULLSCREEN)
        self.configuracion.pantalla_width = self.pantalla.get_rect().width
        self..configuracion.pantalla_height = self.pantalla.get_rect().height
        
        pygame.display.set_caption("Invasión Alienígena")
        self.pantalla= pygame.display.set_mode(
                          (self.configuracion.pantalla_width,
                           self.configuracion.pantalla_height))
        pygame.display.set_caption("Invasión Alienígena")
        self.fondo = pygame.image.load(
                       self.configuracion.fondo_image_filename).convert()
        self.nave07 = Nave(self)
        --recorte-- 
\end{lstlisting}

Se usaron los atributos de ancho y alto del objeto {\ttfamily rect()} de la pantalla ({\ttfamily self.pantalla.get\_rect().width}) para actualizar la configuración de pantalla. Al crear la superficie de la pantalla, le pasamos un tamaño de (0, 0) y el parámetro 
{\ttfamily pygame.FULLSCREEN}. Esto le dice a PyGame que busque una  tamaño de ventana que llenará la pantalla. Porque no conocemos el ancho y el alto de la pantalla con anticipación, actualizamos esta configuración después de crear la pantalla.

\vspace{0.2cm}
\begin{minipage}[t]{0.2\textwidth}
	\centering\raisebox{\dimexpr\topskip-\height}{%
		\includegraphics[scale=0.20]{../A Logos de Python/Advertencia1}}  %[width=\textwidth]
	
\end{minipage}%\hfill
\begin{minipage}[t]{0.70\textwidth} 
	{\scriptsize Asegúrate de poder salir presionando Q antes de ejecutar el juego en modo de pantalla completa;
		PyGame no ofrece una forma predeterminada de salir de un juego en modo de pantalla completa.}
\end{minipage}
\vspace{0.2cm}

En este caso al ejecutar y desplegaremos el juego en modo de pantalla completa no tenemos la opción de hacer clic en la esquina superior derecha de la pantalla por que no se ve; pero con la  opción  {\ttfamily Q} salimos del juego.\\


Algunos juegos se ven mejor en modo de pantalla completa, y en algunos sistemas, el juego puede funcionar mejor en general en
modo de pantalla completa. Si te gusta cómo se ve o como se comporta el juego en modo de pantalla completa, puede sigue
con estos ajustes, y si no puedes volver al enfoque original donde establecimos un tamaño de pantalla específico para
el juego.


\section{Resumiendo el trabajo}

En la siguiente sección, agregaremos la capacidad de disparar balas, lo que implica agregar un nuevo archivo llamado {\ttfamily balas.py} y realizando algunas modificaciones en algunos de los archivos que ya estamos usando. 

Hasta ahora tenemos tres archivos que contienen un número de clases y métodos. Para tener claro cómo está organizado el proyecto,
revisemos cada uno de estos archivos antes de agregar más funciones.


\subsection{invasión\_alien.py}

El archivo principal {\ttfamily invasion\_alien.py}, contiene la clase  {\ttfamily InvasionAlien}. Esta clase crea una 
serie de atributos importantes que se utilizan a lo largo del juego: como las configuraciones se asignan
de la superficie principal de visualización de la pantalla, también se crea una instancia de la nave en este archivo. 
El bucle principal del juego, un bucle while, también se almacena en este módulo. El bucle while llama 
a {\ttfamily \_chequear\_eventos()},  y {\ttfamily \_actualizar\_pantalla()}. También marca el reloj en cada paso por el bucle.

El método {\ttfamily \_chequear\_eventos()} detecta eventos relevantes, como pulsaciones y liberaciones de teclas,  procesa cada uno de estos tipos de eventos a través de los métodos  { }{\ttfamily \_chequear\_keydown\_eventos()} { } y  { } {\ttfamily \_chequear\\\_keyup\_eventos()}. Por ahora, estos métodos gestionan el movimiento de la nave. 

La clase {\ttfamily InvasionAlien()} también contiene {\ttfamily \_actualizar\_pantalla()}, que  dibujar la pantalla en cada paso por el bucle principal.

El archivo {\ttfamily invasion\_alien.py} es el único archivo que necesitas ejecutar cuando quieras jugar {\ttfamily Invasión Alienígena}. 

Los otros archivos, {\ttfamily configurar.py} y {\ttfamily nave.py}, contienen código que se importa a este archivo.

\subsection{configurar.py}

El archivo {\ttfamily configurar.py} contiene la clase {\ttfamily Configuracion()}. Esta clase sólo tiene un método {\ttfamily\_\_init\_\_()}, que inicializa los atributos que controlan la apariencia del juego y la velocidad de la nave.
	
\subsection{nave.py}

El archivo  {\ttfamily nave.py} contiene la clase  {\ttfamily Nave()}, y esta contiene un método {\ttfamily\_\_init\_\_()}, el método   {\ttfamily actualizar()} para gestionar la posición de la nave,  y  el método {\ttfamily  blitme()} para dibujar la nave en la pantalla. La imagen de la nave es almacena en {\ttfamily nave.bmp}, que se encuentra en la carpeta de imágenes.


\section{INTÉNTALO TÚ MISMO}

12-3. Documentación de Pygame: 
Ya hemos avanzado lo suficiente en el juego como para que desees consultar parte de la documentación de Pygame. 
La página de inicio de Pygame está en https://pygame.org y la página de inicio de la documentación está en https://pygame.org/docs. Simplemente lea la documentación por ahora. No lo necesitarás para completar este proyecto, pero te ayudará si luego quieres modificar Alien Invasion o crear tu propio juego.

12-4. Cohete: 
Crea un juego que comience con un cohete en el centro de la pantalla. Permita que el jugador mueva el cohete hacia arriba, abajo, izquierda o derecha usando las cuatro teclas de flecha. Asegúrate de que el cohete nunca se mueva más allá de ningún borde de la pantalla.

12-5. Teclas: 
Cree un archivo Pygame que crea una pantalla vacía. En el bucle de eventos, imprima el atributo event.key cada vez que se detecte un evento pygame.KEYDOWN. Ejecute el programa y presione varias teclas para ver cómo responde Pygame.



\section{Disparar balas}
%  Escribiremos código que dispare una bala, que está representado por un pequeño rectángulo, cuando el jugador presiona la barra espaciadora.

Ahora le agregaremos la capacidad de disparar balas a la nave. Para esta tarea crearemos un módulo {\ttfamily balas.py}, donde se construye las balas, tamaño, posición, y disparo de balas. 

Se usara  el módulo {\ttfamily  pygame.Rect()} para crear un rectángulo que representara las balas;  las balas serán disparadas desde la nave hacia arriba; es decir, las balas viajarán desde la nave hacia arriba hasta que desaparezcan la parte superior de la pantalla.

Y en el módulo  {\ttfamily configurar.py} se agregan las configuración para las balas al final del método {\ttfamily \_\_init\_\_()},  aquí  incluimos los valores que necesitaremos para la nueva clase {\ttfamily Balas}, velocidad, tamaño y color:

\hspace{-2.75cm}{\ttfamily \scriptsize configurar.py}  \vspace{-0.5cm}
\begin{lstlisting}	
--recorte--	
    def __init__(self):
        --recorte--
        # configuraciones para las balas
        self.velocidad_balas = 6.0
        self.ancho_balas = 10  #  2
        self.alto_balas = 16
        self.color_balas = (0, 0, 0)
\end{lstlisting}
Estas configuraciones crean balas de color negro con un ancho de 3 píxeles y un altura de 16 píxeles. Las balas viajarán un poco más rápido que ella nave.

Ahora debemos crear la clase  {\ttfamily Balas()}.
\subsubsection*{Clase Balas}

Debe crear un módulo {\ttfamily balas.py} para almacenar la clase {\ttfamily Balas()}, y guardarlo en la carpeta principal. En este módulo  escribimos el siguiente código:

\hspace{-2.3cm}{\ttfamily \scriptsize  balas.py} \vspace{-0.5cm}
\begin{lstlisting}		
import pygame
from pygame.sprite import Sprite

class  Balas(Sprite):
    """Una clase para manejar las balas disparadas desde la nave"""
    def __init__(self, ia_jugar):
        """crea un objeto bala desde la posicion  actual de la nave"""
        super().__init__()
        self.pantalla = ia_jugar.pantalla
        self.configuraciones = ia_jugar.configuraciones
        self.color = self.configuraciones.color_balas
        # Crea un rectangulo en (0, 0) y se configura la posicion correcta
        self.rectang = pygame.Rect(0, 0, self.configuraciones.ancho_balas,
                                         self.configuraciones.alto_balas)
        self.rectang.midtop = ai_jugar.nave.rectang.midtop
        # Almacena la posicion de las balas como un flotante
        self.y = float(self.rectang.y)
\end{lstlisting}
Aquí estamos usando una clase {\ttfamily Sprite} del módulo {\ttfamily pygame.sprite}  como argumento  de la clase {\ttfamily Balas()}; de esta forma nuestra clase es una clase hija,  heredera de la clase  {\ttfamily Sprite}.\\ 

Se importa del módulo  {\ttfamily pygame.sprite} la clase {\ttfamily Sprite}.  Cuando usas  el módulo {\ttfamily pygame.sprite}, puedes agrupar elementos relacionados en tu juego y actúa sobre todos los elementos agrupados a la vez.
 
Para crear una instancia de {\ttfamily Balas()}, se necesita la instancia actual de la clase {\ttfamily InvasionAlien}, y llamamos
a {\ttfamily super().\_\_init\_\_()} para heredar correctamente de la clase {\ttfamily Sprite}.  También se configuran los 
atributos para la pantalla, objetos de configuración, y  el color de la balas.\\


Ahora creamos el atributo {\ttfamily self.rectang}  para la balas. Las balas no se basa en una imagen, por lo que tenemos que 
construir un rectángulo, para ello se usara la clase {\ttfamily pygame.Rect()}. Esta clase requiere de las coordenadas
$x$ e $y$ de la esquina superior izquierda del rectángulo, el ancho y alto del rectángulo. 

Inicializamos el rectángulo en (0, 0), pero lo moveremos a la ubicación correcta en la siguiente línea, porque 
la posición de la bala depende de la posición de la nave. Obtenemos el ancho y el alto de la bala a partir de los valores almacenados en {\ttfamily self.configuraciones}. Configuramos el atributo{\ttfamily  .midtop} de la bala para que coincida con el atributo 
{\ttfamily .midtop} de la nave. Esto hará que la bala emerja desde la parte superior de la nave, haciendo que parezca que la bala se dispara desde la nave. Usamos un flotador para la coordenada $y$ de la bala para poder hacer ajustes a la velocidad de la bala.

Agregamos una segunda parte al archivo  {\ttfamily balas.py} que consiste en dos métodos, {\ttfamily update()} y {\ttfamily dibuja\_balas()}:

\hspace{-2.3cm}{\ttfamily \scriptsize  balas.py} \vspace{-0.5cm}
\begin{lstlisting}
--recorte--		
    def update(self):
        """Movimiento de bala hacia arriba por la pantalla"""
        # Actualiza  exacta de la posicion de las balas.
        self.y -= self.configuraciones.velocidad_balas
        # Actualiza la posicion del rectangulo
        self.rectang.y = self.y
    
    def dibuja\_balas(self):
        """Dibuja las balas en la pantalla"""
        pygame.draw.rect(self.pantalla, self.color, self.rectang)
\end{lstlisting}

El método {\ttfamily actualiza()} gestiona la posición de la bala. Cuando se dispara una bala, sube por la pantalla, 
lo que corresponde a un valor decreciente de la coordenada $y$. 

Para actualizar la posición, restamos la cantidad almacenada en la configuración {\ttfamily .velocidad\_balas} de {\ttfamily self.y }. Luego usamos el valor de {\ttfamily self.y} para establecer el valor de {\ttfamily self.rectang.y = self.y}.

La configuración {\ttfamily .velocidad\_balas}  nos permite aumentar la velocidad de las balas a medida que avanza el juego o según sea necesario para refinar el comportamiento del juego.

Una vez que se dispara una bala, nunca cambiamos el valor de su coordenada $x$, por lo que viajará verticalmente en línea recta incluso si la nave se mueve.

Cuando queremos dibujar una bala, llamamos a {\ttfamily dibuja\_balas()}. La función {\ttfamily pygame.draw.rect()} llena la parte de la pantalla definida por el rectángulo de la bala con el color almacenado en {\ttfamily self.color}.

\subsubsection{Almacenamiento de balas en un grupo}

Ahora que tenemos la clase  {\ttfamily Balas()}  y las configuraciones necesarias definidas, podemos escribir código para
disparar una bala cada vez que el jugador presiona la barra espaciadora.

Creamos un grupo en la clase {\ttfamily InvasionAlien} para almacenar todas las balas activas que se han disparado.
Este grupo será una instancia de la clase  {\ttfamily pygame.sprite.Group}, que se comporta
como una lista con algunas funciones adicionales que son útiles al crear juegos. Usaremos este grupo para dibujar bala
en la pantalla en cada paso por el bucle principal y para actualizar la posición de cada bala.\vspace{0.2cm}

Actualizamos el archivo {\ttfamily invasion\_alien.py}. 

Se  importa la clase  { } {\ttfamily Balas()}:

\hspace{-2.4cm}{\ttfamily \scriptsize invasión}

\vspace{-0.2cm}
\hspace{-2.1cm}{\ttfamily \scriptsize \_alien.py}\vspace{-0.84cm}
\begin{lstlisting}		
--recorte-- 
from nave import Nave
from balas import Balas
\end{lstlisting}
\noindent Creamos el grupo que contiene las balas en método {\ttfamily \_\_init\_\_()}:

\hspace{-2.4cm}{\ttfamily \scriptsize invasión}

\vspace{-0.2cm}
\hspace{-2.1cm}{\ttfamily \scriptsize \_alien.py}\vspace{-0.84cm}
\begin{lstlisting}		
--recorte--
    def __init__(self):
         --recorte-- 
        self.nave = Nave(self)
        self.balas = pygame.sprite.Group()
\end{lstlisting}
\noindent Luego necesitamos actualizar la posición de las balas en cada paso el bucle while:

\hspace{-2.4cm}{\ttfamily \scriptsize invasión}

\vspace{-0.2cm}
\hspace{-2.1cm}{\ttfamily \scriptsize \_alien.py}\vspace{-0.84cm}
\begin{lstlisting}		
    def run_jugar(self):
        """Inicia el bucle principal del juego"""
        while True:
            self._check_events()
            self.nave.actualizar()
            self.balas.update()
            self._actualiza_pantalla()
            self.clock.tick(60)
\end{lstlisting}
\noindent Cuando se llama a  {\ttfamily update()} en un grupo, el grupo  actualiza  automáticamente  cada sprite del grupo.  En este caso, la línea {\ttfamily self.balas.update()}, actualiza  automáticamente para cada bala que colocamos en el grupo.

\subsection{Disparar balas}	
	
En la clase  {\ttfamily InvasionAlien}, es necesario modificar el método {\ttfamily \_chequear\\ \_keydown\_eventos()} 
para disparar una bala cuando el jugador presiona la barra espaciadora. Sin embargo, en el método  {\ttfamily \_chequear\_keyup\_eventos()}  no es necesario modificaciones, porque no pasa nada cuando al liberar la barra espaciadora.

Y también  necesitamos   modificar el método  {\ttfamily \_actualizar\_pantalla()} para asegurarse de que cada bala se dibuje en la
pantalla antes de llamar a {\ttfamily flip()}.

Y para disparar una bala se escribe un nuevo método, {\ttfamily \_disparar\_balas()}, que hace el trabajo de disparar balas:


\hspace{-2.4cm}{\ttfamily \scriptsize invasión}

\vspace{-0.2cm}
\hspace{-2.1cm}{\ttfamily \scriptsize \_alien.py}\vspace{-0.84cm}
\begin{lstlisting}
--recorte--
    def _chequear_keydown_eventos(self, event):
        --recorte--
        elif event.key == pygame.K_q:
            pygame.quit()
            sys.exit()
        elif event.key == pygame.K_SPACE:
            self._dispara_bala()
    def _check_keyup_events(self, event):
        --recorte---
    def _dispara_bala(self):
        """Crea una nueva bala y la añade al grupo de balas"""
        nueva_bala = Balas(self)
        self.balas.add(nueva_bala)

    def _actualiza_pantalla(self):
        ""Actualiza inagen sobre pantalla, y  cambia  la nueva pantalla"""
        self.pantalla.fill(self.configuracion.color_fondo)
        for bala in self.balas.sprites():
            bala.dibuja_bala()
    self.nave07.blitme()
    pygame.display.flip()
--recorte--
\end{lstlisting}

Al presiona la barra espaciadora se activa el llamado a {\ttfamily\_dispara\_bala()}. Se crea un objeto  de  la clase {\ttfamily Balas()}, llamado  {\ttfamily nueva\_bala}. Luego lo agregamos al grupo de balas usando el método {\ttfamily  self.balas.add(nueva\_bala)}. El método {\ttfamily add()} es similar a método {\ttfamily append()}, pero está escrito específicamente para grupos de PyGame ({\ttfamily self.balas = pygame.sprite.Group()}).



\begin{figure}[h!]
	\centering
	\includegraphics[scale=0.40]{../A Logos de Python/dispara\_balas}
	\caption{La nave después de disparar una serie de balas.}
	\label{dispara}
\end{figure}


El método {\ttfamily balas.sprites()} devuelve una lista de todos los sprites del grupo de balas. 
Para dibujar todas las balas disparadas en la pantalla, recorremos los sprites en balas y llamamos a draw\_bullet() en cada una 4. Colocamos este bucle antes de la línea que dibuja ella nave, para que las balas no comiencen encima de lala nave.

Cuando ejecutes {\ttfamily invasion\_alien.py} ahora, deberías poder mover la nave hacia la derecha y hacia la izquierda y disparar tantas balas como quieras. Las balas suben por la pantalla y desaparecen cuando llegan a la cima, 
como se muestra en la Figura \ref{dispara}. \\

Como tarea Usted puede modificar el tamaño, el color y la velocidad de las balas en {\ttfamily configurar.py}.



\subsubsection{Eliminar balas fuera del alcance}

Por el momento, las balas desaparecen cuando llegan a la parte superior, pero sólo porque PyGame no puede dibujarlas
por encima de la parte superior de la pantalla. En realidad, las balas siguen existiendo; sus valores de coordenadas $y$ se vuelven cada vez más grandes pero negativas. Esto es un problema porque continúan consumiendo memoria y potencia de procesamiento.

Necesitamos deshacernos de estas viejas balas o el juego se ralentizará haciendo tanto trabajo innecesario.
Para resolver este inconveniente, necesitamos detectar el valor inferior del objeto rectángulo de la bala cuando alcance un valor igual a 0, lo que indica que la bala  pasó  parte superior de la pantalla:

\hspace{-2.4cm}{\ttfamily \scriptsize invasión}

\vspace{-0.2cm}
\hspace{-2.1cm}{\ttfamily \scriptsize \_alien.py}\vspace{-0.84cm}
\begin{lstlisting}
--recorte--
    def run_jugar(self):
    """Inicia el principal bucle del juego"""
    while True:
        self._chequear_eventos()
        self.nave.actualiza()
        self.balas.update()
        # Obtener rid de balas que han disparado.
        for balas in self.balas.copy():
            if balas.rectang.bottom <= 0:
                self.balas.remove(balas)
                print("Una",len(self.balas))
        self._actualiza_pantalla()
        self.clock.tick(60)
\end{lstlisting}

Cuando se usa un bucle {\ttfamily for} con una lista o un grupo en PyGame, Python espera que la lista mantenga 
la misma longitud mientras se ejecute el bucle. Eso significa que no puedes eliminar elementos de una lista o grupo dentro de un bucle  {\ttfamily for}, entonces tenemos que optar por una copia del grupo. Para ello usamos el método {\ttfamily copy()} para configurar
 el bucle  {\ttfamily for}, lo que nos deja libres para modificar el grupo de balas original dentro del bucle. 
Comprobamos cada bala para ver si ha desaparecido en la parte superior de la pantalla. 
Si lo tiene, lo eliminamos de las balas. 
Insertamos un llamada {\ttfamily print()} para mostrar cuántas balas existen actualmente en el juego y verificar
se eliminan cuando llegan a la parte superior de la pantalla.

Si este código funciona correctamente, podemos observar la salida en la terminal mientras disparamos
balas y ver que el número de balas disminuye a cero después de cada serie de balas se ha despejado en la parte superior de la pantalla. 

Después de ejecutar el juego y verifiquar que las bala se eliminen correctamente, elimine el {\ttfamily print()}.
Si lo dejas, el juego se ralentizará significativamente porque requiere más tiempo para escribir la salida en la terminal
que para dibujar gráficos en el ventana de juego.

\subsection{Limitar el número de balas}

En muchos juegos de disparos limitan la cantidad de balas que un jugador puede tener en el
pantalla al mismo tiempo; esta limitación hace que el jugado dispare con precisión. Incorporaremos este restricción   
en Invasión Alienígena.   

Primero, almacenamos la cantidad de balas permitidas en pantalla en la variable {\ttfamily self.balas\_permitidas = 10} dentro del archivo  {\ttfamily configurar.py}:

\hspace{-2.75cm}{\ttfamily \scriptsize configurar.py}  \vspace{-0.5cm}
\begin{lstlisting}	
--recorte--	
        # Configurando Balas
        --recorte--
        self.balas_color = (0, 0, 0)
        self.balas_permitidas = 10
\end{lstlisting}


Esta restricción limita al jugador a diez balas a la vez en la pantalla. Usaremos esta configuración en {\ttfamily InvasionAlien} para verificar cuántas balas existen antes de crear una nueva bala en  {\ttfamily \_disparer\_balas()}:



\hspace{-2.4cm}{\ttfamily \scriptsize invasión}

\vspace{-0.2cm}
\hspace{-2.1cm}{\ttfamily \scriptsize \_alien.py}\vspace{-0.84cm}
\begin{lstlisting}
--recorte--	
    def _dispara_bala(self):
        """Crea una nueva bala y la añade al guupo de balas"""
        if len(self.balas) < self.configuracion.balas_permitidas:
            nueva_bala = Balas(self)
            self.balas.add(nueva_bala)
\end{lstlisting}

Cuando el jugador presiona la barra espaciadora, verificamos la longitud del grupo de balas (número de balas). Si  {\ttfamily len(self.balas)} es menor que diez, creamos una nueva bala. Pero si ya hay diez balas activas, no pasa nada cuando se presiona la barra espaciadora. \\

Al ejecutar el juego, solo deberías poder disparar balas en grupos de diez.


\subsection{Organizando la clase InvasionAlien()}    %  {\ttfamily \_actualizar\_balas()}

Es importante mantener la clase principal {\ttfamily InvasionAlien} razonablemente bien organizada, por lo que ahora que hemos escrito y verificado el código para la administración de balas, podemos moverlo a un método específico para actualizar balas. 

Crearemos un nuevo método llamado  {\ttfamily \_actualizar\_balas()} y lo agregaremos justo antes de  {\ttfamily \_actualizar\_pantalla()}:


\hspace{-2.4cm}{\ttfamily \scriptsize invasión}

\vspace{-0.2cm}
\hspace{-2.1cm}{\ttfamily \scriptsize \_alien.py}\vspace{-0.84cm}
\begin{lstlisting}
--recorte--	
    def _actualizar_balas(self):
        """Actualizar la posicion de balas y obtiene rid 
        de las balas fuera del alcance."""
        # Actualizar posicion de balas.
        self.balas.update()
        # Obtenine rid de balas que han disparado
        for balas  in self.balas.copy():
            if balas.rectang.bottom <= 0:
                self.balas.remove(bala)
    --recorte--	
\end{lstlisting}


El código para {\ttfamily \_actualizar\_balas()} se corta y pega desde {\ttfamily run\_jugar()}; todo lo que hemos hecho aquí 
es aclarar los comentarios. 

Ahora el bucle while en {\ttfamily run\_jugar()} vuelve a quedar  simple y claro:


\hspace{-2.4cm}{\ttfamily \scriptsize invasión}

\vspace{-0.2cm}
\hspace{-2.1cm}{\ttfamily \scriptsize \_alien.py}\vspace{-0.84cm}
\begin{lstlisting}
    def run_jugar(self):
        """Inicia el bucle principal del juego."""
        while True:
            self._chequear_eventos()            
            self.nave07.actualizar()            

            self.balas.update()           
            self._actualizar_balas()
            self._actualizar_pantalla()            

            self.nave07.blitme()
            # Actualiza la pantalla, dibujada la más recientemente
            pygame.display.flip()
            self.clock.tick(60)
\end{lstlisting}


Ahora nuestro bucle principal contiene solo un código mínimo, por lo que podemos leer rápidamente los nombres de los métodos
y comprender lo que sucede en el juego. El bucle principal comprueba la entrada del jugador y luego actualiza la posición de la nave y cualquier bala que haya sido disparada. Luego usamos las posiciones actualizadas para dibujar una nueva pantalla y marcar el reloj al final de cada paso por el bucle.


Al ejecutar {\ttfamily invasion\_alien.py} una vez más y asegúrese de que aún pueda disparar balas sin errores.




\section{INTÉNTALO TÚ MISMO}

12-6. Tirador lateral: escribe un juego que coloque una nave en el lado izquierdo del
pantalla y permite al jugador mover ella nave hacia arriba y hacia abajo. 
Haz que  la nave se dispare  bala que viaja a través de la pantalla cuando el jugador presiona la barra espaciadora.
Asegúrese de que las balas se eliminen una vez que desaparezcan de la pantalla.

\section{Resumen}

%  Enumerar lo aprendido:

En este capítulo, aprendiste a hacer un plan para un juego y aprendiste la estructura básica de un juego escrito en Pygame. 
Aprendiste a establecer un fondo.
Configuraciones de color y organizar en una clase separada donde puedes ajustar  más fácilmente. 
Aprendiste  cómo dibujar una imagen en la pantalla y darle  control al jugador sobre el movimiento de
los elementos del juego.
Creaste elementos que se mueven solos, como balas volando por una pantalla, y eliminaste
objetos que ya no son necesarios. También aprendiste a refactorizar código en un
proyecto de forma regular para facilitar el desarrollo continuo.\\


En el Capítulo \ref{cap13}, agregaremos las naves alienígenas  en el juego Invasión Alienígena. 
Por el final de la capítulo, podrás derribar extraterrestres, con suerte antes de que lleguen a la nave.






\chapter{Alienígenas en el cielo} \label{cap13}

\begin{wrapfigure}{l}{0.5\textwidth} 
	\vspace{-25pt}\begin{center} \includegraphics[width=0.50\textwidth]{../A Logos de Python/Alienigenas}
	\end{center}\vspace{-30.5pt} %\caption{Python 3} 
\end{wrapfigure}
En este capítulo, agregaremos los alienígena o extraterrestres al juego Invasión Alienígena.
Agregaremos los alienígenas cerca de la parte superior de la pantalla y luego generamos una flota completa
de alienígenas. Haremos que la flota avance de lado y hacia abajo, y nos desharemos de cualquier Alienígena que sea alcanzado por una bala.
Finalmente, limitaremos la cantidad de la naves que tiene un jugador y finalizaremos el juego cuando el jugador se quede sin naves.\\



A medida que avance en este capítulo, aprenderá más sobre Pygame y sobre la gestión de un gran proyecto. También aprenderás a detectar colisiones entre los objetos del juego, como balas y extraterrestres. Detectar colisiones te ayuda a define interacciones 
entre elementos en los juegos. 

Por ejemplo, puedes confinar un personaje dentro de las paredes de un laberinto o pasar una pelota entre dos paredes. Continuaremos trabajando a partir de un plan que revisamos ocasionalmente para mantener el enfoque de nuestras sesiones de escritura de códigos.

Antes de comenzar a escribir código nuevo para agregar una flota de alienígenas a la pantalla, echemos un vistazo al proyecto y actualicemos nuestro plan.



\section{Revisar el proyecto}


Cuando estás comenzando una nueva fase de desarrollo en un proyecto grande, siempre es buena idea revisar su plan y aclarar lo que quiere
lograr con el código que estás a punto de escribir. 

En este capítulo haremos la siguiente:\\

• Agregamos un único alienígena en la esquina superior izquierda de la pantalla, con espacio suficiente  a su alrededor.

• Llena la parte superior de la pantalla con tantos alienígenas como podamos colocar  horizontalmente. Luego crearemos filas adicionales de alienígenas hasta que tengamos una flota completa.

• Haremos que la flota se mueva hacia los lados y hacia abajo hasta que toda la flota llegue abajo. Si toda la flota es derribada, crearemos una nueva flota.  Si un extraterrestre golpea la nave o llega al suelo destruiremos su nave terrestre y finaliza esta fase del juego.  

• Se limita el número de la naves que el jugador puede usar y finaliza el juego cuando el jugador haya agotado el número de naves asignadas.\\ 


Perfeccionaremos este plan a medida que implementemos funciones, pero esto es lo específico y suficiente para empezar a escribir código.
También debe revisar su código existente cuando comience a trabajar en un nueva serie de características para  proyecto. 
Porque cada nueva fase normalmente hace un proyecto más complejo, es mejor limpiar cualquier desordenado o 
código ineficiente. Hemos estado refactorizando a medida que avanzamos, por lo que no necesitamos ningún código
para refactorizar en este punto.

\section{El primer alienígena}

Dibujar un alienígena en la pantalla es un proceso semejante a de dibujar la nave en la pantalla.  El comportamiento de 
los alienígenas será controlado por una clase llamada {\ttfamily Alien()}, que estructuraremos como la clase {\ttfamily Nave()}. 

Continuaremos usando imágenes de mapa de bits para simplificar (.bmp). Puedes encontrar tu propia imagen para un 
extraterrestre o usar la que se muestra en la Figura \ref{alienigena}.

\begin{figure}[h!]
	\centering
	\includegraphics[scale=0.40]{../A Logos de Python/alienigena}
	\caption{El extraterrestre que usaremos.}
	\label{alienigena}
\end{figure}
que está disponible en los recursos del libro en  la web linis.  % https://ehmatthes.github.io/pcc\_3e.
Esta imagen tiene un fondo gris, que coincide con el color de fondo de pantalla. 
Asegúrese de guardar el archivo de imagen que elija en la carpeta de imágenes.

\subsection{La clase alienígena}

Ahora escribiremos la clase {\ttfamily Alien} y la guardaremos como {\ttfamily alien.py}:

\hspace{-1.6cm}{\ttfamily \scriptsize alien.py}  \vspace{-0.5cm}
\begin{lstlisting}		
import pygame
from pygame.sprite import Sprite
class Alien(Sprite):
    """Una clase para representar un alienígena del aflota."""
    def __init__(self, ia_jugar):
        """Inicializa el alienígena y se configura su posicion
        inicial."""
        super().__init__()
        self.pantalla = ia_jugar.pantalla
        # carga la imagen y configura el atributo rectang
        self.imagen = pygame.image.load('imagenes/alien.bmp')
        self.rectang = self.imagen.get_rect()
        # Cada alienígena aparece cerca de la esquina superior izquierda
        # de la pantalla
        self.rectang.x = self.rectang.width/3
        self.rectang.y = self.rectang.height/3
        # Almacena el alienígena en la posicion exacta
        self.x = float(self.rectang.x)
\end{lstlisting}

La mayor parte de esta clase es como la clase de la nave, excepto por la ubicación del alienígena en la pantalla. 
Inicialmente dibujamos a cada alienígena cerca de la esquina superior izquierda de la pantalla;  agregamos un espacio 
a la izquierda que es igual al ancho del alienígena y un espacio por encima es igual a su altura, por lo que es fácil de ver. 
Nos preocupamos principalmente la velocidad horizontal de los extraterrestres, por lo que rastrearemos la posición horizontal de cada uno
extraterrestre precisamente.

Esta clase Alien no necesita un método para dibujarla en la pantalla; en su lugar, usaremos un método de grupo de Pygame que dibuja automáticamente todos los elementos de un grupo a la pantalla.

\subsection{Una instancia de la clase Alien}

Debemos crear una instancia de Alien para que podamos ver el primer alienígena en la pantalla. Debido a que es parte de nuestro trabajo de configuración, agregaremos el código para la instancia al final del método {\ttfamily \_\_init\_\_()} en la clase {\ttfamily InvasionAlien}. 

Luego debemos crear una flota completa de alienígenas, lo que supondrá bastante trabajo, así que para ello 
creamos  un nuevo método auxiliar llamado  {\ttfamily \_crear\_flota()}.

El orden de los métodos en una clase no importa, siempre que haya algo coherencia en la forma en que se colocan.
Colocaré {\ttfamily \_crear\_flota()} justo antes del método {\ttfamily \_actualizar\_pantalla()}, pero en cualquier
lugar funcionará. \vspace{0.3cm}


\noindent Actualizamos los cambios en  {\ttfamily invasión\_alien.py}:

\noindent Importamos la clase {\ttfamily Alien()}:


\hspace{-2.4cm}{\ttfamily \scriptsize invasión}

\vspace{-0.2cm}
\hspace{-2.1cm}{\ttfamily \scriptsize \_alien.py}\vspace{-0.84cm}
\begin{lstlisting}
from balas import Balas
from alien import Alien	
\end{lstlisting}
\noindent Y actualizamos el método {\ttfamily \_\_init\_\_()}:

\hspace{-2.4cm}{\ttfamily \scriptsize invasión}

\vspace{-0.2cm}
\hspace{-2.1cm}{\ttfamily \scriptsize \_alien.py}\vspace{-0.84cm}
\begin{lstlisting}
--recorte--			
    def __init__(self):
    --recorte--
        self.nave = Nave(self)
        self.balas = pygame.sprite.Group()
        self.aliens = pygame.sprite.Group()
        self._crear_flota()
\end{lstlisting}
\noindent Creamos un grupo para mantener la flota de extraterrestres y llamamos al  método  nuevo  {\ttfamily \_crear\_flota()}:

\hspace{-2.4cm}{\ttfamily \scriptsize invasión}

\vspace{-0.2cm}
\hspace{-2.1cm}{\ttfamily \scriptsize \_alien.py}\vspace{-0.84cm}
\begin{lstlisting}		
--recorte--
    def _crear_flota()(self):
        """Crear la  flota de alienígena."""
        # Hacer un alienígena
        alien = Alien(self)
        self.aliens.add(alien)
\end{lstlisting}
En este método, creamos una instancia de la clase  {\ttfamily Alien} y luego lo agregamos al grupo que contendrá la flota. 
El extraterrestre se colocará en el área superior izquierda predeterminada de la pantalla. 
Para que aparezca el alienígena, necesitamos llamar al método {\ttfamily draw()} del grupo en {\ttfamily \_actualizar\_pantalla()}:

\hspace{-2.4cm}{\ttfamily \scriptsize invasión}

\vspace{-0.2cm}
\hspace{-2.1cm}{\ttfamily \scriptsize \_alien.py}\vspace{-0.84cm}
\begin{lstlisting}		
    def _actualizar_pantalla(self):
        --recorte--
        self.nave.blitme()
        self.aliens.draw(self.pantalla)
        pygame.display.flip()
\end{lstlisting}


Al  llamar al  método {\ttfamily draw()} en un grupo, Pygame dibuja cada elemento del grupo en la posición definida por su atributo {\ttfamily rect}. El método  {\ttfamily draw()}  requiere un argumento: una superficie sobre la cual dibujar los elementos del grupo. 
{\ttfamily self.aliens.draw(self.pantalla)}.


La Figura \ref{aliePantall} muestra el primer alienígena en la pantalla.

\begin{figure}[h!]
	\centering
	\includegraphics[scale=0.35]{../A Logos de Python/alien\_pantalla}
	\caption{Aparece el primer alienígena\\ en la pantalla.}
	\label{aliePantall}
\end{figure}


Ahora que el primer alienígena aparece correctamente, escribiremos el código para dibujar
toda una flota.




\subsection{Construyendo la flota alienígena}

Para dibujar una flota, necesitamos descubrir cómo llenar la parte superior de la pantalla con extraterrestres,
sin saturar la ventana del juego. Hay un número de maneras para lograr este objetivo. Lo abordaremos agregando extraterrestres a través
la parte superior de la pantalla, hasta que no quede espacio para un nuevo extraterrestre. Repitiendo el proceso, siempre y cuando tengamos suficiente espacio vertical para agregar una nueva fila.


\subsubsection{Una fila de extraterrestres}

Ahora estamos listos para generar una fila completa de extraterrestres. Para crear una fila completa,
primero haremos un solo alienígena que tengamos acceso al ancho del alienígena,  colocaremos el extraterrestre en el lado izquierdo de la pantalla y luego seguir agregando extraterrestres hasta  quedarse sin espacio:

Obtenemos el ancho del alienígena del primer alienígena que creamos y luego definimos
una variable llamada {\ttfamily actual\_x}. Esto se refiere a la posición horizontal del
próximo extraterrestre que pretendemos colocar en la pantalla. Inicialmente configuramos esto para el ancho de 
un alienígena, para luego desplazar del primer alienígena de la flota desde el borde izquierdo de la pantalla.


\hspace{-2.4cm}{\ttfamily \scriptsize invasión}

\vspace{-0.2cm}
\hspace{-2.1cm}{\ttfamily \scriptsize \_alien.py}\vspace{-0.84cm}
\begin{lstlisting}
    def _crear_flota(self):
        """Crear la  flota de alienígena."""
        # Hacer un alienígena
        # Espacio entre alienigenas igual un tercio del ancho del alienigena
        alien = Alien(self)
        ancho_alien = alien.rect.width
        actual_x = alien_width
        while actual_x < (self.configuracion.pantalla_width - 
                             1.3 * ancho_alien):       
            nuevo_alien = Alien(self)
            nuevo_alien.x = actual_x
            nuevo_alien.rect.x = actual_x
            self.aliens.add(nuevo_alien)
            actual_x += 1.5 * ancho_alien
\end{lstlisting}



A continuación, comenzamos el ciclo while; vamos a seguir agregando extraterrestres
mientras haya suficiente espacio para colocar uno. Para determinar si hay espacio
para colocar otro alienígena, compararemos {\ttfamily  while actual\_x < (self.configuracion.ancho\_pantalla} con algún valor máximo.  El primer intento de definir este bucle podría verse así:

%\hspace{-2.5cm}{\ttfamily \scriptsize curso.py}  \vspace{-0.5cm}
\begin{lstlisting}	
while actual_x < self.configuracion.ancho_pantalla:	
\end{lstlisting}


Esto parece que podría funcionar, pero colocaría al alienígena final de la fila en el extremo derecho de la pantalla.
Entonces agregamos un pequeño margen en el lado derecho de la pantalla. Siempre y cuando haya al menos dos anchos alienígenas
de espacio en el borde derecho de la pantalla, ingresamos al bucle y agregamos otro  a la flota.





Siempre que haya suficiente espacio horizontal para continuar el bucle, queremos hacer dos cosas: crear un extraterrestre en la posición correcta y definir la posición horizontal del siguiente extraterrestre en la fila. Creamos un alienígena y lo asignamos a {\ttfamily nuevo\_alien}. 

Luego establecemos la posición horizontal precisa en el valor actual de {\ttfamily actual\_x}. También posicionamos el {\ttfamily rect} del alienígena en este mismo valor $x$ y agregamos el nuevo alienígena al grupo  {\ttfamily self.aliens} .

Finalmente, incrementamos el valor de  {\ttfamily actual\_x}. Agregamos dos anchos de alienígenas a la posición horizontal, para pasar el alienígena que acabamos de agregar y dejar algo de espacio entre los alienígenas también.

Python reevaluará la condición al comienzo del ciclo while y decidirá si hay espacio para otro extraterrestre. Cuando no quede espacio, el bucle terminará y deberíamos tener una fila completa de alienígenas.


Al ejecutar {\ttfamily invasión\_alien.py}, deberías ver aparecer la primera fila de alienígenas, como en la Figura \ref{flota01}.


\begin{figure}[h!]
	\centering
	\includegraphics[scale=0.35]{../A Logos de Python/flota01}
	\caption{Aparece la primera fila de \\ extraterrestres en pantalla.}
	\label{flota01}
\end{figure}





NOTA 


\vspace{0.2cm}
\begin{minipage}[t]{0.2\textwidth}
	\centering\raisebox{\dimexpr\topskip-\height}{%
		\includegraphics[scale=0.20]{../A Logos de Python/Sugerencia}}  %[width=\textwidth]
	
\end{minipage}%\hfill
\begin{minipage}[t]{0.70\textwidth} 
	{\scriptsize No debe ser motivo de para desanimarse al construir un bucle como el que se muestra en esta sección, y colocar 
		demasiado alienígenas a la derecha.  Una cosa buena de la programación es que sus ideas iniciales sobre cómo abordar
		un problema como este no tienen por qué ser correctas. Es perfectamente razonable iniciar un bucle como este con los alienígenas colocados demasiado a la derecha y luego modificar el bucle hasta que tengas una cantidad adecuada de espacio en la pantalla.}
\end{minipage}
\vspace{0.2cm}


\subsubsection{Refactorización del método  \_crear\_flota()}

Si el código que hemos escrito hasta ahora fuera todo lo que necesitábamos para crear una flota, probablemente dejaríamos
 {\ttfamily\_crear\_flota()} como está. Pero tenemos más trabajo por hacer, así que limpiemos un poco el método. Agregaremos un nuevo método auxiliar,  {\ttfamily \_crear\_alien()}, y lo llamaremos desde {\ttfamily\_crear\_flota()}:

\hspace{-2.4cm}{\ttfamily \scriptsize invasión}

\vspace{-0.2cm}
\hspace{-2.1cm}{\ttfamily \scriptsize \_alien.py}\vspace{-0.84cm}
\begin{lstlisting}
    def _creare_flota(self): 
        --recorte--
        while actual_x < (self.configuracion.ancho_pantalla - 
                              1.3 * ancho_alien):
            self._crear_alien(actual_x)
            actual_x+= 2 * ancho_alien
            
    def _crear_alien(self, x_position):
        """Crear un alienigena y colocar en la fila."""
        nuevo_alien = Alien(self)
        nuevo_alien.x = x_posicion
        nuevo_alien.rect.x = x_posicion
        self.aliens.add(nuevo_alien)
	
\end{lstlisting}

El método {\ttfamily \_crear\_alien()} requiere otro parámetro además de {\ttfamily self}: el valor $x$ que especifica dónde se 
debe colocar el extraterrestre. El código en el método {\ttfamily \_crear\_alien()}  es el mismo código que estaba 
en  {\ttfamily\_crear\_flota()}: excepto que usamos el parámetro  {\ttfamily x\_posicion} en lugar de  
{\ttfamily actual\_x}. La refactorización facilitará la adición de nuevas filas y la creación de una flota completa.

\subsubsection{Agregar filas de alienígenas}

Para terminar la flota, seguiremos agregando más filas hasta que nos quedemos sin espacio.

Usaremos un bucle anidado; envolveremos otro bucle  {\ttfamily while} alrededor del bucle actual.
El bucle interior colocará a los extraterrestres horizontalmente en una fila enfocándose en
los valores $x$ de los extraterrestres. 

El bucle exterior colocará a los extraterrestres verticalmente enfocándose en los valores de $y$. 
Dejaremos de agregar filas cuando nos acerquemos al final del pantalla, dejando suficiente espacio para la nave y algo de espacio para comenzar a disparar los alienígenas.

Aquí se explica cómo anidar los dos bucles while en  {\ttfamily\_crear\_flota()}:



%\hspace{-2.5cm}{\ttfamily \scriptsize curso.py}  \vspace{-0.5cm}
\begin{lstlisting}	
    def _crear_flota(self):
        """Create the fleet of aliens."""
        # Crear, guardar y anadir alienigenas hasta llenar la pantalla
        # Espacio entre alienigenas es igual al ancho y alto del alienigena.
        alien = Alien(self)
        ancho_alien, alto_alien = alien.rect.size
        actual_x, actual_y =  ancho_alien, alto_alien

        while actual_y < (self.configuracion.pantalla_height - 
                             6 * alto_alien):
            while actual_x <(self.configuracion.pantalla_width - 
                                1.3 * ancho_alien):
                self._crear_alien(actual_x, actual_y)
                actual_x += 1.3 * ancho_alien

            # finaliza una fila; resetea el valor de $x$,
            # e incrementa el valor de $y$.
            actual_x = ancho_alien
            actual_y += 1.5 * alto_alien 
\end{lstlisting}


Necesitaremos saber la altura del alienígena para poder colocar filas, así que tomamos
el ancho y alto del extraterrestre usando el atributo de tamaño de un extraterrestre recto 1. A

El atributo de tamaño de {\ttfamily rect} es una tupla que contiene su ancho y alto.

A continuación, establecemos los valores iniciales  $x$ e $y$   para la ubicación del primer alienígena.
en la flota. Lo colocamos a un ancho alienígena desde la izquierda y un alienígena
altura hacia abajo desde la parte superior. 

Luego definimos el bucle while que controla cómo Se colocan muchas filas en la pantalla. Siempre que el valor de y para el siguiente
fila es menor que la altura de la pantalla, menos tres alturas alienígenas, mantendremos añadiendo filas. (Si esto no deja la cantidad adecuada de espacio, podemos ajústelo más tarde.)





Llamamos a {\ttfamily \_crear\_alien()} y le pasamos el valor y así como su posición x 4. Modificaremos  {\ttfamily \_crear\_alien()} en un momento.
Observe la sangría de las dos últimas líneas del código 5. Están dentro del bucle while externo, pero fuera del bucle while interno. Este bloque se ejecuta una vez finalizado el bucle interno; se ejecuta una vez después de crear cada fila. Después de agregar cada fila, restablecemos el valor de {\ttfamily actual\_x} para que el primer alienígena de la siguiente fila se coloque en la misma posición que el primer alienígena de las filas anteriores. Luego agregamos dos alturas alienígenas al valor actual de {\ttfamily actual\_y}, por lo que la siguiente fila se colocará más abajo en la pantalla. La sangría es realmente importante aquí; si no ves la flota correcta cuando ejecutas {\ttfamily invasion\_alien.py}

Al final de esta sección, verifique la sangría de todas las líneas en estos bucles anidados.

Necesitamos modificar  {\ttfamily \_crear\_alien()} para establecer correctamente la posición vertical del extraterrestre:




\hspace{-2.5cm}{\ttfamily \scriptsize curso.py}  \vspace{-0.5cm}
\begin{lstlisting}	
    def _crear_alien(self, x_posicion, y_posicion):
        """Crear un alienigena y colocar en la flota."""
        nuevo_alien = Alien(self)
        nuevo_alien.x = x_posicion
        nuevo_alien.rect.x = x_posicion
        nuevo_alien.rect.y = y_posicion
        self.aliens.add(nuevo_alien)
\end{lstlisting}


Modificamos la definición del método para aceptar el valor y para el nuevo alienígena y configuramos la posición vertical de rect en el cuerpo del método.


Cuando ejecutes el juego ahora, deberías ver una flota completa de alienígenas, como se muestra en la Figura \ref{flotaComplet}.



\begin{figure}[h!]
	\centering
	\includegraphics[scale=0.35]{../A Logos de Python/flotaComplet}
	\caption{Aparece la flota completa.}
	\label{flotaComplet}
\end{figure}


En la siguiente sección, haremos que la flota se mueva.


\section{INTÉNTALO TÚ MISMO}

13-1. Estrellas: Encuentra una imagen de una estrella. Haz que aparezca una cuadrícula de estrellas en la pantalla.

13-2. Mejores estrellas: puedes crear un patrón de estrellas más realista introduciendo
aleatoriedad al colocar cada estrella. Recuerde del Capítulo 9 que puede obtener
un número aleatorio como este:


\section{La flota se mueve}

Ahora hagamos que la flota de alienígenas se mueva hacia la derecha a través de la pantalla hasta
golpea el borde, y luego del choque se hace caer una cantidad determinada y se mueve en el otra
dirección. Continuaremos este movimiento hasta que todos los alienígenas hayan sido fusilados,
uno choca con la nave o llegan al final de la pantalla.\\

Comencemos haciendo que la flota se mueva hacia la derecha.

\subsubsection{Los extraterrestres se mueven hacia la derecha}

Para mover a los alienígenas, usaremos un método {\ttfamily actualizar()} en {\ttfamily alien.py}, que
Llama a cada extraterrestre del grupo de extraterrestres. Primero, agregue una configuración para controlar el
velocidad de cada alienígena:


\hspace{-2.75cm}{\ttfamily \scriptsize configurar.py}  \vspace{-0.5cm}
\begin{lstlisting}		
    def __init__(self):
        --recorte--
        # Configurancion Alienigena
        self.velocidad_alien = 1.0
\end{lstlisting}

Luego usaremos esta configuración en el método {\ttfamily actualizar()} del módulo {\ttfamily alien.py}:


\hspace{-2.5cm}{\ttfamily \scriptsize alien.py}  \vspace{-0.5cm}
\begin{lstlisting}		
    def __init__(self, ia_jugar):
        """Inicializa el alienígena y configuramos su posicion inicial."""
         super().__init__() 
         --recorte--
         
    def actualizar(self):
        """Mover el alienígena hacia la derecha."""
        self.x += self.configurar.velocidad_alien
        self.rect.x = self.x
\end{lstlisting}




Creamos un parámetro de configuración en {\ttfamily \_\_init\_\_()} para que podamos acceder a la 
velocidad  de alienígena  en el método {\ttfamily actualizar()}. Cada vez que actualizamos la posición de un alienígena, 
lo movemos la cantidad almacenada en {\ttfamily velocidad\_alien}. Realizamos un seguimiento de la posición exacta del alienígena con el atributo {\ttfamily self.x}, que puede contener valores flotantes. Luego usamos el valor de {\ttfamily self.x} para actualizar la posición del {\ttfamily rect} del alienígena. 

En el ciclo while principal, ya tenemos llamadas para actualizar las posiciones de la nave y las balas. Ahora también agregaremos una llamada para actualizar la posición de cada alienígena:

\hspace{-2.4cm}{\ttfamily \scriptsize invasión}

\vspace{-0.2cm}
\hspace{-2.1cm}{\ttfamily \scriptsize \_alien.py}\vspace{-0.84cm}
\begin{lstlisting}	
        while True:
            self._chequear_eventos()            
            self.nave.actualizar()            

            self.balas.update()           
            self._actualizar_balas()
            self._actualizar_aliens()
            self._actualizar_pantalla()            

            self.nave.blitme()
            # Actualiza. Hace visible la pantalla dibujada más recientemente
            pygame.display.flip()
            self.clock.tick(60)
\end{lstlisting}

Estamos a punto de escribir un código para gestionar el movimiento de la flota, por lo que creamos un nuevo método llamado {\ttfamily actualizar\_alien()}. 
Actualizamos las posiciones de los alienígenas después de que se hayan actualizado las balas, porque pronto comprobaremos si alguna bala alcanzó a algún alienígena.

El lugar donde se coloque este método en el módulo no es crítico. Pero para mantener el código organizado y claro, lo colocaré justo después de {\ttfamily actualizar\_balas()} para que coincida con el orden de las llamadas a métodos en el bucle  {\ttfamily while}. Aquí está la primera versión de {\ttfamily actualizar\_alien()}:

\hspace{-2.4cm}{\ttfamily \scriptsize invasión}

\vspace{-0.2cm}
\hspace{-2.1cm}{\ttfamily \scriptsize \_alien.py}\vspace{-0.84cm}
\begin{lstlisting}
    def _actualizar_aliens(self):
        """Actualiza la posiciones de todos las alienígena de la flota."""
        self.aliens.update()
\end{lstlisting}


Usamos el método {\ttfamily update()}  en el grupo de extraterrestres, que llama al método {\ttfamily update()} de cada extraterrestre. 

Al  ejecutar {\ttfamily Invasión\_Alien} , deberías ver que la flota se mueve hacia la derecha y desaparece del costado de la pantalla.

\subsubsection{Crear configuraciones para la dirección de la flota}
Ahora crearemos la configuración que hará que la flota se mueva hacia abajo y hacia la izquierda cuando llegue al borde derecho de la pantalla. 


Agregamos la configuración {\ttfamily flota\_cae\_velocidad} que controla la rapidez con la que la flota desciende de la pantalla cada vez que un extraterrestre llega a cualquiera de los bordes. 

\hspace{-2.5cm}{\ttfamily \scriptsize alien.py}  \vspace{-0.5cm}
\begin{lstlisting}        
    # Configurancion Alienigena
    self.velocidad_alien = 1.0 
    self.flota_cae_velocidad = 10   # velocidad_flota_cae
    # Dirección de flota: 1 representa derecha;
    # -1 representa izquierda.
    self.direccion_flota = 1  
\end{lstlisting}

Es útil separar esta velocidad de la velocidad horizontal de los alienígenas para que puedas ajustar las dos velocidades de forma independiente.

Para implementar la configuración de {\ttfamily direccion\_flota}, podríamos usar un valor tipo texto como 'izquierda' o 'derecha',
pero terminaríamos con declaraciones {\ttfamily if - elif} que prueban la dirección de la flota. 
En cambio, debido a que solo tenemos  dos direcciones con las que lidiar, usemos los valores 1 y -1,  cambiemos entre ellos cada vez que la flota cambia de dirección. 
(Usar números también tiene sentido porque moverse hacia la derecha implica agregar al valor de la coordenada $x$ de cada alienígena, y moverse hacia la izquierda implica restando del valor de la coordenada $x$ de cada extraterrestre.)

\subsubsection{Alienígena ha llegado al borde}

Necesitamos un método para comprobar si hay un extraterrestre en alguno de los bordes, y necesitamos
modificar  {\ttfamily update()} para permitir que cada alienígena se mueva en la dirección apropiada.

Este código es parte de la clase {\ttfamily Alien}:



\hspace{-1.6cm}{\ttfamily \scriptsize alien.py}  \vspace{-0.5cm}
\begin{lstlisting}	
    def chequear_bordes(self):
        """Returna True si el alienígena esta en el borde de la pantalla (ventana)."""
        pantalla_rect = self.pantalla.get_rect()
        return (self.rect.right >= pantalla_rect.right) or (self.rect.left <= 0)
    
    def update(self):
        """Mueve el alienígena a la izquierda o derecha."""
        self.x += (self.configuracion.velocidad_alien * 
                    self.configuracion.direccion_flota)
        self.rect.x = self.x
\end{lstlisting}


Podemos llamar al nuevo método {\ttfamily chequear\_bordes()} en cualquier extraterrestre para ver si
está en el borde izquierdo o derecho. 
El alienígena está en el borde derecho si el atributo del rectángulo de su {\ttfamily rect} es mayor o igual que el atributo derecho {\ttfamily rect}  de la pantalla. Está en el borde izquierdo si su valor izquierdo es menor o igual a 0.

Más bien, en lugar de poner esta prueba condicional en un bloque {\ttfamily if}, ponemos la prueba directamente en el
declaración de devolución. 

Este método devolverá {\ttfamily True} si el alienígena está a la derecha o borde izquierdo y {\ttfamily False} si no está en ninguno de los bordes.

Modificamos el método  {\ttfamily actualizar()} para permitir movimiento hacia la izquierda o hacia la derecha 
multiplicando la velocidad del alienígena por el valor de {\ttfamily direccion\_flota}. 
Si {\ttfamily direccion\_flota} es 1, el valor de  {\ttfamily  velocidad\_alien} se agregará al valor posición actual
del alienígena, moviendo al extraterrestre hacia la derecha; si {\ttfamily direccion\_flota} es -1, el valor se le restará
de la posición delalienígena, moviéndolo hacia la izquierda.

\subsubsection{Abandonar la flota y cambiar de dirección}

Cuando un extraterrestre llega al borde, toda la flota debe descender y
cambia la dirección. Por lo tanto, necesitamos agregar algo de código a la clase {\ttfamily InvasionAlien} 
porque ahí es donde comprobaremos si hay extraterrestres a la izquierda o a la derecha del borde. 

Haremos que esto suceda escribiendo los métodos {\ttfamily \_chequear\_borde\_flota()}
y {\ttfamily \_cambiar\_direccion\_flota()}, y luego modificando {\ttfamily \_actualizar\_aliens()}.

Estos nuevos métodos lo ubicare  después de {\ttfamily \_crear\_alien()}, pero nuevamente, la ubicación de estos
dos métodos en la clase no son críticos.


\hspace{-2.4cm}{\ttfamily \scriptsize invasión}

\vspace{-0.2cm}
\hspace{-2.1cm}{\ttfamily \scriptsize \_alien.py}\vspace{-0.84cm}
\begin{lstlisting}
    def _chequear_borde_flota(self):
        """Responde appropriadamente si cualquier
        alienigena alcanza el borde."""
        for alien in self.aliens.sprites():
            if alien.chequear_borde():
                self._cambiar_direccion_flota()
                break
    def _cambiar_direccion_flota(self):
        """Cae la  flota y cambia de direccion."""
        for alien in self.aliens.sprites():
            alien.rect.y += self.configuracion.flota_cae_velocidad
            self.configuracion.flota_cae_velocidad *= -1
\end{lstlisting}

En {\ttfamily  \_chequear\_borde\_flota()}, recorremos la flota y llamamos a {\ttfamily chequear\_borde()} en cada alienígena.
Si {\ttfamily chequear\_borde()} devuelve  {\ttfamily True},  sabemos que un alienígena está en un borde y toda la flota necesita cambiar de dirección; entonces llamamos a {\ttfamily \_cambiar\_direccion\_flota()} y salimos del bucle.
En {\ttfamily \_cambiar\_direccion\_flota()}, recorremos todos los alienígenas y soltamos a cada uno usando la configuración {\ttfamily flota\_cae\_velocidad}; luego cambiamos el valor de {\ttfamily \_cambiar\_direccion\_flota()} multiplicando su valor actual por  $-1$. 
La línea que cambia la dirección de la flota no forma parte del bucle {\ttfamily for}. 
Queremos cambiar la posición vertical de cada  alienígena, pero sólo queremos cambiar la dirección de la flota una vez.

Veamos los cambios en el método {\ttfamily \_actualizar\_aliens()}:

\hspace{-2.4cm}{\ttfamily \scriptsize invasión}

\vspace{-0.2cm}
\hspace{-2.1cm}{\ttfamily \scriptsize \_alien.py}\vspace{-0.84cm}
\begin{lstlisting}
    def _actualizar_aliens(self):
        """Actualiza la posiciones de todos los alienígenas de la flota."""
        self._chequear_borde_flota
        self.aliens.update()
\end{lstlisting}


Modificamos el método llamando a {\ttfamily  \_chequear\_borde\_flota()} antes de actualizar la posición de cada alienígena.


Cuando ejecutas el juego ahora, la flota debería moverse hacia adelante y hacia atrás entre los bordes de la pantalla y descender cada vez que llegue a un borde. Ahora podemos empezar a derribar alienígenas y estar atentos a cualquier alienígena que golpee la nave o llegue al final de la pantalla.

\section{INTÉNTALO TÚ MISMO}

13-3. Gotas de lluvia: busque una imagen de una gota de lluvia y cree una cuadrícula de gotas de lluvia. Haz que las gotas de lluvia caigan hacia la parte inferior de la pantalla hasta que desaparezcan.

13-4. Lluvia constante: Modifique su código en el Ejercicio 13-3 para que cuando una fila de gotas de lluvia desaparezca de la parte inferior de la pantalla, aparezca una nueva fila en la parte superior de la pantalla y comience a caer.


\subsubsection{Disparar extraterrestres}


Hemos construido nuestra nave y una flota de extraterrestres, pero cuando las balas alcanzan a los extraterrestres, simplemente pasan porque no estamos comprobando si hay colisiones. En la programación de juegos, las colisiones ocurren cuando los elementos del juego se superponen. Para hacer que las balas derriben a los extraterrestres, usaremos la función {\ttfamily sprite.groupcollide()} para buscar colisiones entre miembros de dos grupos.

\section{Detección de colisiones de alienígenas con las balas}

Necesitamos saber cuándo una bala alcanza a un alienígenas para poder tomar la decisión de desaparece el  alienígena tan 
pronto como sea golpeado por una bala (alcanzado). Para hacer esto, buscaremos colisiones inmediatamente después de actualizar 
la posición de todas las balas.

La función  {\ttfamily sprite.groupcollide()} compara los rectángulos ({\ttfamily objetos Rect}) de cada elemento en un grupo con los rectángulos de cada elemento del otro grupo. 
En este caso se compara el rectángulo de cada bala con el rectángulo de cada alienígena y devuelve un diccionario que contiene
las balas y los alienígenas que han chocado. Cada clave en el diccionario será una bala, y el valor correspondiente será el alienígena que fue alcanzado. (Bueno, también utilizaremos este diccionario cuando implementemos un sistema de puntuación en el Capítulo \ref{cap14}.)


Agregue el siguiente código al final de {\ttfamily \_actualizar\_balas()} para verificar si hay colisiones
entre balas y extraterrestres:

\hspace{-2.4cm}{\ttfamily \scriptsize invasión}

\vspace{-0.2cm}
\hspace{-2.1cm}{\ttfamily \scriptsize \_alien.py}\vspace{-0.84cm}
\begin{lstlisting}
    def _actualizar_balas()(self):
        """Actualizamos posicion de balas y o btenemos rid de la balas viejas."""
        --recorte--
        # Chequea para cuaquier bala que tiene hit alienigena.
        # Si es asi, obtiene  rid de bala y alienigena.
        colisiones = pygame.sprite.groupcollide(
                         self.balas, self.aliens, True, True)
\end{lstlisting}

El nuevo código que agregamos compara las posiciones de todas las balas  en
{\ttfamily self.balas} y todos los alienígenas en {\ttfamily  self.aliens}, e identifica aquellos que se superponen.
Siempre que los rectángulos  de una bala y un alienígena se superponen, {\ttfamily  groupcollide()} agrega un valor clave
par al diccionario que devuelve. Los dos argumentos verdaderos le indican a PyGame  borrar las balas y los alienígenas que han chocado. 
(Para hacer un potente bala que puede viajar a la parte superior de la pantalla, destruyendo a todos los alienígenas en su
ruta, puede establecer el primer argumento booleano en Falso y mantener el segundo Argumento booleano establecido en Verdadero. 
Los alienígenas impactados desaparecerían, pero todas las balas permanecería activo hasta que desaparecieran de la parte superior de la pantalla).

Cuando ejecutes Alien Invasion ahora, los alienígenas que golpees deberían desaparecer.



\begin{figure}[h!]
	\centering
	\includegraphics[scale=0.35]{../A Logos de Python/eliminaParcial}
	\caption{La flota ha sido parcialmente derribada.}
	\label{eliminaParcial}
\end{figure}

La Figura \ref{eliminaParcial} muestra un flota que  ha sido parcialmente derribada .


\section{Hacer balas más poderosas para realizar pruebas}

Puedes probar muchas funciones de Alien Invasion simplemente ejecutando el juego, pero algunas funciones son tediosas
de probar en la versión normal del juego. 
Por ejemplo, es mucho trabajo derribar a cada alienígena en la pantalla varias veces
para probar si su código responde correctamente a una flota vacía.

Para probar funciones particulares, puedes cambiar ciertas configuraciones del juego para enfocarte en un área particular. 
Por ejemplo, puedes reducir la pantalla para que haya menos alienígenas a los que derribar o aumentar la velocidad de las balas y recibir muchas balas a la vez.

Mi cambio favorito para probar Alien Invasion es usar balas muy anchas que permanecen activas incluso después de haber golpeado a un alienígena (ver Figura \ref{poderBalas}). Intenta configurar {\ttfamily ancho\_Balas} en 100, o incluso en 300, para ver qué tan rápido puedes derribar la flota\\


\begin{figure}[h!]
	\centering
	\includegraphics[scale=0.35]{../A Logos de Python/balaspoder}
	\caption{Balas más poderosas.}
	\label{poderBalas}
\end{figure}

Las balas  más poderosas hacen algunos aspectos del juego más fácil de ganar (Figura \ref{poderBalas}).

Cambios como estos te ayudarán a probar el juego de manera más eficiente y posiblemente generarán ideas para otorgar poderes
adicionales a los jugadores. Sólo recuerde restaurar la configuración a la normalidad cuando haya terminado de probar una función.

\subsection{Reposición de la flota}

Una característica clave en el juego Invasión Alienígena,  es que los alienígenas son implacables, cada vez que se destruye una flota, 
aparecer una nueva flota.

Para hacer que aparezca una nueva flota de alienígenas después de que una flota haya sido destruida, primero comprobamos si el grupo de alienígenas está vacío. Si es así, hacemos una llamada a {\ttfamily \_crear\_flota()}. 

Realizaremos esta verificación al final de {\ttfamily \_actualizar\_balas()}, porque ahí es donde se destruyen los extraterrestres individuales.

\hspace{-2.4cm}{\ttfamily \scriptsize invasión}

\vspace{-0.2cm}
\hspace{-2.1cm}{\ttfamily \scriptsize \_alien.py}\vspace{-0.84cm}
\begin{lstlisting}
    def _actualizar_balas(self):
        --recorte--
        if not self.aliens:
            # Destruye la existencia de las balas y crea una nueva flota.
            self.balas.empty()
            self._crear_flota()
\end{lstlisting}

Comprobamos si el grupo de extraterrestres está vacío. Se evalúa un grupo vacío
a False, por lo que esta es una forma sencilla de comprobar si el grupo está vacío. Si es así, nosotros
deshazte de las balas existentes utilizando el método vacío(), que elimina
todos los sprites restantes de un grupo. También llamamos \_crear\_flota(), que
vuelve a llenar la pantalla con extraterrestres.

Ahora aparece una nueva flota tan pronto como destruyes la flota actual.

{\ttfamily Acelerando las balas}

Si has intentado disparar a los alienígenas en el estado actual del juego, es posible que encuentres
que las balas no viajan a la mejor velocidad para el juego. Podrían ser un poco demasiado lento o un poco demasiado rápido. En este punto, puede modificar la configuración, para hacer el juego más interesante. 

Tenga en cuenta que el juego se volverá progresivamente más rápido, así que no hagas que el juego sea demasiado rápido al
comienzo.

Modificamos la velocidad de las balas ajustando el valor de {\ttfamily velocidad\_balas}
en {\ttfamily configurar.py}. En mi sistema, ajustaré el valor de bullet\_speed a 2,5, para que el
las balas viajan un poco más rápido:


\hspace{-2.5cm}{\ttfamily \scriptsize configurar.py}  \vspace{-0.5cm}
\begin{lstlisting}	
    # Configuracion de balas
    self.bullet_speed = 2.5
    self.bullet_width = 3
    --recorte--
\end{lstlisting}


El mejor valor para esta configuración depende de tu experiencia en el juego, así que encuentre un valor que funcione para usted. 
También puedes ajustar otras configuraciones.

\subsection{ Refactorización \_actualizar\_balas()}

Refactoricemos { \ttfamily\_actualizar\_balas() } para que no realice tantas tareas diferentes.

Moveremos el código para lidiar con colisiones entre balas y extraterrestres a una sección separada.
método:


\hspace{-2.4cm}{\ttfamily \scriptsize invasión}

\vspace{-0.2cm}
\hspace{-2.1cm}{\ttfamily \scriptsize \_alien.py}\vspace{-0.84cm}
\begin{lstlisting}
    def _actualizar_balas(self):
        --recorte--
        # obtener rid de las balas que se han disparado.
        for balas in self.balas.copy():
            if balas.rect.bottom <= 0:
                self.balas.remove(balas)
        self._chequear_collisions_balas_alien()

    def _chequear_colisiones_balas_alien(self):
        """Responder a las colisiones balas-alienigena."""
        # Remover cualquier bala y alienigena que han colisionado.
        colisiones = pygame.sprite.groupcollide( self.balas, self.aliens, True, True)
        if not self.aliens:
            # Destruye la existencia de las balas y crea una nueva flota.
            self.balas.empty()
            self._crear_flota()
\end{lstlisting}

Hemos creado un nuevo método,{ \ttfamily \_chequear\_colisiones\_balas\_alien}, para buscar colisiones entre balas y extraterrestres y responder adecuadamente si toda la flota ha sido destruida. Hacerlo evita que { \ttfamily  \_actualizar\_balas()} crezca demasiado y simplifica el desarrollo posterior.

\section{INTÉNTALO TÚ MISMO}

13-5. Tirador lateral Parte 2: 
Hemos recorrido un largo camino desde el Ejercicio 12-6, Tirador lateral. 
Para este ejercicio, intenta desarrollar Sideways Shooter hasta el mismo punto al que llevamos Alien Invasion. 

Añade una flota de alienígenas y haz que se muevan de lado hacia la nave o escriba un código que coloque extraterrestres en posiciones aleatorias a lo largo del lado derecho de la pantalla y luego los envíe hacia la nave. Además, escribe código que haga que los alienígenas desaparezcan cuando sean golpeados.

\section{Terminando el juego}

¿Cuál es la diversión y el desafío de jugar un juego que no puedes perder? 

Si el jugador no derriba la flota lo suficientemente rápido, haremos que los alienígenas destruyan la nave cuando hagan contacto. 
Al mismo tiempo, limitaremos la cantidad de naves que un jugador puede usar y destruiremos la nave cuando un extraterrestre llegue al final de la pantalla. El juego terminará cuando el jugador haya agotado todos sus la naves.

\section{Detección de colisiones}

Comenzaremos verificando colisiones entre alienígenas y la nave, para que podamos responder adecuadamente cuando un 
alienígena  golpee la nave. Comprobaremos si hay colisiones entre naves y alienígenas inmediatamente después de actualizar la posición de cada alienígena en {\ttfamily InvasionAlien}:


\hspace{-2.4cm}{\ttfamily \scriptsize invasión}

\vspace{-0.2cm}
\hspace{-2.1cm}{\ttfamily \scriptsize \_alien.py}\vspace{-0.84cm}
\begin{lstlisting}
\begin{lstlisting}
    def _actualiar_aliens(self):
        --recorte--
        self.aliens.update()
        # Ver colisiones entre alienigena y nave.
        if pygame.sprite.spritecollideany(self.nave07, self.aliens):
            print("¡¡¡Nave golpeada!!!")
\end{lstlisting}

La función {\ttfamily spritecollideany()} toma dos argumentos: un sprite (nave) y un grupo (alien). La función busca cualquier miembro del grupo que haya  chocado con el sprite y deja de recorrer el grupo tan pronto como encuentra un miembro que ha chocado con el sprite. 

Aquí, recorre el grupo de alienígenas y devuelve al primer alienígena que encuentra que ha chocado con la nave.


Si no ocurren colisiones, {\ttfamily spritecollideany()} devuelve Ninguno y el bloque {\ttfamily if } no ejecutará. 
Si encuentra un alienígena que ha chocado con la nave, regresa ese alienígena y el bloque {\ttfamily if } se ejecuta:
 imprime ¡¡¡Nave golpeada!!!. 
 
Cuando un extraterrestre golpea la nave, tendremos que realizar una serie de tareas:
eliminar a todos los alienígenas restantes y balas, volver a centrar la nave y crear una nueva flota. 


Para  escribir código ese código debemos saber si nuestro enfoque para detectar colisiones de naves alienígenas
funciona correctamente. Escribir un {\ttfamily print()} como un indicador, es una forma sencilla y muy practica de
asegurarnos de que estamos detectar adecuadamente estas colisiones.

Ahora, cuando ejecutas  Invasión Alienígena, debe aparecer  el mensaje  ¡¡¡Nave golpeada!!! 
en la terminal cada vez que un extraterrestre golpea la nave. 

Cuando estás probando esta característica, establezca  {\ttfamily cae\_velocidad\_flota}  en un valor más alto, como 50 o 100,
para que el Los extraterrestres llegan a tu nave más rápido.

\subsubsection{Respondiendo a las colisiones entre naves alienígenas}

Ahora necesitamos descubrir exactamente qué sucederá cuando un extraterrestre colisione con la nave.
En lugar de destruir la nave y crear una nueva, contaremos cuántas veces la nave ha sido impactado mediante las estadísticas
de seguimiento para el juego. Las estadísticas de seguimiento también serán útiles para las puntuaciones.

Escribamos una nueva clase,  {\ttfamily Estadisticas\_Juego()}, para realizar un seguimiento de las estadísticas del juego y guárdela.
como  {\ttfamily stats\_juego.py}:


\hspace{-2.4cm}{\ttfamily \scriptsize stats}

\vspace{-0.2cm}
\hspace{-2.1cm}{\ttfamily \scriptsize \_juego.py}\vspace{-0.84cm}
\begin{lstlisting}
class Estadisticas_Juego():
    """Track estadisticas parael juego Invasión Alienígena."""
    def __init__(self, ia_jugar):
        """Inicializa estadisticas."""
        self.configuracion = ia_jugar.configuracion
        self.resetear_stats()
   
    def resetear_stats(self):
        """Inicializa estadisticas que cambian durante el juego."""
        self.nave_permitidas = self.cofiguracion.limite_nave
\end{lstlisting}	



Creamos una instancia de {\ttfamily Estadisticas\_Juego()} durante todo el tiempo que dure Invasión Alienígena ejecutándose, pero necesitaremos restablecer algunas estadísticas cada vez que el jugador inicie un nuevo juego.  Para hacer esto, inicializaremos 
la mayoría de las estadísticas en el  método {\ttfamily resetear\_stats()}, en lugar de escribirlas en {\ttfamily \_\_init\_\_()}. 

Llamaremos a este método desde {\ttfamily \_\_init\_\_()} por lo que las estadísticas se configuran correctamente
cuando se crea por primera vez la instancia de {\ttfamily Estadisticas\_Juego()}.
Pero también podremos llamar a  {\ttfamily resetear\_stats()} cada vez que el jugador comience un nuevo juego.

En este momento solo tenemos una estadística, la {\ttfamily nave\_permitidas}, cuyo valor cambiará
durante todo el juego.  La cantidad de naves que dispone el jugador se almacenan en {\ttfamily configurar.py}
como {\ttfamily limite\_nave}:

\hspace{-2.75cm}{\ttfamily \scriptsize configurar.py}  \vspace{-0.5cm}
\begin{lstlisting}
    # Configurar Nave
    self.velocidad_nave = 10
    self.limite_nave = 3
\end{lstlisting}

También necesitamos realizar algunos cambios en {\ttfamily invasión\_alien.py} para crear una instancia de {\ttfamily Estadisticas\_Juego()}. Primero, actualizaremos las declaraciones para importar en la parte superior del archivo:


\hspace{-2.4cm}{\ttfamily \scriptsize invasión}

\vspace{-0.2cm}
\hspace{-2.1cm}{\ttfamily \scriptsize \_alien.py}\vspace{-0.84cm}
\begin{lstlisting}
import sys
from time import sleep
import pygame
from configurar import Configuracion
from stats_juego import Estadisticas_Juego
from Nave import Nave
--recorte--
\end{lstlisting}	



Importamos la función {\ttfamily sleep()} del módulo de {\ttfamily time}  de la biblioteca estándar de Python, 
para que podamos pausar el juego por un momento cuando la nave es golpeado. También importamos {\ttfamily Estadisticas\_Juego()}. Y crearemos una instancia de {\ttfamily Estadisticas\_Juego()} en  {\ttfamily\_\_init\_\_()}:



\hspace{-2.4cm}{\ttfamily \scriptsize invasión}

\vspace{-0.2cm}
\hspace{-2.1cm}{\ttfamily \scriptsize \_alien.py}\vspace{-0.84cm}
\begin{lstlisting}
    def __init__(self):
        --recorte--
        self.pantalla = pygame.pantalla.set_mode(
                         (self.configuracion.pantalla_width,
                           self.configuracion.pantalla_height))
        pygame.display.set_caption("Invasión Alienígena")

        # Crea una instancia para almacenar estadísticas del juego.
        self.estadisticas = Estadisticas_Juego(self)
        self.nave = Nave(self)
         --recorte--
\end{lstlisting}


Creamos la instancia después de crear la ventana del juego, pero antes de definir otros elementos del juego, como por ejemplo la nave.
Cuando un extraterrestre golpee la nave, restaremos 1 del número de naves que quedan, y destruimos todos los extraterrestres y las balas existentes, crearemos una nueva flota y reponemos la nave la posición inicial en la pantalla. 

También pausaremos el juego por un momento para que el jugador pueda notar la colisión y reagruparse antes de que 
aparezca una nueva flota.

Crearemos un nuevo método llamado {\ttfamily \_nave\_golpeada()}, y escribimos la mayor parte de este código en el. 
En el método {\ttfamily \_actualizar\_aliens()} se hace el llamado al nuevo método, cuando un extraterrestre golpee la nave:


\hspace{-2.4cm}{\ttfamily \scriptsize invasión}

\vspace{-0.2cm}
\hspace{-2.1cm}{\ttfamily \scriptsize \_alien.py}\vspace{-0.84cm}
\begin{lstlisting}
    def _nave_golpeada(self):
        """Responde a la nave es golpeada por un alienigena."""
        # Decrementa nave_
        self.estadisticas.nave_permitidas -= 1
        # Obtiene rid de cualquier remaining bala y alienigena
        self.balas.empty()
        self.aliens.empty()
        # Crea una nueva flota y una nave en el centro
        self._crear_flota()
        self.nave.nave_centro()
        # Pause.
        sleep(0.5)
\end{lstlisting}


El nuevo método {\ttfamily \_nave\_golpeada()} coordina la respuesta cuando un extraterrestre
choca contra la nave. Dentro de  {\ttfamily \_nave\_golpeada()}  se reduce en 1, el número de la naves que quedan
y se procede a vaciar los grupos de balas y alienígenas.

Se continua creando  una nueva flota y colocando  la nave en su posición inicial. (Agregaremos el
método {\ttfamily nave\_centro()} para iniciar la nave). Luego agregamos una pausa después de realizar actualizaciones
en todos los elementos del juego, pero antes de cualquier cambio en la pantalla, para permitirle al jugador  ver que su nave ha sido golpeada. 
 
La llamada  {\ttfamily sleep()} detiene la ejecución del programa durante medio segundo, tiempo suficiente para que el 
jugador vea que el extraterrestre ha golpeado la nave.

Cuando la función {\ttfamily sleep()} finaliza, la ejecución del código pasa al método  {\ttfamily \_actualizar\_pantalla()},
que dibuja la nueva flota en la pantalla.

En {\ttfamily\_actualizar\_aliens()}, reemplazamos la llamada {\ttfamily print()} con una llamada a {\ttfamily \_nave\_golpeada()}
cuando un extraterrestre choca con otra la nave:


\hspace{-2.4cm}{\ttfamily \scriptsize invasión}

\vspace{-0.2cm}
\hspace{-2.1cm}{\ttfamily \scriptsize \_alien.py}\vspace{-0.84cm}
\begin{lstlisting}
    def _actualizar_aliens(self):
      --recorte--
        if pygame.sprite.spritecollideany(self.nave, self.aliens):
           self._nave_golpeada()
\end{lstlisting}

El nuevo método  {\ttfamily nave\_centro()} que centra la nave en la pantalla, se agrega  en el archivo {\ttfamily nave.py}:

\hspace{-2.5cm}{\ttfamily \scriptsize nave.py}  \vspace{-0.5cm}
\begin{lstlisting}
    def nave_centro(self):
        """Centra la nave en la pantalla."""
        self.rect.midbottom = self.pantalla_rect.midbottom
        self.x = float(self.rect.x)
\end{lstlisting}


Centramos la nave de la misma manera que lo hicimos en {\ttfamily\_\_init\_\_()}. Después de centrarlo,
restablecemos el atributo {\ttfamily self.x}, que nos permite rastrear la posición exacta de la nave.



Tenga en cuenta que nunca se crea más de una nave; hacemos solo una instancia de la nave para el
todo el juego y se vuelve a centrar cada vez que la nave haya sido golpeado.
La estadística se  {\ttfamily nave\_permitidas}
nos avisará cuando el jugador se haya quedado sin naves.

Al ejecutar el juego, dispare a algunos extraterrestres y deje que un extraterrestre golpee la nave. 

El juego debería hacer una pausa y debería aparecer una nueva flota con ella nave centrado en el
parte inferior de la pantalla nuevamente.

\subsubsection{Aliens que llegan al final de la pantalla} 

Si un extraterrestre llega al final de la pantalla, el juego responderá.
de la misma manera que ocurre cuando un extraterrestre choca contra la nave. Para comprobar cuándo sucede esto,
agregue un nuevo método en {\ttfamily invasion\_alien.py}:


\hspace{-2.4cm}{\ttfamily \scriptsize invasión}

\vspace{-0.2cm}
\hspace{-2.1cm}{\ttfamily \scriptsize \_alien.py}\vspace{-0.84cm}
\begin{lstlisting}
    def _chequear_aliens_piso(self):
        """Verifica si algún alienígena alcanza el piso de la pantalla."""
        for alien in self.aliens.sprites():
            if alien.rect.bottom >= self.configuraciones.pantalla_height:
                # Tratado de la misma forma como si la nave es golpeada.
                self._nave_golpeada()
                break
\end{lstlisting}


El método \_check\_aliens\_bottom() comprueba si algún extraterrestre ha llegado al final de la pantalla. Un extraterrestre llega al fondo cuando su valor rect.bottom es mayor o igual a la altura de la pantalla 1. Si un extraterrestre llega al fondo, llamamos \_ship\_hit(). Si un alienígena toca el fondo, no hay necesidad de comprobar el resto, por lo que salimos del bucle después de llamar a \_ship\_hit().
Llamaremos a este método desde \_update\_aliens():


\hspace{-2.4cm}{\ttfamily \scriptsize invasión}

\vspace{-0.2cm}
\hspace{-2.1cm}{\ttfamily \scriptsize \_alien.py}\vspace{-0.84cm}
\begin{lstlisting}
    def _actualizar_aliens(self):
        --recorte--
        # ver colisiones entre alien-nave.
        if pygame.sprite.spritecollideany(self.nave, self.aliens):
            self._nave_golpeada()
        # Ver alienigena golpea el piso de la pantalla.
        self._chequear_aliens_piso()
\end{lstlisting}


Llamamos a {\ttfamily self.\_chequear\_aliens\_piso()} después de actualizar las posiciones de todos los alienígenas y después de buscar colisiones entre naves alienígenas. Ahora aparecerá una nueva flota cada vez que la nave sea golpeada por un extraterrestre o un extraterrestre llegue al final de la pantalla.

¡Juego terminado!


Invasión se siente más completo ahora, pero el juego nunca termina. 
El valor de las naves dejados se vuelve cada vez más negativo.
 
Agreguemos un indicador {\ttfamily juego\_activo}, para que podamos finalizar el juego cuando el jugador se quede sin  naves. 
Estableceremos esta indicador al final del método \_\_init\_\_() en InvasionAlien:

\hspace{-2.4cm}{\ttfamily \scriptsize invasión}

\vspace{-0.2cm}
\hspace{-2.1cm}{\ttfamily \scriptsize \_alien.py}\vspace{-0.84cm}
\begin{lstlisting}
    def __init__(self):
        --recorte--
        # Comienza   Invasion Alienen un estado activo.
        self.juego_activo = True


\end{lstlisting}

En el método  {\ttfamily \_nave\_golpeada()} agregamos el indicador {\ttfamily juego\_activo  =  False}, para señala que el jugador
ha agotado todos sus naves:



\hspace{-2.4cm}{\ttfamily \scriptsize invasión}

\vspace{-0.2cm}
\hspace{-2.1cm}{\ttfamily \scriptsize \_alien.py}\vspace{-0.84cm}
\begin{lstlisting}
    def _nave_golpeada(self):
        """Responde cuando la nave ha sido golpeada por un alienigena."""
        if self.estadisticas.nave_permitidas > 0:
            # Restamos 1 a nave_permitidas
            self.estadisticas.nave_permitidas  -= 1
            --recorte--
            # Pause.
            sleep(0.5)
        else:
            self.juego_activo = False
\end{lstlisting}

La mayor parte del método  {\ttfamily \_nave\_golpeada()} no ha cambiado. Hemos movido todo el código existente
en un bloque {\ttfamily if}, que prueba para asegurarse de que el jugador tenga al menos una nave
restante. Si es así, creamos una nueva flota, hacemos una pausa y seguimos adelante. Si al jugador no quedan naves, configuramos  {\ttfamily juego\_activo} en  {\ttfamily False}.

\section{Identificar cuándo deben ejecutarse las partes del juego}

Necesitamos identificar las partes del juego que siempre deben ejecutarse y las
partes que deberían ejecutarse sólo cuando el juego está activo:


\hspace{-2.4cm}{\ttfamily \scriptsize invasión}

\vspace{-0.2cm}
\hspace{-2.1cm}{\ttfamily \scriptsize \_alien.py}\vspace{-0.84cm}
\begin{lstlisting}
    def run_jugar(self):
        """Start the main loop for the game."""
        while True:
            self._check_events()
            if self.juego_activo:
                self.nave.update()
            self._actualizar_balas()
            self._actualizar_aliens()
            self._actualizar_pantalla()
            self.clock.tick(60)
\end{lstlisting}


En el bucle principal, siempre es necesario llamar al método  {\ttfamily \_chequear\\ \_eventos()}, incluso si
el juego está inactivo. Por ejemplo, necesitamos saber si el usuario presiona  {\ttfamily Q} para
salir del juego o hace clic en el botón para cerrar la ventana. También continuamos
actualizando la pantalla para que podamos realizar cambios en la pantalla mientras esperamos
ver si el jugador elige comenzar un nuevo juego.

El resto de las funciones  deben se llamadas cuando el juego está activo, porque cuando el juego
está inactivo, no necesitamos actualizar las posiciones de los elementos del juego.
Ahora, cuando juegues Alien Invasion, el juego debería congelarse cuando hayas consumiste todos las  naves permitidas.

\section{INTÉNTALO TÚ MISMO}

13-6. Fin del juego: en Sideways Shooter (dispara de un lado), lleva la cuenta del número de veces que
la nave es golpeado y el número de veces que  la nave golpea a un extraterrestre.
Decidir sobre una condición apropiada para terminar el juego, y detener el juego cuando esta situación
ocurre.

\section{Resumen}

En este capítulo, aprendiste cómo agregar una gran cantidad de elementos idénticos en un juego creando una flota de alienígenas.
Se uso  bucles anidados para crear un cuadrícula de elementos, y se hizo que un gran conjunto de elementos del juego se movieran llamando
con el  método {\ttfamily update()} de cada elemento. Aprendiste a controlar la dirección de
objetos en la pantalla y para responder a situaciones específicas, como cuando la flota llega al borde de la pantalla. Detectaste y respondiste a colisiones cuando las balas alcanzaron a los extraterrestres y los extraterrestres impactaron en la nave. También aprendió a realizar un seguimiento de las estadísticas de un juego y a utilizar el indicador {\ttfamily juego\_activo} para determinar cuándo 
termina el juego.\\


En el próximo y último capítulo de este proyecto, agregaremos un botón  {\ttfamily Jugar} para que el jugador pueda elegir cuándo
comenzar su primer juego y si desea volver a jugar cuando finalice el juego. 
Aceleraremos el juego cada vez que el jugador derribe a toda la flota y agregaremos un sistema de puntuación.\\ 

¡El resultado final será un juego totalmente jugable!




%%%%%%%%%%%%%%%%%%%    %%%%%%%%%%%%%%%      %%%%%%%%%%%%%%%%%%%%%



%%%%%%%%%%%%%%%%%%%    %%%%%%%%%%%%%%%      %%%%%%%%%%%%%%%%%%%%%



%%%%%%%%%%%%%%%%%%%    %%%%%%%%%%%%%%%      %%%%%%%%%%%%%%%%%%%%%

%%%%%%%%%%%%%%%%%%%    %%%%%%%%%%%%%%%      %%%%%%%%%%%%%%%%%%%%%


\chapter{PUNTUACIÓN}

En este capítulo, terminaremos de construir Invasión Alienígena. Agregaremos un botón para iniciar el juego a solicitud del jugador y reiniciar el juego una vez que termine. Le asignaremos un sistema de niveles cuando el jugador elimina una flota; el nivel del juego esta asociado con la rapidez de los alienígenas para bajar;  también implementar un sistema de puntuación. xxx\\


Al final del capítulo, habrás aprendido  lo suficiente como para empezar a escribir juegos que
aumenta en dificultad a medida que el jugador progresa y que cuentan con sistemas completos de puntuación.


\section{Agregar el botón Play}


En esta sección, agregaremos un botón {\ttfamily Play} (Jugar o reproducir) que aparece antes de que comience un juego y reaparece
cuando termina el juego, para que el jugador pueda volver a jugar si lo desea.

Hasta este momento, el juego comienza tan pronto como ejecutas el archivo {\ttfamily invasion\_alien.py},  ahora  vamos a iniciar el juego en un estado inactivo, y solo después de que el jugador haga clic en el botón  {\ttfamily Play}  comenzará el juego. 


Para hacer esto, debemos modificar el método {\ttfamily \_\_init\_\_()} de la clase  {\ttfamily InvasionAlien}:

\hspace{-2.4cm}{\ttfamily \scriptsize invasión}

\vspace{-0.2cm}
\hspace{-2.1cm}{\ttfamily \scriptsize \_alien.py}\vspace{-0.84cm}
\begin{lstlisting}		
    def __init__(self):
        """Inicializa el juego y crea recursos para el juego."""
        pygame.init()
        --recorte--
        # Inicia Invasión Alienígena en un estado inactivo.
        self.juego_activo = Falso
\end{lstlisting}

Ahora el juego debería comenzar en un estado inactivo, sin posibilidad de iniciarlo (hasta que hagamos clic en un botón de Play, pero por ahora no lo hemos creado).

\subsubsection{Crear una clase de botón}

Debido a que Pygame no tiene un método integrado para crear botones, debemos escribir una clase botón para crear un rectángulo relleno con una etiqueta. Puedes usar este código para hacer cualquier botón en un juego.

Creamos un módulo {\ttfamily boton.py}  para desarrollar la de la clase {\ttfamily Boton},  y lo guardamos en la carpeta principal; es escribimos el siguiente código:

\hspace{-2.5cm}{\ttfamily \scriptsize boton.py}  \vspace{-0.5cm}
\begin{lstlisting}
import pygame.font
class  Boton:
    """Una clase para contruir  botón en el juego."""
    def __init__(self, ia_jugar, msg):
        """Inicializa atributos de  Botón."""
        self.pantalla = ia_jugar.pantalla
        self.pantalla_rect = self.pantalla.get_rect()
        # Configuramos las dimensione y propiedade del botón.
        self.ancho, self.alto = 200, 50
        self.color_boton = (0, 135, 0)
        self.color_texto = (255, 255, 255)
        self.font = pygame.font.SysFont(None, 48)
        # Construción del objeto rectangulo del botón y cálculo del centro.
        self.rect = pygame.Rect(0, 0, self.ancho, self.alto)
        self.rect.center = self.pantalla_rect.center
        # El mensaje necesario para presionar una sola vez.
        self._prep_msg(mensaje)
\end{lstlisting} 
% msj = mensaje
Primero, importamos el módulo  {\ttfamily pygame.font}, que permite a PyGame representar texto 
a la pantalla. El método  {\ttfamily \_\_init\_\_()} toma los parámetros  {\ttfamily self} y el objeto {\ttfamily  ia\_jugar}
y el {\ttfamily mensaje}, que contiene el texto del botón. Configuramos las dimensiones ancho 200 y alto 50, el color  verde oscuro del objeto {\ttfamily rect}, y configure {\ttfamily texto\_color} para representar el texto en blanco.


A continuación, preparamos un atributo de fuente para representar el texto. El argumento {\ttfamily None}, 
le dice a PyGame que use la fuente predeterminada, y {\ttfamily 48} especifica el tamaño del texto. Para centrar el botón en la pantalla, creamos un rectángulo para el botón y establezca su atributo central para que coincida con el de la pantalla.
Pygame trabaja con texto representando la cadena que desea mostrar una imagen. Finalmente, llamamos a {\ttfamily \_prep\_msg()} para manejar esta representación.

Aquí está el código para \_prep\_msg():

\hspace{-2.5cm}{\ttfamily \scriptsize boton.py}  \vspace{-0.5cm}
\begin{lstlisting}		
    def _prep_msg(self, mensaje):
        """Retorna mensaje dentro una imagen renderizada y centra 
        el texto en el botón."""
        self.msg_image = self.font.render(mensaje, True, 
                             self.color_texto,self.color_boton)
        self.mensaje_image_rect = self.mensaje_image.get_rect()
        self.mensaje_image_rect.center = self.rect.center
\end{lstlisting}


El método {\ttfamily\_prep\_msg()} necesita un parámetro propio y el texto que se va a representar 
como imagen (mensaje). La llamada a {\ttfamily font.render()} convierte el texto almacenado en una
una imagen, que luego almacenamos en {\ttfamily self.mensaje\_image}. El {\ttfamily font.render()}

El método también toma un valor booleano para activar o desactivar el antialiasing (antialiasing
suaviza los bordes del texto). Los argumentos restantes, el  color de fuente y color de fondo especificados. 
Establecemos antialiasing en True y configuramos el fondo del texto al mismo color que el botón. Si no incluye
un color de fondo, PyGame intentará representar la fuente con un color de fondo transparente.
Centramos la imagen del texto en el botón creando un rectángulo  a partir de la imagen y configurando su atributo 
central para que coincida con el del botón.


Finalmente, creamos un método {\ttfamily dibuja\_boton()} al que podemos llamar para mostrar el botón en pantalla:


\hspace{-2.5cm}{\ttfamily \scriptsize boton.py}  \vspace{-0.5cm}
\begin{lstlisting}
    def dibuja_boton(self):
        """Dibuja boton y dibuja."""
        self.pantalla.fill(self.color_boton, self.rect)
        self.pantalla.blit(self.mensaje_image, self.mensaje_image_rect)
\end{lstlisting}

Llamamos a {\ttfamily pantalla.fill()} para dibujar la parte rectangular del botón, y
 {\ttfamily  pantalla.blit()} para dibujar la imagen del texto en la pantalla, pasándole una imagen
y el objeto {\ttfamily rectángulo } asociado con la imagen. Esto completa la clase Botón.

\subsubsection{Dibujar el botón en la pantalla}

Usaremos la clase Botón para crear un botón Play en la clase {\ttfamily AlienInvasion}. Primero, vamos a
actualice las declaraciones de importación:


\hspace{-2.4cm}{\ttfamily \scriptsize invasión}

\vspace{-0.2cm}
\hspace{-2.1cm}{\ttfamily \scriptsize \_alien.py}\vspace{-0.84cm}
\begin{lstlisting}		
--recorte--
from stats_juego import Estadisticas_Juego
from boton import Boton
\end{lstlisting}


Como solo necesitamos un botón para reproducir  (Play),  para ello creamos el botón dentro del método {\ttfamily\_\_init\_\_()} de la clase  {\ttfamily AlienInvasion} pero lo colocamos al final:

\hspace{-2.4cm}{\ttfamily \scriptsize invasión}

\vspace{-0.2cm}
\hspace{-2.1cm}{\ttfamily \scriptsize \_alien.py}\vspace{-0.84cm}
\begin{lstlisting}	
    def __init__(self):
        --recorte--
        self.juego_activo = False
        # Hace el boton  Play 
        self.play_boton = Boton(self, "Play")
\end{lstlisting}


Este código crea una instancia de Botón con la etiqueta {\ttfamily Play}, pero no dibuja el botón en la pantalla.
Para dibujarlo, llamaremos al método {\ttfamily dibujar\_boton()}  en {\ttfamily \_actualizar\_pantalla()}:

\hspace{-2.4cm}{\ttfamily \scriptsize invasión}

\vspace{-0.2cm}
\hspace{-2.1cm}{\ttfamily \scriptsize \_alien.py}\vspace{-0.84cm}
\begin{lstlisting}		
    def _actualizar_pantalla(self):
        --recorte--
        self.aliens.draw(self.pantalla)
        # Dibuja el botón Play si el juego esta inactivo 
        if not self.jugar_activo:
            self.play_boton.dibuja_boton()
        pygame.display.flip()
\end{lstlisting}

Para que el botón {\ttfamily Play} sea visible sobre todos los demás elementos de la pantalla, lo dibujamos después de que se hayan dibujado todos los demás elementos, pero antes de pasar a una nueva pantalla. Lo incluimos en un bloque {\ttfamily if}, por lo que el botón solo aparece cuando el juego está inactivo.

Ahora, al ejecutar  Invasión Alienígena, debe aparecer el botón Play en el centro de la pantalla, como se muestra en la Figura \ref{play01}.


\begin{figure}[h!]
	\centering
	\includegraphics[scale=0.40]{../A Logos de Python/play01}
	\caption{Aparecer el botón Play.}
	\label{play01}
\end{figure}



\subsubsection{Comenzando el juego}
 
Para comenzar un juego nuevo cuando  el jugador hace clic en Jugar,  se agregue el siguiente bloque  {\ttfamily elif}
al final del método  {\ttfamily \_chequear\_eventos()} para monitorear los eventos del mouse sobre el botón:

\hspace{-2.4cm}{\ttfamily \scriptsize invasión}

\vspace{-0.2cm}
\hspace{-2.1cm}{\ttfamily \scriptsize \_alien.py}\vspace{-0.84cm}
\begin{lstlisting}	
    def _chequear_eventos(self):
        """Respuesta a teclas presionads y eventos del muuse"""
        --recorte--
            elif event.type == pygame.MOUSEBUTTONDOWN:
                mouse_pos = pygame.mouse.get_pos()
                self._chequear_play_boton(mouse_pos)
\end{lstlisting}

Pygame detecta un evento {\ttfamily MOUSEBUTTONDOWN} cuando el jugador hace clic en cualquier lugar
en la pantalla, pero queremos restringir el juego para que responda sólo a los clics del mouse
en el botón Play. Para lograr esto, usamos {\ttfamily pygame.mouse.get\_pos()},
que devuelve una tupla que contiene las coordenadas $x$ e $y$ del cursor del mouse
al hace clic. Para escribir estos códigos creamos un nuevo método  {\ttfamily\_chequear\_play\_button()}, y lo ubicamos después de {\ttfamily \_chequear\_eventos()}:

\hspace{-2.4cm}{\ttfamily \scriptsize invasión}

\vspace{-0.2cm}
\hspace{-2.1cm}{\ttfamily \scriptsize \_alien.py}\vspace{-0.84cm}
\begin{lstlisting}	
    def _chequear_play_boton(self, mouse_pos):
        """Comenzar un nuevo juego cuando hace clic en Play."""
        if self.play_boton.rect.collidepoint(mouse_pos):
            self.juego_activo= True
\end{lstlisting}

Usamos el método {\ttfamily rect.collidepoint()} se usa  para verificar si el punto del clic del mouse se superpone a la región definida por el rectángulo del botón Play.  Entonces si configuramos {\ttfamily juego\_activo}  en  True, haciendo clic en Play:   ¡el juego comienza!


En este punto, deberías poder comenzar y jugar un juego completo. 

Cuando el el juego termina, el valor de {\ttfamily juego\_activo} debe volver al estado inactivo  {\ttfamily False } y el botón Play
debería reaparecer. 


\subsubsection{Reiniciar el juego}

El código del botón Play que acabamos de escribir funciona la primera vez que el jugador hace clic Jugar. 
Pero no funciona después de que termina el primer juego, porque las condiciones que provocaron que el juego 
terminara no se han reiniciado.

Para reiniciar el juego cada vez que el jugador hace clic en Play, debemos resetear las
estadísticas del juego, eliminar viejos alienígenas y balas, construir una nueva flota y
centrar la nave, como se muestra aquí:


\hspace{-2.4cm}{\ttfamily \scriptsize invasión}

\vspace{-0.2cm}
\hspace{-2.1cm}{\ttfamily \scriptsize \_alien.py}\vspace{-0.84cm}
\begin{lstlisting}	
    def _chequear_play_boton(self, mouse_pos):
        """Comienza un nuevo juego cuando el jugador hace clic en Play."""
        if self.play_boton.rect.collidepoint(mouse_pos):
            # Restablecer las estadísticas del juego.
            self.estadisticas.resetear_stats()
            self.juego_activo = True

            # Obtiene las rid de los balas y alienigenas  remaining.
            self.balas.empty()
            self.aliens.empty()
            # Crea una nueva flota y centra la nave
            self._crear_flota()
            self.nave.nave_centro()
\end{lstlisting}



Restablecemos las estadísticas del juego, lo que le da al jugador tres naves nuevos. Luego configuramos {\ttfamily juego\_activo}
 en True para que el juego comience tan pronto como termine de ejecutarse el código de esta función. Vaciamos los alienígenas y 
 los grupos de balas, y luego creamos una nueva flota y centramos la nave.

Ahora el juego se reiniciará correctamente cada vez que hagas clic en Jugar, ¡lo que te permitirá jugar tantas veces como quieras!


\subsubsection{Desactivar el botón Play}

Un problema con nuestro botón Play es que la región del botón en la pantalla seguirá respondiendo a los clics incluso cuando el botón Play no esté visible. Si haces clic en el área del botón Play  por accidente después de que comience un juego, ¡el juego se reiniciará!. 
Para solucionar este problema, configura el juego para que se inicie solo cuando {\ttfamily juego\_activo} sea False:


\hspace{-2.4cm}{\ttfamily \scriptsize invasión}

\vspace{-0.2cm}
\hspace{-2.1cm}{\ttfamily \scriptsize \_alien.py}\vspace{-0.84cm}
\begin{lstlisting}		
    def _chequear_play_boton(self, mouse_pos):
        """Comienza un nuevo juego cuando el jugador hace clic en Play."""
        boton_clicked = self.play_boton.rect.collidepoint(mouse_pos)
        if boton_clicked and not self.juego_activo: 
            # Restablecer las estadísticas del juego.
            self.estadisticas.resetear_stats()
            self.juego_activo = True
         --recorte--
\end{lstlisting}

El indicador {\ttfamily boton\_clicked} almacena un valor True o Falso, y el juego se reiniciará solo si se hace clic en Play y el juego no está activo actualmente. Para probar este comportamiento, inicie un nuevo juego y haga clic repetidamente donde debería estar el botón Play. 

Si todo funciona como se esperaba, al hacer clic en el área del botón Play no debe  tener ningún efecto en el juego.

\subsubsection{Ocultar el cursor del mouse}

Queremos que el cursor del mouse esté visible cuando el juego esté inactivo, pero una vez que comienza el juego, simplemente estorba. Para solucionar este problema, lo haremos invisible cuando el juego se active. Podemos hacer esto al final del bloque {\ttfamily if}  en {\ttfamily \_chequear\_play\_boton()}:


\hspace{-2.4cm}{\ttfamily \scriptsize invasión}

\vspace{-0.2cm}
\hspace{-2.1cm}{\ttfamily \scriptsize \_alien.py}\vspace{-0.84cm}
\begin{lstlisting}
    def _chequear_play_boton(self, mouse_pos):
        """Comienza un nuevo juego cuando el jugador hace clic en Play."""
        boton_clicked = self.play_boton.rect.collidepoint(mouse_pos)
        if boton_clicked and not self.juego_activo: 
            --recorte--
            # Ocultar el cursor del mouse
            pygame.mouse.set_visible(False)
\end{lstlisting}


Pasar {\ttfamily False} a  {\ttfamily set\_visible()} le dice a PyGame que oculte el cursor cuando
el mouse está sobre la ventana del juego.

Haremos que el cursor reaparezca una vez que termine el juego para que el jugador pueda
Haz clic en Jugar nuevamente para comenzar un nuevo juego. Aquí está el código para hacer eso:

\hspace{-2.4cm}{\ttfamily \scriptsize invasión}

\vspace{-0.2cm}
\hspace{-2.1cm}{\ttfamily \scriptsize \_alien.py}\vspace{-0.84cm}
\begin{lstlisting}
    def _nave_golpeada(self):
        """Responde a cuando la nave ha sido golpeada por un alienigena."""
        if self.estadisticas.nave_izquierda > 0:
        --recorte--
        else:
            self.juego_activo = False
            pygame.mouse.set_visible(True)
\end{lstlisting}

Hacemos que el cursor vuelva a ser visible tan pronto como el juego se vuelve inactivo,
lo que sucede en {\ttfamily \_nave\_golpeada()}. La atención a detalles como este hace que su juego tenga un 
aspecto más profesional y permite al jugador concentrarse en jugar, en lugar de que descubrir la interfaz de usuario.

\section{INTÉNTALO TÚ MISMO}


14-1. Presione P para reproducir: debido a que Alien Invasion usa la entrada del teclado para controlar el
nave, sería útil iniciar el juego presionando una tecla. Agregue código que permita al
El jugador presiona P para comenzar. Podría ser útil mover algún código de {\ttfamily \_chequear\_play\_boton()}
a un método {\ttfamily  \_start\_game()} que se puede llamar desde {\ttfamily  \_chequear\_play\_boton()} y
{\ttfamily \_chequear\_keydown\_eventos()}.



14-2. Práctica de tiro: cree un rectángulo en el borde derecho de la pantalla que
sube y baja a un ritmo constante. Luego, en el lado izquierdo de la pantalla, cree
un nave que el jugador puede mover hacia arriba y hacia abajo mientras dispara balas al
objetivo rectangular. Añade un botón Jugar que inicia el juego y cuando el jugador
falla el objetivo tres veces, finaliza el juego y haz que reaparezca el botón Jugar.

Deje que el jugador reinicie el juego con este botón Jugar.


\section{Subir de nivel}



En nuestro juego actual, una vez que un jugador derriba a toda la flota alienígena, el
El jugador alcanza un nuevo nivel, pero la dificultad del juego no cambia, debemos agregarle algo de dificultad para darle  
un poco emoción y hacer que el juego sea más desafiante; si aumentamos la velocidad de los alienígenas cada vez que un jugador elimina una flota, el juego se complica un poco  más en cada nivel.


\subsubsection{Modificar la configuración de velocidad}


Primero reorganizaremos la clase Configuración para agrupar la configuración del juego en
los estáticos y los dinámicos. También nos aseguraremos de que cualquier configuración cambie
durante el reinicio del juego cuando al comenzar un juego nuevo:



\hspace{-2.75cm}{\ttfamily \scriptsize configurar.py}  \vspace{-0.5cm}
\begin{lstlisting}	
    def __init__(self):
        """Inicializamos la configuracion de estatica del juego."""
        #  Configuración de pantalla    
        self.fondo_image_filename = 'pygame.jpg' 
        self.pantalla_width = 800
        self.pantalla_height = 600
        self.color_fondo = (170,238,187)#(230, 230, 230)

        # Configuraciones de Nave        
        self.limite_nave = 3

        # configuraciones para las balas
        self.velocidad_balas = 1.0
        self.ancho_balas = 3
        self.alto_balas = 10
        self.color_balas = (0, 0, 0)
        self.balas_permitidas = 10

        # Configurancion Alienigena
        self.velocidad_alien = 1.0 
        self.flota_cae_velocidad = 5
        
        # Dirección de flota: 1 representa derecha;
        # -1 representa izquierda.
        self.direccion_flota = 10  
       
        # Que tan rapido la sube velocidad del juego
        self.velocidad_sube_escala = 1.1

        self.inicialize_configuracion_dinamica()

\end{lstlisting}


Las configuraciones que permanecen constantes la ubicamos en el método {\ttfamily \_\_init\_\_()}. 
Agregamos una configuración {\ttfamily  self.velocida\_sube\_scale = 1.1} para controlar qué tan rápido se acelera el juego: un valor de 2 duplicará la velocidad del juego cada vez que el jugador alcance un nuevo nivel; un valor de 1 mantendrá la velocidad constante. Un valor como 1,1 debería aumentar la velocidad lo suficiente como para que el juego sea desafiante pero no imposible. 

Finalmente, para inicializar los valores de los atributos que cambian a lo largo del juego llamamos  {\ttfamily inicialize\_configuracion\_dinamica()}:



\hspace{-2.75cm}{\ttfamily \scriptsize configurar.py}  \vspace{-0.5cm}
\begin{lstlisting}	
    # configuracion  dinámica
    def inicialize_configuracion_dinamica(self):
        """Inicializa la configuración de dinámica del juego."""
        self.velocidad_nave = 6
        self.velocidad_balas = 1.0
        self.velocidad_alien = 1.0 
        # Dirección de flota: 1 representa derecha;
        # -1 representa izquierda.
        self.direccion_flota = 1
\end{lstlisting}

Este método establece los valores iniciales para las velocidades de la nave, la bala y el alienígena. Aumentaremos estas 
velocidades a medida que el jugador avance en el juego y las restableceremos cada vez que el jugador comience un nuevo juego. 
Incluimos {\ttfamily direccion\_flota } en este método para que los alienígenas siempre se muevan justo al comienzo de un nuevo juego. 

No necesitamos aumentar el valor de {\ttfamily self.flota\_cae\_velocidad = 5}, porque cuando los extraterrestres se muevan más rápido por la pantalla, también vendrán bajar la pantalla más rápido.

Para aumentar la velocidad de la nave, de las balas y de los alienígenas cada vez que el
jugador alcanza un nuevo nivel, escribiremos un nuevo método llamado  {\ttfamily incrementa\_velocidad()}:

\hspace{-2.75cm}{\ttfamily \scriptsize configurar.py}  \vspace{-0.5cm}
\begin{lstlisting}	
    def incrementar_velocidad(self):
        """Configuaracion para incremento de velocidad."""
        self.velociodad_nave *= self.velocidad_sube_escala
        self.velocidad_balas *= self.velocidad_sube_escala
        self.velcidad_alien *=  self.velocidad_sube_escala
\end{lstlisting}

Para aumentar la velocidad de estos elementos, multiplicamos cada uno  de ellos por el valor 
de {\ttfamily self.velocidad\_sube\_escala}. 
Aumentamos el tiempo del juego llamando a  {\ttfamily incrementa\_velocidad()} en {\ttfamily \_chequear\_colision\_balas\_alien()} cuando el último alienígena de una flota ha sido derribado:


\hspace{-2.4cm}{\ttfamily \scriptsize invasión}

\vspace{-0.2cm}
\hspace{-2.1cm}{\ttfamily \scriptsize \_alien.py}\vspace{-0.84cm}
\begin{lstlisting}	
    def _chequear_colisiones_balas_alien(self):
        --recorte--
        if not self.aliens:
            # Se destruye existencia de balas y se crea una nueva flota.
            self.balas.empty()
            self._creear_flota()
            self.configuracion.incrementar.velocidad()
\end{lstlisting}

cambiando los valores de velocidad  de la nave, del alienígena, y de las balas,  es suficiente para subir el nivel del juego. 

\subsubsection{Resetear la velocidad}

Ahora necesitamos retornar cualquier cambio en la configuración de sus valores iniciales cada vez que comience un nuevo juego; 
por otro lado en cada tiempo de juego; por otro lado, cada juego nuevo  debe comenzar con un incremento de velocidad configurado previamente:

\hspace{-2.4cm}{\ttfamily \scriptsize invasión}

\vspace{-0.2cm}
\hspace{-2.1cm}{\ttfamily \scriptsize \_alien.py}\vspace{-0.84cm}
\begin{lstlisting}
    def _chequear_play_boton(self, mouse_pos):
        """Comienza un nuevo juego cuando el jugador hace clic en Play."""
        boton_clicked = self.play_boton.rect.collidepoint(mouse_pos)
        if boton_clicked and not self.juego_activo: 
            # resetea la configuracion del juego
            self.configuracion.inicialize_configuracion_dinamica()
        --recorte--
\end{lstlisting}



Y ahora jugar  Invasión Alienígena debería ser más divertido y desafiante.
Cada vez que inicia un nuevo juego, debe acelerarse y volverse rápido y por supuesto más difícil.
Si el juego se vuelve demasiado difícil, muy  rápido, disminuya los valores de {\ttfamily self.velocidad\_sube\_escala}, o si el juego no es lo suficientemente desafiante, aumente ligeramente el valor. 

Encuentra un punto óptimo aumentando la dificultad en una cantidad de tiempo razonable. Los primeros dos juegos deben ser fáciles, los siguientes deberían ser desafiantes pero factibles, y luego increíblemente difícil.


\section{INTÉNTALO TÚ MISMO}

14-3. Práctica de tiro desafiante: comience con su trabajo del ejercicio 14-2 (página 283). Haz que el objetivo se mueva más rápido a medida que avanza el juego y reinicia el objetivo a la velocidad original cuando el jugador hace clic en Reproducir.\\

14-4. Niveles de dificultad: crea un conjunto de botones para Alien Invasion que permitan al jugador seleccionar un nivel de dificultad inicial apropiado para el juego. Cada botón debe asignar los valores apropiados para los atributos en Configuración necesarios para crear diferentes niveles de dificultad.


\section{Puntuación} 

Implementemos un sistema de puntuación para realizar un seguimiento de la puntuación del juego en tiempo real y mostrar 
la puntuación más alta, el nivel y la cantidad de naves restantes. La puntuación es una estadística del juego, por lo que 
agregaremos un atributo de puntuación en el módulo  {\ttfamily Stats\_Juego}:



\hspace{-2.4cm}{\ttfamily \scriptsize stats}

\vspace{-0.2cm}
\hspace{-2.1cm}{\ttfamily \scriptsize \_juego.py}\vspace{-0.84cm}
\begin{lstlisting}	
class Estadisticas_Juego():
    """Track estadísticas parael juego Invasión Alienígena."""
    def __init__(self, ia_jugar):
        --recorte:
        
    def resetear_stats(self):
        """Inicializa los estadísticas  que cambian durante el juego .""" 
        self.nave_permitidas = self.configuracion.limite_naves
        self.puntuacion = 0
\end{lstlisting}

Para restablecer la puntuación cada vez que comienza un nuevo juego, inicializamos la puntuación en  {\ttfamily  resetear\_stats()}
en lugar de {\ttfamily \_\_init\_\_()}. Visualización de la puntuación Para mostrar la puntuación en la pantalla, primero creamos
 una nueva clase, Marcador. Por ahora, esta clase solo mostrará la puntuación actual. Con el tiempo, también lo
  usaremos para informar la puntuación más alta, el nivel y la cantidad de naves restantes. Aquí está la primera
   parte de la clase; guárdelo como marcador.py:



\hspace{-2.9cm}{\ttfamily \scriptsize puntuacion.py}  \vspace{-0.5cm}
\begin{lstlisting}
import pygame.font
class Puntuacion: 
    """Una clase para reportar el informacion de la puntuacion."""
    def __init__(self, ia_jugar): 
        """Inicializa atributo de puntaje.""" 
        self.pantalla = ia_jugar.pantalla
        self.pantalla_rectang = self.pantalla.get_rect() 
        self.configuracion = ia_jugar.configuracion 
        self.estadisticas =  ia_jugar.stats 
        # Configuarcion de fuente para informar puntaje scoring
        self.color_texto =  (30, 30, 30) 
        self.fuente = pygame.font.SysFont(None, 48)
        # Prepara la imagen inicial de la puntuacion
        preparacion_puntuacion()
\end{lstlisting}

Debido a que {\ttfamily puntuacion.py} escribe texto en la pantalla, comenzamos importando el
módulo {\ttfamily pygame.font}. 

A continuación, le damos a {\ttfamily \_\_init\_\_()} el parámetro {\ttfamily ia\_jugar} para que pueda
acceder a los objetos de configuración, pantalla y estadísticas, que necesitará para informar el
valores que estamos rastreando. Luego configuramos un color de texto y creamos una instancia de una fuente
objeto.

Para convertir el texto que se mostrará en una imagen, llamamos a {\ttfamily preparacion\_puntuacion()},
que definimos aquí:

\hspace{-2.9cm}{\ttfamily \scriptsize puntuacion.py}  \vspace{-0.5cm}
\begin{lstlisting}		
    def preparacion_puntuacion(self):
        """retorna la puntuacion dentro de la imagen renderizada."""
        puntuacion_str = str(self.estadisticos.puntuacion)
        self.puntuacion_imagen = self.fuente.render(puntuacion_str, True,
                                    self.color_texto, 
                                    self.configuracion.color_fondo)
        # Muestra los puntos en la parte superior izquierda de la pantalla.
        self.rectang_puntuacion = self.puntuacion_imagen.get_rect()
        self.rectang_puntuacion.right = self.pantalla_rectang.right - 20
        self.rectang_puntuacion.top = 20
\end{lstlisting}


En el método {\ttfamily preparacion\_puntuacion()}, convertimos el valor numérico de {\ttfamily stats.puntaje} en una cadena
y luego pasa esta cadena a  la función {\ttfamily render()}, para que crear la imagen.
Para mostrar la puntuación clara en pantalla, pasamos el color de fondo de pantalla y el color del texto a {\ttfamily renderizar()}.

Colocaremos la información  en la esquina superior derecha de la pantalla y  que se expanda hacia la izquierda a medida que aumenta la puntuación y el crece ancho del número. 
Para asegurarse de que la puntuación siempre se alinee con el lado derecho de
la pantalla, creamos un {\ttfamily rectángulo} llamado {\ttfamily rectang\_puntuacion} y configuramos su borde derecho en 20 píxeles
desde el borde derecho de la pantalla 4. Luego colocamos el borde superior a 20 píxeles
hacia abajo desde la parte superior de la pantalla 5.

Luego creamos un método {\ttfamily mostrar\_puntuacion()} para mostrar la puntuación  renderizada.
imagen:

\hspace{-2.9cm}{\ttfamily \scriptsize puntuacion.py}  \vspace{-0.5cm}
\begin{lstlisting}
    def mostrar_puntuacion(self):
        """Dibujar puntuacion en pantalla."""
        self.pantalla.blit(self.puntuacion_imagen, self.rectang_puntuacion)
\end{lstlisting}

Este método dibuja la imagen de la puntuación en la pantalla en la ubicación{\ttfamily mostrar\_puntuacion()}
especifica.

\subsubsection{El marcador en pantalla}

Para mostrar la puntuación, crearemos una instancia de Marcador en {\ttfamily AlienInvasion}. Primero,
actualicemos las declaraciones de importación:

\hspace{-2.4cm}{\ttfamily \scriptsize invasión}

\vspace{-0.2cm}
\hspace{-2.1cm}{\ttfamily \scriptsize \_alien.py}\vspace{-0.84cm}
\begin{lstlisting}		
--recorte--
from stats_juegos import Estadisticas_Juego
from puntuacion import Puntuacion
--snip--
\end{lstlisting}

 

Ahora, creamos una instancia de {\ttfamily Puntuacion} en {\ttfamily \_\_init\_\_()} en la clase {\ttfamily AlienInvasion}:



\hspace{-2.4cm}{\ttfamily \scriptsize invasión}

\vspace{-0.2cm}
\hspace{-2.1cm}{\ttfamily \scriptsize \_alien.py}\vspace{-0.84cm}
\begin{lstlisting}
    def __init__(self): 
        --recorte-- 
        pygame.display.set_caption("Invasión Alienígena") 
        # Crea una instancia para almacenar estadisticos del juego 
        # y crea un puntuacion en  pantalla
        self.estadisticas = Estadisticas_Juego(self) 
        self.Pp = Puntuacion(self)
        --recorte--
\end{lstlisting}

Y  el dibujo de puntuación sobre pantalla en {\tt actualizar\_pantalla}:


\hspace{-2.4cm}{\ttfamily \scriptsize invasión}

\vspace{-0.2cm}
\hspace{-2.1cm}{\ttfamily \scriptsize \_alien.py}\vspace{-0.84cm}
\begin{lstlisting}
    def actualizar_pantalla(self): 
        --recorte-- 
        self.aliens.draw(self.pantalla) 
        # Dibuja la puntuacion en pantalla Pp
        self.Pp.mostrar_puntuacion() 
\end{lstlisting}


Llamamos a {\ttfamily mostrar\_puntuacion()} justo antes de dibujar el botón Play. 

Cuando ejecutes Invasión Alienígena debe aparecer un 0 en la parte superior derecha de la pantalla 
(en este punto, solo queremos asegurarnos de que la puntuación aparezca en el lugar correcto antes de continuar  desarrollar el sistema de puntuación). 


La Figura \ref{puntuacion01} muestra la puntuación tal como aparece antes de que comience el juego. A continuación, asignaremos valores de puntos a cada extraterrestre!





\begin{figure}[h!]
	\centering
	\includegraphics[scale=0.35]{../A Logos de Python/puntuacion01}
	\caption{La puntuación en la esquina superior derecha.}
	\label{puntuacion01}
\end{figure}



\subsubsection{Actualización de la puntuación los alienígenas}

Para escribir una puntuación en vivo en pantalla, actualizamos el valor de stats.score cada vez que
se golpea a un extraterrestre y luego se llama a {\ttfamily preparacion\_puntuacion()} para actualizar la imagen de la puntuación. Pero
Primero, determinemos cuántos puntos obtiene un jugador cada vez que dispara.
derribar a un extraterrestre:

\hspace{-2.75cm}{\ttfamily \scriptsize configurar.py}  \vspace{-0.5cm}
\begin{lstlisting}	
    def inicialize_configuracion_dinamica(self):
        """Inicializa la configuración de dinámica del juego."""
        --recorte--
        # Configuracion de puntuacion
        self.puntos_alien = 50
\end{lstlisting}

Aumentaremos el valor de puntuación cada alienígena a medida que avance el juego. Para asegurar que éste el valor de puntuación  se restablezca cada vez que comienza un nuevo juego, configuramos el valor en el método {\ttfamily inicialize\_configuracion\_dinamica()}.

Actualicemos la puntuación en {\ttfamily \_chequear\_colisiones\_balas\_alien()} cada vez que
alienígena es derribado:

\hspace{-2.4cm}{\ttfamily \scriptsize invasión}

\vspace{-0.2cm}
\hspace{-2.1cm}{\ttfamily \scriptsize \_alien.py}\vspace{-0.84cm}
\begin{lstlisting}       
    def _chequear_colisiones_balas_alien(self):
        """Responder a las colisiones balas-alienigena."""
        # Remover cualquier bala y alienigena que han colisionado.
        colisiones = pygame.sprite.groupcollide(
                        self.balas, self.aliens, True, True)

        if colisiones:
            self.estadisticas.puntuacion += self.configuracion.puntos_alien 
            self.Pp.preparacion_puntuacion()
            --recorte--  
\end{lstlisting}

Cuando una bala alcanza a un extraterrestre, Pygame devuelve un diccionario de colisiones.
Comprobamos si el diccionario existe, y si existe, el valor del alienígena es
añadido a la puntuación. Luego llamamos al método {\ttfamily preparacion\_puntuacion()} para crear una nueva imagen de la
puntuación actualizada.

¡Ahora, cuando juegues Invasión  Alienígena, debe acumular puntos por cada alienígena que derribe.

\subsubsection{Restablecer la puntuación}

En este momento, solo estamos preparando una nueva puntuación después de que un extraterrestre haya sido derribado
que funciona durante la mayor parte del juego. Pero cuando comenzamos un nuevo juego, todavía
se ve puntuación del juego anterior hasta que derribes al primer alienígena.

Podemos solucionar este problema preparando la puntuación al iniciar un nuevo juego:

\hspace{-2.4cm}{\ttfamily \scriptsize invasión}

\vspace{-0.2cm}
\hspace{-2.1cm}{\ttfamily \scriptsize \_alien.py}\vspace{-0.84cm} 
\begin{lstlisting}
    def _chequear_play_boton(self, mouse_pos):
        """Comienza un nuevo juego cuando el jugador hace clic en Play."""
        boton_clicked = self.play_boton.rect.collidepoint(mouse_pos)
		--recorte--
        # Restablecer las estadísticas del juego.
        self.estadisticas.resetear_stats()
        self.Pp.preparacion_puntuacion()
        --recorte--
\end{lstlisting}



Llamamos el método {\ttfamily preparacion\_puntuacion()} después de resetear las estadísticas del juego al iniciar un juego nuevo. 
Esto prepara el marcador con una puntuación de 0. 
Asegurándose de anotar todos los aciertos tal como está escrito actualmente,
nuestro código podría no puntuar para algunos extraterrestres. 
Por ejemplo, si dos balas chocan con alienígenas durante el mismo paso por el bucle o si hacemos una bala extra ancha para impactar a varios alienígenas, el jugador solo recibirá puntos por impactar a uno de los alienígenas.  Para solucionar este problema, reorientamos  la forma en que se detectan las colisiones entre balas y extraterrestres. 

En  {\ttfamily \_chequear\_colisiones\_balas\_alien(self)}, cualquier bala que colisione con un extraterrestre se convierte en una clave en el diccionario de colisiones. El valor asociado con cada bala es una lista de extraterrestres con los que chocó. Revisamos los valores en el diccionario de colisiones para asegurarnos de otorgar puntos por cada impacto alienígena:


\hspace{-2.4cm}{\ttfamily \scriptsize invasión}

\vspace{-0.2cm}
\hspace{-2.1cm}{\ttfamily \scriptsize \_alien.py}\vspace{-0.84cm} 
\begin{lstlisting}		
    def _check_bullet_alien_collisions(self):
        --recorte--
        if colisiones: 
            for aliens in colisiones.values(): 
                self.estadisticas.puntuacion += 
                    self.configuracion.puntos_alien * len(aliens)
                self.Pp.preparacion_puntuacion()     
        --recorte--
\end{lstlisting}

Si se ha definido el diccionario de colisiones, recuerda que cada valor del diccionario es una lista de alienígenas alcanzados 
por una sola bala. Multiplicamos el valor de cada alienígenas por el número de  alienígenas en cada lista y sumamos esta cantidad a la puntuación actual. Para probar esto, cambia el ancho de una bala a 300 píxeles y verifica que recibes puntos por cada alienígena que 
derribes con tus balas extra anchas; luego regresa el ancho de la bala a su valor normal. 

Aumento del valor de los puntos debido a que el juego se vuelve más difícil cada vez que un jugador alcanza un nuevo nivel,
los alienígenas en niveles posteriores deberían valer más puntos. Para implementar esta funcionalidad, agregaremos código para aumentar el valor de los puntos cuando aumente la velocidad del juego: settings.pyclass Configuración:


\hspace{-2.75cm}{\ttfamily \scriptsize configurar.py}  \vspace{-0.5cm}
\begin{lstlisting}	
class Configuracion:
    """Una clase para almacenar todas las configuraciones de 
    Invasión Alienígena.""" 
    def __init__(self):
        """Inicializamos la configuracion de estatica del juego."""
        #  Configuración de pantalla    

        --recorte--
        # Incremento de la velocidad del juego
        self.velocidad_sube_escala = 1.1
    
        # Incremento de los valores puntuacion de los alienigenas  
        self.puntuacion_escale = 1.5 
    
        self.initialize_dynamic_settings() 
    
        def initialize_dynamic_settings(self): --snip--
            --recorte--	

    def incrementar_velocidad(self):
        """Configuaracion para incremento de velocidad, y la puntuacion por
        alienigena."""
        self.velocidad_nave *= self.velocidad_sube_escala
        self.velocidad_balas *= self.velocidad_sube_escala
        self.velocidad_alien *=  self.velocidad_sube_escala
        self.puntos_alien = int(self.puntos_alien  = 50 * 
                                   self.velocidad_sube_escala)	
\end{lstlisting}


Definimos una tasa {\ttfamily puntuación\_escala}, para aumentar la puntuación. Un pequeño aumento en la 
velocidad (1.1) hace que el juego sea un poco más rápido. Sin embargo, para ver una diferencia  notable usamos el valor de la tasa   en 1,5. 
 
Ahora cuando aumentamos la velocidad del juego, también aumentamos el valor en puntos de cada nave derribada.
Usamos la función {\ttfamily int()} para aumentar el valor del punto en números enteros. 
Para ver el valor de cada alienígena, agregue una llamada {\ttfamily print()} al incremento\_speed()
método en Configuración:

\hspace{-2.75cm}{\ttfamily \scriptsize configurar.py}  \vspace{-0.5cm}
\begin{lstlisting}
    def increase_speed(self):
        --recorte--	
       self.puntos_alien = int(self.puntos_alien  = 50 * 
                                   self.velocidad_sube_escala)	
       print(self.puntos_alien)
\end{lstlisting}


El nuevo valor de puntos  aparece en la terminal cada vez que alcanzar un nuevo nivel.


Después de verificar el valor de puntos está aumentando, debe eliminar la llamada {\ttfamily print()}, que puede afectar el 
rendimiento del juego y distraer al jugador.

\subsubsection{Redondeando la puntuación}

La mayoría de los juegos del estilo de disparos  reportan puntuaciones como múltiplos de 10, así que sigamos
con nuestras puntuaciones. Además, formateemos la información  para incluir comas separadores en grandes cantidades. 
Haremos este cambio en el Marcador:


\hspace{-2.9cm}{\ttfamily \scriptsize puntuacion.py}  \vspace{-0.5cm}
\begin{lstlisting}		
    def preparacion_puntuacion(self):
        """retorna la puntuacion dentro de la imagen renderizada."""
        redondeo_puntuacion = round(self.estadisticas.puntuacion, -1)
        puntuacion_str = f"{redondeo__puntuacion:,}"
        self.puntuacion_imagen = self.fuente.render(puntuacion_str, True,
                                    self.color_texto, 
                                    self.configuracion.color_fondo) 
        --recorte--	
\end{lstlisting}



La función  {\ttfamily round()} normalmente redondea un flotante a un número determinado de decimales
dados como segundo argumento. Sin embargo, cuando pasas un negativo número como segundo argumento, 
función  {\ttfamily round()} redondeará el valor al 10, 100, 1000, etc. más cercanos. Este código le dice a Python que redondee el valor
de {\ttfamily estadisticas.puntuacion} al 10 más cercano y asígnelo a {\ttfamily redondeo\_puntuacion}.




\begin{figure}[h!]
	\centering
	\includegraphics[scale=0.35]{../A Logos de Python/comas}
	\caption{Puntuación redondeada con separadores de coma.}
	\label{comas}
\end{figure}



Luego usamos el formato de cadena {\ttfamily f-string} para la puntuación. Un formato
especificador es una secuencia especial de caracteres que modifica la forma en que una variable
se presenta el valor. En este caso la secuencia: le dice a Python que inserte comas en los lugares apropiados 
del valor numérico proporcionado. Esto da como resultado cadenas como 1,000,000 en lugar de 1000000. 

Ahora, cuando ejecutes el juego, debes ver una puntuación redondeada y perfectamente formateada incluso cuando 
acumule muchos puntos, como se muestra en la Figura Figura \ref{comas}.   



\subsubsection{Puntuaciones altas}

Cada jugador quiere superar la puntuación más alta de un juego, así que hagamos un seguimiento e informemos las puntuaciones altas para darles a los jugadores algo por lo que trabajar.

Almacenaremos puntuaciones altas en {\ttfamily stats\_juego.py}:
%\hspace{-2.5cm}{\ttfamily \scriptsize stats\_juego.py}  \vspace{-0.5cm}

\hspace{-2.4cm}{\ttfamily \scriptsize stats}

\vspace{-0.2cm}
\hspace{-2.1cm}{\ttfamily \scriptsize \_juego.py}\vspace{-0.84cm}
\begin{lstlisting}
    def __init__(self, ia_jugar):
        --recorte--	
        # Puntuacione altas nunca debe ser eliminadas. 
        self.puntuacion_alta = 0	
\end{lstlisting}
Debido a que la puntuación más alta nunca debe ser eliminada,  inicializamos {\ttfamily alta\_puntuacion} en {\ttfamily\_\_init\_\_()} en lugar de {\ttfamily resetear\_stats()}. 


A continuación, modificaremos el marcador para mostrar la puntuación más alta. 


\hspace{-2.9cm}{\ttfamily \scriptsize puntuacion.py}  \vspace{-0.5cm}
\begin{lstlisting}		
    def __init__(self, ia_jugar):
        """Inicializa atributo de puntaje.""" 
        --recorte--	
	    
	    # Prepara la imagen de puntuacion inicial. 
	    self.preparacion_puntuacion() 
	    self.preparacion_alta_puntuaciones()	
\end{lstlisting}

La puntuación más alta se mostrará por separado de la puntuación, por lo que necesitamos un nuevo método, 
{\ttfamily preparacion\_alta\_puntuacion}, para preparar la imagen de puntuación más alta.
	  
	  
\hspace{-2.9cm}{\ttfamily \scriptsize puntuacion.py}  \vspace{-0.5cm}
\begin{lstlisting}		
    def preparacion_alta_puntuaciones(self):
        """Retorna la puntuacion alta dentro de magen renderizada."""
        puntuacion_alta = round(self.estadisticas.puntuacion_alta, -1)
        puntuacion_alta_str = f"{puntuacion_alta:,}"
        self.puntuacion_alta_imagen = self.font.render(puntuacion_alta_str,
                                          True,
                                          self.color_texto,
                                          self.configuracion.color_fondo)
        # puntuacion_alta al centro y en la parte superior de la pantalla.
        self.puntuacion_alta_rectang = self.puntuacion_alta_imagen.get_rect()
        self.puntuacion_alta_rectang.centerx = self.pantalla_rectang.centerx
        self.puntuacion_alta_rectang.top = pantalla_rectang.top
\end{lstlisting}

Redondeamos la puntuación más alta a la decena más cercana y la formateamos con comas.
Luego generamos una imagen a partir de la puntuación más alta, centramos la
puntuación más alta en el rectángulo  horizontalmente y establece su atributo superior para que coincida con el superior
de la imagen de la puntuación.

El método {\ttfamily mostrar\_puntuacion()} ahora dibuja la puntuación actual en la parte superior derecha
y la puntuación más alta en la parte superior central de la pantalla:

\hspace{-2.9cm}{\ttfamily \scriptsize puntuacion.py}  \vspace{-0.5cm}
\begin{lstlisting}		
    def mostrar_puntuacion(self):  # dibuja o muestra
        """Dibuja puntuacion en pantalla."""
        self.pantalla.blit(self.puntuacion_imagen, self.rectang_puntuacion)
        self.screen.blit(self.alta_puntuacion_imagen, 
                             self.rectang_puntuacion_alta)
\end{lstlisting}
Para comprobar puntuaciones altas, escribiremos en {\ttfamily puntuacion.py} un nuevo método, {\ttfamily chequear\_altas\_puntuaciones()}:  %(marcador de puntuación)

\hspace{-2.9cm}{\ttfamily \scriptsize puntuacion.py}  \vspace{-0.5cm}
\begin{lstlisting}	
    def chequear_altas_puntuaciones(self):
        """Compruebe para ver si hay una puntuaciones altas nuevas"""
        if self.estadisticas.puntuacion > self.estadisticas.puntuacion_alta:
            self.estadisticas.puntuacion_alta = self.estadisticas.puntuacion	
            self.preparacion_alta_puntuaciones()
\end{lstlisting}
% puntuacion_alta



El método {\ttfamily chequear\_altas\_puntuaciones()} compara la puntuación actual con la
puntuación más alta. Si la puntuación actual es mayor, actualizamos el valor de {\ttfamily puntuacion\_alta} 
y llame  {\ttfamily preparacion\_alta\_puntuacion()} para actualizar la imagen de la puntuación más alta.

Necesitamos llamar a  {\ttfamily chequear\_altas\_puntuaciones()}  cada vez que un extraterrestre sea derribado  después de la actualización.
la puntuación en  {\ttfamily \_chequear\_colisiones\_balas\_alien()}:


\hspace{-2.4cm}{\ttfamily \scriptsize invasión}

\vspace{-0.2cm}
\hspace{-2.1cm}{\ttfamily \scriptsize \_alien.py}\vspace{-0.84cm} 
\begin{lstlisting}		
    def _chequear_colisiones_balas_aliens(self):	
        --recorte--
        if colisiones:
            for aliens in colisiones.values():
                self.estadisticas.puntuacion +=  self.\
                         configuracion.puntos_alien * len(aliens)
                self.Pp.preparacion_puntuacion()     
                self.Pp.chequear_altas_puntuaciones()
                print(self.Pp.chequear_altas_puntuaciones() )    
        --recorte--
\end{lstlisting}


Llamamos a  {\ttfamily chequear\_altas\_puntuaciones} cuando el diccionario de colisiones está presente, y
lo hacemos después de actualizar la puntuación de todos los alienígenas que han sido alcanzados.

La primera vez que juegues  Invasión Alienígena, tu puntuación será la puntuación más alta, por lo que se mostrará como la puntuación 
actual como la puntuación más alta. Pero cuando comienzas un segundo juego, tu puntuación más alta debe aparecer en el medio y tu puntuación actual debería aparecer a la derecha, como se muestra en la Figura \ref{altapuntuacion}.



\begin{figure}[h!]
	\centering
	\includegraphics[scale=0.35]{../A Logos de Python/altapuntuacion}
	\caption{En la parte superior central de la}\vspace{-0.35cm}
    \caption*{pantalla se muestra la puntuación más alta.}
	\label{altapuntuacion}
\end{figure}





\subsubsection{Visualización del nivel}

Para mostrar el nivel del jugador, necesitamos definir un atributo en {\ttfamily Estadisticas\_Juego()}, que represente el nivel actual. 
Se inicializa en {\ttfamily resetear\_stats()} para restablecer el nivel al inicio de cada juego nuevo:




\hspace{-2.4cm}{\ttfamily \scriptsize stats}

\vspace{-0.2cm}
\hspace{-2.1cm}{\ttfamily \scriptsize \_juego.py}\vspace{-0.84cm}
\begin{lstlisting}
    def resetear_stats(self):
        """Inicializa estadisticas que cambian durante el juego"""
        self.nave_permitidas = self.configuracion.limite_naves
        self.puntuacion = 0
        self.nivel = 1
\end{lstlisting}
Para mostrar el nivel actual, creamos un nuevo método {\ttfamily  preparando\_nivel()}, en {\ttfamily puntuacion.py}:



 
\hspace{-2.9cm}{\ttfamily \scriptsize puntuacion.py}  \vspace{-0.5cm}
\begin{lstlisting}		
    def preparando_nivel(self): 
        """Retorna el nivel en una imagen renderizada.""" 
        nivel_str = str(self.estadisticas.nivel)
        self.nivel_imagen = self.fuente.render(nivel_str,
                                True, self.color_texto, 
                                self.configuracion.color_fondo)
        
        # Posicion del nivel bajo la puntuacion
        self.nivel_rectang = self.nivel_imagen.get_rect()
        self.nivel_rectang.right = self.rectang_puntuacion.right
        self.nivel_rectang.top =   self.rectang_puntuacion.bottom + 10	
\end{lstlisting}

Y lo llamamos en {\ttfamily  \_\_init\_\_(self, ia\_jugar)} 

\hspace{-2.9cm}{\ttfamily \scriptsize puntuacion.py}  \vspace{-0.5cm}
\begin{lstlisting}	
	def __init__(self, ia_jugar):
        --recorte-- 
	    self.preparando_puntuacion_alta() 
	    self.preparando_nivel()	
\end{lstlisting}

El método {\ttfamily  preparando\_nivel()} crea una imagen a partir del valor almacenado en
{\ttfamily estadisticas.nivel} y establece el atributo {\ttfamily rectangulo } de la imagen para que coincida con el atributo {\ttfamily rect} de la puntuación.  Luego establece el atributo superior a 10 píxeles debajo de la parte inferior de
la imagen de la puntuación  para dejar espacio entre la  puntuación y el nivel.
Y actualizamos el método  {\ttfamily mostrar\_puntuacion()}:


\hspace{-2.9cm}{\ttfamily \scriptsize puntuacion.py}  \vspace{-0.5cm}
\begin{lstlisting}	
    def mostrar_puntuacion(self):
        """Dibuja puntuacion y nivel en la pantalla."""
        self.pantalla.blit(self.puntuacion_imagen, self.rectang_puntuacion)
        self.pantalla.blit(self.puntuacion_alta_imagen, 
                            self.puntuacion_alta_rectang )
        self.pantalla.blit(self.nivel_imagen,  self.nivel_rectang)
\end{lstlisting}

Esta nueva línea de código dibuja la imagen del nivel en la pantalla.


Incrementamos {\ttfamily estadisticas.nivel} y actualizamos la imagen del nivel en {\ttfamily \_chequear\_colisiones\_balas\_alien()}:


\hspace{-2.4cm}{\ttfamily \scriptsize invasión}

\vspace{-0.2cm}
\hspace{-2.1cm}{\ttfamily \scriptsize \_alien.py}\vspace{-0.84cm}
\begin{lstlisting}		
    def _chequear_colisiones_balas_alien(self):
        --recorte--
        if not self.aliens:
            # Destruye la existencia de las balas y crea una nueva flota.
            self.balas.empty()
            self._crear_flota()
            self.configuracion.incrementar_velocidad()
            # Increar nivel
            self.estadisticas.nivel += 1
            self.Pp.preparando_nivel()	
\end{lstlisting}

Si una flota es destruida, incrementamos el valor de stats.level y llamamos
prep\_level() para asegurarse de que el nuevo nivel se muestre correctamente.
Para garantizar que la imagen del nivel se actualice correctamente al comienzo de un nuevo juego,
También llamamos a prep\_level() cuando el jugador hace clic en el botón Reproducir:


\hspace{-2.4cm}{\ttfamily \scriptsize invasión}

\vspace{-0.2cm}
\hspace{-2.1cm}{\ttfamily \scriptsize \_alien.py}\vspace{-0.84cm}
\begin{lstlisting}		
    def _chequear_play_boton(self, mouse_pos):
        --recorte--
        if boton_clicked and not self.juego_activo: 
            --recorte--
            self.Pp.preparacion_puntuacion()
            self.Pp.preparando_nivel()            
            --recorte--
\end{lstlisting}

Llamamos a {\ttfamily self.Pp.preparando\_nivel()}  justo después de llamar a {\ttfamily self.Pp.preparacion\_puntuacion()}.

Ahora puedes ver cuántos niveles has completado; el nivel actual aparece justo debajo de la puntuación actual,  como se muestra en Figura \ref{nivel}. 

\begin{figure}[h!]
	\centering
	\includegraphics[scale=0.35]{../A Logos de Python/nivel}
	\caption{El nivel actual aparece justo }\vspace{-0.35cm}
	\caption*{debajo de la puntuación actual.}
	\label{nivel}
\end{figure}




NOTA 

En algunos juegos clásicos, las puntuaciones tienen etiquetas, como Puntuación, Puntuación alta y Nivel.
Hemos omitido estas etiquetas porque el significado de cada número queda claro una vez que has jugado. 
Para incluir estas etiquetas, agréguelas a las cadenas de puntuación justo antes de las llamadas a font.render() en el marcador. 




\subsection{Visualización del número de naves}

Finalmente, mostremos el número de naves que le quedan al jugador, pero esta vez usemos un gráfico. Para hacerlo, dibujaremos naves en la esquina superior izquierda de la pantalla para representar cuántos naves quedan, tal como lo hacen muchos juegos de arcade clásicos.
Primero, debemos hacer que la nave herede de {\tt  Sprite} para poder crear un grupo:


\hspace{-2.5cm}{\ttfamily \scriptsize nave.py}  \vspace{-0.5cm}
\begin{lstlisting}		
import pygame 
from pygame.sprite import Sprite

class Nave(Sprite):
    """Una clase que maneja el comportamiento de la nave"""
    def __init__(self, ia_jugar):
        """Inicia la nave y configura la posicion inicial."""
         super().__init__()
         --recorte--
\end{lstlisting}

Aquí importamos  {\ttfamily  Sprite}, nos aseguramos de que la nave herede de {\ttfamily Sprite} y llamamos a {\ttfamily super()} al comienzo de {\ttfamily \_\_init\_\_()}.


Y ahora debemos modificar la puntuación para crear un grupo de naves que podamos mostrar:

\hspace{-2.9cm}{\ttfamily \scriptsize puntuacion.py}  \vspace{-0.5cm}
\begin{lstlisting}		
import pygame.font
from pygame.sprite import Group
from nave import Nave	
\end{lstlisting}


Debido a que estamos creando un grupo de naves, importamos el grupo y la nave.

\hspace{-2.9cm}{\ttfamily \scriptsize puntuacion.py}  \vspace{-0.5cm}
\begin{lstlisting}		
def __init__(self, ia_jugar):
    """Initializamos los atributos de la puntuacion guardada."""
    self.ia_jugar = ia_jugar
    self.pantalla = ia_jugar.pantalla
    --recorte--
    self.preparando_nivel()
    self.preparando_nave()	
\end{lstlisting}

Asignamos la instancia del juego a un atributo, porque la necesitaremos para
crear algunas naves. Llamamos a {\ttfamily preparando\_nave()} después de la llamada a {\tt preparando\_nivel() }.

Creamos    {\ttfamily preparando\_nave()} 

  
\hspace{-2.9cm}{\ttfamily \scriptsize puntuacion.py}  \vspace{-0.5cm}
\begin{lstlisting}		
    def preparando_nave(self):
        """Muestra numero de naves permitidas."""
        self.naves = Group()
        for numero_nave in range(self.estadisticas.naves_permitidas):
            nave = Nave(self.ia_jugar)
            nave.rect.x = 10 + numero_naves * nave.rect.width
            nave.rect.y = 10
            self.nave.add(nave)	
\end{lstlisting}

El método {\ttfamily preparando\_nave()} crea un grupo vacío, {\ttfamily self.naves}, para almacenar las naves 
(instancia de Nave). Para llenar este grupo, se ejecuta un bucle una vez por cada nave que el jugador ha dejado. 
Dentro del bucle, creamos una nueva nave y configuramos el valor de la coordenada $x$  de las naves para que  aparezcan uno al lado del otro con un espacio de 10 píxeles de margen en el lado izquierdo del grupo de naves3. Establecemos la coordenada $y$ con 
un valor de 10 píxeles hacia abajo desde la parte superior de la pantalla para que las naves aparezcan en la
esquina superior izquierda de la pantalla. Luego agregamos cada nave nueva al grupo.


Ahora necesitamos dibujar las naves en la pantalla:

\hspace{-2.9cm}{\ttfamily \scriptsize puntuacion.py}  \vspace{-0.5cm}
\begin{lstlisting}		
    def mostrar_puntuacion(self):
        """Dibuja puntuaciones, niveles y naves en la pantalla."""
        self.pantalla.blit(self.puntuacion_imagen, self.rectang_puntuacion)
        self.pantalla.blit(self.puntuacion_alta_imagen, self.puntuacion_alta_rectang )
        self.pantalla.blit(self.nivel_imagen,  self.nivel_rectang)        
        self.nave.draw(self.pantalla)	
\end{lstlisting}

Para mostrar las naves en la pantalla, llamamos a {\ttfamily draw()} en el grupo, y Pygame dibuja cada nave. 
Para mostrarle al jugador con cuántos naves tiene para comenzar, llamamos a  {\ttfamily preparando\_nave()} cuando 
comienza un nuevo juego. Hacemos esto en {\ttfamily \_chequear\_play\_boton()} en {\ttfamily AlienInvasion}:

\hspace{-2.4cm}{\ttfamily \scriptsize invasión}

\vspace{-0.2cm}
\hspace{-2.1cm}{\ttfamily \scriptsize \_alien.py}\vspace{-0.84cm}
\begin{lstlisting}	
    def _chequear_play_boton(self, mouse_pos):
        --recorte-- 
        if boton_clicked and not self.juego_activo: 
            --recorte--  
            self.Pp.preparando_nivel() 
            self.Pp.preparando_nave() 
            --recorte--
\end{lstlisting}


También llamamos a {\ttfamily preparando\_nave()} cuando una nave es impactada, para actualizar la visualización de las imágenes
de la naves el jugador pierde una nave:


\hspace{-2.4cm}{\ttfamily \scriptsize invasión}

\vspace{-0.2cm}
\hspace{-2.1cm}{\ttfamily \scriptsize \_alien.py}\vspace{-0.84cm}
\begin{lstlisting}	
    def _nave_golpeada(self):
        """Responde a cuando la nave ha sido golpeada por un alienigena."""
        if self.estadisticas.nave_permitidas > 0:
            # Resta 1 a nave_permitidas y actualiza puntuacion
            self.estadisticas.nave_permitidas -= 1
            self.Pp.preparando_nave()
            --recorte--
\end{lstlisting}

Llamamos al método  {\ttfamily preparando\_nave()} después de disminuir el valor de {\ttfamily nave\_permitidas}, de modo que se
muestre el número correcto de naves restantes cada vez que se destruye un nave. 

La Figura \ref{numeroNaves} muestra el sistema de 
puntuación completo, con las naves restantes mostrados en la parte superior izquierda de la pantalla.


\begin{figure}[h!]
	\centering
	\includegraphics[scale=0.35]{../A Logos de Python/numeroNaves}
	\caption{Las naves disponibles aparecen}\vspace{-0.35cm}
	\caption*{en la parte superior izquierda de la pantalla.}
	\label{numeroNaves}
\end{figure}


\section{INTÉNTALO TÚ MISMO}

14-5. Puntuación más alta de todos los tiempos: la puntuación más alta se restablece cada vez que un jugador cierra
y reinicia Alien Invasion. Solucione este problema escribiendo la puntuación más alta en un archivo antes
llamando a sys.exit() y leyendo la puntuación más alta al inicializar su valor en
Estadísticas del juego.

14-6. Refactorización: busque métodos que realicen más de una tarea y
Refactorícelos para organizar su código y hacerlo eficiente. Por ejemplo, mover
parte del código en \_check\_bullet\_alien\_collisions(), que inicia una nueva
nivel cuando la flota de alienígenas ha sido destruida, a una función llamada inicio
\_nuevo nivel(). Además, mueva las cuatro llamadas a métodos separados en el método \_\_init\_\_()
en Scoreboard a un método llamado prep\_images() para acortar \_\_init\_\_(). la preparación
El método \_images() también podría ayudar a simplificar \_check\_play\_button() o start\_game()
si ya has refactorizado \_check\_play\_button().
NOTA Antes de intentar refactorizar el proyecto, consulte el Apéndice D para obtener más información.
cómo restaurar el proyecto a un estado de funcionamiento si introduce errores
mientras se refactoriza.

14-7. Ampliando el juego: piensa en una manera de expandir Alien Invasion. Para
Por ejemplo, puedes programar a los extraterrestres para que disparen balas a tu nave. Tú
También puedes agregar escudos para que tu nave se esconda detrás, los cuales pueden ser destruidos por
balas de ambos lados. O puedes usar algo como el módulo pygame.mixer
para agregar efectos de sonido, como explosiones y sonidos de disparos.


14-8. Sideways Shooter, versión final: continúa desarrollando Sideways Shooter,
utilizando todo lo que hemos hecho en este proyecto. Agrega un botón Jugar, crea el juego.
acelerar en los puntos apropiados y desarrollar un sistema de puntuación. Asegúrate de refactorizar
su código mientras trabaja y busca oportunidades

\section{Resumen}

En este capítulo, aprendió cómo implementar un botón Reproducir para iniciar una nueva
juego. También aprendió a detectar eventos del mouse y ocultar el cursor en
juegos activos. Puedes usar lo que has aprendido para crear otros botones, como un
Botón de ayuda para mostrar instrucciones sobre cómo jugar a tus juegos. tu tambien aprendiste
cómo modificar la velocidad de un juego a medida que avanza, implementar un progresivo
sistema de puntuación y mostrar información en forma textual y no textual.

\chapter{Empaquetando el juego, y creando un archivo EXE}              %   Harrison Kinsley


Si ha llegado hasta aquí, seguro logro escribir un juego con PyGame, y ahora quieres compartir y distribuir su obra maestra; la forma más sencilla de distribuir su juego es agrupar el código Python y con sus datos en un
archivo comprimido, como {\ttfamily ZIP},  {\ttfamily TAR} o  {\ttfamily GZIP}, y cargarlo en su sitio web o envíelo por correo electrónico. 
Sin embargo, el problema con este enfoque es que Python y cualquier módulo externo que utilice deben ser instalado antes de correr el juego, lo que genera dificultades para el publico común que no maneja esas técnicas.

Para distribuir su juego a una audiencia más amplia y no técnica, necesitarás empaquetar tu juego en un formato
sencillo y amigable  con  las plataformas elegidas, por ejemplo en Windows en un archivo  {\ttfamily EXE}.

En ésta sección se cubre el tema referente a empaquetar el juego y construir un archivo ejecutable en un formato que permita a  cualquier usuario instalarlo y jugarlo.

\subsection{Creando paquetes en Windows}         

Instalar un juego en Windows generalmente implica hacer doble clic en un archivo {\ttfamily EXE}, 
y éste inicia la instalación en forma automática. 

El instalador suele tener la forma de un asistente con varias páginas que muestran la licencia, acuerdo y preguntar al usuario
dónde copiar los archivos del juego y qué iconos instalar; y finaliza con un botón en la página final para iniciar la copiar de archivos y crear iconos.

Hay dos pasos necesarios para crear un instalador fácil de usar para su juego PyGame en la Plataforma Window:

Primero debe convertir su archivo principal de Python en un ejecutable, que se ejecute sin necesidad de instalar Python.

Segundo, debe utilizar el software de instalación para crear un único archivo {\ttfamily EXE} que contenga sus archivos del juego.\\


Con un instalador EXE, hace que su juego sea accesible para una audiencia más amplia, y  esencial si su juego
está destinado a ser comercial. 

Puede omitir el segundo paso si  va dirigido a un público es lo suficientemente técnico como para
descomprima un archivo {\ttfamily ZIP} y haga doble clic en un archivo {\ttfamily EXE}.


\subsection{Usando cx\_Freeze}

Para convertir un archivo Python en un archivo ejecutable, puede usar el comando {\ttfamily cx\_Freeze}, 
que es en sí mismo un módulo de Python {\ttfamily cx\_Freeze}, no forma parte de la biblioteca estándar de Python, pero se puede 
instalar fácilmente con el siguiente comando en cmd.exe/bash/terminal:  \footnote{cmd.exe, bash y la terminal son términos relacionados con la línea de comandos y la interacción con sistemas operativos:	
	El cmd.exe es el ejecutable del Símbolo del Sistema (Command Prompt) de Windows. 	
	El bash es una shell (interfaz de línea de comandos) popular en sistemas operativos como Linux y macOS.	
	La terminal es un término genérico que se refiere a una aplicación de línea de comandos que se utiliza para interactuar con un sistema operativo. En Windows, el Terminal se conoce también como cmd o Símbolo del Sistema. En Linux y macOS, la terminal es una aplicación que se utiliza para ejecutar comandos y scripts en un entorno de texto.}\\
	
{\bf{Características importantes de cx\_Freeze}}:

\begin{itemize}
	\item[$-$]  Crea ejecutables independientes a partir de scripts de Python con el mismo rendimiento. 
    \item[$-$] 	Admite versiones de Python de 3.8 a 3.12.
	\item[$-$]  Es multiplataforma y debería funcionar en cualquier plataforma en la que funcione Python.
	\item[$-$]  Utiliza una licencia derivada de  la Fundación de Software de Python.
	\item[$-$]  Tiene una documentación disponible en el repositorio oficial.
	\item[$-$]  Tiene un canal de discusión para ayuda y soporte.
\end{itemize}

La última versión estable de {\ttfamily cx\_Freeze} es la 7.1.0, lanzada el 26 de mayo de 2024. 
Esta versión incluye nuevas características como la opción {\ttfamily --zip-filename}  en {\ttfamily build\_exe}, 
soporte para {\ttfamily pyproject.toml}, creación de formatos {\ttfamily AppImage} y {\ttfamily DEB} para Linux, mejoras en {\ttfamily bdist\_mac}, nuevos {\ttfamily hooks} y soporte para Python 3.12. Muy importante  revisar la documentación oficial para obtener más información sobre la instalación y actualizaciones.



\subsubsection*{Instalación en un Entorno Virtual}

Activar el Entorno Virtual: Asegúrate de estar en el entorno virtual donde deseas instalar {\ttfamily cx\_Freeze}.
Instalar {\ttfamily cx\_Freeze} utilizando pip:
\begin{lstlisting}
pip install --upgrade cx_Freeze
\end{lstlisting}	


\subsubsection*{Instalación en un Entorno no Virtual}

Instalar {\ttfamily cx\_Freeze} utilizando {\ttfamily pip}:
\begin{lstlisting}
python -m pip install --upgrade cx_Freeze}
\end{lstlisting}	

\subsubsection*{Instalación utilizando Pipenv}

Instalar o actualizar {\ttfamily cx\_Freeze}  utilizando {\ttfamily pipenv}:
\begin{lstlisting}	
pipenv install cx_Freeze
\end{lstlisting}


Antes de crear un ejecutable para su código Python, debe escribir un archivo {\ttfamily setup.py} por convenio, y  contiene toda la 
información de su proyecto y ejecuta {\ttfamily cx\_Freeze}. \\


En mi caso instalé el módulo con el siguiente comando:   
\begin{lstlisting}
python -m pip install --upgrade cx_Freeze}
\end{lstlisting} 



Ahora creamos el ejecutable, que por convenio debe llamarse  {\ttfamily setup.py}.
Primero importamos de {\ttfamily cx\_Freeze},  {\ttfamily setup}, y   {\ttfamily Executable}, luego especificamos qué vamos a crear un ejecutable;


\hspace{-2.2cm}{\ttfamily \scriptsize  setup.py}  \vspace{-0.5cm}
\begin{lstlisting}			
import sys
from cx_Freeze import setup, Executable

ejecutables = [Executable("Invasión_Alienígena.py")]

setup(
    name="Invasión Alienígena",
    version="0.1",
    description="Un Juego",
    options={"build_exe": {"packages": ["pygame"],
                           "include_files":[ "configurar.py",
                           "stats_juego.py","puntuacion.py",
                           "nave.py",
                           "boton.py", "balas.py","alien.py",
                           "pygame.jpg","alienigena.bmp",
                           "nave.bmp"]}},
    executables=[Executable("Invasión_Alienígena.py", base="Win32GUI")]
)
\end{lstlisting}	
Y configuramos la compilación {\ttfamily setup( )}, con los parámetros de nombre, opciones y ejecutables.
Dentro de las  {\ttfamily options} (un diccionario)  tiene una clave llamada  {\ttfamily build\_exe} que se refiere a un sub-diccionario que contiene opciones específicas para el proceso de compilación, como  {\ttfamily include\_files} la clave que se refiere a los módulos  o librería de objetos del juego que construimos, y los incluimos como una lista. 

Muy importante el parámetro  {\ttfamily executables} debe ser una lista de objetos {\ttfamily  Executable}, y la base  base= {\ttfamily  "{}Win32GUI"{}}\footnote{En el contexto de {\ttfamily cx\_Freeze} el parámetro {\ttfamily  Win32GUI} se refiere al subsystem de Windows que se utiliza para crear aplicaciones de escritorio que se ejecutan en modo de ventana (GUI). Este subsystem es compatible con ambas arquitecturas, 32 bits y 64 bits.}, y todos los archivos en la misma carpeta o directorio.   Aquí puede colocar música, imágenes y cualquier otro archivo, como los módulos {\ttfamily py} que necesita su programa.\\






Bien, una vez creado el archivo {\ttfamily script setup.py},  debe ser  ejecutado  en {\ttfamily cmd.exe/bash/terminal} con el siguiente comando:
\begin{lstlisting}		
python setup.py build	
\end{lstlisting}

Si recibe el error, como que Python no es un comando reconocido, puede hacer lo siguiente:

\begin{lstlisting}		
C:/Python34/python setup.py build
\end{lstlisting}

Esta línea de código  busca y copia todos los archivos utilizados por el {\ttfamily script} en una carpeta llamada {\ttfamily build}. 
La construcción interior es otra carpeta con información de la máquina, y luego allí encontrará el archivo  {\ttfamily Invasión\_Alienígena.exe}, que inicia la simulación de la máquina en estado del Juego, así como otros archivos
necesarios para que se ejecute sin instalar Python. 

Si comprimes esta carpeta, puedes compartirla con cualquier persona, tenga o no Python o los distintos
módulos instalados y pueden ejecutar el juego, pero para darle un toque profesional también debes crear un instalador.




% balas.py: Contiene la clase Balas que representa las balas disparadas por la nave.
% boton.py: Contiene la clase Boton que representa el botón de "Play".
% configurar03.py: Contiene la clase Configuracion que almacena configuraciones del juego.
% estadisticas\_juego.py: Contiene la clase Estadisticas\_Juego que almacena estadísticas del juego.
% nave07.py: Contiene la clase Nave que representa la nave del jugador.
% puntuacion.py: Contiene la clase Puntuacion que maneja la puntuación del juego
% alien13.py: Contiene la clase Alien que representa a los alienígenas.






\subsubsection{Construyendo el instalador}

Para un estilo un poco más profesional,  puedes crear un instalador, que proporcionará un único instalador para facilitar el uso del juego de manera muy sencilla. Sin necesidad de lidiar con la descompresión con algún extractor de archivos,  simplemente haciendo doble clic en la aplicación inicia el proceso de instalación. Para ello, se mantiene todo igual excepto cuando vas a construir el archivo, debe usar el siguiente código:

\begin{lstlisting}		
python setup.py bdist_msi	
\end{lstlisting}	

%	python setup.py install

Esto crea un directorio {\ttfamily dist}; y dentro  está el instalador.  De esta forma puedes distribuir sólo este instalador a tu amigos que luego podrán instalarlo y jugar. 


Tenga en cuenta que si tiene una versión de Windows de 32 bits, esto supone un
instalador de 32 bits para el juego. Windows de 32 o 64 bits puede instalarlo y ejecutarlo. 
Si está operando con un sistema de 64 bits versión de Windows, su distribución será de 64 bits y
 no podrá ejecutarse en máquinas de 32 bits.\\

	



Puede parecer bastante simple, pero se sabe que construir  un ejecutable causa muchos dolores de cabeza.
La documentación para obtener más información sobre cómo funciona el módulo  {\ttfamily cx\_Freeze}, la encuentra en  
\href{https://cx-freeze.readthedocs.io/en/stable/}{\color{vino}https://cx-freeze.readthedocs.io/en/stable/}.
Un conjunto de  módulos y funciones para crear ejecutables compatibles con varias plataformas, incluyendo Windows, Mac OS X y Linux.
La documentación también cubre cómo utilizar  {\ttfamily cx\_Freeze} con otras herramientas y tecnologías, como PyInstaller, Setuptools y Virtualenv. 

En general, la documentación de {\ttfamily cx\_Freeze} es una guía detallada sobre cómo utilizar el herramienta para congelar scripts de Python en ejecutables y es un recurso valioso para cualquier persona que busque crear ejecutables de Python utilizando {\ttfamily cx\_Freeze}.

La documentación proporciona información detallada sobre cómo utilizar  {\ttfamily cx\_Freeze}, incluyendo:

Iniciación: Esta sección proporciona una visión general de cómo utilizar  {\ttfamily cx\_Freeze}, incluyendo cómo instalarlo y cómo crear un ejecutable básico.

Uso básico: Esta sección explica cómo utilizar  {\ttfamily cx\_Freeze} para congelar un script de Python en un ejecutable.

Uso avanzado: Esta sección cubre cómo utilizar  {\ttfamily cx\_Freeze} para congelar un script de Python en un ejecutable con opciones avanzadas, como personalizar el icono del ejecutable y crear un archivo setup.exe.


Solución de problemas: Esta sección proporciona información sobre cómo solucionar problemas comunes con {\ttfamily cx\_Freeze}, como errores y problemas de compatibilidad.

Preguntas frecuentes: Esta sección responde a preguntas frecuentes sobre  {\ttfamily cx\_Freeze}.\\




	
La siguiente tabla muestra una lista de  algunas  métodos y funciones de {\ttfamily cx\_Freeze} que podrían ayudarte para crear ejecutables en Windows:


\begin{center}
	\noindent Table \ref{cx_Freeze}. Lista de métodos del módulo  {\ttfamily cx\_Freeze}:\\
	{\footnotesize{{%\scriptsize 
				\hspace{-0.7cm}	{\scriptsize
					\begin{tabular}{l l l }	% white
						\multicolumn{1}{|>{\columncolor{vino}} l }{\color{white}\bf  }& 
						\multicolumn{1}{|>{\columncolor{vino}} l }{\color{white}\bf  Descripción}\\ 
						pip install cx\_Freeze  &Comando para instalar cx\_Freeze usando   pip.\\
					  	                        &Requiere tener Python agregado al PATH durante la instalación\\
						setup()	                & Función principal para configurar el proceso de empaquetado.\\
					                            & Recibe varios argumentos como name, version, description, author, url, etc\\
					    Executable()            & Clase para definir un ejecutable a partir de un script Python.\\
					                            & Recibe argumentos como script, base, icon, etc.\\
				     	build\_exe              &Comando para construir el ejecutable a partir de la configuración definida \\
					                            &en setup.py. Define el directorio de salida para el ejecutable. \\					   
					    include\_files          & Permite incluir archivos adicionales en el ejecutable, como imágenes\\
                                                & o archivos de configuración\\
                        include\_msvcr          &Incluye la biblioteca de runtime de Microsoft Visual C++ (MSVCR)\\
                                                & en el ejecutable\\                     
                        msi\_name&	Define el nombre del instalador MSI\\
                        msi\_version	&Define la versión del instalador MSI\\
                        msi\_description&	Define la descripción del instalador MSI\\
                        msi\_author&	Define el autor del instalador MSI\\
                        msi\_url&	Define la URL del autor del instalador MSI\\                   	                  
						\arrayrulecolor{vino}\hline    
					\end{tabular}\label{cx_Freeze} 		
					%  	\vspace{0.2cm} 
	}	}}} 
\end{center}

Estos métodos y funciones  permiten personalizar el proceso de empaquetado y crear un instalador   MSI
("{}Microsoft Software Installer{}") más avanzado.\\




Como una técnica alternativa  para generar un ejecutable de Python  se tiene a {\ttfamily PyInstaller}, que se puede usar como:
\begin{lstlisting}		
pyinstaller --windowed --onefile --icon= nombre.ico script.py		
\end{lstlisting}	
{\ttfamily --windowed} oculta la consola,
{\ttfamily --onefile} crea un archivo único,
{\ttfamily --icon} agrega un ícono al ejecutable.






%    Troubleshooting: 

\subsubsection*{Ventajas de  cx\_Freeze}

Facilidad de uso:  {\ttfamily cx\_Freeze} es fácil de instalar y utilizar, especialmente para usuarios con experiencia en Python.
Compatibilidad multiplataforma: Puede generar ejecutables para Windows, Linux y Mac.

Flexibilidad: Permite personalizar la configuración de los ejecutables, como la inclusión de iconos y la ocultación de la consola.
Soporte para instaladores: Permite crear instaladores de Windows (MSI) para facilitar la distribución de aplicaciones.


\subsubsection*{Ventajas  PyInstaller}

Soporte amplio: {\ttfamily PyInstaller} es ampliamente utilizado y tiene una gran comunidad que lo mantiene actualizado y mejorado.

Flexibilidad: Permite generar ejecutables que incluyan todas las dependencias necesarias para su funcionamiento.

Opción de crear un único archivo ejecutable: Permite crear un archivo ejecutable único que incluya todas las dependencias necesarias.
Soporte para iconos: Permite agregar iconos a los ejecutables.

Desventajas de {\ttfamily cx\_Freeze}:

No tan ampliamente utilizado como {\ttfamily PyInstaller}: Esto puede significar que no hay tanto apoyo y actualizaciones como PyInstaller.
Puede requerir configuración adicional: Para generar ejecutables que incluyan todas las dependencias necesarias, puede ser necesario realizar configuración adicional.

{\ttfamily PyInstaller} puede generar ejecutables grandes, si se incluyen todas las dependencias necesarias.

Puede requerir configuración adicional: Para generar ejecutables que incluyan todas las dependencias necesarias,
puede ser necesario realizar configuración adicional.\\

Ambas herramientas son útiles para empaquetar scripts de Python en ejecutables, pero PyInstaller es más ampliamente utilizado y tiene una mayor comunidad detrás.



